%%% POST-RECORD v1 BEGIN
% format: hexagon-post-format/1
% record_id: S6_07_PAPER
% version: 1
% date: 2026-11-16
% title: "The (3,4,infinity) family of complex structures on S^6: periods, moduli and the ends of the disc"
% kind: theorem
% claim_status: "T, K and A per result; see the concordance in Appendix C"
% statement_plain: "The cusp-normalized period has maximum -(3/pi)log(6), only at t=1; the known punctured-disc parameter w has exact radius 216 and Albanese nome iw/216."
% statement_tex: "M=-(3/\pi)\log 6; 0<|w|<216; q_2=iw/216."
% corrects: "Engel Remark 8.1; AC26 Remarks 1.3 and 9.21"
% refutes: "Engel Conjecture 8.2 as worded"
% answers: "Engel Remark 8.4, Klingler's question"
% known_before: "Construction: AC26 and E26; standard-data 2Z, intrinsic pencil, tangent cohomology, Kuranishi germ and full automorphisms: arXiv:2610.00250v1."
% new_here: "Closed periods and exact radius, all discrete data, the specified cohomology computations, real structures, monodromy and ends; prior overlaps credited in Appendix C."
% main_statements: ["FAM C.1.1", "FAM C.2.3", "ARI C.8.5", "INV A3.1", "FAM F.1.1", "PUSH D.1.2", "PUSH A.3.2"]
% hard_step: "FAM C.2.3 and PUSH E.4: interval certificates; ARI C.8.1-C.8.5: ladder"
% external_results: "Exact locators and opened-source checks in the Verification record and Appendix C"
% produced_by:
%   - {role: source proofs, model: "GPT-6 Astra", effort: "ChatGPT Pro mode", date: "2026-11-12/2026-11-16"}
%   - {role: workbench proof, model: "GPT-5.6 Sol and GPT-6 Astra", effort: "Codex reasoning effort ultra", date: "2026-09"}
%   - {role: assembly and source audit, model: "GPT-6 Astra", effort: "ChatGPT Pro mode", date: "2026-11-16"}
% checks:
%   - {by: "Claude Fable 5.2", effort: high, type: "F,D,C,S,R", scope: "Exactly the supplied review and intake records"}
%   - {by: "Claude Fable 5.2", effort: max, type: "R,D,S", scope: "Exactly RESPONSE_EXPERT.md"}
%   - {by: "Claude Opus 5.5", effort: max, type: "S,D", scope: "Checks reported in the supplied status files"}
%   - {by: "GPT-6 Astra", effort: "ChatGPT Pro mode", type: S, scope: "Opened citations, three-source credit search and source-to-print comparison"}
% human_understanding: none
% text_origin: primarily AI-generated
% formalization: none
% open_points: ["algebraic dimension zero", "real S3 knot", "replacement for the zero-end limit", "other families' monodromy"]
% not_claimed: ["human verification", "fresh execution of missing interval data", "proofs of the A-items in this paper"]
% literature_search: {date: "2026-11-16", scope: "supplied primary sources and opened primary web sources", outcome: "Section-by-section overlaps and limitations recorded in Appendix C"}
% licence: CC0-1.0
% provenance:
%   prompt: "PROMPT(20261116-143923).md (supplied separately)"
%   return: "S6_07_PAPER_paper.tex"
%   reports: ["S6_07_PAPER_report.md", "records/"]
%   register: "S6_07_PAPER_source_table.csv"
%   repository: "https://github.com/KokunoYumeto/s6-complex-structure-workbench"
% hexagon:
%   subject: math.CV
%   cross_lists: [math.AG, math.DG, math.NT]
%   msc: [32G05, 32J17, 14K25, 11F03]
%   tags: [periods, moduli, complex-six-sphere, cohomology, monodromy]
%   tools: ["GPT-6 Astra, ChatGPT Pro mode", "Claude Fable 5.2, effort high", "Claude Fable 5.2, effort max", "Claude Opus 5.5, effort max", "GPT-5.6 Sol and GPT-6 Astra, Codex reasoning effort ultra", "LuaLaTeX"]
%   comments: "AI-produced; no human checking and no formalization. Status differs by result."
%   provenance_statement: "Source-preserving assembly of six supplied returns, with reviews and source audits."
%%% POST-RECORD v1 END
\documentclass[11pt,a4paper]{article}
\usepackage[T1]{fontenc}
\usepackage{lmodern}
\usepackage[margin=22mm]{geometry}
\usepackage{amsmath,amssymb,amsthm,mathtools,mathrsfs}
\usepackage{booktabs,tabularx,longtable,array,multirow}
\usepackage{enumitem}
\usepackage{graphicx,xcolor}
\usepackage{microtype}
\PassOptionsToPackage{hyphens}{url}
\usepackage[hidelinks,unicode]{hyperref}
\usepackage[nameinlink,noabbrev]{cleveref}
\usepackage{newunicodechar}
\newunicodechar{−}{\ensuremath{-}}
\newunicodechar{∞}{\ensuremath{\infty}}
\newunicodechar{Θ}{\ensuremath{\Theta}}
\newunicodechar{χ}{\ensuremath{\chi}}
\newunicodechar{μ}{\ensuremath{\mu}}
\newunicodechar{β}{\ensuremath{\beta}}
\newunicodechar{η}{\ensuremath{\eta}}
\newunicodechar{τ}{\ensuremath{\tau}}
\newunicodechar{λ}{\ensuremath{\lambda}}
\newunicodechar{Δ}{\ensuremath{\Delta}}
\newunicodechar{π}{\ensuremath{\pi}}
\newunicodechar{²}{\ensuremath{{}^2}}
\newunicodechar{³}{\ensuremath{{}^3}}
\newunicodechar{⁶}{\ensuremath{{}^6}}
\newunicodechar{⁷}{\ensuremath{{}^7}}
\newunicodechar{¹}{\ensuremath{{}^1}}
\newunicodechar{₁}{\ensuremath{{}_1}}
\newunicodechar{₂}{\ensuremath{{}_2}}
\newunicodechar{₃}{\ensuremath{{}_3}}
\newunicodechar{→}{\ensuremath{\to}}
\newunicodechar{≠}{\ensuremath{\ne}}
\newunicodechar{≤}{\ensuremath{\le}}
\newunicodechar{≥}{\ensuremath{\ge}}
\newunicodechar{∪}{\ensuremath{\cup}}
\newunicodechar{∩}{\ensuremath{\cap}}
\newunicodechar{⊕}{\ensuremath{\oplus}}
\newunicodechar{⊗}{\ensuremath{\otimes}}
\newunicodechar{∂}{\ensuremath{\partial}}
\newunicodechar{∈}{\ensuremath{\in}}
\newunicodechar{ℤ}{\ensuremath{\mathbb Z}}
\newunicodechar{ℂ}{\ensuremath{\mathbb C}}
\newunicodechar{ℙ}{\ensuremath{\mathbb P}}
\newunicodechar{ℝ}{\ensuremath{\mathbb R}}
\newunicodechar{⋅}{\ensuremath{\cdot}}
\newunicodechar{√}{\ensuremath{\surd}}
\newunicodechar{±}{\ensuremath{\pm}}
\newunicodechar{≅}{\ensuremath{\simeq}}
\newunicodechar{∗}{\ensuremath{*}}
\newunicodechar{‖}{\ensuremath{\Vert}}
\newunicodechar{Λ}{\ensuremath{\Lambda}}
\newunicodechar{′}{\ensuremath{\prime}}
\newunicodechar{℘}{\ensuremath{\wp}}
\newunicodechar{₀}{\ensuremath{{}_0}}
\newunicodechar{₄}{\ensuremath{{}_4}}
\newunicodechar{⁴}{\ensuremath{{}^4}}
\newunicodechar{⁸}{\ensuremath{{}^8}}
\newunicodechar{ⁿ}{\ensuremath{{}^n}}
\newunicodechar{ᶜ}{\ensuremath{{}^c}}
\newunicodechar{α}{\ensuremath{\alpha}}
\newunicodechar{ψ}{\ensuremath{\psi}}
\newunicodechar{ρ}{\ensuremath{\rho}}
\newunicodechar{κ}{\ensuremath{\kappa}}
\newunicodechar{𝒩}{\ensuremath{\mathcal N}}
\newunicodechar{𝓛}{\ensuremath{\mathcal L}}
\newunicodechar{𝒜}{\ensuremath{\mathcal A}}
\newunicodechar{𝒲}{\ensuremath{\mathcal W}}
\newunicodechar{𝓔}{\ensuremath{\mathcal E}}
\newunicodechar{𝒮}{\ensuremath{\mathcal S}}
\newunicodechar{ℛ}{\ensuremath{\mathcal R}}
\newunicodechar{∏}{\ensuremath{\prod}}
\newunicodechar{∧}{\ensuremath{\wedge}}
\newunicodechar{̄}{\ensuremath{\bar{\vphantom x}}}
\emergencystretch=2em
\allowdisplaybreaks[2]
\setlength{\parskip}{3pt}
\setlist{topsep=3pt,itemsep=2pt,parsep=0pt}
\theoremstyle{plain}
\newtheorem{theorem}{Theorem (T)}[section]
\newtheorem{proposition}[theorem]{Proposition (T)}
\newtheorem{lemma}[theorem]{Lemma (T)}
\newtheorem{corollary}[theorem]{Corollary (T)}
\theoremstyle{definition}
\newtheorem{cited}[theorem]{Known result (K)}
\newtheorem{announcement}[theorem]{Announcement (A)}
\newtheorem{computed}[theorem]{Certified computation (T)}
\newtheorem{refuted}[theorem]{Refuted statement (false)}
\newtheorem{definition}[theorem]{Definition}
\newtheorem{remark}[theorem]{Remark}
\newtheorem{example}[theorem]{Example}
\newtheorem{openq}[theorem]{Open question}
\newcommand{\uses}[1]{\par\noindent{\small\textsf{Uses:} #1}\par\smallskip}
\newcommand{\step}[2]{\par\smallskip\noindent\textbf{Step #1} (#2).\ \ignorespaces}
\newcolumntype{L}[1]{>{\raggedright\arraybackslash}p{#1}}
\newcolumntype{Y}{>{\raggedright\arraybackslash}X}
\newcommand{\C}{\mathbb C}
\newcommand{\R}{\mathbb R}
\newcommand{\Q}{\mathbb Q}
\newcommand{\Z}{\mathbb Z}
\newcommand{\N}{\mathbb N}
\newcommand{\Pp}{\mathbb P}
\newcommand{\HH}{\mathbb H}
\newcommand{\OO}{\mathcal O}
\newcommand{\partt}{\mathrm{part}}
\newcommand{\gh}{\widehat\gamma}
\newcommand{\uh}{\widehat u}
\newcommand{\wh}{\widehat w}
\newcommand{\dhat}{\widehat\delta}
\newcommand{\RP}{\mathbb R\mathbb P}
\newcommand{\eps}{\varepsilon}
\newcommand{\dbar}{\bar\partial}
\newcommand{\dd}{\mathrm d}
\newcommand{\id}{\mathrm{id}}
\newcommand{\Hh}{\mathbb H}
\newcommand{\cO}{\mathcal O}
\newcommand{\cV}{\mathcal V}
\newcommand{\cL}{\mathcal L}
\newcommand{\cH}{\mathcal H}
\newcommand{\cB}{\mathfrak B}
\DeclareMathOperator{\Aut}{Aut}
\DeclareMathOperator{\Alb}{Alb}
\DeclareMathOperator{\NS}{NS}
\DeclareMathOperator{\Pic}{Pic}
\DeclareMathOperator{\End}{End}
\DeclareMathOperator{\Hom}{Hom}
\DeclareMathOperator{\Ext}{Ext}
\DeclareMathOperator{\Tor}{Tor}
\DeclareMathOperator{\Sym}{Sym}
\DeclareMathOperator{\Spec}{Spec}
\DeclareMathOperator{\Proj}{Proj}
\DeclareMathOperator{\rank}{rank}
\DeclareMathOperator{\coker}{coker}
\DeclareMathOperator{\im}{im}
\DeclareMathOperator{\tr}{tr}
\DeclareMathOperator{\Pf}{Pf}
\DeclareMathOperator{\GL}{GL}
\DeclareMathOperator{\SL}{SL}
\DeclareMathOperator{\Sp}{Sp}
\DeclareMathOperator{\MW}{MW}
\DeclareMathOperator{\Diff}{Diff}
\DeclareMathOperator{\Supp}{Supp}
\DeclareMathOperator{\Res}{Res}
\DeclareMathOperator{\Div}{Div}
\DeclareMathOperator{\ord}{ord}
\DeclareMathOperator{\lcm}{lcm}
\DeclareMathOperator{\sgn}{sgn}
\DeclareMathOperator{\Log}{Log}
\DeclareMathOperator{\diag}{diag}
\DeclareMathOperator{\Vol}{Vol}
\DeclareMathOperator{\Rea}{Re}
\DeclareMathOperator{\Ima}{Im}
\title{The $(3,4,\infty)$ family of complex structures on $S^6$:\\periods, moduli and the ends of the disc}
\author{AI-produced by GPT-6 Astra (OpenAI, ChatGPT Pro mode)}
\date{16 November 2026\qquad Version 1\qquad CC0 1.0}
\begin{document}
\maketitle

\begin{abstract}
For the cusp-normalized $(3,4,\infty)$ family on $S^6$, the particular period is $\frac{6i}{\pi}\operatorname{Log}\frac{\vartheta_1(\mu,\tau)}{2\sqrt3\,\eta_D(\tau)}$; its real determinant function has maximum $-(3/\pi)\log6$, attained only at the multiplicity-four fibre. The known punctured-disc parameter $w=e^{-\pi ic}$ therefore has exact radius $216$, and the Albanese nomes are $iw/216$ and $-w^2/48^3$. The period construction also gives the stated all-data invariants, Hodge and tangent-twist calculations, real structures, index-$12$ monodromy group, and descriptions of the special fibres. At the outer circle the family has no extension of the specified kind; Engel's proposed central space has virtual Euler number $0$ instead of $2$. Theorems are assembled from proofs by GPT-6 Astra (ChatGPT Pro mode), with the bounded readings and computations of Claude Fable 5.2 (efforts high and max) and Claude Opus 5.5 (effort max) specified per result; no human has checked them. Announcements retain repository proofs, including the zero-end non-filling, trivial smooth monodromy and component claims. The construction is credited to Claude (effort not recorded) and Alp\"oge and to Engel; the intrinsic pencil, standard-data $2\mathbb Z$ theorem, tangent cohomology, Kuranishi germ and full automorphisms are credited to Du, Guo, Kusakabe and Xie, arXiv:2610.00250v1.
\end{abstract}

\section*{Result}
\subsection*{Statement}
\noindent The rows below locate the main groups of results. Each individual statement, its status, its proof and its kind have a separate row in Appendix~\ref{app:audits}; its run identifier is retained in its heading. Here T means a theorem with a printed proof, K a known theorem re-derived in Appendix~\ref{app:known}, and A an announcement with proof in the repository.
\begin{center}\small
\begin{tabularx}{\linewidth}{@{}lYll@{}}
\toprule Identifier & Statement & Status & Proof / kind\\\midrule
FAM B.1--B.3 & Dictionary $u=w^2$; all discrete data; $2kl\mathbb Z$ in the stated model & T & \S\S1--2 / theorem\\
FAM C; ARI C & Closed periods, division values, maximum and radius $216$ & T & \S\S3--4 / theorem, certificate\\
INV A--D; CDP C5b & Hodge, Fr\"olicher, Bott--Chern, Aeppli and tangent twists & T & \S\S5--6 / theorem\\
FAM B.4, D, F & Real structures, monodromy, special fibres and subvarieties & T & \S\S7--9 / theorem\\
PUSH D.1, A.3 & Outer non-extension; conjectured centre has the wrong Euler number & T & \S10 / counterexample\\
DGKX v1 & Intrinsic pencil, $2\mathbb Z$, deformations, $\Aut=\mathbb C^*$ & K & App. B / new proof\\
Remaining selected rows & Exact announcement statements and their mechanisms & A & \S\S2--7,10--11 / theorem\\
\bottomrule
\end{tabularx}
\end{center}
For example, $c=0$ gives $w=1$ and Albanese nome $i/216$. The component and zero-end assertions remain announcements; the dictionary for higher multiplicities is not included in the theorem for the stated model.

\subsection*{Status}
The same model, GPT-6 Astra (ChatGPT Pro mode), produced the six mathematical returns and assembled this paper. Claude Fable 5.2 (effort high) reviewed FAM and read the explicitly listed parts of the intakes; Claude Fable 5.2 (effort max) supplied the expert readings and derivations, and Claude Opus 5.5 (effort max) supplied the main-session checks recorded there. The historical program checks include the exact recomputations, the full interval re-execution and the Arb certification. This assembly adds source comparisons and opened-citation checks, not a second independent mathematical review. No human has checked any of it.
\begin{center}\small
\begin{tabularx}{\linewidth}{@{}lY@{}}
\toprule Kind & Theorem, computation, correction, counterexample and new proof, assigned per result.\\
Claim & T: proved as in the supplied text; K: known and re-derived; A: announced.\\
Checks & One row per recorded check in Appendix~\ref{app:verification}; scope and limitations are explicit.\\
Human understanding & none.\\
Formalization & none.\\
Open points & Four question groups in \S\ref{sec:open}; source discrepancies in Limits and open points.\\\bottomrule
\end{tabularx}
\end{center}

\subsection*{Known and new}
The section-by-section credit table in Appendix~\ref{app:credits} starts from the expert's supplied table and incorporates the three-source audit. We credit \cite[Main Theorem; Sections 3--9]{AC26} for the construction, canonical bundle and the stated cohomological, fibre and automorphism results, \cite[Sections 2--8]{E26} for the geometric construction, and \cite[Theorems 1.3, 1.4, 7.1 and 7.2]{DGKXv1} for the known deformation and period-identification results. In particular, the N\'eron--Severi equation of \cite[Lemma 9.2]{AC26} and the elliptic-logarithm reading of \cite[Remark 3.12]{AC26} are credited there, although their source-run proofs are retained. The source audit used the supplied full TeX and transcription and the primary references opened on 16 November 2026; negative searches are bounded searches, not assertions of absence from all literature.

The constructed $X_c$ has algebraic dimension $1$ and $c_3(X_c)=2$. This contradicts Corollary~2.3 of the corrigendum of Campana, Demailly and Peternell \cite[Corollary 2.3]{CDP20}, which asserts algebraic dimension $0$ for a complex structure on $S^6$. The counterexample is due to \cite[Main Theorem; Section 10]{AC26}, whose Section~10 locates two failing steps for $X$. Where the argument fails in general, and what is true in its place, is the subject of a separate note.

\subsection*{Proof map}
\textbf{Idea.} The marked period lattice carries the affine deck data, the theta norm and the root-base direct images. Its norm determines the admissible disc; its two finite fibres retain the Albanese nomes; its cusp supplies the non-normal fibre used in the cohomology calculations. The exact period identification is a known input.
\begin{enumerate}[leftmargin=*,label=\arabic*.]
\item Fix the marked construction and dictionary; use the norm obstruction for affine data (\S\S1--2, Appendix B).
\item Compute the theta expression and division values; evaluate the maximum by the certificates and ladder (\S\S3--4).
\item Calculate direct images on the root base and pass to the requested cohomology (\S\S5--6).
\item Read real forms, monodromy and fibres from the same lattice, and test the boundary and central model (\S\S7--10).
\end{enumerate}
\textbf{Hard step.} The exact maximum uses the two interval certifications as proof of record. The supplied review records a full re-execution of the $19\,701$ norm boxes and a derivation of the local estimate; the ladder's displayed inequalities were checked, but ARI C.8.1--C.8.3 were not read line by line in that intake. The certificate files were not supplied to this assembly.

\textbf{Reading routes.} The results table and credit table give the claims and their scope. A reader using the known period-identification theorem can read \S\S1--4 and then the relevant application. A full check follows the dependency rows in Appendix~\ref{app:dependencies}, including the printed re-derivations in Appendices A--B. The principal external inputs, with exact locators and S-checks, are the cited construction, the known DGKX theorems, the Mordell--Weil height formula, proper base change and relative duality, the structure theorem for double complexes, and the extension results used in \S10.
\clearpage


% ===== root source-preserving fragment =====

\section{The family and its geometric parameter}\label{sec:family}
\subsection{Conventions}
\begin{definition}[FAM conventions]\label{def:conventions}
We use $B=\Pp^1_t$, $p_1=0$, $p_2=1$, $p_0=\infty$, and
$B^\circ=B\setminus\{0,1,\infty\}$. Base meridians are counterclockwise
in a local holomorphic coordinate. The analytic generators are clockwise
root deck transformations, $s_j(g_jz)=e^{-2\pi i/m_j}s_j(z)$, with
$g_1g_2g_0=1$. A positive base loop ends at the inverse deck translate;
identifying its endpoint with the initial fibre applies $g_j$, giving
homology monodromy $A_j$. The retained multiplicities are $m_1=3,m_2=4$;
$f^*p_j=m_jS_j$ includes the nonreduced structure, whereas $S_j$ is reduced.
The cusp fibre $W=f^{-1}(\infty)$ is reduced.
\[
 \Lambda=\Z\langle\gh,\uh,\wh,\dhat\rangle,\qquad
 \Pi_c=\begin{pmatrix}6\mu&\tau&1&0\\\beta^{\partt}+c&\mu&0&1\end{pmatrix},
 \qquad w=e^{-\pi ic}.
\]
We fix $(\beta^{\partt}+\tau)(p_0)=0$. Write
$D_{\partt}=\Im\beta^{\partt}-6(\Im\mu)^2/\Im\tau$ and
$M=\sup D_{\partt}$. Engel's line bundle is
$M_S=\OO_S(P-O)$; his indices $0,1,\infty$ correspond to our $1,2,0$.
In his common integral basis, the columns of the following matrix are
our ordered basis vectors:
\[
 P_{\rm mark}=\begin{pmatrix}
 0&-2&-1&-2\\2&-1&0&-1\\0&2&1&3\\-1&0&0&0
 \end{pmatrix},\quad
 T_0P_{\rm mark}=P_{\rm mark}A_1,\quad
 T_1P_{\rm mark}=P_{\rm mark}A_2,\quad
 \delta_{\rm E}=P_{\rm mark}\dhat.
\]
For our finite twist vectors $v_1,v_2$, his positive-root affine vectors
are $a_0=-P_{\rm mark}v_1/3$, $a_1=-P_{\rm mark}v_2/4$.
In additive fibre coordinates $(\xi,\eta)$, put $v=e^{2\pi i\eta}$.
The first period acts by
$(\xi,v)\mapsto(\xi+6\mu,e^{2\pi i\beta_c}v)$; Engel's contracting
generator is its inverse, over $-6P$ with multiplier $e^{-2\pi i\beta_c}$.
We choose his reference linearization to be this transported inverse
period lift at $c=0$. Then $\lambda_0=1$ and $u=w^2$.
The comparison maps and the effect of changing this reference are
proved below. Homology has integral coefficients unless stated
otherwise; tori and the threefold have their complex orientations.

\end{definition}
\uses{The period construction of \cite[Sections 2--6]{AC26}; \cite[Sections 2--5]{E26}.}

Throughout this section $\mathcal P=\{c:\operatorname{Im}c<-M\}$, and $X_c$ means the actual glued space with $v_1=\varepsilon=(1,2,-4,0)^t$, $v_2=-\varepsilon'=(-1,-3,3,0)^t$ and the untwisted cusp collar. We write $d=\widehat\delta$, $e=(0,1)^t$, $E=d\gamma$, and use column vectors. Our peripheral generators are clockwise: $s_j(g_jz)=e^{-2\pi i/m_j}s_j(z)$ and $g_1g_2g_0=1$. Thus $s(g_0z)=s(z)-1$ at the cusp. All period comparisons use
\[
 A_1=\begin{pmatrix}1&0&0&0\\6&0&1&0\\-6&-1&-1&0\\-2&1&0&1\end{pmatrix},\quad
 A_2=\begin{pmatrix}1&0&0&0\\0&0&-1&0\\-6&1&0&0\\3&0&1&1\end{pmatrix},\qquad
 R_g(z)\Pi_c(z)=\Pi_c(gz)A_g.
\]
Here $R_1=\left(\begin{smallmatrix}-1/\tau&0\\(1-\mu)/\tau&1\end{smallmatrix}\right)$, $R_2=\left(\begin{smallmatrix}1/\tau&0\\-\mu/\tau&1\end{smallmatrix}\right)$, and $R_0=I$. The distinction between source and target markings in each equation below is retained by its displayed period equation.


The construction, the smooth manifold $S^6$, its algebraic reduction, its canonical bundle and the connected automorphism group are due to \cite[Main Theorem; Sections 3--9]{AC26}. The rational elliptic surface, its line bundle and inverse-translation quotient are those of \cite[Section 2]{E26}. The equality $u=w^2$ uses the transported reference in Definition~\ref{def:conventions}; Engel's parameter has no absolute scale.

\subsection{The intrinsic pencil (known)}
The following result is printed with status K as prescribed. The intrinsic-pencil and base-preservation assertions for the standard data are \cite[Theorem 7.1]{DGKXv1}; the canonical bundle is \cite[Main Theorem (4); Theorem 9.1]{AC26}. The all-affine-data scope and direct-image formulation are retained from FAM A.1.1; the scope of the cited theorem is distinguished in the source table.
\begin{cited}[FAM A.1.1; K]\label{fam:A.1.1}
For the constructed holomorphic map $f_c:X_c\to\mathbb P^1$, and also for every admissible choice of its affine translations and cusp collar, we have
\[
 \omega_{X_c}\simeq\mathcal O_{X_c}(-2S_2),\qquad
 f_{c*}\mathcal O_{X_c}(kS_2)=\mathcal O_B(\lfloor k/4\rfloor p_2)\quad(k\in\mathbb Z).
\]
The latter equality is an equality of subsheaves of meromorphic functions under pullback. The anticanonical section ring is $\mathbb C[U,V]$, with $\deg U=1$, $\deg V=2$. The complete degree-two map is $f_c$ followed by a projective coordinate isomorphism of its base. Every biholomorphism between two such threefolds is over the identity of $B$.
\end{cited}
\uses{Definition~\ref{def:conventions}; \cite[Theorem 7.1]{DGKXv1}; \cite[Theorem 9.1]{AC26}.}
\begin{cited}[DGKX Theorem 7.2; K]\label{dgkx:all-maps-factor}
For every $c\in\mathcal P$, every holomorphic map
$h:X_c\to\mathbb P^1$ factors uniquely as
$h=\varphi\circ f_c$ with a holomorphic map
$\varphi:\mathbb P^1\to\mathbb P^1$.
\end{cited}
\uses{\cite[Theorem 7.2, \texttt{thm:all-maps-factor}]{DGKXv1}.}

\subsection{The elliptic surface and its reference lift}
\begin{proposition}[FAM B.1.1; T]\label{fam:B.1.1}
Let $E_\tau=\mathbb C/(\mathbb Z+\tau\mathbb Z)$, $P=[\mu]$, and $L_\mu=\mathcal O_{E_\tau}(P-O)$. The holomorphic map
\[
 [\xi,\eta]\longmapsto[\xi,e^{2\pi i\eta}]
\]
identifies $\mathbb C^2/\Pi_c\Lambda$ with the quotient of $L_\mu^\times$ by the lift
$A(\xi,v)=(\xi+6\mu,e^{2\pi i\beta_c}v)$. Engel's generator, in Definition 2.3 of the attached paper, is its inverse, with translation $-6P$ and multiplier $e^{-2\pi i\beta_c}$. The section $P=[\mu]$ is a generator of the Mordell--Weil group of the stated pointed Weierstrass surface.
\end{proposition}
\uses{Definition~\ref{def:conventions}; FAM A.3.1 and preceding dictionary identities; \cite[Remark 3.12]{AC26}; \cite[Lemma 2.1; Corollary 2.2; Definition 2.3]{E26}.}
\begin{proof}
\par\smallskip\noindent\textbf{Step 1.} First quotient by columns two, three and four. Exponentiation kills column four; columns three and two then give
\[
 (\xi,v)\sim(\xi+1,v)\sim(\xi+\tau,e^{2\pi i\mu}v).
\]
For the odd theta function with periods $1,\tau$, the quotient $\theta(\xi-\mu)/\theta(\xi)$ is a meromorphic section of this line bundle: its factors are $1,e^{2\pi i\mu}$ and its divisor is $P-O$. Sending its divisor section to this quotient is the asserted line-bundle isomorphism. Quotienting the remaining first column gives exactly $A$. On sufficiently small simply connected charts the logarithm of $v$ is a holomorphic inverse; changes of logarithm are column-four periods. Thus the displayed map is a biholomorphism of the two quotient complex spaces, with its inverse explicitly defined by local logarithms. Choosing $A^{-1}$ changes the generator of the same infinite cyclic group and gives precisely Engel's declared inverse-translation convention.

\par\smallskip\noindent\textbf{Step 2.} For clarity the elliptic surface in this dictionary is the displayed cubic, rather than only a curve with the same $j$-value. Let $f_{\rm sc}$ be the branch satisfying
\[
 f_{\rm sc}^4=\frac{12t^3(t-1)}{g_2(\tau)},\qquad
 f_{\rm sc}^6=-\frac{8t^4(t-1)^2}{g_3(\tau)},\qquad
 f_{\rm sc}/t\longrightarrow i\sqrt3/\pi.
\]
The identities $g_2^3/(g_2^3-27g_3^2)=t$ make these equations compatible. Their zero and pole orders give a holomorphic branch on the period cover; the indicated cusp sign fixes it. The map $\xi\mapsto(f_{\rm sc}^2\wp_\tau(\xi),\frac12f_{\rm sc}^3\wp'_\tau(\xi))$ is a pointed isomorphism to
$y^2=x^3-3t^3(t-1)x+2t^4(t-1)^2$. Its factors are $f_{\rm sc}(g_1z)=-f_{\rm sc}/\tau$, $f_{\rm sc}(g_2z)=f_{\rm sc}/\tau$; therefore this isomorphism descends over $B^\circ$.

\par\smallskip\noindent\textbf{Step 3.} We first verify that $P=[\mu]$ gives a rational section of the resolved elliptic surface. On $B^\circ$ its cubic coordinates descend by the displayed scale and affine laws. At the two finite centres, $\mu_j$ is not a lattice point, while $f_{\rm sc}$ has orders two and one in the root parameter. Its cubic coordinates are therefore holomorphic and invariant, hence holomorphic in the base coordinate. At the cusp $\mu$ is holomorphic and bounded, and the convergent multiplicative expansion of $\wp$ gives meromorphic cubic coordinates with finite pole orders. Even when its exponential meets one, the pole is controlled by its finite vanishing order; it cannot be identically one because the two finite values of $\mu$ differ. The coordinates are thus meromorphic on $\mathbb P^1$ and rational. The closure of its graph on the proper resolved surface is an irreducible curve; its normalization maps finitely with degree one to the smooth base, so is the base itself. This gives the unique section of that surface extending $P$.

\par\smallskip\noindent\textbf{Step 4.} We also verify that the specified section is a Mordell--Weil generator using its full logarithm. A rational section of this cubic pulls back to a holomorphic section of the good elliptic family at both finite covering points: multiply its projective coordinates by a common power to remove poles, and use the nonsingular limiting Weierstrass equation. At the cusp put $r_{\rm cusp}=1/t$, $X=r_{\rm cusp}^2x,Y=r_{\rm cusp}^3y$. Its equation is $Y^2=X^3-3(1-r_{\rm cusp})X+2(1-r_{\rm cusp})^2$. A section cannot pass through the node $(1,0)$: substitution of $X=1+O(r_{\rm cusp})$, $Y=O(r_{\rm cusp})$ leaves the coefficient $+1$ of $r_{\rm cusp}$ on the left-hand side of the defining equation. At every other central point the multiplicative uniformization is a local isomorphism. Explicitly, at $r_{\rm cusp}=0$ it is
\[
 X=1+\frac{12U}{(U-1)^2},\qquad
 Y=\frac{12\sqrt3 U(U+1)}{(U-1)^3},\qquad U=e^{2\pi i\xi},
\]
with inverse $U=(Y/(X-1)+\sqrt3)/(Y/(X-1)-\sqrt3)$ away from $O$; a local parameter at $O$ is $X/Y$. The convergent lattice expansion of $\wp$ gives the same holomorphic map for small $r_{\rm cusp}$. Hence a section has a holomorphic unit $U(r_{\rm cusp})$ and a bounded logarithm at the cusp. Lifting the section on the simply connected cover gives a global logarithm $\nu$, with this cusp branch. Its only remaining ambiguity is an integer, since a nonzero multiple of $\tau$ is unbounded.

\par\smallskip\noindent\textbf{Step 5.} Its two affine laws have integral defects:
\[
 \nu(g_1z)=-\nu/\tau+\alpha_1+\beta_1\tau(g_1z),\quad
 \nu(g_2z)=\nu/\tau+\alpha_2+\beta_2\tau(g_2z).
\]
Composing them and using cusp invariance gives $\alpha_2=-\beta_1$, $\beta_2=\alpha_1+\beta_1$. Put $n=-\alpha_1-2\beta_1$ and $a=\alpha_1+\beta_1$. Substitution shows that $n\mu+a$ has exactly these defects. Their difference is a global section of the first-coordinate sheaf $\mathcal O_B(-p_0)$ from A.3.1, so is zero. Every rational section is therefore $nP$. The same cusp argument and the distinct fixed values of $\mu$ show that $P$ has infinite order. This proves $\operatorname{MW}(S)=\mathbb ZP$.
\end{proof}
\begin{proposition}[FAM B.1.2; T]\label{fam:B.1.2}
The exact set of nonzero parameters for which every ordinary fibre and both normalized finite root fibres of Engel's period-reference quotient are compact complex two-tori is
\[
 0<|u|<e^{-2\pi M}.
\]
This is the parameter domain for the compact family with both smooth finite root models. The condition on ordinary fibres alone at equality is determined by whether a maximizing point of $D_{\partt}$ is ordinary; this boundary issue is treated in Section~\ref{sec:periods}. For an independent reference, the radius is $|\lambda_0|e^{-2\pi M}$.
\end{proposition}
\uses{Definition~\ref{def:conventions}; FAM A.3.1 and preceding dictionary identities; \cite[Remark 3.12]{AC26}; \cite[Lemma 2.1; Corollary 2.2; Definition 2.3]{E26}.}
\begin{proof}
\par\smallskip\noindent\textbf{Step 1.} The function $N(\xi,v)=|v|\exp(2\pi\operatorname{Im}\mu\,\operatorname{Im}\xi/\operatorname{Im}\tau)$ is invariant under both line-bundle period identifications. The inverse lift multiplies it by $e^{2\pi D_c}$. More directly the real determinant of the four period columns is a nonzero sign times $(\operatorname{Im}\tau)D_c$. Thus the four columns generate a full lattice exactly when $D_c\ne0$, and then the quotient is a compact torus by a closed fundamental parallelepiped. At $D_c=0$ their real span has dimension three; a rank-four lattice in $\mathbb R^4$ cannot result, so the quotient cannot be a compact complex two-torus. The function $D_{\partt}$ descends continuously to the connected base with the cusp removed and tends to $-\infty$ at the cusp. Its image is therefore $(-\infty,M]$: the maximum is attained, and the intermediate value theorem fills all smaller values. Consequently $D_c$ is nonzero at every base point exactly when $\operatorname{Im}c<-M$. Since $|u|=e^{2\pi\operatorname{Im}c}$, the assertion follows, with the indicated reference scaling.
\end{proof}
The geometric identification in this proposition is already implicit in \cite[Remark 3.12]{AC26}. Its intersection calculation re-derives the height computation of \cite[Main Theorem, entry 49; (3.2)]{OS91} and \cite[Theorem 8.6; (8.16)]{Shi90} in the given charts.
\begin{proposition}[FAM B.1.3; T]\label{fam:B.1.3}
For the original affine logarithm $\mu$, we have $\wp_\tau(\mu)=0$ identically. Equivalently the generator $P=[\mu]$ of B.1.1 is one of the two sections $P_*=(0,\sqrt2t^2(t-1))$ and $-P_*$. This assertion does not use a numerical evaluation of $\mu$ or of $\beta$.
\end{proposition}
\uses{Definition~\ref{def:conventions}; FAM A.3.1 and preceding dictionary identities; \cite[Remark 3.12]{AC26}; \cite[Lemma 2.1; Corollary 2.2; Definition 2.3]{E26}.}
\begin{proof}
\par\smallskip\noindent\textbf{Step 1.} We give the intersection calculation that identifies the generator. At $t=1$, put $v_t=t-1$ and blow up the origin. In the $v_t$-chart $x=v_t X,y=v_t Y$, the strict equation is
$Y^2=v_t X^3-3(1+v_t)^3X+2(1+v_t)^4$. Its exceptional curve is smooth, and $P_*$ meets it at $(X,Y)=(0,\sqrt2)$. In the other chart $v_t=xU,y=xY$, the equation is $Y^2=x-3(1+xU)^3U+2(1+xU)^4U^2$. The implicit function theorem solves for $x$ at the common point; the two fibre branches $U=0,x=0$ meet with intersection multiplicity two. These charts cover the tangent conic and resolve the singularity. The two components have multiplicity one; the one met by $P_*$ is the nonidentity component $D$, with $D^2=-2$ from intersection with the full fibre.

\par\smallskip\noindent\textbf{Step 2.} At $t=0$, two blow-ups in the successive $t$-charts give
\[
 Y^2=tX^3+3t^2(1-t)X+2t^2(1-t)^2,
\quad
 Y_2^2=t^2X_2^3+3t(1-t)X_2+2(1-t)^2.
\]
The second equation has two smooth multiplicity-one components $E_{2,\pm}$ at $Y_2=\pm\sqrt2$, and $P_*$ meets one of them. The remaining singular point, in the second $X$-chart, has
\[
 V^2=H(X,T),\quad H=X^2T+(2+3X)T^2-(4X+3X^2)T^3+2X^2T^4,
 \qquad t=XT.
\]
Let $T_0(X)$ solve $H_T(X,T_0)=0$. It has $T_0=-X^2/4+O(X^3)$. Taylor's formula gives
$H-H(X,T_0)=(T-T_0)^2h_0$, $h_0(0,0)=2$, and $H(X,T_0)=X^4e(X)$, $e(0)=-1/8$. These identities follow by substitution in the displayed polynomial, not by discarding its higher terms. With $Z=(T-T_0)\sqrt{h_0}$, $a=V-Z$, $b=(V+Z)/e(X)$, the equation becomes exactly $ab=X^4$ through a local analytic isomorphism. Its resolution has charts
\[
 a=U_i^{i+1}V_i^i,\quad b=U_i^{3-i}V_i^{4-i},\quad X=U_iV_i
 \quad(0\le i\le3),
\]
with changes $U_{i+1}=1/V_i,V_{i+1}=U_iV_i^2$. They give a chain of three $(-2)$-curves. These are the charts of the two successive point blow-ups resolving $ab=X^4$, so the resolution is proper. Retaining $t=XT$ in the comparison gives multiplicities $2,3,2$ on that chain. The two new multiplicity-one components meet its ends; the first exceptional curve meets its middle component, and the original strict fibre meets the other end of that first exceptional curve. All these statements follow by substituting the displayed monomials and their orders; the only fibre components of multiplicity one are the strict component and $E_{2,\pm}$.

\par\smallskip\noindent\textbf{Step 3.} Order the six nonidentity components by the graph
$\Theta_1-\Theta_2-\Theta_3-\Theta_4$ and $\Theta_2-\Theta_5-\Theta_6$, with $E_{2,\pm}=\Theta_4,\Theta_6$. Their intersection matrix is $-C$, where
\[
 C=\begin{pmatrix}2&-1&0&0&0&0\\-1&2&-1&0&-1&0\\0&-1&2&-1&0&0\\0&0&-1&2&0&0\\0&-1&0&0&2&-1\\0&0&0&0&-1&2\end{pmatrix},\quad
 C^{-1}=\frac13\begin{pmatrix}6&9&6&3&6&3\\9&18&12&6&12&6\\6&12&10&5&8&4\\3&6&5&4&4&2\\6&12&8&4&10&5\\3&6&4&2&5&4\end{pmatrix}.
\]
Multiplication verifies this inverse. Self-intersections follow as well from the complete multiplicity vector and $F\cdot\Theta_i=0$.

\par\smallskip\noindent\textbf{Step 4.} The form $dx\wedge dt/y$ is holomorphic and nonvanishing on the finite resolved charts. In the last charts its expression is a unit times $dX\wedge da/a=-dU_i\wedge dV_i$; the earlier blow-up equations cancel the apparent denominators at $y=0$. At infinity it is $-r_{\rm cusp}^{-1}dX\wedge d r_{\rm cusp}/Y$ and has one fibre pole. Thus $K_S=-F$. Adjunction gives $Q^2=-1$ for every section $Q$, since $F\cdot Q=1$. A section meets a unique multiplicity-one component transversely at a smooth fibre point: the pullback of the full base equation to the section has order one, excluding all other possibilities.

\par\smallskip\noindent\textbf{Step 5.} Let $v_Q$ be the intersection vector with $\Theta_i$ and let $d_Q=Q\cdot D$. The possibilities are $v_Q=0,e_4,e_6$ and $d_Q=0,1$. Define the rational divisor class
\[
 \varphi(Q)=Q-O-(Q\cdot O+1)F+\sum\Theta_i(C^{-1}v_Q)_i+\tfrac12d_QD.
\]
It is orthogonal to $O,F,\Theta_i,D$ by direct multiplication. The chord-and-tangent rational function on the generic cubic has divisor $(Q)+(R)-(Q+R)-(O)$; extending it to the surface adds only vertical components. Applying the displayed orthogonal projection therefore proves $\varphi(Q+R)=\varphi(Q)+\varphi(R)$. Expansion gives
\[
 h(Q)=-\varphi(Q)^2=2+2Q\cdot O-v_Q^tC^{-1}v_Q-\tfrac12d_Q.
\]
For $Q\ne O$, the nonnegative intersection $Q\cdot O$ gives $h(Q)\ge2-4/3-1/2=1/6$. The visible $P_*$ is disjoint from $O$ everywhere, and the explicit blow-up charts place it on both nonidentity components. Hence $h(P_*)=1/6$. The logarithm argument in B.1.1 already proved $P_*=nP$ for an integer $n\ne0$, without identifying the visible section. Additivity gives $1/6=n^2h(P)\ge n^2/6$, so $n=\pm1$. The pointed cubic map in B.1.1 has $x(P)=f_{\rm sc}^2\wp_\tau(\mu)$ with $f_{\rm sc}\ne0$ on the ordinary locus. This proves the asserted zero identity.
\end{proof}
\begin{proposition}[FAM B.1.4; T]\label{fam:B.1.4}
The line-bundle comparison of the period family with Engel's family over $B^\circ$ transports its inverse period lift to a global section of $\operatorname{Hom}(t_{6P}^*M_S,M_S)$ with divisor $F_\infty$. Take this transported section for $c=0$ as the reference linearization. Then $\lambda_0=1$ exactly and $u=e^{-2\pi ic}=w^2$. For another independently fixed reference the precise constant is its scalar ratio with this section. With the cusp branch fixed in Section~\ref{sec:periods} and $f_{\rm sc}/t\to i\sqrt3/\pi$, the section is $P=(0,-\sqrt2t^2(t-1))$.
\end{proposition}
\uses{Definition~\ref{def:conventions}; FAM A.3.1 and preceding dictionary identities; \cite[Remark 3.12]{AC26}; \cite[Lemma 2.1; Corollary 2.2; Definition 2.3]{E26}.}
\begin{proof}
\par\smallskip\noindent\textbf{Step 1.} The zero identity $\wp_\tau(\mu)=0$, proved without numerical input in B.1.3, identifies $P$ with the signed visible section. With the cusp branch in Section~\ref{sec:periods} and the specified scale, substitution of $U=-(5+2\sqrt6)$ in the nodal uniformization of B.1.1 gives $Y=-\sqrt2$, so $P=(0,-\sqrt2t^2(t-1))$.

\par\smallskip\noindent\textbf{Step 2.} To compare the line bundles as families, use their common divisor $P-O$ and rigidify at the zero section. Put $F_\theta(\xi)=\theta(\xi-\mu)/\theta(\xi)$, $r_\theta=\theta(-\mu)/\theta'(0)$, and let $s_{P-O}$ be the canonical meromorphic divisor section. The map
\[
 [\xi,v]\longmapsto\frac{2r_\theta}{f_{\rm sc}}
        v\frac{s_{P-O}(\xi)}{F_\theta(\xi)}
\]
is a holomorphic line-bundle isomorphism, including at $P,O$ after cancellation of the equal divisors. At $O$ it sends the chosen frame to $dx/y$, since $F_\theta\sim r_\theta/\xi$ and $x/y\sim-\xi/f_{\rm sc}$ with the corresponding canonical divisor frame. Two such maps that agree at $O$ differ by a holomorphic scalar on each compact elliptic fibre and hence agree everywhere. This proves descent under the actual period identifications.

\par\smallskip\noindent\textbf{Step 3.} Engel's \cite[Lemma 2.1]{E26}, which we have read together with its proof, identifies $\operatorname{Hom}(t_{6P}^*M_S,M_S)$ with $\pi_S^*\mathcal O_B(1)$. The lemma applies because the preceding pointed cubic map and Mordell--Weil calculation identify his $(S,O,P)$ with ours, including the chosen sign. To determine which section is transported, it suffices to compute its order at one generic smooth point of each special fibre, since a section of this pulled-back line is a base section. At a finite centre write $6\mu=p+r\tau+\delta$ with $(p,r)=(4,-2),(3,-3)$ and $\delta/f_{\rm sc}$ holomorphic; this divisibility follows by linearizing the affine laws, whose homogeneous orders are two and one. Evaluate at $\xi=f_{\rm sc}z$ for generic fixed $z$. The factor $(2r_\theta/f_{\rm sc})/F_\theta(f_{\rm sc}z)$ tends to $2z$, a unit. Restoring the lattice shift $p+r\tau$ at the target changes the fibre multiplier to $e^{-2\pi i(\beta_c-r\mu)}$, whose finite limits are $e^{-2\pi i\kappa_1}$ and $e^{-2\pi i\kappa_2}$. The source and target theta ratios remain units. Therefore the base section has order zero at both finite points. At the cusp $f_{\rm sc}$ has a simple pole and $r_\theta$ has a nonzero limit; the two factors $2r_\theta/f_{\rm sc}$ cancel between source and target. The remaining multiplier is
$e^{-2\pi i\beta_c}=q_\tau e^{-2\pi i(\beta_c+\tau)}$, of exact order one. Thus the transported section has divisor $F_\infty$. Engel's \cite[Corollary 2.2]{E26} now identifies it with a nonzero scalar multiple of his reference, since two base sections with the same divisor have constant ratio. Changing $\beta$ by $c$ multiplies it exactly by $e^{-2\pi ic}$. With our stated period-reference normalization, this scalar is one. For any independently chosen reference $\phi_{\rm ref}$, the scalar is instead the precisely defined ratio $\lambda_0=\phi_{\rm part}/\phi_{\rm ref}$, evaluated on any ordinary fibre. Scaling the reference by $a\in\mathbb C^*$ changes $\lambda_0$ to $\lambda_0/a$; no scalar value is invariant under that declared choice.
\end{proof}
\subsection{Correction to Engel's Remark 8.1}
\noindent\textbf{Printed sentence.} In \cite[Remark 8.1]{E26}: ``We may thus build a proper, holomorphic submersion $\mathcal Y^*\to\Delta^*$ whose fiber over $u$ is $X(u,a_0,a_1)$''. The fault is the global descent for fixed sphere logarithmic data; the local construction is retained. No contact with the source's author is recorded. The finite order-two obstruction below identifies the precise qualification.
\begin{proposition}[FAM B.2.3; T]\label{fam:B.2.3}
The assertion in \cite[Remark 8.1]{E26} of a global proper family over the original punctured $u$-disc for fixed sphere logarithmic data requires a double-cover qualification. The retained family with such data does not descend through $u=w^2$. It does give a proper holomorphic family on each simply connected $u$-open set, and a global proper family over the punctured $w$-disc as in A.3.1.
\end{proposition}
\uses{FAM A.2.1, A.3.1, A.4.1 and B.2.1; \cite[Remark 8.1]{E26}.}
\begin{proof}
\par\smallskip\noindent\textbf{Step 1.} Descent through $u=w^2$ would furnish an isomorphism between the two pulled-back fibres at $w$ and $-w$, namely $X_c$ and $X_{c+1}$. A.2.1 excludes that isomorphism by the nonzero class of $d$ in $\Lambda^{A_2}/N_2\Lambda$. This is an exact order-two obstruction on the normalized finite cover. It applies to all sphere data by B.2.1 because $r_2$ is odd.

\par\smallskip\noindent\textbf{Step 2.} There is also a direct analytic witness near zero. The relative Albanese invariant of the multiplicity-four reduced fibre in a holomorphic descended family would be a single-valued holomorphic function of $u$. After pullback it would be even in $w$. But A.4.1 gives $j_2(w)=C^{-1}w^{-1}+744+O(w)$, where $C=e^{-\pi i\kappa_2^{\rm part}}\ne0$; its Laurent coefficient of $w^{-1}$ is nonzero. Thus it is not even on any punctured neighbourhood. This verifies the obstruction independently of a choice of isomorphism lift. Local inverse logarithms of $u$ give the proper local families already constructed; coherent even shifts give their global double-cover family. The untwisted middle quotient itself has the integer shear identification, so this refutation concerns the fixed sphere logarithmic data in the stated remark.
\end{proof}

\section{Discrete data and parameter identifications}\label{sec:discrete}
\subsection{The standard data (known)}
The $2\mathbb Z$ equivalence is Theorem 1.4 of Du, Guo, Kusakabe and Xie, arXiv:2610.00250, version 1 \cite[Theorem 1.4]{DGKXv1}. Version 2 does not contain it. The same statement was obtained independently in the workbench from which this paper grew, and Appendix A prints that proof in the present notation.
\begin{cited}[FAM A.5.1; K]\label{fam:A.5.1}
Fix the original cusp-regular period solution, the clockwise peripheral marking, the prescribed periodic triangular cusp fan, the affine vectors $(\varepsilon,-\varepsilon')$, and the untwisted cusp collar. For every $c,c'\in\mathcal P$, the constructed compact complex threefolds satisfy
\[
 X_c\simeq X_{c'}\quad\Longleftrightarrow\quad c'-c\in2\mathbb Z.
\]
Every biholomorphism is over the identity of the original $\mathbb P^1$. The proper holomorphic family of their isomorphism classes has parameter $0<|w|<e^{-\pi M}$ and nonzero Kodaira--Spencer class at every point.
\end{cited}
\uses{FAM A.1.1, A.2.1, A.3.1 and A.4.2; \cite[Theorems 1.3, 1.4 and 7.1; punctured-disc corollary]{DGKXv1}.}

\subsection{All affine vectors and cusp sections}
\begin{proposition}[FAM B.2.1; T]\label{fam:B.2.1}
Let
\[
 w_1=r_1\varepsilon+k_1d,\qquad w_2=r_2\varepsilon'+k_2d,
 \qquad 3\nmid r_1,\quad 2\nmid r_2,
\]
and use an arbitrary holomorphic cusp collar section $\sigma$ in the construction of \cite[Section 6.4]{AC26}. Write $l_0=\gamma(c(\sigma))$ in its positive cusp-winding convention, and set $p=12l_0-4r_1-3r_2$. Let primes denote a second such datum with the same period solution and cusp fan. Then
\[
 X(c;w_1,w_2,\sigma)\simeq X(c';w'_1,w'_2,\sigma')
\]
holds exactly for the following arithmetic conditions:
\[
 n=c'-c\in\mathbb Z,\qquad
 p'=a p\text{ for some }a\in\{1,-1\},\qquad
 n\equiv k_2-k'_2\pmod2.
\]
All maps are over $B$. A map with the indicated $a$ has fibre-linear part $aI_2$. In particular the complete invariants are
\[
 \bigl(|p|,\ e^{-\pi i(c+k_2)}\bigr),\qquad
 |p|\ge1,\quad\gcd(|p|,6)=1.
\]
Every value of this positive integer occurs. The full collar section enters the construction, but its additional analytic part and its second integral winding contribute no further invariant.
\end{proposition}
\uses{Definition~\ref{def:conventions}; FAM A.1.1--A.3.1 and the preceding discrete-data calculation; \cite[Section 6.4]{AC26}; \cite[Sections 3 and 5]{E26}.}
\begin{proof}
\par\smallskip\noindent\textbf{Step 1.} The determinant calculation in A.1.1 was independent of translations, so the base is still intrinsic. The centralizer proof gives $c'=c+n$ and $aI_2$ exactly as before. Put $S_n=I+nE$ and compare all target periods in source coordinates. On the finite root cover the exact deck equation is
\[
 B_j(g_jz)-R_jB_j(z)=\Pi_c(g_jz)\left(\frac{d_j}{m_j}+\ell_j\right),
 \quad d_j=S_nw'_j-a w_j,
 \quad N_j\ell_j=-d_j.                                      \tag{B2.1}
\]
Necessity follows by the normalized-base-change lift already proved. Conversely the norm equation has a rational solution $r_j^*$ of
$(I-A_j)r_j^*=d_j/m_j+\ell_j$, since $\ker N_j$ over $\mathbb Q$ equals the image of $I-A_j$. We can choose $\gamma(r_j^*)=0$: the columns $\widehat u,\widehat w$ already span that image over $\mathbb Q$. Then $B_j=\Pi_cr_j^*$ is holomorphic on the whole root disc and solves the exact equation. Thus the norm condition is sufficient for a local isomorphism, not only for conjugacy on the central real torus.

\par\smallskip\noindent\textbf{Step 2.} The computed norm images translate (B2.1) into
\[
 r'_1\equiv ar_1\pmod3,\quad r'_2\equiv ar_2\pmod4,
 \quad ak_2-k'_2-nr'_2\equiv0\pmod2.                       \tag{B2.2}
\]
There is no condition on $k_1$. Since $r'_2$ is odd, the last condition is the displayed parity. On the finite collar the middle translation is
$A_j=B_j+(\log s_j)/(2\pi i)\,\Pi_cd_j$ and has integral defect $\ell_j$. Its $\gamma$ defect is
$\alpha_j=(ar_j-r'_j)/m_j$.

\par\smallskip\noindent\textbf{Step 3.} We prove the precise cusp necessity used next. A holomorphic isomorphism between cusp neighbourhoods with fibre-linear part $aI_2$ has, after lifting near a generic smooth central point to the toric covers, the form
\[
 (x,t_c)\longmapsto(t_c^q v_{\rm unit}(t_c)x^a,t_c),\qquad q\in\mathbb Z^2,
 \quad v_{\rm unit}(0)\in(\mathbb C^*)^2.                           \tag{B2.3}
\]
Here the assertion concerns this local lift near a generic central point. To prove it, take a source chart near the generic point of its component and a target chart near its image. The etale toric quotient permits a local lift over a simply connected neighbourhood. The torus characters in a chart at the height-one ray $v$ are $x_i=t_c^{v_i}X_i$, with $X_i$ invertible. On $t_c\ne0$ the affine fibre equation gives $x'_i/x_i^a$ independent of the $X$ coordinates. Its expression in the two charts is therefore $t_c^{v'_i-av_i}$ times a holomorphic unit of $t_c$, including zero. This proves (B2.3). Taking logarithms shows that its winding is in $\Lambda_{\rm tor}=\mathbb Z\widehat w+\mathbb Zd$. Changing a lift adds $(I-A_0)\Lambda=\Lambda_{\rm tor}$, so the vanishing of its class in $\bar\Lambda=\Lambda/\Lambda_{\rm tor}$ is intrinsic. This proof supplies the necessary extension test; no assertion about arbitrary punctured translations is used.

\par\smallskip\noindent\textbf{Step 4.} We retain the actual collar convention. The map from the cusp piece to the middle family is $T_{-\sigma}G_0^{-1}$. Consequently the middle translation induced by (B2.3) is
\[
 A_0=sq+a\sigma-\sigma'+\eta_0,
\]
with $\eta_0$ holomorphic on the cusp disc after a logarithm of the unit. For the clockwise generator $s\mapsto s-1$, the class of its defect in $\bar\Lambda$ is $-a c(\sigma)+c(\sigma')$. On the other hand every middle integral cocycle satisfies
$\ell_0=-A_0(\ell_1+A_1\ell_2)$, whose $\gamma$ coordinate is $-\alpha_1-\alpha_2$. Equality of these coordinates is exactly
\[
 l'_0-a l_0=-(ar_1-r'_1)/3-(ar_2-r'_2)/4,
\]
or $p'=ap$. This proves the remaining necessity. Notice that $p\equiv-r_1\pmod3$ and $p\equiv r_2\pmod2$, so $p$ is nonzero and coprime to six.

\par\smallskip\noindent\textbf{Step 5.} For sufficiency, $p'=ap$ implies the first two congruences of (B2.2) by reduction modulo three and four. The parity supplies the remaining one. Choose integral solutions $\ell_1,\ell_2$ of (B2.1). Their $\gamma$ coordinates already have the required values. Adding any element of $\ker N_1=(I-A_1)\Lambda$ preserves (B2.1). In particular adding an arbitrary integer multiple of $e_1=(0,1,-2,0)^t$ changes the $\widehat u$ coordinate of $\ell_0$ by an arbitrary integer while leaving its $\gamma$ coordinate fixed. We can therefore arrange
$[\ell_0]=-a c(\sigma)+c(\sigma')$ in the entire rank-two group $\bar\Lambda$. Choose lifts of $\sigma,\sigma'$ on the cusp half-plane. Their remaining defect differs from $\ell_0$ by an actual toric lattice vector; an integral $q$ makes $sq+a\sigma-\sigma'$ have exactly defect $\ell_0$.

\par\smallskip\noindent\textbf{Step 6.} The finite translations and this cusp reference now define a locally nonempty affine torsor with linear difference sheaf precisely $\mathcal V$ of A.3.1. The terms $a\sigma-\sigma'$ are retained, even when their chosen lifts have essential singularities at the puncture; they belong to the specified collar reference, while differences of allowed extensions are holomorphic $\eta_0$. The proof of $H^1(B,\mathcal V)=0$ gives a global middle translation. The finite maps descend by (B2.1), and the cusp map is $x\mapsto t_c^qv_{\rm unit}(t_c)x^a$. The fan is preserved by all integer translations and by $y\mapsto-y$; exponentials of the integer period shift agree. Hence this cusp map and its inverse extend over every stratum. The finite overlap is the equality defining $A_j$, and the cusp overlap is the equality defining $A_0$. Thus the maps glue globally.

\par\smallskip\noindent\textbf{Step 7.} Finally parity is equivalent to equality of the two displayed exponentials. Passing from $p$ to $-p$ is effected by $a=-1$ and the same proof; for fixed sign the congruences have just been shown sufficient. Every positive integer coprime to six is obtained by choosing $r_1\in\{1,2\}$ and odd $r_2$ in the appropriate residue classes modulo three and four, and then taking $l_0=(p+4r_1+3r_2)/12$. This also proves that the extra collar data have been accounted for, rather than omitted from the statement.
\end{proof}
\begin{proposition}[FAM B.2.2; T]\label{fam:B.2.2}
For multiplicities $(3,4)$ and $|p|=1$, every admissible discrete member is biholomorphic to one member of A.5.1, uniquely modulo $2\mathbb Z$. At fixed nonzero admissible Engel parameter $u$ there are exactly two isomorphism classes, exchanged by a positive circuit of $u$ about zero. In the convention that Engel's root generator is the inverse of our clockwise one and his circle class is $d$, the identity of the middle quotient sends
\[
 X_{c+1}\longleftrightarrow X(u;a_0-d/3,\ a_1+d/4).
\]
Reversing both declared root generators reverses these two displayed signs.
\end{proposition}
\uses{Definition~\ref{def:conventions}; FAM A.1.1--A.3.1 and the preceding discrete-data calculation; \cite[Section 6.4]{AC26}; \cite[Sections 3 and 5]{E26}.}
\begin{proof}
\par\smallskip\noindent\textbf{Step 1.} Take the canonical target data $(r'_1,r'_2,k'_1,k'_2,l'_0)=(1,-1,0,0,0)$, for which $p'=-1$. Choose $a$ so that $p'=ap$, and an integer $n\equiv k_2\pmod2$. B.2.1 supplies the isomorphism to $X_{c+n}$. The uniqueness is A.5.1. At fixed $u=e^{-2\pi ic}$, its logarithms differ by integers; the invariant $e^{-\pi i(c+k_2)}$ takes exactly the two square roots of $u$. Both occur, already as $X_c,X_{c+1}$, and A.2.1 distinguishes them.

\par\smallskip\noindent\textbf{Step 2.} We spell out the integral comparison to Engel's common basis $B_0=(e_1,e_2,e_3,e_4)$, retaining his $C f_4=e_1+e_4$:
\[
 P_{\rm mark}=\begin{pmatrix}0&-2&-1&-2\\2&-1&0&-1\\0&2&1&3\\-1&0&0&0\end{pmatrix},\qquad
 T_0P_{\rm mark}=P_{\rm mark}A_1,\quad T_1P_{\rm mark}=P_{\rm mark}A_2.
\]
Its determinant is $-1$, $P_{\rm mark}d=\delta_0$, $\psi P_{\rm mark}=-\gamma$, $P_{\rm mark}\varepsilon=-e_4$ and $P_{\rm mark}\varepsilon'=\delta_0-Cf_4$. These are exact lattice homomorphism identities, verified independently through Viaclovsky's intermediate basis in \texttt{ab\_geometry.py}. They also agree with the complex periods: at $z_1$ the complex-linear map
$(\zeta_1,\zeta_2)\mapsto(-\zeta_1/(\rho+1),\zeta_2+\zeta_1/3)$ carries $\Pi_cP_{\rm mark}^{-1}$ to Engel's columns $(1,0),(\rho-1,0),((1+\rho)/3,1/3),(0,-\kappa_1)$. At $z_2$ his $Cf_4$ has preimage $d-\varepsilon'$ and fixed period $1-\kappa_2$. Thus the comparison fixes the positive circle class and the actual local complex structures, in addition to their monodromy matrices.

\par\smallskip\noindent\textbf{Step 3.} For the sign assertion, inverting the concrete clockwise affine generator changes its invariant translation $\Pi_cv_j/m_j$ to $-\Pi_cv_j/m_j$. Thus our dictionary is $a_{j-1}=-P_{\rm mark}v_j/m_j$ in Engel's lattice marking (and its period image in the torus). Integer shift gives $\Pi_{c+1}v_j=\Pi_c(v_j+\gamma(v_j)d)$, with $\gamma(v_1)=1,\gamma(v_2)=-1$, proving the two signs. A positive $u$ circuit changes $\log u$ by $2\pi i$ and hence changes $c$ by $-1$, exchanging the square roots. This gives the claimed permutation with its orientation.
\end{proof}

\subsection{Coprime higher multiplicities}
The theorem concerns the stated model and marking. The identification with the coordinates of \cite[Section 2]{VL26} is separately announced below; this distinction retains the gap noted by the supplied review.
\begin{proposition}[FAM B.3.1; T]\label{fam:B.3.1}
Let $k,l$ be positive integers with $\gcd(3k,4l)=1$, and put $m_1=3k,m_2=4l$. Consider Viaclovsky's sphere members expressed in the same clockwise marking, with $w_j=r_j\varepsilon_j+s_jd$, $\gcd(r_j,m_j)=1$, cusp winding $l_0$, and $p=m_1m_2l_0-m_2r_1-m_1r_2=\pm1$, where $\varepsilon_1=\varepsilon,\varepsilon_2=\varepsilon'$. Two such members are biholomorphic precisely when their multiplicity pair agrees and, for some $a\in\{1,-1\}$ and integer $n=c'-c$,
\[
 p'=ap,\quad r'_j\equiv ar_j\pmod{m_j},\quad
 nr'_1\equiv as_1-s'_1\pmod{k},\quad
 nr'_2\equiv as_2-s'_2\pmod{2l}.
\]
The congruences on $r'_j$ already follow from $p'=ap$. For fixed discrete vectors and cusp collar the identification subgroup is $2kl\mathbb Z$. For fixed admissible $u$ and this multiplicity pair, all the stated sphere translation data give exactly $2kl$ isomorphism classes, permuted in one cycle by a circuit of $u$ about zero.
\end{proposition}
\uses{FAM A.1.2, A.2.1, B.2.1; \cite[Proposition 4.7]{VL26}.}
\begin{proof}
\par\smallskip\noindent\textbf{Step 1.} We first recover the base for these sphere members without using the degree-two canonical pencil. A biholomorphism $F:X\to X'$ gives $g=f'F$. Since $H^2(S^6,\mathbb R)=0$, the closed form $g^*\omega_{\rm FS}$ is exact. Restrict to a regular torus $T$ and wedge with a flat Kahler form on $T$. Stokes' theorem gives integral zero, while the integrand is nonnegative and strictly positive wherever $dg|_T\ne0$. Thus $g$ is constant on each regular fibre. The image of $(f,g)$ is a compact analytic curve whose projection to $B$ has degree one; its normalization gives a holomorphic $h:B\to B$ with $g=hf$. Apply the same construction to $F^{-1}$ to see that $h$ is an automorphism. The three homological local monodromy orders are three, four and infinity, independently of $k,l$; their intertwining maps under $F$ therefore fix each special point. Hence $h=1$. The nonreduced fibre multiplicities now force equal $k,l$. This is the full reduction in the read Proposition 4.7 of Viaclovsky's paper, written here with its exact maps.

\par\smallskip\noindent\textbf{Step 2.} The root model of degree $m_1$ is the pullback of the degree-three model by $s\mapsto s^k$, and likewise at the other point. Its linear action remains $A_1$ or $A_2$, while its root generator has order $m_j$. The primitive invariant coordinate $r_j/m_j$ proves freeness of every nonidentity power. The normalization and deck-commutation arguments of A.1.2 and A.2.1 apply to these root covers without changing a generator. Their norms are
\[
 N_{3k}=kN_1,\qquad N_{4l}=lN_2,
\]
so the fixed-coordinate norm images are $3k\mathbb Z\varepsilon+k\mathbb Zd$ and $4l\mathbb Z\varepsilon'+2l\mathbb Zd$. The norm equation (B2.1), with $m_j$ replaced by these actual covering degrees, gives exactly the displayed congruences. The centralizer and first-column comparison remain unchanged.

\par\smallskip\noindent\textbf{Step 3.} For completeness the integer $p$ records the primitive sphere relation in this marking. The common torus coinvariants are detected by $\gamma$, since $(A_1-I)\Lambda+(A_2-I)\Lambda=\ker\gamma$ by their explicit saturated bases. The finite meridian relations are $x^{m_1}=c^{r_1}$, $y^{m_2}=c^{r_2}$, and the cusp relation is $xy=c^{l_0}$ with $c$ central. Eliminating $y$ gives an abelian relation matrix with determinant $-p$ and primitive first elementary divisor because $\gcd(m_1,m_2)=1$. Thus the sphere data have $p=\pm1$. Reducing $p'=ap$ modulo $m_1,m_2$ gives the two congruences for $r_j$.

\par\smallskip\noindent\textbf{Step 4.} To prove sufficiency, use the local rational solutions of the norm equation and the identical cusp comparison in B.2.1. The first-coordinate extension on the new root discs has vanishing orders $2k$ and $l$, precisely the orders of the pulled-back $\Theta$. Dividing by it again gives $\mathcal O_B(-p_0)$, and finite averaging again gives the same rank-two extension sheaf. Thus $H^1(B,\mathcal V)=0$ still removes the analytic cocycle. The images of $I-A_j$ and their arbitrary $\widehat u$ adjustments are unchanged, so $p'=ap$ is exactly the remaining cusp winding condition. Every local map and overlap is the one already written in B.2.1, with the new denominator $m_j$; this proves global sufficiency.

\par\smallskip\noindent\textbf{Step 5.} For identical data, $p\ne0$ excludes $a=-1$. The congruences reduce to $k\mid nr_1$ and $2l\mid nr_2$. The two primitivity assumptions imply $k\mid n$ and $2l\mid n$. Since $\gcd(k,2l)=1$, their least positive common multiple is $2kl$. For different sphere translation data, $r'_1$ is invertible modulo $k$, $r'_2$ modulo $2l$, and the Chinese remainder theorem gives exactly one integer shift class modulo $2kl$. The values of $s_1$ modulo $k$ and $s_2$ modulo $2l$ all occur without affecting freeness or $p$. Thus there are exactly $2kl$ classes at fixed $u$. The change $c\mapsto c-1$ on a positive $u$ circuit acts by a generator of this cyclic residue group. Program \texttt{ab\_exact.py} verifies the repeated norms and shift steps for five coprime pairs; the preceding symbolic identities prove all pairs.
\end{proof}
\begin{announcement}[PUSH E.1.1; A]\label{push:E.1.1}
Let $(\delta,a,b,d)$ be the marking in
\cite[Section 2, equations (2.1), (2.2), (2.7), (2.11)--(2.12)]{VL26},
and put $u_0=d+2a+2b$, $u_1=d+3a+3b$.
The homomorphism from the marking of the question to this marking is
\[
 P_V=\begin{pmatrix}
 -2&1&0&1\\-6&0&-1&0\\0&-1&0&0\\-1&0&0&0
 \end{pmatrix},\qquad \det P_V=-1.
\]
It intertwines the two finite monodromies and satisfies
\[
 P_V\widehat\delta=\delta,\qquad
 P_V\varepsilon=-u_0,\qquad
 P_V\varepsilon'=\delta-u_1.
\]
With the counterclockwise deck generators of that paper, its data
$m_0=3k$, $m_1=4l$ and
$a_i=[v_i/m_i]$ with $v_i=\rho_i\delta+N_i u_i$ correspond to
\[
 w_1=N_0\varepsilon-\rho_0\widehat\delta,
 \qquad
 w_2=N_1\varepsilon'-(\rho_1+N_1)\widehat\delta.
\]
For the section-preserving cusp comparison used there, $l_0=0$ and
$p=-(nN_0+mN_1)$.  Thus its sphere condition is
$b:=nN_0+mN_1=\pm1$.
The brackets in $a_i$ mean the torsion point in the smooth
root torus, using its section-preserving origin and the
transported period vector.  Replacing $v_i$ by an integral
multiple $m_i$ times a period does not change that torsion
point, although its chosen lift fixes the collar bookkeeping.
\end{announcement}
\uses{FAM B.3.1; PUSH E.1.1--E.1.2; \cite[Section 2]{VL26}.}
\noindent\textit{Mechanism.} The integral change of basis transports the positive-root translations; repeated norm images give the displayed congruences. Proof in the repository, S6\_02\_PUSH, E.1.1 \cite[run S6\_02\_PUSH, E.1.1]{Repo}.
\begin{announcement}[PUSH E.1.2; A]\label{push:E.1.2}
Use a local logarithm $c$ of the transported linearization,
$u=e^{-2\pi i c}$, and the same reference section and cusp comparison
on both sides.  Two sphere members with data
$(m,\rho_0,N_0;n,\rho_1,N_1;c)$ and primed data are biholomorphic
if and only if their multiplicities agree and there are
$\alpha\in\{1,-1\}$ and $h=c'-c\in\mathbb Z$ such that
\[
 b'=\alpha b,\qquad
 hN'_0\equiv\rho'_0-\alpha\rho_0\pmod{k},\qquad
 hN'_1\equiv\rho'_1-\alpha\rho_1\pmod{2l}.
\]
Every such biholomorphism covers the identity of $B$.
For fixed discrete data the period identification subgroup is
$2kl\mathbb Z$.  At fixed $u$ all sphere translation data give
exactly $2kl$ classes.  A positive $u$-circuit acts on them by a
single cycle.
\end{announcement}
\uses{FAM B.3.1; PUSH E.1.1--E.1.2; \cite[Section 2]{VL26}.}
\noindent\textit{Mechanism.} The integral change of basis transports the positive-root translations; repeated norm images give the displayed congruences. Proof in the repository, S6\_02\_PUSH, E.1.2 \cite[run S6\_02\_PUSH, E.1.2]{Repo}.


% ===== periods source-preserving fragment =====

\section{Periods in closed form}\label{sec:periods}
\subsection{The normalized theta expression}
We use the odd Jacobi function and the Dedekind function in FAM C.1:
\[
\vartheta_1(v,\tau)=2\sum_{n\ge0}(-1)^n e^{\pi i\tau(n+1/2)^2}\sin((2n+1)\pi v),\qquad
\eta_D(\tau)=e^{\pi i\tau/12}\prod_{n\ge1}(1-e^{2\pi in\tau}).
\]
The cusp branch is reached from $\tau=iT$, $T>1$, on which $\mu=1/2-iy(T)$ and $0<y(T)<T/2$.
Set $a_{\rm cusp}=5+2\sqrt6$. The elliptic-logarithm interpretation is due to \cite[Remark~3.12]{AC26}; the zero formula is due to Eichler and Zagier \cite[unnumbered theorem, p.~399]{EZ82}, also stated in \cite[§2, equation~(6)]{DI10}.

\begin{proposition}[FAM C.1.1; T]\label{fam:C.1.1}
For the period section identified in Statement~\ref{fam:B.1.3}, the particular solution
with \((\beta^{\rm part}+\tau)(p_0)=0\) is
\begin{equation}\label{fam:eq:cd-beta}
 \beta^{\rm part}=\frac{6i}{\pi}\operatorname{Log}
 \frac{\vartheta_1(\mu,\tau)}{2\sqrt3\,\eta_D(\tau)}.
\end{equation}
Here the logarithm is continued on the simply connected uniformizing
half-plane from its cusp expansion. Its special values are
\[
 \mu_0=\frac12-\frac{i\log a_{\rm cusp}}{2\pi},\qquad
 \mu(z_1)=\frac{2-\rho}{3},\qquad \mu(z_2)=\frac{1-i}{2},
\]
\[
 \beta^{\rm part}(z_2)=i\left(\frac32-\frac3\pi\log6\right),
 \qquad \operatorname{Re}\beta^{\rm part}(z_1)=-\frac12.
\]
In particular \(e^{2\pi i\mu_0}=-a_{\rm cusp}\).
\end{proposition}
\noindent\textit{Uses:} FAM B.1.3; the cusp normalization; the Jacobi transformation laws.\par
\begin{proof}
\noindent\textbf{Step 1.} Statement~\ref{fam:B.1.3} identifies the elliptic point represented by \(\mu\) with
one of the sections \((0,\pm\sqrt2t^2(t-1))\). The uniformization map
\((v,\tau)\mapsto(f_{\rm sc}^2\wp_\tau(v),f_{\rm sc}^3\wp'_\tau(v)/2)\) therefore gives
\(\wp_\tau(\mu)=0\). No zero of \(\wp\) is a lattice point. Thus the numerator
of \eqref{fam:eq:cd-beta} is nonzero; the denominator is nonzero as well.
A holomorphic logarithm exists on the simply connected cover.


\noindent\textbf{Step 2.} For completeness the cusp value follows directly from the normally
convergent Fourier expansion, rather than from a quoted Eichler--Zagier
constant:
\begin{equation}\label{fam:eq:cd-wp}
 \frac{\wp_\tau(v)}{\pi^2}
 =\csc^2(\pi v)-\frac13+
 8\sum_{n\geq1}\frac{nq_\tau^n}{1-q_\tau^n}(1-\cos(2\pi nv)),
 \qquad q_\tau=e^{2\pi i\tau}.
\end{equation}
This identity follows by differentiating the logarithm of the Jacobi
product twice and using
\(E_2=1-24\sum nq_\tau^n/(1-q_\tau^n)\). Boundedness of the chosen elliptic logarithm
at the cusp permits passage to the limit. We obtain
\(\sin^2(\pi\mu_0)=3\), equivalently
\(v_{\rm cusp}^2+10v_{\rm cusp}+1=0\), \(v_{\rm cusp}=e^{2\pi i\mu_0}\).
On the chosen imaginary-axis branch \(\operatorname{Im}\mu<0\), so \(v_{\rm cusp}=-a_{\rm cusp}\).
One can see the branch without a numerical convention: on the rectangular
elliptic curve with \(\tau=iT\), \(T>1\), \(\wp\) is real on the segment
\(1/2-iy\), \(0\leq y\leq T/2\). Its endpoint values are the largest and
middle real roots of its cubic, respectively positive and negative.
Its derivative has no zero in the open segment, since the zeros of
\(\wp'\) are the three nonzero half-periods. There is exactly one zero
there; this gives the indicated continuation. Consequently
\(2\sin(\pi\mu_0)=2\sqrt3\).


\noindent\textbf{Step 3.} Put \(R=\vartheta_1/\eta_D\). Poisson summation for the defining Gaussian
series and its index shift give
\[
 \begin{aligned}
 R(v,\tau+1)&=e^{\pi i/6}R(v,\tau),\\
 R(v/\tau,-1/\tau)&=-i e^{\pi iv^2/\tau}R(v,\tau),\\
 R(1-v,\tau)&=R(v,\tau),\qquad R(v+1,\tau)=-R(v,\tau).
 \end{aligned}
\]
Inserting the actual transformed arguments gives
\[
 \begin{aligned}
 R(\mu\circ g_1,\tau\circ g_1)
 &=e^{-\pi i/3}e^{\pi i(1-\mu)^2/\tau}R(\mu,\tau),\\
 R(\mu\circ g_2,\tau\circ g_2)
 &=i e^{\pi i\mu^2/\tau}R(\mu,\tau).
 \end{aligned}
\]
Multiplication of the logarithms by \(6i/\pi\) gives exactly
\(2-6(1-\mu)^2/\tau\) and \(-3-6\mu^2/\tau\), modulo constants in
\(12\mathbb Z\). These constants vanish: each defect is constant on the
connected cover, and evaluation at the corresponding fixed point makes
both its two sides zero. The fixed-point equations themselves give
\(\mu(z_1)=(2-\rho)/3\) and \(\mu(z_2)=(1-i)/2\).
The Jacobi product has cusp expansion
\(R(\mu,\tau)=2q_\tau^{1/12}\sin(\pi\mu)(1+O(q_\tau))\).
Equation \eqref{fam:eq:cd-beta} consequently gives
\(\beta^{\rm part}+\tau\to0\). The difference from the specified particular
solution is an invariant holomorphic function extending across all three
points of \(B\), hence a constant, and the cusp normalization makes it zero.


\noindent\textbf{Step 4.} At \(\tau=i\), the half-period translation formula gives
\(\vartheta_1((1-i)/2,i)=e^{\pi/4}\vartheta_3(0,i)\).
The identities \(2\eta_D^3=\vartheta_2\vartheta_3\vartheta_4\),
\(\vartheta_2(0,i)=\vartheta_4(0,i)=2^{-1/4}\vartheta_3(0,i)\)
give \(R=e^{\pi/4}\sqrt2\), proving the displayed value of \(\beta\).
At \(\rho\), \(q_\tau\) is negative real, \(\mu=1/2-i/(2\sqrt3)\), and the
product for \(R\) has phase \(\pi/12\), with all its remaining factors
positive real. This gives \(\operatorname{Re}\beta=-1/2\).
\end{proof}

\begin{computed}[FAM C.1.2; T]\label{fam:C.1.2}
The following decimal evaluations of the exact expressions in C.1.1 agree
in two independent evaluations through the displayed twenty decimal places:
\[
\begin{array}{c|l}
\mu_0&0.5-0.36485183191106784792\,i\\
\mu(z_1)&0.5-0.28867513459481288225\,i\\
\mu(z_2)&0.5-0.5\,i\\
\beta^{\rm part}(z_1)&-0.5-1.27101206057537346287\,i\\
\beta^{\rm part}(z_2)&-0.21100425815613411418\,i
\end{array}
\]
These are evaluations of explicit convergent expressions, with working
precision ninety decimal digits. The formula \eqref{fam:eq:cd-beta}, together
with \(j(\tau)=1728t\) and \(\wp_\tau(\mu)=0\), determines the periods
throughout a fundamental domain.
\end{computed}
\noindent\textit{Uses:} FAM C.1.1; the theta product and its displayed tail bound.\par
\begin{proof}
\noindent\textbf{Step 1.} The source run records that the executable \texttt{cd\_numerics.py} evaluates \(\mu_0\) first by its
algebraic logarithm and then by solving
\(\pi^2\csc^2(\pi v)-\pi^2/3=0\). For \(\beta\), its first calculation uses
the Jacobi theta series and the eta product; its second uses
\[
 R(v,\tau)=2e^{\pi i\tau/6}\sin(\pi v)
 \prod_{n\geq1}(1-2q_\tau^n\cos(2\pi v)+q_\tau^{2n}).
\]
The product is truncated after 120 factors. At the two specified points
\(|q_\tau|\leq e^{-\pi\sqrt3}\) and \(|\operatorname{Im}v|\leq1/2\), so the
omitted logarithm is bounded by
\(\sum_{n>120}(2|q_\tau|^n\cosh\pi+|q_\tau|^{2n})/(1-2|q_\tau|^n\cosh\pi-|q_\tau|^{2n})\),
which is less than \(10^{-280}\). The measured discrepancies of the
ninety-digit calculations are below \(10^{-88}\). We do not promote floating-point rounding to an interval certificate.


\noindent\textbf{Step 2.} 
A further sixty-digit interval computation, \texttt{cd\_special\_intervals.py},
evaluates the real positive product at $\rho$, including its forty-factor
tail, and the exact logarithmic expression at $i$. It encloses both
imaginary parts within intervals narrower than $3\cdot10^{-59}$ and
therefore certifies the displayed twenty decimal places. It also proves
$-0.2<\operatorname{Im}(\kappa_1^{\rm part}-\kappa_2^{\rm part})<-0.1$.
The exact formulas, product-tail bounds, independent algorithms, and
full outputs are reported by the source run as accompanying the decimals. The special values of \(\mu\) are also checked directly in
the fixed-point equations, an exact algebraic calculation.
\end{proof}

The computation in FAM C.1.2 is the recorded source computation. Neither \path{cd_numerics.py} nor \path{cd_special_intervals.py}, nor their full outputs, is among the supplied attachments; no current execution is claimed. The latter interval is the input used by FAM A.4.1 for its generic Albanese cross-check. All program and output-file references in the period and maximum proofs below likewise describe the source runs; they do not assert a current execution or the presence of those artifacts in the supplied attachments.

\subsection{The elliptic norm and its special values}
The elliptic lattice in this subsection is denoted $\Lambda_{\rm ell}$ to distinguish it from the marked rank-four lattice. Put
\[
\mathcal N(v,\tau)=e^{-\pi(\operatorname{Im}v)^2/\operatorname{Im}\tau}
\left|\frac{\vartheta_1(v,\tau)}{\eta_D(\tau)}\right|,
\qquad U_\tau(v)=\log\mathcal N(v,\tau),
\qquad \mathcal L(t)=U_\tau(\mu).
\]
For $\Lambda_{\rm ell}=\omega(\mathbb Z+\tau\mathbb Z)$ use $U_{\Lambda_{\rm ell}}(z)=U_\tau(z/\omega)$.
The theta formula gives the source identity (FAM C.2.1, equation~\texttt{cd-norm})
\begin{equation}\label{fam:eq:cd-norm}
D_{\rm part}=\frac6\pi\log\frac{\mathcal N(\mu,\tau)}{2\sqrt3}.
\end{equation}

\begin{proposition}[ARI C.1.1; T; distribution under isogenies]\label{ari:C.1.1}
Let $\Lambda_{\mathrm{ell}},\Lambda_{\mathrm{ell}}'\subset\mathbb C$ be lattices and let
$\alpha\ne0$ satisfy $\alpha\Lambda_{\mathrm{ell}}\subset\Lambda_{\mathrm{ell}}'$. Write
$\phi:\mathbb C/\Lambda_{\mathrm{ell}}\to\mathbb C/\Lambda_{\mathrm{ell}}'$ for $z\mapsto\alpha z$,
$K=\alpha^{-1}\Lambda_{\mathrm{ell}}'/\Lambda_{\mathrm{ell}}$, and
$d=[\Lambda_{\mathrm{ell}}':\alpha\Lambda_{\mathrm{ell}}]$. Then
\begin{equation}\label{ari:eq:cs-isogeny}
 U_{\Lambda_{\mathrm{ell}}'}(\alpha z)=\sum_{v\in K}U_{\Lambda_{\mathrm{ell}}}(z+v).
\end{equation}
The equality holds away from the zeros and as an equality of locally
integrable functions at them. If $\Delta_{\Lambda_{\mathrm{ell}}}=g_2(\Lambda_{\mathrm{ell}})^3-
27g_3(\Lambda_{\mathrm{ell}})^2$, then
\[
 \prod_{0\ne v\in K}\mathcal N_{\Lambda_{\mathrm{ell}}}(v)
 =|\alpha|\left|\frac{\Delta_{\Lambda_{\mathrm{ell}}'}}{\Delta_{\Lambda_{\mathrm{ell}}}}\right|^{1/12}.
\]
In particular, for every nonzero endomorphism of one elliptic curve,
$\prod_{v\in\ker\alpha,\ v\ne0}\mathcal N(v)=|\alpha|$.
\end{proposition}
\noindent\textit{Uses:} the theta product and Haar probability measure.\par
\begin{proof}
\noindent\textbf{Step 1.} The Jacobi product, with $q_\tau=e^{2\pi i\tau}$ and
$X=e^{2\pi iz}$, gives
\[
 \vartheta _1(z,\tau)/\eta_D(\tau)
 =i q_\tau^{1/12}X^{-1/2}(1-X)
     \prod_{m\ge1}(1-q_\tau^mX)(1-q_\tau^mX^{-1}),
\]
where a harmless fixed unit changes with the theta convention. For
$0<y<T$, integration of the logarithm over $z=x+iy$, $0\le x\le1$,
gives $-\pi T/6+\pi y-\pi y^2/T$. Each product factor has mean logarithm
zero by the absolutely convergent logarithmic series; the endpoint
singularities are integrable. Integrating in $y$ proves
$\int U_{\Lambda_{\mathrm{ell}}}\,d\mu_{\Lambda_{\mathrm{ell}}}=0$ for Haar probability measure.
Moreover
\[
 \Delta_{\rm Eucl}U_{\Lambda_{\mathrm{ell}}}
   =2\pi\delta_0-\frac{2\pi}{\operatorname{area}\Lambda_{\mathrm{ell}}},\qquad
 U_{\Lambda_{\mathrm{ell}}}(z)=\log|z|+\frac1{12}\log|\Delta_{\Lambda_{\mathrm{ell}}}|+o(1).
\]
The second formula follows from
$\vartheta _1'(0,\tau)=2\pi\eta_D(\tau)^3$ and
$\Delta_{\mathbb Z+\tau\mathbb Z}=(2\pi)^{12}\eta_D(\tau)^{24}$,
then scaling by $\omega$. The two sides of
\eqref{ari:eq:cs-isogeny} have the same logarithmic masses at $K$ and the
same constant Laplacian, since
$d=|\alpha|^2\operatorname{area}\Lambda_{\mathrm{ell}}/
\operatorname{area}\Lambda_{\mathrm{ell}}'$. Their difference is a smooth harmonic
function on the compact torus, hence constant. Both means vanish:
translations preserve Haar measure and an isogeny sends Haar
probability measure to Haar probability measure. The constant is zero.
Taking $z\to0$ and cancelling the logarithmic singularity gives the
product. The kernel, Laplacian and mean-zero argument fix the normalization.
\end{proof}

\begin{proposition}[ARI C.1.2; T; the phases and the two periods]\label{ari:C.1.2}
With the branch continued from the specified cusp, the exact values are
\begin{align*}
 \frac{\vartheta _1((2-\rho)/3,\rho)}{\eta_D(\rho)}
   &=e^{\pi i/12}\,3^{1/4}e^{\pi/(6\sqrt3)},\\
 \frac{\vartheta _1((1-i)/2,i)}{\eta_D(i)}
   &=\sqrt2\,e^{\pi/4},\\
 \beta^{\rm part}(z_1)
   &=-\frac12+i\left(\frac1{\sqrt3}-\frac3{2\pi}\log48\right),\\
 \beta^{\rm part}(z_2)
   &=i\left(\frac32-\frac3\pi\log6\right).
\end{align*}
Consequently $\mathcal N(z_1)=3^{1/4}$, $\mathcal N(z_2)=\sqrt2$,
$D_{\rm part}(z_1)=-3\log48/(2\pi)$, and
$D_{\rm part}(z_2)=-3\log6/\pi$.
\end{proposition}
\noindent\textit{Uses:} ARI C.1.1; FAM C.1.1; the fixed cusp branch.\par
\begin{proof}
\noindent\textbf{Step 1.} By ARI C.1.1, multiplication by $1+\rho$ has matrix
$\left(\begin{smallmatrix}1&-1\\1&2\end{smallmatrix}\right)$,
its kernel is $\{0,\mu,-\mu\}$, and $(1+\rho)\mu=1$ for
$\mu=(2-\rho)/3$. Statement ARI C.1.1 gives $\mathcal N(\mu)^2=\sqrt3$.
For the requested parallel calculation at $i$, multiplication by $1+i$
has kernel $\{0,(1-i)/2\}$ and absolute value $\sqrt2$; the same
identity gives $\mathcal N((1-i)/2,i)=\sqrt2$ without a separate theta-constant
evaluation.


\noindent\textbf{Step 2.} At $\rho$, $\mu=1/2-i/(2\sqrt3)$ and $q_\tau=-e^{-\pi\sqrt3}$. In
\[
 R(z,\tau)=2e^{\pi i\tau/6}\sin(\pi z)
 \prod_{m\ge1}\bigl(1-2q_\tau^m\cos(2\pi z)+q_\tau^{2m}\bigr),
\]
$\sin(\pi\mu)$ is positive real. The factors are
$(1+q_\tau^m e^{\pi/\sqrt3})(1+q_\tau^m e^{-\pi/\sqrt3})$, both positive:
when $m$ is odd the larger subtracted number is at most
$e^{-2\pi/\sqrt3}<1$. Thus the phase is exactly $\pi/12$ modulo
$2\pi$. There is no winding ambiguity in the continued logarithm.
Indeed continue from the cusp along $\tau=1/2+iT$,
$T\ge\sqrt3/2$. Real conjugation of the rhombic lattice keeps the
chosen zero on $\mu=1/2-iy$, $0<y<T$. It cannot leave that component:
a collision of the two zeros would require $g_3=0$, which does not
occur on this ray. Indeed, writing $q_\tau=-r$, $0<r<1/200$,
the tail after the first term of the $E_6$ series has absolute value
at most $504\cdot32r^2/((1-r^2)(1-8r))<1/2$, while its
first term is positive. Thus $E_6>1/2$. Hitting $y=0$ would again give a half-period zero, and
hitting $y=T$ would give a pole. Each displayed product factor remains
positive because $e^{-2\pi T}e^{2\pi y}<1$. The cusp phase tends to
$\pi/12$, so the continued argument has that value. On the imaginary
ray $\tau=iT$, $T\ge1$, the same branch is
$\mu=1/2-iy$ with $0<y\le T/2$; every product factor is positive,
so its argument at $i$ is zero. These checks supply the branch step
used in FAM C.1.1. The Gaussian factors are respectively
$e^{\pi/(6\sqrt3)}$ and $e^{\pi/4}$. Substitution in the definition of
$\beta^{\rm part}$ proves all four identities.
\end{proof}

\subsection{The Albanese nomes}

\begin{proposition}[ARI C.2.1; T; exact nomes and continuation radius]\label{ari:C.2.1}
For $w=e^{-\pi ic}$, the fixed branches give
\begin{align*}
 \kappa_1&=c+\tfrac12-\frac{3i}{2\pi}\log48,&
 q_1&=-w^2/48^3,\\
 \kappa_2&=c+\tfrac32-\frac{3i}{\pi}\log6,&
 q_2&=iw/216,\qquad q_1=(27/64)q_2^2.
\end{align*}
The family is parameterized by $0<|q_2|<1$ and exists precisely in
that domain. Writing $J(q_\tau)=q_\tau^{-1}+744+\cdots$, one has
\[
 j(\operatorname{Alb}S_1)=J(-w^2/48^3).
\]
This function is holomorphic on $0<|w|<48^{3/2}$, has a pole of order
two at zero, and has a natural boundary at $|w|=48^{3/2}$.
In particular it continues holomorphically through every point of the
family boundary $|w|=216$.
\end{proposition}
\noindent\textit{Uses:} FAM A.4.1 and A.5.1; ARI C.1.2; the maximum theorem stated in the Result.\par
\begin{proof}
\noindent\textbf{Step 1.} The precise Albanese lattices in FAM A.4.1 and A.5.1 are verified by projecting all
translations on the product covers. At the order-three fibre the
projection is $(1/3)(\mathbb Z+\kappa_1\mathbb Z)$; multiplication by
three gives the first stated lattice. At the order-four fibre it is
$\tfrac12\mathbb Z+\tfrac14\kappa_2\mathbb Z$; multiplication by two
gives $\mathbb Z+(\kappa_2/2)\mathbb Z$. Thus the positive-imaginary
period representatives are $-\kappa_1$ and $-\kappa_2/2$.
Adding $2\mu_1=1-i/\sqrt3$ and
$3\mu_2=3/2-3i/2$ to C.1.2 proves the two $\kappa$ identities,
including their real parts. Exponentiation gives the signs in $q_1,q_2$;
the ratio is $216^2/48^3=27/64$.


\noindent\textbf{Step 2.} The factor $1/2$ in $q_2=iw/216$ is exactly the Albanese projection just computed. The natural-boundary argument uses modular invariance at rational boundary points. The existence criterion and exact $2\mathbb Z$
identification are FAM C.2.3 and FAM A.4.1 and A.5.1; the latter is also Theorem~1.4
(label \texttt{p2:thm:exact-period}) of
\cite[Theorem~1.4 (\texttt{p2:thm:exact-period})]{DGKXv1}. Its hypotheses specify precisely the
retained period family, affine finite twists and cusp collar. They
coincide with the family in this prompt; the normalization differs
only by the fixed choice of the particular solution, which does not
alter parameter differences. Version~2 \cite[Introduction]{DGKXv2} omits that theorem; we infer
no mathematical conclusion from this omission.


\noindent\textbf{Step 3.} The Fourier expansion for $J$ converges for $0<|q_\tau|<1$ and has a simple
pole at $q_\tau=0$. The inverse image under $q_\tau=-w^2/48^3$ is the claimed
punctured disc and the pole has order two. The inequality
$216<48^{3/2}$ is equivalent to $27<64$. At every nonzero point on the
larger boundary the map to $q_\tau$ is locally biholomorphic. A continuation
there would continue $J$ through the unit circle, contradicting the
same rational-cusp argument. The phrase ``up to the
radius'' therefore means the open punctured disc, not holomorphy on
its bounding circle or at its central pole.
\end{proof}

\subsection{Division values}
For division-polynomial calculations we use the rational model of ARI C.4:
\[
E^{\rm div}_t:\quad y_{\rm div}^2=x_{\rm div}^3+A(t)x_{\rm div}+B(t),\quad
A=-\tfrac34t^3(t-1),\quad B=\tfrac14t^4(t-1)^2,
\]
\[
P^{\rm div}=(0,t^2(t-1)/2),\qquad \Delta_{\rm div}(t)=27t^8(t-1)^3.
\]
Its coordinates are distinguished from the cubic coordinates in FAM B.1.1. The notation sheet records the homothety and the sign of its distinguished section; the invariant norm is the same. We use $\lambda_\infty(Q)=-\log\mathcal N(Q)$ as the normalized local-height notation of ARI C.4.1.

\begin{proposition}[ARI C.4.2; T; the full division identity]\label{ari:C.4.2}
Let $\psi_m$ be the standard division polynomial of a smooth
Weierstrass model, normalized by $\psi_1=1$, $\psi_2=2y$.
For $m\ge2$ and $Q,mQ\ne O$,
\begin{equation}\label{ari:eq:cs-division}
 \mathcal N(Q)^{m^2}
       =\mathcal N(mQ)\frac{|\Delta_E|^{(m^2-1)/12}}{|\psi_m(Q)|}.
\end{equation}
For the displayed family this specializes to
\[
 \mathcal N(P^{\mathrm{div}})^4=3^{3/4}|t-1|^{-1/4}\mathcal N(2P^{\mathrm{div}}),\qquad
 \mathcal N(P^{\mathrm{div}})^9=16|t|^{-2/3}\mathcal N(3P^{\mathrm{div}}).
\]
If $P^{\mathrm{div}}_t$ has exact order $n$ and $m\ge2$ satisfies
$m\equiv1$ or $-1\pmod n$, then
\begin{equation}\label{ari:eq:cs-torsionvalue}
 \mathcal N(P^{\mathrm{div}}_t)^{12}
       =\frac{|\Delta_{\mathrm{div}}(t)|}{|\psi_m(P^{\mathrm{div}}_t)|^{12/(m^2-1)}}.
\end{equation}
The positive value on the right is independent of $m$.
\end{proposition}
\noindent\textit{Uses:} ARI C.1.1; the displayed rational Weierstrass model and standard division-polynomial normalization.\par
\begin{proof}
\noindent\textbf{Step 1.} Use the complex uniformization $x=\wp_{\Lambda_{\mathrm{ell}}}(z)$,
$y=\wp'_{\Lambda_{\mathrm{ell}}}(z)/2$. The rational function $\psi_m$ has divisor
$\sum_{0\ne Q\in E[m]}(Q)-(m^2-1)(O)$ and leading Laurent term
$(-1)^{m-1}m z^{1-m^2}$. One way to verify both normalizations without
a model convention is to form
$(-1)^{m-1}\sigma_{\Lambda_{\mathrm{ell}}}(mz)/\sigma_{\Lambda_{\mathrm{ell}}}(z)^{m^2}$.
Its lattice factors cancel, its zeros are exactly the nonzero
$m$-division points, and its expansion at zero is the displayed
one. The addition formula for $\sigma$ gives the standard
division-polynomial recurrence with initial values $1,2y$, so it is
the stated rational function.


\noindent\textbf{Step 2.} The function
\[
 m^2U_{\Lambda_{\mathrm{ell}}}(z)-U_{\Lambda_{\mathrm{ell}}}(mz)
                    +\log|\psi_m(Q(z))|
\]
is harmonic after its logarithmic singularities cancel. Its
Laplacian has zero constant term as well, so it is constant.
The expansion in ARI C.1.1 and the leading term of $\psi_m$ make that
constant $(m^2-1)\log|\Delta_E|/12$. This proves
\eqref{ari:eq:cs-division}. Direct substitution of
\[
 \psi_2(P^{\mathrm{div}})=t^2(t-1),\qquad
 \psi_3(P^{\mathrm{div}})=-9t^6(t-1)^2/16
\]
gives the two special identities. For $m\equiv\pm1\pmod n$,
$mP^{\mathrm{div}}=\pm P^{\mathrm{div}}$ and $\mathcal N$ is even. Cancelling $\mathcal N(P^{\mathrm{div}})$ in
\eqref{ari:eq:cs-division} gives \eqref{ari:eq:cs-torsionvalue}.
Every right side is the same positive number $\mathcal N(P^{\mathrm{div}})^{12}$, which
proves independence of $m$, not merely numerical agreement.
\end{proof}

\begin{proposition}[ARI C.4.3; T; algebraicity and explicit fields for every order]\label{ari:C.4.3}
Let $t_0\in\mathbb C\setminus\{0,1\}$ and suppose $P^{\mathrm{div}}_{t_0}$ has
exact order $n\ge2$. Then $t_0$ and $\mathcal N(P^{\mathrm{div}}_{t_0})$ are algebraic.
Put $K_0=\mathbb Q(t_0,\bar t_0)$,
\[
 W_{\rm rad}=\Delta_{\mathrm{div}}(t_0)\Delta_{\mathrm{div}}(\bar t_0)>0,\qquad
 V_{\rm rad}=\psi_{n+1}(P^{\mathrm{div}}_{t_0})\psi_{n+1}(P^{\mathrm{div}}_{\bar t_0})>0.
\]
With positive real radicals,
\[
 \mathcal N(P^{\mathrm{div}}_{t_0})^{12n}=W_{\rm rad}^{n/2}V_{\rm rad}^{-6/(n+2)}
 \ \in\ K_0\bigl(W_{\rm rad}^{1/2},V_{\rm rad}^{1/(n+2)}\bigr).
\]
This supplies an explicit field, without a claim that it is minimal.
More generally, if $m\equiv\pm1\pmod n$, then
\[
 \lambda_\infty(P^{\mathrm{div}}_{t_0})
 =\frac{\log|\psi_m(P^{\mathrm{div}}_{t_0})|}{m^2-1}
       -\frac1{12}\log|\Delta_{\mathrm{div}}(t_0)|.
\]
These formulas determine the height and $\mathcal N^{12}$ on the entire
torsion locus.
\end{proposition}
\noindent\textit{Uses:} ARI C.4.2; the nonreal cusp value of the chosen elliptic logarithm.\par
\begin{proof}
\noindent\textbf{Step 1.} The division polynomial $\psi_n(P^{\mathrm{div}}_t)$ is a nonzero rational
polynomial in $t$. To check nonvanishing for every $n$, a hypothetical
torsion identity would express the elliptic logarithm as
$a\tau+b$ with constant rational $a,b$. Its bounded cusp limit
forces $a=0$, but that limit has nonzero imaginary part, a
contradiction. Thus $t_0$ is a zero of a nonzero polynomial over
$\mathbb Q$. Formula \eqref{ari:eq:cs-torsionvalue} shows algebraicity
of $\mathcal N$ using only positive rational powers of nonzero algebraic
absolute values. Taking $m=n+1$, raising to the $n$th power, and
using $|\Delta|=W_{\rm rad}^{1/2}$, $|\psi_{n+1}|=V_{\rm rad}^{1/2}$ gives the exact
field statement. The local-height formula is its logarithm.


\noindent\textbf{Step 2.} Here is a polynomial relation involving only $t$ and $r=\mathcal N^{12}$,
if one prefers that form to positive radicals. Remove the factors
at $0,1$ from $\psi_n(P^{\mathrm{div}}_t)$ and divide successively by the
corresponding polynomials for proper divisors of $n$; let $F_n$ be
the resulting primitive polynomial. At a point of exact order
$d\mid n$, the formal group gives
$z([n/d]Q)=(n/d)z(Q)+O(z(Q)^2)$, so the vanishing multiplicities for
$\psi_d$ and $\psi_n$ agree. Hence this cancellation leaves exactly
the exact-order-$n$ divisor, with its multiplicities. Set
$m=n+1$ and $e=m^2-1$. Clearing rational denominators in
\[
 \operatorname{Res}_S\left(
 F_n(S),\
 r^{2e}\{\psi_m(P^{\mathrm{div}}_T)\psi_m(P^{\mathrm{div}}_S)\}^{12}
              -\{\Delta_{\mathrm{div}}(T)\Delta_{\mathrm{div}}(S)\}^{e}\right)
\]
gives a polynomial in $\mathbb Q[T,r]$ that vanishes at
$(t_0,\mathcal N(P^{\mathrm{div}}_{t_0})^{12})$, since $S=\bar t_0$ is a common root.
It is nonzero: at an exact-order-$n$ root no $\psi_m$ vanishes,
so the leading coefficient in $r$ of the resultant is nonzero.
This is an explicit recurrence-and-resultant prescription for
every $n$, including nonreal torsion parameters.
\end{proof}

\begin{proposition}[ARI C.4.4; T; the five values and all orders through eight]\label{ari:C.4.4}
The exact special values are
\[
 \begin{array}{c|c|c}
 t&\operatorname{ord}(P^{\mathrm{div}})&\mathcal N(P^{\mathrm{div}})^{12}\\ \hline
 0&3&27\\
 1&2&64\\
 32/5&4&10\\
 64(47-27\sqrt3)/11&6&17+10\sqrt3\\
 64(47+27\sqrt3)/11&6&10\sqrt3-17 .
 \end{array}
\]
At $t=0,1$ the orders refer to the smooth models after the specified
base changes. On the smooth base, $F_2=F_3=1$ and
\[
 \begin{aligned}
 F_4&=5t-32,\\
 F_5&=89t^2+4736t-4096,\\
 F_6&=11t^2-6016t+8192,\\
 F_7&=881t^4+4745344t^3-12079104t^2\\
    &\hspace{13mm}+24641536t-16777216,\\
 F_8&=49t^4-2614336t^3-7864320t^2\\
    &\hspace{13mm}+19398656t-8388608.
 \end{aligned}
\]
All these polynomials are squarefree. At each root of exact order
$n=4,5,6,7,8$, the positive number $r=\mathcal N^{12}$ is respectively
determined by
\[
 \begin{array}{c|l}
 n&\text{equation for }r\\ \hline
 4&r=64/|t|\\
 5&r^5=432^5/\bigl((1-t)(5t-32)^4\bigr)\\
 6&r^2=1728^2/F_5(t)\\
 7&r^{35}=2^{192}3^{81}/
       \bigl(|t|^8|t-1|^3|F_6(t)|^{12}\bigr)\\
 8&r^4=1728^4/|F_7(t)|.
 \end{array}
\]
The order-five roots are real and below one; the order-six roots
are real and above one. The two complex order-seven roots are
included in the formula.
\end{proposition}
\noindent\textit{Uses:} ARI C.1.2, C.4.2 and C.4.3; the division-polynomial recurrence.\par
\begin{proof}
\noindent\textbf{Step 1.} The division-polynomial recurrence, evaluated at $x=0$,
$y=t^2(t-1)/2$, gives
\[
 \begin{aligned}
 \psi_4(P^{\mathrm{div}})&=-t^{10}(t-1)^4F_4/32,\\
 \psi_5(P^{\mathrm{div}})&=t^{16}(t-1)^6F_5/4096,\\
 \psi_6(P^{\mathrm{div}})&=9t^{24}(t-1)^9F_6/65536,\\
 \psi_7(P^{\mathrm{div}})&=-t^{32}(t-1)^{12}F_7/16777216,\\
 \psi_8(P^{\mathrm{div}})&=-t^{42}(t-1)^{16}F_4F_8/268435456.
 \end{aligned}
\]
The recursion and all coefficients are checked over $\mathbb Q[t]$
by \path{Cspecial_division.py}. The same program computes
$\gcd(F_n,F_n')=1$ and excludes common factors for proper orders.
Thus no point or multiplicity is omitted. Applying
\eqref{ari:eq:cs-torsionvalue} with $m=n-1$ gives every displayed
formula. The signs for $F_5,F_6$ follow directly from their quadratic
formula; in particular $F_5$ is positive at the two roots of $F_6$.
At order six, substituting
$t=64(47\mp27\sqrt3)/11$ gives the exact identities
\[
 F_5(t)(17\pm10\sqrt3)^2=1728^2.
\]
The positive roots are $17+10\sqrt3$ and $10\sqrt3-17$; their
product is $11$. The substitution is checked in
\path{Cspecial_relations.py}. The order-four value follows from
$64/(32/5)=10$, and the two central values are C.1.2.


\noindent\textbf{Step 2.} For clarity, the recurrence used by the exact program is
\[
 \begin{aligned}
 \psi_{2m+1}&=\psi_{m+2}\psi_m^3-\psi_{m-1}\psi_{m+1}^3,\\
 \psi_{2m}&=\frac{\psi_m}{\psi_2}
   (\psi_{m+2}\psi_{m-1}^2-\psi_{m-2}\psi_{m+1}^2).
 \end{aligned}
\]
Its initial values at $P^{\mathrm{div}}$ are $\psi_0=0,\psi_1=1,\psi_2=2y$,
$\psi_3=-A^2$, $\psi_4=4y(-8B^2-A^3)$. No approximate algebra
enters this calculation. The same code records $\psi_9$, and the
separate file \texttt{Cspecial\_psi16.out} records the factored
$\psi_{16}$ for the duplication ladder. A rational polynomial
relation for the nonreal order-seven values is explicitly computed
by \path{Cspecial_relations.py}, by the resultant prescription
of C.4.3.
\end{proof}

\subsection{Integral identities}

\begin{proposition}[ARI C.8.6; T; the Mahler-measure identities]\label{ari:C.8.6}
Let $\Lambda_{\mathrm{ell}}\subset\C$ be a lattice of area $A$, let $v$ be a zero
of $\wp_{\Lambda_{\mathrm{ell}}}$, and let $d\mu=A^{-1}dA_z$ be Haar probability
measure on $\C/\Lambda_{\mathrm{ell}}$. Put $U(z)=\log \mathcal N(z)$ and
$\mathcal L=U(v)$. Then
\begin{align}
 6\mathcal L&=\int_{\C/\Lambda_{\mathrm{ell}}}\log\left|\frac{\wp'(z)^2}{\wp(z)^3}\right|d\mu(z),
 \label{ari:eq:ca-mahler}\\
 \mathcal N(v)^2&=|\Delta(\Lambda_{\mathrm{ell}})|^{1/6}
      \exp\left(-\int_{\C/\Lambda_{\mathrm{ell}}}\log|\wp(z)|d\mu(z)\right).
 \label{ari:eq:ca-wpmean}
\end{align}
If $C(v_{\wp})=1-(g_2/4)v_{\wp}^2-(g_3/4)v_{\wp}^3$, the pushforward of $d\mu$
under $v_{\wp}=1/\wp(z)$ is
\begin{equation}\label{ari:eq:ca-nu}
 d\nu(v_{\wp})=\frac1{2A}\frac{dA_{v_{\wp}}}{|v_{\wp}|\,|C(v_{\wp})|}.
\end{equation}
In particular its unnormalized integral is exactly $2A$, and
\[
 \mathcal N(v)^6=4\exp\left(\int_{\C}\log|C(v_{\wp})|d\nu(v_{\wp})\right).
\]
Writing $e_1,e_2,e_3$ for the three branch values and
$p(x)=\int\log|x-\wp(z)|d\mu(z)$, the maximum theorem is equivalently
\[
 \int\log|C|d\nu\leq\log2,\qquad
 p(0)\geq\frac{p(e_1)+p(e_2)+p(e_3)}3-\frac{\log2}3.
\]
Equality in either inequality occurs exactly when $g_3=0$.
\end{proposition}
\noindent\textit{Uses:} ARI C.1.1; the maximum theorem stated in the Result.\par
\begin{proof}
\noindent\textbf{Step 1.} The mean of $U$ is zero by the theta product calculation in C.1,
and near zero
$U(z)=\log|z|+\frac1{12}\log|\Delta|+o(1)$.
The divisor of $\wp$ is $(v)+(-v)-2(O)$, with multiplicity
two when $v=-v$. Subtract its logarithmic divisor from the
corresponding combination of $U$'s. The result is harmonic
on the compact torus and hence constant. Its expansion at
zero gives
\[
 \log|\wp(z)|=U(z-v)+U(z+v)-2U(z)
                  +\frac16\log|\Delta|-2\mathcal L.
\]
Integration proves \eqref{ari:eq:ca-wpmean}.
Similarly the divisor of $\wp'$ is the sum of the three
nonzero half-periods minus $3(O)$. The multiplication-by-two
identity of C.1 gives the product of $\mathcal N$ at those half-periods
as $2$. Since $\wp'(z)\sim-2z^{-3}$, comparison at zero yields
\[
 \log|\wp'(z)|=\sum_{w\in E[2]\setminus\{O\}}U(z-w)-3U(z)
                         +\frac14\log|\Delta|.
\]
Its mean is $\frac14\log|\Delta|$. Twice this identity minus
three times the preceding mean proves \eqref{ari:eq:ca-mahler}.
All logarithmic singularities are locally integrable, so these
integrations also cover a zero at a half-period.


\noindent\textbf{Step 2.} The map $z\mapsto x=\wp(z)$ has degree two. Its two inverse
sheets each have Jacobian $|\wp'(z)|^{-2}$; branch values form
a measure-zero set. Thus its pushforward density is
$2/(A|4x^3-g_2x-g_3|)$. Under $x=1/v_{\wp}$ the Jacobian is
$|v_{\wp}|^{-4}$ and the denominator is $4|v_{\wp}|^{-3}|C(v_{\wp})|$.
This proves \eqref{ari:eq:ca-nu} and its normalization. The identity
$\wp'^2/\wp^3=4C(v_{\wp})$ gives the next formula.


\noindent\textbf{Step 3.} Replacing $v$ above by a half-period gives
$p(e_j)=\frac16\log|\Delta|-2U(w_j)$; summing and using
$\sum_jU(w_j)=\log2$ proves the centroid formulation.
The equality characterization is that of C.8.5,
since $j=1728$ is equivalent to $g_3=0$ for a nonsingular lattice.
\end{proof}

\begin{proposition}[ARI C.8.7; T; the two complex beta integrals]\label{ari:C.8.7}
For the square and equianharmonic lattices respectively, the
integrals in C.8.6 are
\[
 \int\log|C|\,d\nu=\log2,\qquad
 \int\log|C|\,d\nu=\frac32\log3-2\log2.
\]
They recover $\mathcal N^2=2$ and $\mathcal N^4=3$ with the same normalization.
\end{proposition}
\noindent\textit{Uses:} ARI C.8.6; the gamma duplication and triplication formulas.\par
\begin{proof}
\noindent\textbf{Step 1.} By homogeneity it suffices to take $C(v_{\wp})=1-v_{\wp}^m$, with $m=2$
or $m=3$. For $-1<\Re s<1-1/m$ put
\[
 I_m(s)=\int_{\C} |v_{\wp}|^{-1}|1-v_{\wp}^m|^{s-1}dA_{v_{\wp}}.
\]
The substitution $w=v_{\wp}^m$, with all $m$ sheets counted, gives
\[
 I_m(s)=\frac1m B_{\C}\left(\frac1{2m},\frac{s+1}2\right),
 \quad
 B_{\C}(a,b)=\pi\frac{\Gamma(a)\Gamma(b)\Gamma(1-a-b)}
                   {\Gamma(1-a)\Gamma(1-b)\Gamma(a+b)}.
\]
We include a proof of the integral evaluation. For
$0<a,b$ and $a+b<1$, insert
\[
 |w|^{2a-2}=\frac1{\Gamma(1-a)}
       \int_0^\infty t^{-a}e^{-t|w|^2}\,dt
\]
and the corresponding formula for $|1-w|^{2b-2}$.
Tonelli's theorem applies to the positive integrand. The
Gaussian integral over $w$ is
$\pi(t+s)^{-1}\exp(-ts/(t+s))$.
Putting $t=rv$, $s=r(1-v)$ reduces the remaining integrals
to $\Gamma(1-a-b)$ and the beta integral $B(a,b)$, proving
the displayed formula. The same integrable bounds justify
differentiation in a neighbourhood of $s=0$.


\noindent\textbf{Step 2.} Consequently, with $a=1/(2m)$ and $\psi=\Gamma'/\Gamma$,
\[
 \frac{I_m'(0)}{I_m(0)}
 =\frac12\left(2\psi(1/2)-\psi(1/2-a)-\psi(1/2+a)\right).
\]
The gamma duplication and triplication formulas, differentiated
logarithmically, evaluate this as $\log2$ for $m=2$ and
$\frac32\log3-2\log2$ for $m=3$. These can also be checked
directly from
$\psi(1/2)=-\gamma-2\log2$,
$\psi(1/4)+\psi(3/4)=-2\gamma-6\log2$ and
$\psi(1/3)+\psi(2/3)=-2\gamma-3\log3$.
Equation \eqref{ari:eq:ca-mahler} completes the normalization test.
\end{proof}

\subsection{Announcements on the norm and the Betti measure}

\begin{announcement}[ARI C.8.8; A; the signed Gaussian identity]\label{ari:C.8.8}
For $\Lambda_{\mathrm{ell}}=\Z+\tau\Z$, $T=\Im\tau$, and every $z\in\C$,
\[
 \sqrt{2T}\,|\eta_D(\tau)|^2\mathcal N(z,\tau)^2
 =\sum_{m\in\frac12\Lambda_{\mathrm{ell}}}\epsilon(m)
                         e^{-2\pi|z-m|^2/T},
\]
where $\epsilon=-1$ on $\Lambda_{\mathrm{ell}}$ and $+1$ on the three other
cosets of $\Lambda_{\mathrm{ell}}$ in $\frac12\Lambda_{\mathrm{ell}}$.
Both sides and all their derivatives converge locally uniformly.
\end{announcement}
\noindent\textit{Uses:} the theta series and Poisson summation.\par
The mechanism is comparison of Fourier coefficients on the real torus. The signed sum of the four half-period translates of the Gaussian has the same coefficients as the squared theta series with its Gaussian normalization.
\noindent\textit{Proof in the repository:} \href{https://github.com/KokunoYumeto/s6-complex-structure-workbench}{S6 complex-structure workbench}, S6\_05\_ARITHMETIC, Statement C.8.8.

We take $dd^c=(i/\pi)\partial\bar\partial$ and $dA_{z_B}=(i/2)dz_B\wedge d\bar z_B$ (ARI C.8, measure subsection). Here $z_B$ is a local base coordinate; it is unrelated to the global parameter $w$.

\begin{announcement}[ARI C.8.9; A]\label{ari:C.8.9}
On a local period chart write $\mu=a\tau+b$, with $a,b$ real,
and put $T=\Im\tau$. The positive Betti measure is
\[
 m_{P^{\mathrm{div}}}=db\wedge da
 =\frac{|\mu_{z_B}-a\tau_{z_B}|^2}{T}\,dA_{z_B}.
\]
Thus the form $da\wedge db$ in the source request labelled Lead 15(b) has the wrong sign.
The globally single-valued function $\mathcal L(t)=\log \mathcal N(P^{\mathrm{div}}_t)$ satisfies
on $\mathbb{P}^1(\C)$
\begin{equation}\label{ari:eq:ca-current}
 dd^c\mathcal L=\frac1{12}\delta_\infty-m_{P^{\mathrm{div}}},\qquad
               \int_{\mathbb{P}^1}m_{P^{\mathrm{div}}}=\frac1{12}.
\end{equation}
There are no atoms of $m_{P^{\mathrm{div}}}$ at $0$, $1$ or $\infty$.
Let $T_0=\sqrt3/2$, $\rho=(1+i\sqrt3)/2$, and
$p=\wp'_\rho((2-\rho)/3)$. The density in the ordinary
$t$ coordinate satisfies
\begin{align}
 \frac{dm_{P^{\mathrm{div}}}}{dA_t}
 &=\frac{|t|^{-2/3}}{64T_0|p|^{2/3}}
                      \bigl(1+O(|t|^{1/3})\bigr),
                         &&t\longrightarrow0,\label{ari:eq:ca-density0}\\
 \frac{dm_{P^{\mathrm{div}}}}{dA_t}
 &=\frac{|t-1|^{-3/2}}{24\pi^2\sqrt{E_4(i)}}
                      \bigl(1+O(|t-1|^{1/2})\bigr),
                         &&t\longrightarrow1.\label{ari:eq:ca-density1}
\end{align}
\end{announcement}
\noindent\textit{Uses:} the normalized theta norm; local root expansions and the cusp branch.\par
The mechanism is the curvature of the real Gaussian factor and the Jacobian of the two Betti coordinates. The cusp logarithm supplies the point mass, and the two finite root expansions supply the stated integrable densities; we retain the corrected sign and total mass from the source.
\noindent\textit{Proof in the repository:} \href{https://github.com/KokunoYumeto/s6-complex-structure-workbench}{S6 complex-structure workbench}, S6\_05\_ARITHMETIC, Statement C.8.9.

\begin{announcement}[ARI C.8.10; A; the correctly normalized equidistribution]\label{ari:C.8.10}
Let $D_n$ be the divisor, with intersection multiplicity, of
parameters $t\in\mathbb{P}^1\setminus\{0,1,\infty\}$ for which
$nP^{\mathrm{div}}_t=O$. Then
\[
 \deg D_n=\frac{n^2}{12}+O(1),\qquad
 \frac1{n^2}[D_n]\ \longrightarrow\ m_{P^{\mathrm{div}}},\qquad
 \frac1{\deg D_n}[D_n]\ \longrightarrow\ 12m_{P^{\mathrm{div}}}
\]
weakly as $n\to\infty$ through indices for which $D_n\ne0$.
For any nonrepeating sequence of individual algebraic torsion
parameters, their Galois-orbit probability measures likewise
converge to $12m_{P^{\mathrm{div}}}$. Thus the probability-limit assertion
in the source request labelled Lead 15(b) requires the factor $12$.
\end{announcement}
\noindent\textit{Uses:} ARI C.8.9; the height metric and its equidistribution theorem.\par
The mechanism applies the height metric and the small-height equidistribution theorem of DeMarco and Mavraki \cite[Theorem~1.1 and Corollary~1.2]{DM20}, with the normalization of ARI C.8.9. The passage to full division divisors uses bounded multiplicities at a fixed torsion parameter.
\noindent\textit{Proof in the repository:} \href{https://github.com/KokunoYumeto/s6-complex-structure-workbench}{S6 complex-structure workbench}, S6\_05\_ARITHMETIC, Statement C.8.10.

\begin{announcement}[ARI C.8.11; A; the two first-derivative identities]\label{ari:C.8.11}
On an ordinary modulus chart put
\[
 a=\frac{\Im\mu}{T},\qquad
 \mathcal W=\frac{\vartheta_1'(\mu,\tau)}{\vartheta_1(\mu,\tau)}
                         +2\pi ia,\qquad
 \mathcal R_\mu=\frac{d\mu}{d\tau}-a.
\]
Then
\[
 \mathcal L_\tau=\frac{\mathcal W}{8\pi i}(\mathcal W+4\pi i\mathcal R_\mu),\qquad
 \det\frac{d}{d\tau}
 \begin{pmatrix}\tau&\mu\\\mu&\beta^{\rm part}/6\end{pmatrix}
       =\frac{(\mathcal W+2\pi i\mathcal R_\mu)^2}{4\pi^2}.
\]
Moreover
\[
 \mathcal R_\mu=\frac{g_2}{6\pi i\wp'(\mu)}-\frac{\mathcal W}{2\pi i},
 \qquad
 \mathcal L_\tau=\frac{i\mathcal W(3\mathcal W\wp'(\mu)-2g_2)}
                         {24\pi\wp'(\mu)}.
\]
The critical-point factors in the source checkpoint K4 follow with their stated
normalization after multiplying by $6/\pi$.
\end{announcement}
\noindent\textit{Uses:} the theta heat equation and the Weierstrass differential equation.\par
The mechanism differentiates the logarithm of the theta quotient and its real Gaussian factor, and then differentiates the zero equation for the chosen elliptic logarithm. The root chart handles a half-period, where the displayed quotients are interpreted by continuation.
\noindent\textit{Proof in the repository:} \href{https://github.com/KokunoYumeto/s6-complex-structure-workbench}{S6 complex-structure workbench}, S6\_05\_ARITHMETIC, Statement C.8.11.

\section{The exact moduli disc}\label{sec:maximum}

\subsection{The two interval certificates}
The known abstract orbit-disc description is credited to \cite[Corollary~\texttt{p2:cor:bare-period}]{DGKXv1}; it leaves the maximum unevaluated. The following maximum and its arithmetic certification are FAM C.2.3 and PUSH E.4.

\begin{computed}[FAM C.2.3; T]\label{fam:C.2.3}
For the cusp normalization $\beta^{\rm part}+\tau\to0$,
\[
 M=-\frac3\pi\log6=-1.71100425815613411418\ldots,
 \qquad \operatorname*{arg\,max}_{B\setminus\{p_0\}}D_{\rm part}=\{1\}.
\]
The exact parameter discs are
\[
 0<|w|<216,
 \qquad 0<|u|<46656
\]
for the reference linearization $u=w^2$; an arbitrary reference constant
$\lambda_0$ multiplies the latter radius by $|\lambda_0|$.
The proof combines an analytic cusp bound, the quantitative local proof
C.2.2, and a completed outward-rounded interval cover of the remaining
compact domain.
\end{computed}
\noindent\textit{Uses:} FAM C.1.1 and C.2.2; the interval cover described below.\par

\begin{computed}[PUSH E.4.1; T]\label{push:E.4.1}
For all 19701 rational rectangles in the supplied certificate, and for
each $\tau$ in the rectangle, both zeros of $\wp_\tau$ satisfy
$\mathcal N^2<2$. The new largest certified upper bound is less than
\[
 1.999877571767410525.
\]
The 19701 rectangles, 36 local rectangles and 174 excluded rectangles
have disjoint interiors and cover exactly
$[0,1/2]\times[17/20,2]$. Every local rectangle is contained in
$|\tau-i|<3/20000$, and every excluded rectangle satisfies
$|\tau|<1$ throughout.
\end{computed}
\noindent\textit{Uses:} the supplied rational rectangle endpoints; the theta norm; Arb interval arithmetic.\par

\paragraph{Certificate block.}
The source certificate is the JSON file \path{f_interval_parallel_certificate.json} named in the FAM report. Its verifier is \path{f_interval_verify.py}; the interval calculation is \path{f_interval_cert.py}, scheduled by \path{f_interval_parallel.py}. The certificate has 19,911 leaves: 19,701 norm boxes, 36 boxes in the analytic local square, and 174 boxes strictly inside the unit circle. The root rectangle is $[0,1/2]+i[17/20,2]$, its exact area is $23/40$, and its maximum subdivision depth is 38. There are zero pending boxes. At 30 decimal digits the largest squared-norm upper endpoint is
\[
\frac{9999999982782673130994674616307871}
 {5000000000000000000000000000000000}<2.
\]
The norm boxes use a uniform complex contraction enclosure for a zero of $\wp_\tau$, followed by an outward upper bound for $\mathcal N^2$. Sign and lattice invariance give the same norm at both zeros. The verifier reconstructs the partition and checks the strict bounds and rational exclusions. The local square is covered by FAM C.2.2; the Fourier-product estimate for $T\ge2$ gives $\mathcal N<1.082$. These are the three parts of the certified proof in FAM C.2.3.

PUSH E.4.1 reuses only the rational rectangle endpoints and recomputes roots, preconditioners, radii, contraction estimates, values and tails using Arb through python-flint 0.9.0. The two programs are \path{code/e4_arb_certificate.py} (192 bits) and \path{code/e4_cover_and_local.py} (256 bits). The cover is checked at all 438 distinct vertical edges; each open vertical slab is partitioned once. The independent largest norm bound is
\[
1.999877571767410524070536478244687423578792301873142913
\ \pm9.26\cdot10^{-56},
\]
which is strictly below two. The value belongs to the enclosure of an upper bound, not to an evaluation of the maximum. The midpoint-radius representation is described in \cite[§§2--3]{Johansson2016}; the constructor and outward endpoints are documented in \cite[\texttt{arb}, \texttt{abs\_upper}, \texttt{upper}]{PythonFlint2026}.

The supplied review records a re-execution of all 19,701 FAM boxes with identical upper-bound strings (73 seconds, 12 processes), by Claude Fable 5.2 at effort high, using Python 3.13.9, sympy 1.13.1 and mpmath 1.3.0. The source review records a read of the verifier and interval routines. The PUSH certification is a producing-run computation by GPT-6 Astra, ChatGPT Pro mode; the intake did not rerun that code.

The attached material contains the TeX descriptions and the status records, but contains neither certificate file nor its executable verifier. Consequently their byte sizes, SHA-256 digests and current executions cannot be supplied from these attachments. We have not reconstructed missing certificates or claimed a new execution. This limitation concerns the supplied computational artifacts; the full duplication proof follows in the text. Its presence does not change the designated proof of record.

\subsection{Local parameters for the duplication proof}
Put $G_\square=g_2(i)=\Gamma(1/4)^8/(16\pi^2)$ and $D_*=-3\log6/\pi$ (ARI C.3.1).

\begin{proposition}[ARI C.3.1; T; fourth order at the maximum]\label{ari:C.3.1}
On either root branch at $\tau=i$, let
$h=\mu-(1-\tau)/2$ and $\delta=\tau-i$. Then
\begin{align*}
 \tau&=i-\frac{3\pi i}{2}h^2-\frac{9\pi^2i}{8}h^4+O(h^6),\\
 h^2&=\frac{2i}{3\pi}\delta+\frac1{3\pi}\delta^2+O(\delta^3),\\
 \mu&=\frac{1-\tau}{2}
  \mathbin{\pm}\left(\frac{2i\delta}{3\pi}\right)^{1/2}
          \left(1-\frac{i}{4}\delta+O(\delta^2)\right).
\end{align*}
The square root is taken holomorphically on the chosen root chart.
Furthermore
\begin{align}\label{ari:eq:cs-Li}
 \log \mathcal N={}&\tfrac12\log2-\frac\pi2|h|^2
   +\frac{3\pi^2}{8}|h|^4
   -\frac{3\pi^2}{4}|h|^2\operatorname{Re}(h^2)
   -\frac{5G_\square}{192}\operatorname{Re}(h^4)+O(|h|^6),\\
 D_{\rm part}={}&D_* -\frac2\pi|\delta|
  +\frac1\pi\bigl(|\delta|^2+|\delta|\operatorname{Im}\delta\bigr)
  +\frac{5G_\square}{72\pi^3}\operatorname{Re}(\delta^2)
  +O(|\delta|^3).\nonumber
\end{align}
In particular the leading terms of the local expansion in the source prompt are correct.
\end{proposition}
\noindent\textit{Uses:} the theta heat equation; the half-period Weierstrass equation; the cusp-normalized branch.\par
\begin{proof}
\noindent\textbf{Step 1.} Write $e(\tau)=\wp_\tau((1-\tau)/2)$ and
$\mathcal E_2=\pi^2E_2/3$. The half-period Taylor series is
\[
 \wp_\tau((1-\tau)/2+h)
 =e+(3e^2-g_2/4)h^2+(3e^3-g_2e/4)h^4+O(h^6).
\]
The exact coefficients are
$e(i)=0$, $e'(i)=iG_\square/(6\pi)$. This proves the claimed identity
$3e'(i)=(2i/\pi)G_\square/4$. For clarity its closed form follows from
$e=\pi^2(\vartheta_2^4-\vartheta_4^4)/3$ and the heat equation:
if $A=\vartheta_2(i)^4=\vartheta_4(i)^4$, then
$G_\square=4\pi^4A^2$ and $e'(i)=2\pi^3iA^2/3$.
The classical beta integral
$K(1/2)=\Gamma(1/4)^2/(4\sqrt\pi)$ and
$\vartheta_3(i)^2=2K(1/2)/\pi$ give
$G_\square=16K(1/2)^4=\Gamma(1/4)^8/(16\pi^2)$.
The beta evaluation follows by changing variable in
$\int_0^\infty(1+x^4)^{-1/2}\,dx
=\tfrac14 B(1/4,1/4)$; it is the lemniscatic complete period.


\noindent\textbf{Step 2.} Ramanujan's identities give $g_2'(i)=2iG_\square$. Differentiating
$e(-1/\tau)=\tau^2 e(\tau)$ twice at $i$ gives
$e''(i)=3i e'(i)=-G_\square/(2\pi)$. Inserting
$\delta=\alpha_2h^2+\alpha_4h^4+O(h^6)$ into the displayed half-period series
therefore gives
\[
 \alpha_2=-3\pi i/2,\qquad
 \alpha_4=\frac{g_2'\alpha_2/4-e''\alpha_2^2/2}{e'}=-9\pi^2i/8.
\]
The implicit-function theorem applies because $e'(i)\ne0$;
evenness in $h$ gives the stated remainder. Series inversion gives
the formulas for $h^2$ and $\mu$.


\noindent\textbf{Step 3.} Let $H(h,\tau)=\log(\vartheta_3(h,\tau)/\eta_D(\tau))$. The exact
identity from the half-period translation is
$\log \mathcal N=\operatorname{Re}H-\pi(\operatorname{Im}h)^2/
\operatorname{Im}\tau$. At $(0,i)$ its relevant Taylor data are
\[
 \begin{gathered}
 H_0=\tfrac12\log2,\quad H_\tau=0,\quad
 [h^2]H=-\pi/2,\quad [h^4]H=G_\square/48,\\
 \partial_\tau[h^2]H=-\pi i/4-iG_\square/(16\pi),\qquad
 H_{\tau\tau}=-G_\square/(24\pi^2).
 \end{gathered}
\]
Indeed $\partial_h^2\log\vartheta_3=-\wp(h+(1-\tau)/2)-\mathcal E_2$;
the heat equation at $h=0$ gives $H_\tau=ie/(4\pi)$, and
$\mathcal E_2'(i)=\pi i/2-iG_\square/(24\pi)$. These identities prove every entry.
After substitution $\tau=\tau(h)$, the pure $h^4$ coefficient of
$H$ is $-3\pi^2/8-5G_\square/192$. Expanding
$\operatorname{Im}\tau=1-(3\pi/2)\operatorname{Re}(h^2)+O(|h|^4)$
in the Gaussian gives \eqref{ari:eq:cs-Li}. Finally
\[
 |h|^2=\frac{2}{3\pi}|\delta|
          \left(1+\tfrac12\operatorname{Im}\delta+O(|\delta|^2)\right)
\]
and substitution into the quartic expression, followed by
multiplication by $6/\pi$, gives the formula for $D_{\rm part}$.
The source run records the script \path{Cspecial_expansion_checks.py} independently checks
these coefficients against theta evaluations; its output is retained.
\end{proof}

\subsection{The three-step duplication proof}
Throughout this proof $q_\tau=e^{2\pi i\tau}$, $T=\operatorname{Im}\tau$, and
\[
\Delta_{\rm mod}=q_\tau\prod_{n\ge1}(1-q_\tau^n)^{24},\qquad
 t=\frac{E_4^3}{1728\Delta_{\rm mod}},\qquad
 t-1=\frac{E_6^2}{1728\Delta_{\rm mod}}.
\]
This modular discriminant differs from the discriminant of the rational division model.

\begin{proposition}[ARI C.8.1; T; three uniform bounds on the elliptic norm]\label{ari:C.8.1}
For every $z\in\C$ and every $\tau$ in the standard modular
fundamental region, we have
\[
 \begin{array}{c|ccc}
 \text{restriction}&T\le2&T\le7/5&T\le11/10\\ \hline
 \mathcal N(z,\tau)&<17/10&<3/2&<143/100.
 \end{array}
\]
If $\wp_\tau(z)=0$ and $T\ge2$, then $\mathcal N(z,\tau)<541/500$.
\end{proposition}
\noindent\textit{Uses:} the theta product, Poisson summation and Fourier majorants.\par
\begin{proof}
\noindent\textbf{Step 1.} Taking absolute values in the theta series and completing the square
gives
\[
 e^{-\pi(\Im z)^2/T}|\vartheta_1(z,\tau)|
 \le\sum_{n\in\Z}
 e^{-\pi T(n+1/2+\Im z/T)^2}
 \le\sum_{n\in\Z}e^{-\pi Tn^2}.
\]
The last inequality follows from Poisson summation: the Fourier
coefficients of the periodized Gaussian are positive, so its maximum
is at an integer. The eta product and $n^2\ge1+3(n-1)$ for $n\ge1$
therefore imply, on $a\le T\le b$,
\begin{equation}\label{ari:eq:ca-sup}
 \mathcal N(z,\tau)\le
 \exp\!\left(\frac{r_a}{(1-r_a)^2}\right)
 \max\{F(a),F(b)\},\qquad
 F(T)=e^{\pi T/12}
 \left(1+\frac{2e^{-\pi T}}{1-e^{-3\pi T}}\right),
\end{equation}
where $r_a=e^{-2\pi a}$. Indeed,
$-\sum\log(1-r^n)\le r/(1-r)^2$. The function $F$ is convex,
as follows by expanding its denominator as a geometric series;
each summand is an exponential of a real linear function of $T$.
We use $6/7<\sqrt3/2$ and the following rational bounds:
\[
\begin{array}{c|@{\,}c@{\,}c@{\,}c@{\,}c@{\,}c}
T&6/7&1&11/10&7/5&2\\ \hline
F(T)&<14211/10000&<14116/10000&<14180/10000&<14782/10000&<16944/10000.
\end{array}
\]
The eta factors in \eqref{ari:eq:ca-sup} are less than $10047/10000$
for $a=6/7$ and less than $10019/10000$ for $a=1$.
Multiplying these fractions proves the asserted bounds, using
$[6/7,1]$ and $[1,2]$ separately for the first one.


\noindent\textbf{Step 2.} We specify elementary arithmetic for all the scalar estimates in
this proof. One may use $314159/100000<\pi<355/113$ and, for
$0\le x\le20$, the Taylor bounds
\begin{equation}\label{ari:eq:ca-exp}
 S_{80}(x)\le e^x\le
 S_{80}(x)+\frac{x^{81}}{81!}\frac1{1-x/82},
 \qquad S_{80}(x)=\sum_{j=0}^{80}\frac{x^j}{j!}.
\end{equation}
Taking reciprocals bounds negative exponentials. The remainder
bound follows by comparison of consecutive terms, not by numerical
sampling. These inequalities verify each displayed rational bound;
much shorter Taylor sums also suffice. Thus the decimal diagnostics
in the accompanying programs are dispensable.


\noindent\textbf{Step 3.} For the cusp we check directly the analytic step of supplied FAM C.2.3.
Choose $|\Im z|\le T/2$ by translation, put $b_0=43/100$ and
$r_0=e^{-4\pi}$. The Fourier series
\[
 \frac{\wp_\tau(z)}{\pi^2}
 =\csc^2(\pi z)-\frac13
   +8\sum_{n\ge1}\frac{nq_\tau^n}{1-q_\tau^n}
          \bigl(1-\cos(2\pi nz)\bigr)
\]
shows, if $|\Im z|\ge b_0$, that
\begin{align*}
 \left|\frac{\wp_\tau(z)}{\pi^2}+\frac13\right|
 &\le \frac1{\sinh^2(\pi b_0)}
 +8\left\{\frac{r_0}{(1-r_0)^3}
 +\frac1{2(1-r_0)}
 \left(\frac{e^{-2\pi}}{(1-e^{-2\pi})^2}
       +\frac{e^{-6\pi}}{(1-e^{-6\pi})^2}\right)\right\}
 <\frac13.
\end{align*}
The two terms in the last parentheses bound respectively
$r^n e^{2\pi n|\Im z|}$ and $r^n e^{-2\pi n|\Im z|}$.
This inequality excludes a zero. Hence a zero has
$|\Im z|<b_0$. The Jacobi product, $\log(1+x)\le x$, and
the monotonicity of $\cosh(\pi y)e^{-\pi y^2/T}$ on
$0\le y\le b_0$ give
\[
 \mathcal N(z,\tau)\le
 2\cosh(\pi b_0)e^{-\pi/3-\pi b_0^2/2}
 \exp\!\left(
 \frac{2\cosh(2\pi b_0)r_0}{1-r_0}
 +\frac{r_0^2}{1-r_0^2}\right)<\frac{541}{500}.
\]
For the stated monotonicity, the concave function
$\tanh(\pi y)-y$ is zero at zero and positive at $b_0$.
For $T\ge2$ the right side before substituting $T=2$ decreases:
the derivative of $-\pi T/6-\pi b_0^2/T$ is negative and
both product errors decrease. The two strict scalar inequalities
follow from \eqref{ari:eq:ca-exp}. This proves the cusp assertion
without using any certificate from FAM C.2.3.
\end{proof}

\begin{proposition}[ARI C.8.2; T; an explicit local maximum on a larger disc]\label{ari:C.8.2}
If $|\tau-i|\le1/50$ and $\wp_\tau(z)=0$, then
$\mathcal N(z,\tau)\le\sqrt2$, with equality exactly when $\tau=i$.
\end{proposition}
\noindent\textit{Uses:} ARI C.8.1 and C.3.1; the theta heat equation and Cauchy estimates.\par
\begin{proof}
\noindent\textbf{Step 1.} Write $\delta=\tau-i$, $h=z-(1-\tau)/2$ and
\[
 G(h,\tau)=\wp_\tau((1-\tau)/2+h),\qquad
 H(h,\tau)=\log\frac{\vartheta_3(h,\tau)}{\eta_D(\tau)}.
\]
We choose the holomorphic logarithm whose value at $(0,i)$ is
$\log\sqrt2$. The half-period theta identity gives
\begin{equation}\label{ari:eq:ca-H}
 \log \mathcal N(z,\tau)=\Re H(h,\tau)-\frac{\pi(\Im h)^2}{T}.
\end{equation}
Both $G$ and $H$ are even in $h$. On the polydisc
$|h|\le1/4$, $|\delta|\le1/10$, they satisfy
\begin{equation}\label{ari:eq:ca-GHbounds}
 |G|<30,\qquad |H|<2/3.
\end{equation}
Here are sufficient series estimates, including their margins.
Put
\[
 \mathcal S_k=\sum_{n\ge1}n^k e^{-9\pi n^2/10}\cosh(\pi n/2),
 \quad a_0=2\mathcal S_0,\quad a_1=4\pi \mathcal S_1,\quad a_2=8\pi^2\mathcal S_2.
\]
Then $a_0<3/10$, $a_1<19/10$, $a_2<12$.
The first two summands suffice to verify these inequalities: for
$n\ge3$, each of the two exponentials in the hyperbolic cosine
has successive-term ratio at most $e^{-29\pi/5}$, and the
factor $n^k$ increases this ratio by at most $(4/3)^2$.
We use \eqref{ari:eq:ca-exp} for the first terms and this geometric
bound for the tail. With $r=e^{-9\pi/5}$, the identities
\[
 G=-\partial_h^2\log\vartheta_3-\frac{\pi^2}{3}E_2,
 \qquad
 |E_2|\le1+\frac{24r}{(1-r)^3}
\]
give
\[
 |G|\le\frac{a_2}{1-a_0}+\frac{a_1^2}{(1-a_0)^2}
       +\frac{\pi^2}{3}\left(1+\frac{24r}{(1-r)^3}\right)<30.
\]
Also
\[
 |H|\le-\log(1-a_0)+\frac{11\pi}{120}
                    +\frac{r}{(1-r)^2}<\frac23.
\]
The logarithm is well defined since $|\vartheta_3-1|<3/10$;
the displayed eta product fixes its branch. A geometric bound
for $-\log(1-x)=\sum_{n\ge1}x^n/n$ verifies the last estimate.


\noindent\textbf{Step 2.} Let
\[
 e(\tau)=\frac{\pi^2}{3}(\vartheta_2(0,\tau)^4-
                              \vartheta_4(0,\tau)^4),\qquad
 \mathcal E_2(\tau)=\frac{\pi^2}{3}E_2(\tau),\qquad
 g=\frac{g_2(i)}4.
\]
On $|\delta|\le1/50$ the following loose bounds hold:
\begin{equation}\label{ari:eq:ca-derivatives}
 |e'|<14,\quad |e''|<60,\quad |g_2'|<500,
 \quad |\mathcal E_2'|<2,\quad |E_6''|<100,
 \qquad 47<g<95/2.
\end{equation}
For completeness these are scalar Fourier-series bounds, not
interval computations. At $T\ge49/50$, upper bounds for
$|\vartheta_j|,|\vartheta_j'|,|\vartheta_j''|$ are
\[
\begin{array}{c|ccc}
j&0\text{ derivatives}&1\text{ derivative}&2\text{ derivatives}\\ \hline
2&929/1000&742/1000&670/1000\\
4&1093/1000&290/1000&911/1000.
\end{array}
\]
They follow by differentiating the absolutely convergent theta
series, taking absolute values, retaining two terms, and using
the same successive-term bound. Differentiating the fourth
powers proves the first two bounds in \eqref{ari:eq:ca-derivatives}.
For $r\le1/470$, the remaining derivative majorants are
\begin{align*}
 |g_2'|&\le640\pi^5\frac{r(1+11r+11r^2+r^3)}{(1-r)^7}<500,\\
 |\mathcal E_2'|&\le16\pi^3\frac{r(1+r)}{(1-r)^5}<2,\\
 |E_6''|&\le2016\pi^2
 \frac{r(1+r)(1+120r+1191r^2+2416r^3+1191r^4+120r^5+r^6)}
 {(1-r)^{11}}<100.
\end{align*}
These follow by summing $\sum n^k r^n$ and bounding
$(1-r^n)^{-1}$ by $(1-r)^{-1}$. Substitution of $r=1/470$
and $\pi<355/113$ verifies the rational inequalities.
Finally $g=\pi^4E_4(i)/3$ and
$E_4=1+240\sum n^3q_\tau^n/(1-q_\tau^n)$ give $47<g<95/2$.
The identity of ARI C.3.1, checked by differentiating $e$ and using
the theta heat equation, reads
\[
 e(i)=0,\qquad e'(i)=id,\qquad d=\frac{2g}{3\pi}>99/10.
\]


\noindent\textbf{Step 3.} The differential equation $\wp''=6\wp^2-g_2/2$ determines the
first coefficients in the even expansion:
\[
 G(h,\tau)=e+a h^2+ea h^4+\mathcal R_G,
 \qquad a=3e^2-g_2/4.
\]
Consequently, on $|\delta|\le1/50$,
\[
 |e|\le7/25,\quad |a+g|\le\frac52+3(7/25)^2=\frac{3419}{1250},
 \quad |a|\le\frac{62794}{1250}.
\]
Cauchy's estimate from \eqref{ari:eq:ca-GHbounds} yields
\[
 |\mathcal R_G|\le
 \frac{30(4|h|)^6}{1-16|h|^2}.
\]
On $|h|=1/10$ the term $ah^2$ has modulus at least
$55331/125000>0.4426$, whereas the sum of the other terms is at most
\[
 \frac7{25}+\frac7{25}\frac{62794}{1250}\frac1{10000}
       +\frac{30(4/25)^3}{1-4/25}<0.4277.
\]
Rouch\'e's theorem puts two zeros, with multiplicity, in
$|h|<1/10$. The disc injects into the torus and $\wp$ has exactly
two zeros on the torus, so it contains both of them.


\noindent\textbf{Step 4.} Write $G=e+h^2Q$. The preceding estimates give, for $|h|\le1/10$,
\[
 |Q+g|\le
 \frac{3419}{1250}+\frac7{25}\frac{62794}{1250}\frac1{100}
       +\frac{30(4/25)^3}{(1-4/25)/100}<\frac{88}{5}.
\]
Taylor's formula and $|e''|<60$ give
$|e|\ge(99/10-3/5)|\delta|=(93/10)|\delta|$.
At either zero, therefore,
\begin{equation}\label{ari:eq:ca-delta-h}
 |\delta|<7|h|^2,
\end{equation}
because $(95/2+88/5)/(93/10)=7$.


\noindent\textbf{Step 5.} Expand
$H=H_0(\tau)+H_2(\tau)h^2+H_4(\tau)h^4+\mathcal R_H$.
The heat equation and the identity for $G$ give
\[
 H_0'=\frac{ie}{4\pi},\quad
 H_2=-\frac{e+\mathcal E_2}{2},\quad H_4=-\frac a{12}.
\]
Thus $H_0(i)=\log\sqrt2$, $H_0'(i)=0$,
$|H_0''|<6/5$, $H_2(i)=-\pi/2$, $|H_2'|<8$, and
$|H_4|<21/5$. Cauchy's estimate gives
\[
 |\mathcal R_H|\le
 \frac{(2/3)(4|h|)^6}{1-16|h|^2}
 \le\frac{2048}{63}|h|^4
 \quad (|h|\le1/10).
\]
Combining \eqref{ari:eq:ca-H} and \eqref{ari:eq:ca-delta-h}, we obtain
\begin{align*}
 \log \mathcal N-\log\sqrt2
 &\le-\frac\pi2\Re(h^2)-\pi(\Im h)^2+C|h|^4\\
 &=-\frac\pi2|h|^2+C|h|^4,\\
 C&=\frac{147}{5}+56+\frac{21}{5}
                 +\frac{2048}{63}+\frac{1100}{49}<150.
\end{align*}
The final term of $C$ bounds the change of $1/T$ from $1$:
$7\pi/(1-1/50)<1100/49$. Since $150|h|^2<3/2<\pi/2$,
the result is strict unless $h=0$. Equation
\eqref{ari:eq:ca-delta-h} then gives $\delta=0$, where equality
holds by the checked isogeny computation at $i$.
\end{proof}

\begin{proposition}[ARI C.8.3; T; the comparison of the two local discs]\label{ari:C.8.3}
Every complex number $t_0$ with $|t_0-1|\le1/500$ has a
representative $\tau$ with $|\tau-i|<1/50$ and $j(\tau)=1728t_0$.
Consequently every zero of $\wp$ for that modulus satisfies
$\mathcal N\le\sqrt2$, with equality only for $t_0=1$.
\end{proposition}
\noindent\textit{Uses:} ARI C.8.2; the Ramanujan differential identities and Rouche theorem.\par
\begin{proof}
\noindent\textbf{Step 1.} At $i$, $E_6=0$ and the Ramanujan differential identity gives
$|E_6'(i)|=\pi E_4(i)^2>13/2$, since $E_4(i)>29/20$.
On $|\delta|=1/50$, \eqref{ari:eq:ca-derivatives} gives
\[
 |E_6(\tau)|>\frac{13}{100}-\frac1{50}=\frac{11}{100}.
\]
On the whole disc, $|q_\tau|<1/470$ and
$|\Delta_{\mathrm{mod}}|\le |q_\tau|\exp(24|q_\tau|/(1-|q_\tau|))<1/400$.
It follows on the boundary that
\[
 |t-1|=\frac{|E_6|^2}{1728|\Delta_{\mathrm{mod}}|}
 >\frac{121}{43200}>\frac1{500}.
\]
Rouch\'e applied to $E_6'(i)\delta$ shows that $E_6$ has
exactly one zero in the disc. Hence $t-1$ has two zeros there,
with multiplicity, and another application of Rouch\'e shows
that $t-t_0$ has two for each $|t_0-1|\le1/500$.
This argument proves coverage of the $t$-disc, rather than
inferring coverage merely from a boundary bound. Elliptic
curves over $\C$ with the same $j$ are isomorphic. In terms of
lattices the isomorphism is a homothety; it carries the two
zeros of $\wp$ to the two zeros of the scaled $\wp$ and
preserves $\mathcal N$. Statement C.8.2 applies to either zero.
\end{proof}

\begin{proposition}[ARI C.8.4; T; the first three normalized duplication steps]\label{ari:C.8.4}
Put $v_t=t-1$. For $n=2^k$, $k=1,2,3$, and $nP^{\mathrm{div}}\ne O$,
\begin{equation}\label{ari:eq:ca-ladder}
 \left(\frac{\mathcal N(P^{\mathrm{div}})}{\sqrt2}\right)^{n^2}
 =\frac{3^{3/4}}{2n}\,
       \frac{\mathcal N(nP^{\mathrm{div}})|v_t|^{-1/4}}{|\mathcal A_k(v_t)|},
\end{equation}
where
\begin{align*}
 \mathcal A_1(v_t)&=1,\qquad \mathcal A_2(v_t)=1-5v_t/27,\\
 \mathcal A_3(v_t)&=1-\frac{217}{27}v_t-\frac{20486}{729}v_t^2
       +\frac{10910}{19683}v_t^3
       +\frac{161383}{177147}v_t^4
       -\frac{245}{14348907}v_t^5.
\end{align*}
On $|v_t|\le1/28$ we have $|\mathcal A_3(v_t)|\ge1-10|v_t|$.
For every $k\ge1$, not just these first three steps, define
\[
 E_k(t)=\frac{|\Delta_E|^{(4^k-1)/12}|t-1|^{1/4}}
 {2^{4^k/2}|\psi_{2^k}(P^{\mathrm{div}})|}.
\]
Then $(\mathcal N(P^{\mathrm{div}})/\sqrt2)^{4^k}=E_k(t)\mathcal N(2^kP^{\mathrm{div}})|t-1|^{-1/4}$,
and the removable limiting value is
\[
 \lim_{t\to1}E_k(t)=\frac{3^{3/4}}{2^{k+1}}
                 =2^{1-k}\frac{3^{3/4}}4.
\]
\end{proposition}
\noindent\textit{Uses:} ARI C.4.2, C.4.4 and C.3.1.\par
\begin{proof}
\noindent\textbf{Step 1.} Substitute the exact division polynomials from Statement C.4.2
in \eqref{ari:eq:cs-division}. Equivalently, remove
$t^{2(n^2-1)/3}(t-1)^{n^2/4}$ from $\psi_n(P^{\mathrm{div}})$ and divide
the resulting polynomial by its value at $t=1$. For $n=2,4,8$
this gives the displayed polynomials and the constant
$3^{3/4}/(2n)$. All operations are rational polynomial operations;
\path{code/C_ladder_algebra.py} reproduces them and records
the next two steps as well. The general identity is C.4.2.
For its limit write $\mu=(1-\tau)/2+h$ as in C.3.1. Since
$n=2^k$ is even, $n\mu$ is congruent to $nh$ modulo the
period lattice. Hence
\[
 \mathcal N(nP^{\mathrm{div}})=2\pi|\eta_D(i)|^2 n|h|(1+o(1)),\qquad
 t-1=\tfrac94\pi^4E_4(i)h^4(1+o(1)).
\]
The identity $E_4(i)^3=1728\eta_D(i)^{24}$ then gives
$\mathcal N(nP^{\mathrm{div}})|t-1|^{-1/4}\to2n/3^{3/4}$.
Since the left side of the general ladder tends to one,
this proves the stated limit for every $k$.
The triangle inequality gives
$|\mathcal A_3(v_t)|\ge1-|v_t|M_{\rm coeff}$, where
\[
 M_{\rm coeff}\le\frac{81}{10}+\frac{29}{28}
                +\frac1{28^2}+\frac1{28^3}+\frac1{28^4}<10.
\]
Here we bounded respectively the five absolute coefficients by
$81/10,29,1,1,1$. This is a rational inequality on one disc,
not a subdivision or a root search.
\end{proof}

\begin{proposition}[ARI C.8.5; T; FAM C.2.3 by explicit analytic estimates]\label{ari:C.8.5}
For every lattice $\Z+\tau\Z$ and every zero $z$ of its
Weierstrass function,
\[
 \mathcal N(z,\tau)\le\sqrt2.
\]
Equality holds exactly for the modular orbit of $i$.
In the parameter of the family this says
$\max D_{\mathrm{part}}=-(3/\pi)\log6$, only at $t=1$.
\end{proposition}
\noindent\textit{Uses:} ARI C.8.1--C.8.4; the norm values of ARI C.1.2.\par
\begin{proof}
\noindent\textbf{Step 1.} Reduce to the standard modular region. The portion $T\ge2$
is covered by Statement C.8.1. Suppose $T\le2$ and put
$r=|t-1|$. For $r\ge9/10$, the step $n=2$ in
\eqref{ari:eq:ca-ladder} and $\mathcal N(2P^{\mathrm{div}})<17/10$ give strict inequality,
because
\[
 \frac{27(17/10)^4}{256}<\frac9{10}.
\]
There is no $2$-torsion point of the section on the smooth base.


\noindent\textbf{Step 2.} If $r\le9/10$, then $T<7/5$. Indeed, $T\ge7/5$ would give
$|q_\tau|<1/6500$ and
\[
 |E_4|>24/25,\qquad |\Delta_{\mathrm{mod}}|<1/6400,
 \qquad |t|>\frac{6400(24/25)^3}{1728}>19/10,
\]
contrary to $|t|\le19/10$. The estimates follow from
\[
 \sum_{n\ge1}\sigma_3(n)r^n
 \le\frac{r(1+4r+r^2)}{(1-r)^5},\qquad
 |\Delta_{\mathrm{mod}}|\le r\exp\frac{24r}{1-r}.
\]
For $1/28\le r\le9/10$, the $n=4$ step and
$\mathcal N(4P^{\mathrm{div}})<3/2$ now apply. The function
$(1-5r/27)r^{1/4}$ is increasing up to $r=27/25$.
At its left endpoint its fourth power is
\[
 \frac1{28}\left(\frac{751}{756}\right)^4
 >\frac{2187}{65536}
 =\left(\frac{3^{3/4}}8\frac32\right)^4.
\]
This proves strict inequality on the whole annulus. The lower
bound for $\mathcal A_2$ is positive there, so no division by a zero has
occurred.


\noindent\textbf{Step 3.} If $r\le1/28$, then $T<11/10$. For otherwise $|q_\tau|<1/1000$,
and the Lambert series with
\[
 \sum_{n\ge1}\sigma_5(n)r^n
 \le\frac{r(1+26r+66r^2+26r^3+r^4)}{(1-r)^7}
\]
gives $|E_6|>9/20$, $|\Delta_{\mathrm{mod}}|<1/900$, and hence
$|t-1|>900(9/20)^2/1728>1/28$.
On $1/500\le r\le1/28$, the $n=8$ step,
$\mathcal N(8P^{\mathrm{div}})<143/100$, and $|\mathcal A_3|\ge1-10r$ prove the claim.
Indeed $(1-10r)r^{1/4}$ has a single critical point in this
interval, a maximum at $r=1/50$, so its minimum is at an endpoint.
At both endpoints its fourth power exceeds $9/5000$, whereas
\[
 \left(\frac{3^{3/4}}{16}\frac{143}{100}\right)^4
 <\frac7{4000}<\frac9{5000}.
\]
Again the displayed positive bound excludes a division-polynomial
zero. The remaining disc $r\le1/500$ is covered by Statement
C.8.3. The smooth limiting fibre at $t=0$ has
$\mathcal N=3^{1/4}<\sqrt2$; the formula also extends there by
continuity, so the singular rational model causes no omitted
modulus. The only equality occurs at $t=1$.
\end{proof}

\subsection{Announcement on rays and level curves}

\begin{announcement}[ARI C.8.15; A; every critical point and every level set]\label{ari:C.8.15}
For every zero $\mu$ of $\wp_\tau$, the zero sets of the two
critical factors are exactly
\[
 \begin{aligned}
 \mathcal W=0&\quad\Longleftrightarrow\quad
       \tau\in\operatorname{SL}_2(\Z)\{i,\rho\},\\
 3\mathcal W\wp'_\tau(\mu)-2g_2=0&\quad\Longleftrightarrow\quad
       \tau\in\operatorname{SL}_2(\Z)\rho.
 \end{aligned}
\]
On the coarse $t$-plane the function $\mathcal L$ has no critical point
away from $t=1$. The point $t=0$ is a noncritical $C^1$ point,
with $\mathcal L_t(0)=7/384$. The point $t=1$ is the unique global
maximum, of value $\log\sqrt2$, and has the fractional-power
local maximum described in C.3. In particular there is no
saddle in $0<t<1$. On the root charts, $h=0$ over $t=1$
is a nondegenerate maximum, while $\delta=0$ over $t=0$
is a degenerate critical point with nonzero harmonic cubic
leading term and three alternating ascending and descending
sectors. These are all critical points on the root charts.

For every angle $\theta$, $r\mapsto \mathcal L(1+re^{i\theta})$ is
strictly decreasing for $r>0$, from $\log\sqrt2$ to $-\infty$.
For each $c<\log\sqrt2$, the level set $\{\mathcal L=c\}$ is one
Jordan curve, strictly star-shaped about $1$. Its bounded
complement is $\{\mathcal L>c\}$; its unbounded complement is
$\{\mathcal L<c\}$. The curve is real analytic except possibly at
$t=0$ when $c=\frac14\log3$, where it is $C^1$ with nonzero
tangent. The closed superlevel sets form nested topological
discs. At the maximum the level set is the single point $1$,
and above it the level set is empty. The corresponding
statements for $D_{\rm part}$ follow by its increasing affine
change from $\mathcal L$.
\end{announcement}
\noindent\textit{Uses:} ARI C.8.11--C.8.14 and the two local root expansions.\par
The mechanism bounds the two critical factors on a cusp region and on the remaining compact modular region. The root expansions identify their behavior over $t=0,1$; the radial derivative then gives the one-radius parametrization of each lower level.
\noindent\textit{Proof in the repository:} \href{https://github.com/KokunoYumeto/s6-complex-structure-workbench}{S6 complex-structure workbench}, S6\_05\_ARITHMETIC, Statement C.8.15.


% ===== cohomology source-preserving fragment =====

\section{Cotangent cohomology and its consequences}\label{sec:cohomology}
\subsection{Direct images and the decisive period derivative}\label{sec:cotangent-images}
The calculations concern every admissible member of the fixed marked
$(3,4)$ period family.  The structure-sheaf direct images and the previously
known Hodge constraints are credited to \cite[Proposition 9.13 and Theorem 9.1]{AC26};
the root-curve version also occurs in \cite[Proposition 6.10]{DGKXv1}.
The source proof of the structure-sheaf calculation is retained below as a
new proof, in accordance with the prescribed T status of this subsection.


\begin{definition}[inv:notation-cotangent; text; Cotangent notation]\label{inv:notation-cotangent}
We use the marked threefold of the question, with its arbitrary admissible
constant $c$.  Coefficients of coherent cohomology are complex.  In this
section $\omega_f=K_X\otimes f^*K_B^{-1}$,
$\mathcal L=\Omega_B^1(\log\Sigma)=\mathcal O_B(1)$,
$\mathcal A=\Omega_X^1(\log D)$, and
$\mathcal E_{\rm cot}=\mathcal A/f^*\mathcal L$.  We write $\gamma,u_{\Lambda}$ also for
their images in $R^1f_*\mathcal O_X$; this convention concerns cohomology
classes and does not identify them with differential-form representatives.
\end{definition}
\uses{The family and marking fixed in Section~\ref{sec:family}.}

\begin{theorem}[inv:A1.0; T; The structure-sheaf input in the precise form used below]\label{inv:A1.0}
The direct images $R^qf_*\mathcal O_X$, for $q=0,1,2$, are
$\mathcal O$, $G_B=\mathcal O\gamma\oplus\mathcal O(-1)$, and $\mathcal O(-1)$.
The inclusion $\mathcal O\gamma\subset G_B$ is a subbundle.  At the cusp,
$(\gamma,u_{\Lambda})$ is a holomorphic frame, and cup product identifies
$\bigwedge^2G_B$ with $R^2f_*\mathcal O_X$.  In the cyclic covering coordinate
$s$, the action on $G_B$ has characters $(1,\zeta^{-a})$, where
$(m,a)=(3,2)$ or $(4,1)$ and $\zeta=e^{-2\pi i/m}$.
These assertions concern full scheme fibres, including $3S_1$ and $4S_2$.
\end{theorem}
\uses{The cyclic and cusp charts; analytic base change \cite[III.4.10, III.4.12]{BS76}; analytic duality \cite[pp.~261, 264--266]{RRV71} and \cite[p.~91, Proposition~2]{RR70}.}

\begin{proof}
\par\smallskip\noindent\textbf{Step 1.}\ We derive the assertion from the local charts and the discriminant;
the same direct images are stated in Proposition 9.13 of
\cite{AC26}, but that proposition is not an input to this proof.
Flatness holds
because a nonzero local parameter on the smooth base is a nonzero divisor
in the local ring of the smooth total space.  The normalization complex
of the cusp has rational surface and curve strata and the cellular
cochain complex of a torus as its constant part.  It gives
$h^q(\mathcal O_W)=(1,2,1)$; the two degree-one classes $\gamma,u_{\Lambda}$ have nonzero
cup product.  To compute the determinant line without assuming a
canonical-bundle formula, we first retain the two finite stabilizers.
Let $\mathfrak B$ be the root curve of orders three and four over $B$,
with root lines $R_1,R_2$, and put $R=R_1R_2^{-1}$.  The presentation
$3R_1=4R_2=H$ gives $\operatorname{Pic}(\mathfrak B)=\mathbb Z R$,
$R_1=R^4$, $R_2=R^3$ and $\mathcal O_{\mathfrak B}(p_0)=R^{12}$;
one may also verify this by removing a line bundle's two stabilizer
characters and descending it to $\mathbb P^1$.

\par\smallskip\noindent\textbf{Step 2.}\ Write $\bar f:X\to\mathfrak B$.  Translation-invariant vertical
fields on the torus charts, extended as logarithmic fields on the cusp
charts, have transition matrices depending only on the base.  They
give a rank-two bundle $V$ on $\mathfrak B$ whose pullback is the
relative logarithmic tangent bundle.  Its fixed field $\xi$ is a
nowhere-zero subline, including in the logarithmic frame at the cusp.
The quotient line $Q=V/\mathcal O\xi$ has automorphy factors $-1/\tau$
and $1/\tau$ for the two finite generators, read directly from the
linear part of the given affine deck transformations.  Thus
$\Delta_{\rm mod}(\tau)^{-1}$ is a meromorphic section of $Q^{12}$.
The discriminant product has no zero in the upper half-plane and
has a simple zero in the cusp coordinate, since
$\tau=(\log t_c)/(2\pi i)+h(t_c)$ with $h$ holomorphic.  Its divisor
therefore proves $Q^{12}=\mathcal O_{\mathfrak B}(-p_0)=R^{-12}$.
The displayed torsion-free Picard group gives
$\det V=Q=R^{-1}$.

\par\smallskip\noindent\textbf{Step 3.}\ The exact logarithmic tangent sequence has quotient
$\bar f^*T_{\mathfrak B}(-p_0)$.  Its determinant, together with
$W=\bar f^*p_0$, gives
$\omega_{X/\mathfrak B}=\bar f^*(\det V)^{-1}$.
The geometric fibres on the root charts have cohomology dimensions
$(1,2,1)$: they are smooth tori at the two finite points, and the
normalization calculation above gives the same result at the cusp.
Analytic base change makes the three root direct images locally free.
Connectedness gives $\bar f_*\mathcal O_X=\mathcal O_{\mathfrak B}$, and fibrewise
Gorenstein duality with its trace pairing now yields
$R^2\bar f_*\mathcal O_X=\det V=R^{-1}$.  The invariant local exponents
of $R^{-1}$ are two at the order-three point and one at the
order-four point.  Removing these characters leaves
$R^{-1}R_1^{-2}R_2^{-1}=R^{-12}$, so its coarse direct image is
$\mathcal O_B(-1)$.  This supplies the degree used below independently of
the asserted direct-image or canonical-bundle formulas in \cite{AC26}.

\par\smallskip\noindent\textbf{Step 4.}\ The same base-change calculation identifies $\gamma,u_{\Lambda}$ with a
frame through the cusp.  At a multiple fibre finite pushforward and the
Reynolds projector identify cohomology with cyclic invariants.  If $b_0$
and $b_1$ have characters $1,\zeta^{-a}$, their invariant frames are
$b_0,s^ab_1$ and $s^a(b_0\wedge b_1)$.  Their products have exactly
these exponents, so the determinant comparison is an isomorphism,
including after reduction modulo $s^m$.  Thus $\deg G_B=-1$.
The invariant flat class $\gamma$ does not enter the holomorphic
two-plane spanned by
$\sigma_1=6\mu\gamma+\tau u_{\Lambda}+w_{\Lambda}$ and
$\sigma_2=\beta\gamma+\mu u_{\Lambda}+\delta$: comparison of the $w_{\Lambda},\delta$
coefficients proves this even on the smooth cyclic covers at their
centres.  It therefore gives the stated subbundle everywhere.
The quotient is $\mathcal O(-1)$ and the extension splits because
$H^1(\mathcal O(1))=0$.  This also specifies which subline in the otherwise
noncanonical splitting is used below.  Coherent Leray now gives
$(h^q(\mathcal O_X))_{q=0}^3=(1,1,0,0)$, since
$H^0(\mathcal O(-1))=H^1(\mathcal O(-1))=H^1(\mathcal O)=0$.
\end{proof}

\begin{lemma}[inv:A1.1; T; Local freeness and the determinant line]\label{inv:A1.1}
Both $\mathcal A$ and $\mathcal E_{\rm cot}$ are locally free, of ranks three and
two.  The vertical vector $\xi=\partial/\partial\zeta_2$ is nowhere
zero as a section of $\mathcal E_{\rm cot}^\vee$.  There is an exact sequence
of holomorphic vector bundles
\begin{equation}\label{inv:eq:A-E-line}
0\longrightarrow N_{\rm det}\longrightarrow\mathcal E_{\rm cot}
 \xrightarrow{\langle\cdot,\xi\rangle}\mathcal O_X\longrightarrow0,
\qquad N_{\rm det}=\det\mathcal E_{\rm cot}=\omega_f(-2S_1-3S_2).
\end{equation}
\end{lemma}
\uses{The local cyclic and normal-crossing charts; inv:A1.0.}

\begin{proof}
\par\smallskip\noindent\textbf{Step 1.}\ At a smooth point of a reduced cusp branch the assertion is the usual
coordinate calculation.  At a double point write $t_c=xy$ with third
coordinate $z$; then
$\mathcal A=R\,d\log x\oplus R\,d\log y\oplus R\,dz$ and its
base subline is $R(d\log x+d\log y)$.  At either triple point write
$t_c=xyz$; then $\mathcal A$ has the three logarithmic coordinate
basis vectors and its base subline is their sum.  In each case that
subline is a direct summand.  These calculations apply on every one
of the three double curves, including their endpoints; no simplicity
of the global divisor is assumed.

\par\smallskip\noindent\textbf{Step 2.}\ At $S_j$ we use the actual free cyclic quotient
$q:J'_j\to N_j$, which is finite and etale on total spaces, and
$t_j=s^{m_j}$ on the base.  On $J'_j$ the logarithmic conormal is
$m_j\,ds/s$, and its coefficient is a unit over $\mathbb C$.  It is again a
direct summand of the logarithmic cotangent bundle.  The comparison
$q^*\mathcal A=\Omega_{J'_j}^1(\log(s=0))$ identifies the quotient
bundles and their exact sequence; faithful flat descent proves the
claim downstairs.  This is not a replacement of $m_jS_j$ by $S_j$.

\par\smallskip\noindent\textbf{Step 3.}\ On the marked torus family $\xi$ is the constant second vertical
vector.  The linear monodromies fix it, and translations preserve it.
On a cusp toric chart it is a nonzero constant vector in the relative
logarithmic tangent lattice of rank two.  A toric coordinate change
$x_i\mapsto c_i(t_c)t_c^{b_i}x_i$ changes $d\log x_i$ by a multiple
of $d\log t_c$, so this constant relative vector descends and is
nowhere zero in that bundle.  It can vanish as an ordinary tangent
vector at a toric stratum; this does not imply vanishing in the
logarithmic bundle.  Hence evaluation is surjective and its kernel
is the determinant line.  Taking determinants of
$0\to f^*\mathcal L\to\mathcal A\to\mathcal E_{\rm cot}\to0$ gives
\[
\det\mathcal E_{\rm cot}=K_X(D)\otimes f^*\mathcal L^{-1}
=\omega_f\bigl(D-f^*\Sigma\bigr)
=\omega_f(-2S_1-3S_2),
\]
which proves \eqref{inv:eq:A-E-line}.
\end{proof}

\begin{proposition}[inv:A1.2; T; The logarithmic direct images by vertical twisting and duality]\label{inv:A1.2}
One has, as holomorphic vector bundles on $B$,
\[
\begin{array}{c|ccc}
q&0&1&2\\\hline
R^qf_*\mathcal E_{\rm cot}&\mathcal O\oplus\mathcal O(-1)&
 \mathcal O^2\oplus\mathcal O(-1)^2&\mathcal O\oplus\mathcal O(-1).
\end{array}
\]
They commute with base change to every full scheme fibre.
\end{proposition}
\uses{inv:A1.0, inv:A1.1; analytic duality \cite[pp.~261, 264--266]{RRV71}.}

\begin{proof}
\par\smallskip\noindent\textbf{Step 1.}\ Put $H=2S_1+3S_2$.  We first compute the line bundle $\mathcal O_X(H)$,
rather than a tensor product of two separately descended lattices.
On the cover at $p_j$, its sections allow poles of order at most $m_j-1$
in $s$.  A fixed character has invariant exponent zero and gains no
pole, whereas the variable character $\zeta^{-a}$ has least allowed
invariant exponent $a-m$.  Thus the inclusion
$R^1f_*\mathcal O\to R^1f_*\mathcal O(H)$ is the identity on $\mathcal O\gamma$ and
changes $s^ab_1$ to $s^{a-m}b_1=t_j^{-1}s^ab_1$ on its quotient.
It makes exactly one upper modification at each finite point.
The same computation applies to $R^2$.  No modification is made at
the cusp or a regular fibre.  The result is
\begin{equation}\label{inv:eq:A-H-images}
f_*\mathcal O(H)=\mathcal O,\qquad
R^1f_*\mathcal O(H)=\mathcal O\oplus\mathcal O(1),\qquad
R^2f_*\mathcal O(H)=\mathcal O(1).
\end{equation}
Here the middle extension splits because $H^1(\mathcal O(-1))=0$.
The local invariant modules are free.  Tensoring them with
$\mathbb C\{t_j\}/(t_j)$ gives invariants of
$s^{-(m_j-1)}\mathbb C\{s\}/(s^m)$, with the indicated surviving character
basis vectors.  The Reynolds idempotent is linear over the base,
so this calculation also proves the comparison with the nonreduced
fibre.  At the cusp the line bundle is trivial.  Cohomology and base
change, applied successively from the top degree, proves the claimed
local freeness and comparison for \eqref{inv:eq:A-H-images}.

\par\smallskip\noindent\textbf{Step 2.}\ For clarity the relative duality map we use is the cup product followed
by the analytic trace, giving a morphism of bundles
\[
R^qf_*(\omega_f\otimes\mathcal O(-H))\longrightarrow
\bigl(R^{2-q}f_*\mathcal O(H)\bigr)^\vee.
\]
Every fibre is a compact Gorenstein surface: it is the Cartier fibre
in a smooth threefold.  Serre duality on that scheme makes the map
an isomorphism on its residue field, also for $m_jS_j$.  The constant
dimensions just computed imply local freeness on both sides, so
Nakayama's lemma makes the bundle map an isomorphism.  Therefore
\begin{equation}\label{inv:eq:A-N-images}
(f_*N_{\rm det},R^1f_*N_{\rm det},R^2f_*N_{\rm det})
=(\mathcal O(-1),\mathcal O\oplus\mathcal O(-1),\mathcal O).
\end{equation}
This is the explicit reduction of analytic relative duality needed
here; the scheme fibre is retained throughout.

\par\smallskip\noindent\textbf{Step 3.}\ The connecting maps of \eqref{inv:eq:A-E-line} are zero on a regular
fibre: its relative cotangent is a trivial bundle and evaluation on
a constant nonzero vector splits.  Source and target direct images
in \eqref{inv:eq:A-N-images} and Statement inv:A1.0 are vector bundles,
so these holomorphic connecting maps, zero on a dense open set, are
zero everywhere.  We obtain extensions in each degree of the listed
$R^qN$ by $R^q\mathcal O$.  Every resulting pair of line-bundle degrees is
zero or minus one; their extension groups are
$H^1(\mathcal O(-1)),H^1(\mathcal O),H^1(\mathcal O(1))$, all zero.  These extensions
give exactly the three asserted bundles.  Applying the already proved
fibre comparisons to the short exact sequences proves base change
for $R^q\mathcal E_{\rm cot}$ as well.  This argument uses the vertical twist
$\mathcal O(H)$ and duality; it does not use the workbench's filtration of
$f_*\mathcal E_{\rm cot}\otimes G_B$.
\end{proof}

\begin{definition}[inv:notation-period-branch; text; The order-four branch]\label{inv:notation-period-branch}
The scalar depends on the fourth root chosen in $t-1=s^4$.  We make
that ambiguity explicit.  We choose the local branch for which, on
the positive $t>1$ arc leading to the distinguished cusp,
$\tau=iy$ with $y>1$, and
\begin{equation}\label{inv:eq:A-s-branch}
\tau=i+\frac{i}{\pi\sqrt{E_4(i)}}s^2+O(s^4),\qquad
\operatorname{Im}\mu_s(0)>0.
\end{equation}
The nonvanishing proved below makes the second condition meaningful;
it fixes the remaining sign of $s$.  Equivalently it is the lower
cusp-regular zero branch of the Weierstrass function in fam:B.1.3.
Replacing $s$ by $i^js$ multiplies the reported derivative by $i^{-j}$.
The workbench instead normalizes $ds/dz(z_2)=1$; its number is obtained
by the chain rule and is not numerically the same unspecified coordinate.
\end{definition}
\uses{The period functions and the first-run clockwise convention.}

\begin{theorem}[inv:A2.1; T; The complete raw period determinant]\label{inv:A2.1}
On the marked upper half-plane,
\begin{equation}\label{inv:eq:A-modular-K}
6\mu_z^2-\tau_z\beta_z
=\tau_z^2\frac{E_4(\tau)^2}{E_6(\tau)}.
\end{equation}
The right side extends holomorphically through the apparent poles.
Its exact orders are two at $z_1$ and zero at $z_2$; it has no other
zero away from their translates.  The logarithmic connecting map
$\kappa_0:f_*\mathcal E_{\rm cot}\to\mathcal L\otimes G_B$ has divisor
$\operatorname{divisor}(\det\kappa_0)=p_1+p_2$.
\end{theorem}
\uses{inv:A1.0; the period transformations \cite[Definition 3.1 and Theorem 3.4(i)]{AC26}; the modular identities stated in the proof.}

\begin{proof}
\par\smallskip\noindent\textbf{Step 1.}\ Consider the holomorphic quadratic differential
$\mathscr Q=6(d\mu)^2-d\tau\,d\beta$.  Direct differentiation of
the two affine transformation laws in AC26 Definition 3.1 gives
\[
g_j^*\mathscr Q=\tau^{-2}\mathscr Q\quad(j=1,2).
\]
In particular $g_2^2$ fixes $\tau$ and sends
\[
\mu\longmapsto1-\tau-\mu,\qquad
\beta\longmapsto\beta+6\tau+12\mu-6,
\]
which leaves $\mathscr Q$ exactly unchanged.  The producing run reports zero symbolic residuals for all three identities in \texttt{a\_calculations.py}.
The kernel of
$\Delta=\mathbb Z/3*\mathbb Z/4\to\mathbb Z/3*\mathbb Z/2=\operatorname{PSL}_2(\mathbb Z)$
is the normal subgroup generated by $g_2^2$.  The squared multiplier
of $\mathscr Q$ is the ordinary modular multiplier
$(c\tau+d)^{-2}$, so it is one on every kernel element.  Two generic
points of the marked cover with the same $\tau$ have the same $t$
and differ by an element of $\Delta$; its modular image fixes a
generic point, hence is the identity.  It follows that
$k=\mathscr Q/(d\tau)^2$ is a single-valued meromorphic function of
$\tau$, transforming as a modular form of weight two.  This kernel
argument is needed: mere equivariance of the separate period functions
would not prove descent.

\par\smallskip\noindent\textbf{Step 2.}\ At $z_1$, write a linearizing parameter $r$ with
$g_1r=\zeta r$, $\zeta=e^{-2\pi i/3}$.  If $\mathscr Q=K_rdr^2$
has first nonzero coefficient of order $n$, its transformation law
at the fixed point reads $\zeta^{n+2}=\rho^{-2}=\zeta$.
Thus $n\equiv2\pmod3$, and $K_r$ vanishes to order at least two.
The map $\tau$ is unramified there by AC26 Theorem 3.4(i).  At
$z_2$, $\tau-i$ has order two and $K_s$ is holomorphic; therefore
$k$ has at most a simple pole at $i$.  Away from these points
$\tau$ is unramified.  At the cusp $\mu$ and $\beta+\tau$ are
holomorphic in $q$, so $k=1+O(q)$.

\par\smallskip\noindent\textbf{Step 3.}\ The modular function $k E_6/E_4^2$ is therefore holomorphic at
$\rho$, at $i$, and at the cusp, as well as everywhere else.
The usual zero statements for $E_4,E_6$ (simple at $\rho,i$,
respectively) are the modular facts used in AC26 Section 3.1.
This function descends to the compact modular curve $X(1)=\mathbb P^1$,
so it is constant.  Its cusp value is one.  This proves
\eqref{inv:eq:A-modular-K} and the exact two local orders, without the
workbench's degree count for the coarse determinant line.

\par\smallskip\noindent\textbf{Step 4.}\ For completeness we translate to the actual coarse lattices.  Choose
cover eigenframes $e_0,e_1$ of $\mathcal E_{\rm cot}$ with characters
$1,\zeta^a$, and $b_0,b_1$ of $G_B$ with characters
$1,\zeta^{-a}$.  Their invariant bases are
$e_0,s^{m-a}e_1$ and $b_0,s^ab_1$.  The logarithmic base frame is
$\ell=dt_j/t_j=m\,ds/s$.  Changing determinants from these frames
to the marked ones gives
\[
\det\kappa_0=\epsilon(t_j)m^{-2}s^{m+2-2a}K_s,
\qquad \epsilon(0)\ne0.
\]
This equality follows by taking the two factors $s/m$ from replacing
$ds$ by $\ell$, the factor $s^{m-a}$ in the source, and the factor
$s^{-a}$ in the target.  Thus at $(m,a)=(3,2)$ the total order in
$s$ is $1+2=3$, and at $(4,1)$ it is $4+0=4$.
Both coarse zeros are simple.  At the cusp, the frame change has
no elementary divisor and $dZ/d\tau$ tends to
$\left(\begin{smallmatrix}0&1\\-1&0\end{smallmatrix}\right)$,
with determinant one.  Formula \eqref{inv:eq:A-modular-K} excludes any
further zero, proving the asserted divisor.
\end{proof}

\begin{theorem}[inv:A2.2; T; The evaluated scalar and the two local leading terms]\label{inv:A2.2}
In the coordinate \eqref{inv:eq:A-s-branch},
\begin{equation}\label{inv:eq:A-mu-value}
\mu_s(0)=i\,C_4,
\qquad
C_4=\frac{\sqrt{2/3}}{\pi E_4(i)^{1/4}}
=\frac{4\sqrt\pi}{3^{3/4}\Gamma(1/4)^2}>0.
\end{equation}
Numerically,
\[
C_4=0.2366093140807852246877051134504496062943051340858243372192243072\ldots.
\]
For $r^3=t$, with $r>0$ on the arc $0<t<1$ and $\tau$ travelling
from $\rho$ to $i$ along the unit circle, set
\[
\lambda=\frac{3e^{5\pi i/6}}{2\pi E_6(\rho)^{1/3}}.
\]
Then $\tau=\rho+\lambda r+O(r^2)$ and
\[
K_r=-\frac{3i\lambda}{2\pi}r^2+O(r^3).
\]
The leading coefficient has real and imaginary parts
\[
\begin{aligned}
\operatorname{Re}K_2&=0.08010221108753827140763650309934451243\ldots,\\
\operatorname{Im}K_2&=0.13874109940222333968402480721280640441\ldots.
\end{aligned}
\]
\end{theorem}
\uses{inv:A2.1; the already identified elliptic section fam:B.1.3; the displayed Ramanujan identities.}

\begin{proof}
\par\smallskip\noindent\textbf{Step 1.}\ Ramanujan's differential relations give
$E_6'(i)=-\pi i E_4(i)^2$ and
$E_4'(\rho)=-(2\pi i/3)E_6(\rho)$.
Using $j=1728E_4^3/(E_4^3-E_6^2)$, we obtain
\[
j(\tau)-1728=-1728\pi^2E_4(i)(\tau-i)^2+O((\tau-i)^3).
\]
Since $j=1728(1+s^4)$, the chosen branch is precisely
\eqref{inv:eq:A-s-branch}.  Formula \eqref{inv:eq:A-modular-K} now gives
\[
K_s(0)=-\frac4{\pi^2\sqrt{E_4(i)}}=6\mu_s(0)^2,
\]
because $\tau_s(0)=0$.  This proves nonvanishing and the first
formula in \eqref{inv:eq:A-mu-value}; the sign is fixed by the stated
coordinate convention.  Similarly
$j(\tau)=-512i\pi^3E_6(\rho)(\tau-\rho)^3+cdots$.
This gives the stated $\lambda$, and substitution into
\eqref{inv:eq:A-modular-K} gives the coefficient $-3i\lambda/(2\pi)$.

\par\smallskip\noindent\textbf{Step 2.}\ There is also a geometric nonvanishing proof independent of this
modular determinant.  Using the section identification proved in
fam:B.1.3, set $x=s^2X,y=s^3Y,t=1+s^4$.  The smooth family becomes
\[
Y^2=X^3-3(1+s^4)^3X+2(1+s^4)^4s^2,
\qquad P=(0,-\sqrt2(1+s^4)^2s)
\]
on the chosen branch.  At $s=0$ its equation is
$Y^2=X^3-3X$ and the point is $(0,0)$.  The first variation of the
curve itself is zero.  In the local coordinate $Y$ at that point,
$dX/(2Y)=dY/(3X^2-3)$, so the Abel integral of $P$ has derivative
$\sqrt2/3$.  The marked period is $-i\Omega_R$, rather than the
positive real period $\Omega_R$: the positive imaginary branch
$\tau=iy$, $y>1$, has elliptic scale $F/s=i/\Omega_R$ at zero.
Its period has zero first derivative.  Division gives
$\mu_s(0)=i\sqrt2/(3\Omega_R)\ne0$.
Using $g_2(\mathbb Z+i\mathbb Z)=(4\pi^4/3)E_4(i)$ and
$g_2(\Omega_R(\mathbb Z+i\mathbb Z))=12$ gives
$\Omega_R=\pi E_4(i)^{1/4}/\sqrt3$, which identifies this value
with \eqref{inv:eq:A-mu-value}.  This also explains why choosing the
positive real period with a positive real $s$ instead would select
the reciprocal branch $\operatorname{Im}\tau<1$.

\par\smallskip\noindent\textbf{Step 3.}\ The producing run reports the following numerical methods and outputs.  The first evaluates
$E_4(i)=1+240\sum_{n=1}^{100}\sigma_3(n)e^{-2\pi n}$ at 100
decimal working digits and inserts it in \eqref{inv:eq:A-mu-value}.
The second integrates
\[
\Omega_R=\frac{\sqrt2}{3^{1/4}}
\int_0^{\pi/2}(1-\tfrac12\sin^2\theta)^{-1/2}\,d\theta,
\qquad C_4=\frac{\sqrt2}{3\Omega_R}.
\]
This is the real period of $Y^2=X^3-3X$ for $dX/(2Y)$; substituting
the three roots $\sqrt3,0,-\sqrt3$ in the elementary elliptic-integral
change of variable proves the formula.  The beta-integral evaluation
of the displayed complete elliptic integral gives the gamma
expression in \eqref{inv:eq:A-mu-value}.  The Fourier and quadrature
answers agree with absolute discrepancy $1.79\cdot10^{-102}$.
As a third check, Newton iteration on the Weierstrass Fourier series
with $j(\tau)=1728(1+s^4)$, using the two roots adjacent to
$(1-\tau)/2$, gives their symmetric derivative.  At $s=10^{-15}$ it
is
\[
i\,0.23660931408078522468770511345048081721971218\ldots,
\]
with discrepancy $3.13\cdot10^{-32}$ and maximum residual below
$2.9\cdot10^{-101}$.  The order-three coefficient was independently
evaluated using the divisor-sum expansion of $E_6$ and its Jacobi
theta-constant formula, with agreement to the working precision.
The source reports that the code, output, and precision accompanied that run.  These are numerical
corroborations, not interval certificates; nonvanishing and the exact
orders were proved analytically above.  For example the elementary
bound $\sigma_3(n)\le n^4$ and $e^{-2\pi}<1/500$ bounds the omitted
Fourier tail after $N$ terms by
$250(N+1)^4 500^{-(N+1)}$; roundoff is separately controlled by
the working precision and the independent calculation.
\end{proof}

\begin{theorem}[inv:A2.4; T; The global rank and Lead 2]\label{inv:A2.4}
The map
$\kappa_1:R^1f_*\mathcal E_{\rm cot}\to\mathcal L\otimes R^2f_*\mathcal O=\mathcal O$
is surjective, and $H^0(\kappa_1):\mathbb C^2\to\mathbb C$ has rank one.
Let $h=a\,d\zeta_1+d\zeta_2$ be the unique global section of
$f_*\mathcal E_{\rm cot}$ with $h(\xi)=1$.  In a genuine invariant target
frame normalized as below,
\[
\kappa_1(h\gamma)(p_2)=-\tfrac14\mu_s(0)=-\tfrac i4 C_4.
\]
The rank-one condition is equivalent to
$a\tau_z+\mu_z\not\equiv0$.  Under the same direct-image and
sheaf-surjectivity hypotheses a rank-zero global map would instead
give $h^{1,2}=2$, $h^{1,1}=3$; that alternative does not occur here.
\end{theorem}
\uses{inv:A1.2, inv:A2.1, inv:A2.2; for the explicitly counterfactual final clause, the residue calculation in inv:A1.3.}

\begin{proof}
\par\smallskip\noindent\textbf{Step 1.}\ Let $d_{ij}(s)$ be the cover Kodaira--Spencer coefficients in the
eigenframes of Statement inv:A2.1 after removing $ds$.
The invariant coarse matrix and its cup-product row have the forms
\[
[\kappa_0]=\begin{pmatrix}A&t_jB\\C&D_0\end{pmatrix},\qquad
[\kappa_1]=(-C,A,-D_0,B),
\]
in the ordered basis
$e_0b_0,s^ae_0b_1,s^{m-a}e_1b_0,e_1b_1$.  Specifically
$A=sd_{00}/m$, $B=s^{1-a}d_{01}/m$,
$C=s^{1-a}d_{10}/m$, $D_0=s^{m+1-2a}d_{11}/m$.
These are holomorphic invariant functions; the row follows by
wedging the connecting class with $b_0,b_1$.
If the row had no unit at a finite point, every entry would vanish
modulo $t_j$ and $\det\kappa_0$ would vanish to order at least two,
contradicting inv:A2.1.  Away from the finite points, including the cusp,
$\kappa_0$ is invertible, which gives a unit in the cup-product row.
Thus $\kappa_1$ is surjective.

\par\smallskip\noindent\textbf{Step 2.}\ The lift $h$ exists and is unique because $f_*\mathcal E_{\rm cot}$ is an
extension of $\mathcal O$ by $\mathcal O(-1)$.  Its invariance under $g_2$ gives
$a(g_2z)=\tau a+\mu$, hence $a(z_2)=1/2$.
In the quotient frame $(\gamma,u_{\Lambda})$ of $G_B$ the connecting map is
\[
\kappa_0(h)=\frac{s}{4}
\bigl((6a\mu_s+\beta_s)\gamma+(a\tau_s+\mu_s)u\bigr)\ell.
\]
Choose the variable eigenframe $b_1$ by Reynolds averaging so that
$\gamma\wedge b_1=v(s)\gamma\wedge u_{\Lambda}$, $v(0)=1$.
Then $s\,v(s)(\gamma\wedge u_{\Lambda})\ell$ is the actual invariant
target frame; the uncorrected expression with $v=1$ need not be
invariant away from zero.  Cup product with $\gamma$ gives at the
central point $-(a\tau_s+\mu_s)/4=-\mu_s/4$, proving rank one.

\par\smallskip\noindent\textbf{Step 3.}\ To check the stated equivalence, there is a second global section
\[
q_E=\sigma_1\otimes u_{\Lambda}+6\sigma_2\otimes\gamma
\quad\text{of }R^1f_*\mathcal E_{\rm cot}.
\]
It is the negative of the relative $(1,1)$ image of the invariant
flat class $\eta=u_{\Lambda}\wedge w_{\Lambda}+6\gamma\wedge\delta$.
It extends at the cusp because the relative logarithmic forms and
$\gamma,u_{\Lambda}$ extend there.  At a finite point it is an invariant
holomorphic tensor on the smooth cover, hence belongs to the four
explicit invariant generators above.  Its cup-product row is zero,
since the two terms are $6\mu_z$ and $-6\mu_z$.
It is independent of $h\gamma$ on every regular fibre (its
$\sigma_1\otimes u_{\Lambda}$ coefficient is nonzero), so these two sections
form a basis of the two-dimensional global section space.
Consequently the global map has rank one precisely when its value
on $h\gamma$, equivalently $a\tau_z+\mu_z$, is not identically
zero.  If its rank were zero, the long exact cohomology sequence of
$0\to\mathcal K\to\mathcal O^2\oplus\mathcal O(-1)^2\to\mathcal O\to0$
would give $h^0(\mathcal K)=2$, $h^1(\mathcal K)=1$.
The residue calculation in inv:A1.3 and Leray would then give exactly
the two counterfactual Hodge numbers stated in the claim.
\end{proof}

\begin{theorem}[inv:A1.3; T; All absolute cotangent direct images]\label{inv:A1.3}
For every admissible $c$, the six direct images are
\begin{equation}\label{inv:eq:A-absolute-table}
\begin{array}{c|cc}
q&R^qf_*\Omega_X^1&R^qf_*\Omega_X^2\\\hline
0&\mathcal O(-2)&\mathcal O(-2)\oplus\mathcal O(-1)\\
1&\mathcal O\oplus\mathcal O(-1)^2&\mathcal O(-2)\oplus\mathcal O(-1)^2\\
2&\mathcal O\oplus\mathcal O(-1)&\mathcal O.
\end{array}
\end{equation}
In particular every listed sheaf is torsion-free.
\end{theorem}
\uses{inv:A1.0--inv:A1.2, inv:A2.1, inv:A2.2, inv:A2.4; the normalized residue sequence; analytic relative duality \cite{RRV71}.}

\begin{proof}
\par\smallskip\noindent\textbf{Step 1.}\ We use Statements inv:A2.1, inv:A2.2 and inv:A2.4 above, whose proofs do not use
\eqref{inv:eq:A-absolute-table}.  They give
$\coker\kappa_0=\mathbb C_{p_1}\oplus\mathbb C_{p_2}$,
surjectivity of $\kappa_1$, and rank one on its global sections.
Write $\mathcal K=\ker\kappa_1$.  The logarithmic conormal sequence
then gives
\[
f_*\mathcal A=\mathcal O(1),\quad
0\to T\to R^1f_*\mathcal A\to\mathcal K\to0,\quad
R^2f_*\mathcal A=\mathcal O\oplus\mathcal O(-1),
\quad T=\mathbb C_{p_1}\oplus\mathbb C_{p_2}.
\]
Let $\nu:\widetilde D=\mathrm{dP}_6\sqcup S_1\sqcup S_2\to D$
be normalization.  The local bases in Statement inv:A1.1 give the exact
residue sequence
\[
0\to\Omega_X^1\to\mathcal A\to\nu_*\mathcal O_{\widetilde D}\to0.
\]
Its three residue direct images are
$\mathbb C_{p_0}\oplus\mathbb C_{p_1}\oplus\mathbb C_{p_2}$,
$\mathbb C_{p_1}\oplus\mathbb C_{p_2}$, and zero: the normalization of $W$ is
rational, and a bielliptic surface has $(h^0,h^1,h^2)(\mathcal O)=(1,1,0)$.
In degree zero the map from $\mathcal O(1)$ is local evaluation with the
nonzero multipliers $1,3,4$.  Its kernel is $\mathcal O(-2)$.

\par\smallskip\noindent\textbf{Step 2.}\ We check the torsion cancellation, rather than discarding it.  At
$p_j$, in the invariant frame $(b_0,s^ab_1)$ of $G_B$, residue of
$(dt_j/t_j)(x b_0+y s^ab_1)$ in $H^1(S_j,\mathcal O)$ is $m_jx[b_0]$.
The residue factors through the base logarithmic line on the covering
chart.  Naturality of the connecting map consequently gives
$\operatorname{res}_1\kappa_0=0$.  The determinant of $\kappa_0$
has a simple zero, so its reduction has rank one, and its image is
exactly the one-dimensional kernel $\mathbb C s^ab_1$ of this residue.
Hence $\coker\kappa_0\to H^1(S_j,\mathcal O)$ is an isomorphism.
Both torsion summands cancel, giving the canonical identification
$R^1f_*\Omega_X^1=\mathcal K$ and the asserted $R^2$.
The rank-one global map
$\mathcal O^2\oplus\mathcal O(-1)^2\to\mathcal O$ has a nonzero constant component
on a trivial summand.  Eliminating all other components by bundle
automorphisms gives
$\mathcal K=\mathcal O\oplus\mathcal O(-1)^2$.

\par\smallskip\noindent\textbf{Step 3.}\ Finally $\Omega_X^2=(\Omega_X^1)^\vee\otimes K_X$.  Relative
duality, with the locally free direct images now established, gives
\[
R^qf_*\Omega_X^2
=\bigl(R^{2-q}f_*\Omega_X^1\bigr)^\vee\otimes K_B.
\]
Substitution of $K_B=\mathcal O(-2)$ gives every entry of the second column.
No Euler-characteristic identity is used to select a Hodge number.
\end{proof}

\subsection{The Hodge table}\label{sec:hodge}

The values of $h^{0,q}$, the vanishing of $h^{p,0}$ for $p>0$, and $h^{1,1}=h^{1,2}+1$ were already obtained in \cite[Theorem 9.1 and Remark 1.3]{AC26}. The evaluated middle value completes that table.

\begin{theorem}[inv:A3.1; T; The complete Hodge table for every admissible parameter]\label{inv:A3.1}
For every $\operatorname{Im}c<-M$, the Hodge numbers are
\[
\begin{array}{c|rrrr}
q\backslash p&0&1&2&3\\\hline
3&0&0&0&1\\
2&0&1&2&1\\
1&1&2&1&0\\
0&1&0&0&0.
\end{array}
\]
\end{theorem}
\uses{inv:A1.0, inv:A1.3; coherent Leray on $\mathbb P^1$ and Serre duality.}

\begin{proof}
Since $B$ is a curve, coherent Leray has no differential and gives
for $F=\Omega_X^p$ the exact sequences
$0\to H^1(B,R^{q-1}f_*F)\to H^q(X,F)\to H^0(B,R^qf_*F)\to0$.
Substitution of \eqref{inv:eq:A-absolute-table}, with
$h^0(\mathcal O)=h^1(\mathcal O(-2))=1$ and all cohomology of $\mathcal O(-1)$ zero,
gives columns $p=1,2$ directly.  Statement inv:A1.0 gives column zero,
and Serre duality gives column three.  Thus the numbers are outputs
of the sheaf computation.  Replacing $\beta$ by $\beta+c$ leaves
every derivative, finite eigencharacter, invariant lattice and
logarithmic residue used above unchanged.  The toric gluing units
change, but their relative logarithmic derivatives still differ only
by base logarithmic forms.  Hence the same computations apply at
each admissible $c$, including every point arbitrarily close to the
strict boundary; no manifold at the excluded boundary is asserted.
\end{proof}

\subsection{The Fr\"olicher sequence}\label{sec:frolicher}

The earlier conditional study of a complex structure on $S^6$ supplies
general Hodge and spectral-sequence constraints in
\cite[Proposition 3.1, Theorem 3.2 and Corollary 3.3]{Ugarte00}.
McHugh gives Bott--Chern and Aeppli formulas in terms of Dolbeault and
Fr\"olicher data \cite[Section 7.1, pp.~24--25]{McHugh18}; Angella records
the spectral-sequence and double-complex descriptions and the resulting
Bott--Chern relations \cite[Figures 1--2 and Section 4]{Angella17}.
The geometric input here is the evaluated Hodge table of inv:A3.1;
the proof of inv:C2.1 uses the structure theorem of \cite{Stelzig21}.

Nondegeneration at $E_1$ and injectivity of the map from $H^{0,1}$ were already established in \cite[Corollary 9.22]{AC26}. The following source calculation determines all the first-differential ranks and degeneration at $E_2$.

\begin{definition}[inv:notation-hodge-input; text; The formal Hodge parameter]\label{inv:notation-hodge-input}
In this section $X$ is the standard marked $(3,4)$ threefold and
$h=h^{1,2}(X)=1$ by inv:A3.1.  We retain $h$ in the formal calculations
to specify exactly which conclusions depend on its evaluated value.
Every assertion involving $h$ is valid for every value
consistent with the Hodge table
\begin{equation}\label{inv:eq:C-hodge-input}
\begin{array}{c|rrrr}
q\backslash p&0&1&2&3\\\hline
3&0&0&0&1\\
2&0&h&h+1&1\\
1&1&h+1&h&0\\
0&1&0&0&0.
\end{array}
\end{equation}
Thus no conclusion below silently uses $h=1$.  The geometric inputs for
this part are the structure-sheaf calculation, the vanishing of
holomorphic forms, and the sphere topology, which give
\eqref{inv:eq:C-hodge-input}; Statements inv:A1.3 and inv:A3.1 determine $h$.
We use $d=\partial+\bar\partial$ on complex-valued
smooth forms and the complex orientation for integration.
\end{definition}
\uses{inv:A1.0, inv:A1.3, inv:A3.1; sphere topology.}

\begin{theorem}[inv:C1.1; T; The complete Fr\"olicher differential in terms of $h$]\label{inv:C1.1}
For a compact complex threefold with \eqref{inv:eq:C-hodge-input} and the
complex de Rham cohomology of $S^6$, the Fr\"olicher spectral sequence
satisfies $E_1\ne E_2=E_\infty$.  Its only nonzero $d_1$ ranks are
\[
\rank d_1^{0,1}=1,\quad \rank d_1^{1,1}=h,\quad
\rank d_1^{1,2}=h,\quad \rank d_1^{2,2}=1.
\]
In particular its two middle rows are exact.
\end{theorem}
\uses{The Hodge table in inv:notation-hodge-input; convergence of the Fr\"olicher spectral sequence to de Rham cohomology.}

\begin{proof}
Here $E_1^{p,q}=H^q(X,\Omega_X^p)$ and
$d_r:E_r^{p,q}\longrightarrow E_r^{p+r,q-r+1}$ are complex-linear.
For $r=2$, any source in the two middle rows has target in one of
$(2,0),(3,0),(3,1)$ or outside the table, except the source $(0,2)$,
which is zero.  Sources in row $q=3$ outside $(3,3)$ are zero, and
$(3,3)$ has no target.  For $r=3$ the only possible targets inside the
table have sources $(0,2)$ or $(0,3)$, both zero.  For $r>3$ there is
no arrow inside the table.  A zero term remains zero on subsequent
pages.  Hence $E_2=E_\infty$.

The filtration on $H^k_{\mathrm{dR}}(X;\mathbb C)$ has no nonzero associated
graded term for $1\le k\le5$.  Every middle term of $E_2$ therefore
vanishes.  Consequently
\[
0\longrightarrow\mathbb C\longrightarrow\mathbb C^{h+1}
\longrightarrow\mathbb C^h\longrightarrow0,
\qquad
0\longrightarrow\mathbb C^h\longrightarrow\mathbb C^{h+1}
\longrightarrow\mathbb C\longrightarrow0
\]
are exact, giving the stated ranks.  The first rank is one, so $E_1$
cannot equal $E_2$.
\end{proof}

\subsection{Bott--Chern and Aeppli cohomology}\label{sec:bc-aeppli}

\begin{theorem}[inv:C2.1; T; Bott--Chern and Aeppli tables as functions of $h$]\label{inv:C2.1}
With
\[
H^{p,q}_{BC}=
\frac{\ker\partial\cap\ker\bar\partial\subset A^{p,q}}
{\im\partial\bar\partial},\qquad
H^{p,q}_{A}=
\frac{\ker\partial\bar\partial\subset A^{p,q}}
{\im\partial+\im\bar\partial},
\]
the complete tables are
\begin{equation}\label{inv:eq:C-BC-A}
\begin{array}{c|rrrr}
q\backslash p&0&1&2&3\\\hline
3&0&0&1&1\\
2&0&h&2h&1\\
1&0&2&h&0\\
0&1&0&0&0
\end{array}_{BC}
\qquad
\begin{array}{c|rrrr}
q\backslash p&0&1&2&3\\\hline
3&0&0&0&1\\
2&0&h&2&0\\
1&1&2h&h&0\\
0&1&1&0&0.
\end{array}_{A}
\end{equation}
In particular, substitution of the value $h=1$ proved in inv:A3.1 gives
the workbench's numerical tables.  The totals are
$(1,0,2,2,2,2,1)$ and $(1,2,2,2,2,0,1)$, respectively.
\end{theorem}
\uses{inv:C1.1; the bounded-double-complex decomposition \cite[Theorem 3 and Proposition 6(1)]{Stelzig21}.}

\begin{proof}
\par\smallskip\noindent\textbf{Step 1.}\ We apply the structure theorem for bounded double complexes over a
field, \cite[Theorem 3]{Stelzig21}, to the actual algebraic double
complex $(A^{\bullet,\bullet}(X),\partial,\bar\partial)$.
The word bounded refers to the bidegrees, not to the dimensions of the
spaces of smooth forms, so that theorem applies here.  It gives a
direct-sum decomposition by isomorphisms of double complexes into
squares and zigzags.  An odd zigzag has one-dimensional total
cohomology.  All odd zigzags in intermediate total degree are therefore
excluded by $b_1=\cdots=b_5=0$.  An odd zigzag of length greater than
one would have at least two vertices of its cohomological total degree.
The extreme diagonals $p+q=0$ and $p+q=6$ contain just one bidegree
each, so the remaining odd zigzags are exactly one dot at $(0,0)$ and
one at $(3,3)$.

\par\smallskip\noindent\textbf{Step 2.}\ An even zigzag of length $2r$ contributes a nonzero $d_r$ to one of
the two Fr\"olicher spectral sequences.  inv:C1.1 and complex conjugation
exclude $r\ge2$ in both sequences.  Thus the remaining summands are
arrows.  The horizontal arrows, with their multiplicities, are
\begin{equation}\label{inv:eq:C-arrows}
\begin{array}{c|c|r}
\text{source}&\text{target}&\text{multiplicity}\\\hline
(0,1)&(1,1)&1\\
(1,1)&(2,1)&h\\
(1,2)&(2,2)&h\\
(2,2)&(3,2)&1.
\end{array}
\end{equation}
The vertical arrows are their reflections in $p=q$, with the same
multiplicities.  The correspondence with spectral-sequence ranks is
\cite[Proposition 6(1)]{Stelzig21}; alternatively one computes
column cohomology on each horizontal arrow directly.

\par\smallskip\noindent\textbf{Step 3.}\ We now evaluate the displayed definitions, rather than use a formula
for the dimensions.  On an arrow, the source is killed by neither
the outgoing differential nor its total differential, whereas its
target is closed under both.  There is no composite
$\partial\bar\partial$.  Hence the target contributes one dimension
to Bott--Chern cohomology, and the source contributes one dimension
to Aeppli cohomology.  On a square, only the top right corner is closed
under both operators, and it is the image of
$\partial\bar\partial$ from the bottom left corner.  The other
three corners of the square are images of $\partial$ or
$\bar\partial$, while the bottom left corner is not in
$\ker\partial\bar\partial$.  Thus a square contributes to neither
cohomology.  A dot contributes once to each.  These computations
commute with direct sums, including the possibly infinite sum of
squares.  Counting the sources and targets in \eqref{inv:eq:C-arrows}
and its reflection gives \eqref{inv:eq:C-BC-A}.

\par\smallskip\noindent\textbf{Step 4.}\ The producing run reports that \texttt{code/c\_double\_complex.py} performs a separate
kernel-and-image computation over $\mathbb Q$ and a source-and-target
count.  It includes an acyclic square with its signs, checks
$\partial\bar\partial+\bar\partial\partial=0$, and verifies the
two answers for $h=0,1,2,3,4,5$.  Its output is
\texttt{outputs/c\_double\_complex.txt}.  This is exact linear
algebra, not a floating-point check and not a determination of $h$.
\end{proof}

\begin{announcement}[inv:C3.3; A; The complete Bott--Chern algebra, for every $h$]\label{inv:C3.3}
The positive-total-degree ideal $N_{BC}\subset H^{\bullet,\bullet}_{BC}(X)$
satisfies
\[
H^{\bullet,\bullet}_{BC}(X)=\mathbb C\cdot1\oplus N_{BC},\qquad N_{BC}^2=0.
\]
Its nonzero bidegrees and dimensions are those of inv:C2.1.  Thus the
square-zero conclusion is independent of the evaluated value of $h$.
\end{announcement}
\uses{inv:C1.1, inv:C2.1; repository inv:C1.2 and inv:C3.1; the graded Leibniz rule.}

The repository argument represents the positive Bott--Chern classes by the arrow targets in the double-complex decomposition. It combines the curvature identity of repository inv:C3.1, the graded Leibniz rule, the first-differential ranks and integration in top degree to make the products vanish.
\par\noindent\emph{proof in the repository} (\texttt{INV}, C3.3).

\subsection{Higher multiplicities and finite covers}\label{sec:multiplicities-covers}

For the original multiplicities, the canonical bundle is credited to \cite[Proposition 9.11]{AC26}; the fundamental group and rational homology of the admissible finite-fundamental-group members are credited to \cite[Theorems 7.17 and 7.22]{AC26}. The statements below retain the source hypotheses for higher multiplicities and for the finite universal covers.

\begin{lemma}[inv:B1.1; T; The relative cohomology comparison]\label{inv:B1.1}
Fix the period family and cusp filling. Let $X'$ be obtained by any
admissible translation twists and multiplicities $m_1=3k$, $m_2=4l$,
where $k,l$ are positive integers. Write $D'=W+S'_1+S'_2$ and
\[
\mathcal E_{\rm cot}'=\Omega^1_{X'}(\log D')/f'^*\Omega^1_B(\log\Sigma).
\]
The identity on the common smooth family induces isomorphisms on $B$
\[
R^qf_*\mathcal O_X\simeq R^qf'_*\mathcal O_{X'},\qquad
R^qf_*\mathcal E_{\rm cot}\simeq R^qf'_*\mathcal E_{\rm cot}'\quad(0\le q\le2).
\]
These isomorphisms identify the logarithmic connecting maps $\kappa_q$.
They also hold when the cusp collar is changed by an admissible translation.
\end{lemma}
\uses{The fixed period and cusp charts; inv:A1.0--inv:A1.2; smooth proper base change \cite[III.4.12]{BS76}.}

\begin{proof}
\par\smallskip\noindent\textbf{Step 1.}\ We construct the comparison at every point. Over $B\setminus\Sigma$
the two families are identified, and a translation acts trivially on
$H^q(F,\mathcal O_F)$ and $H^q(F,\Omega^1_F)$ for a complex torus $F$.
For the second assertion, translation fixes the invariant holomorphic
one-forms, and the first follows either from invariant Dolbeault forms
or the exterior algebra $H^\bullet(F,\mathcal O_F)=\bigwedge^\bullet H^1(F,\mathcal O_F)$.
These are statements about the induced maps on cohomology, including
their holomorphic variation.

\par\smallskip\noindent\textbf{Step 2.}\ At a point with original order $d=3$ or $4$, let
$\widetilde f:T\to\Delta_z$ be the smooth proper torus model with
coarse coordinate $q=z^d$. The new model is its pullback by $z=s^e$,
where $m=de$, followed by the free cyclic quotient of order $m$.
Smooth proper base change identifies the higher direct images upstairs
with the pullbacks of those for $T$. The relative logarithmic cotangent
bundle upstairs is precisely $\Omega^1_{T/\Delta_z}$ pulled back:
the base logarithmic direction is killed in the quotient defining $\mathcal E_{\rm cot}'$.
The derivative of a translation contributes only a multiple of $ds$;
it consequently acts on this quotient as ordinary translation on relative forms.

\par\smallskip\noindent\textbf{Step 3.}\ Average over the cyclic action and choose holomorphic eigenframes.
If a frame for $H^q(\mathcal O)$ or $H^q(\Omega^1)$ has character
$\zeta_d^{-a}$, $0\le a<d$, its original invariant extension is
$z^a\mathcal O_{\Delta_q}$. On the order-$m$ cover its character is
$\zeta_m^{-ea}$, and its invariant extension is
$s^{ea}\mathcal O_{\Delta_q}=z^a\mathcal O_{\Delta_q}$.
Taking invariants is exact in characteristic zero; hence this equality
is an equality of the extensions inside the common punctured-disc
bundle, not only an equality of their ranks or degrees. A translation
twist induces the identity on the cohomology just used and changes none
of these lattices. This proves the comparison at both finite points,
including the full nonreduced fibres.

\par\smallskip\noindent\textbf{Step 4.}\ At the cusp the relative logarithmic forms have the global toric frame
$d\log x_1,d\log x_2$, modulo $d\log q$.
A toric gluing has the form
$x_i\mapsto c_i(q)q^{b_i}x_i$ and therefore changes $d\log x_i$ by
$(q c_i'/c_i+b_i)d\log q$. It induces the identity on the relative
logarithmic quotient. The structure-sheaf comparison is also the identity
on the normalization cochain complex. This verifies the cusp extension
and the assertion about translated collars.

\par\smallskip\noindent\textbf{Step 5.}\ The logarithmic cotangent exact sequences are isomorphic on the smooth
base, so their connecting maps agree there under these comparisons.
Their source and target sheaves are locally free by the local calculations
in Section~\ref{sec:cotangent-images} and the eigenframe calculation above. A holomorphic map between
vector bundles on a reduced curve which is zero on a dense open set is
zero everywhere. Thus the connecting maps agree globally.
\end{proof}

\begin{theorem}[inv:B2.1-canonical; T; Canonical bundles for all coprime multiplicities]\label{inv:B2.1-canonical}
Let $X=X_{3k,4l}(u)$ with admissible translation data as in
\cite[Theorem 1.1 and Proposition 2.3]{VL26},
$\gcd(3k,4l)=1$, and any nonzero $u$ in its construction domain.
Then
\[
K_X=f^*\mathcal O_B(-1)\otimes
\mathcal O_X\bigl((k-1)S_1+(3l-1)S_2\bigr).
\]
\end{theorem}
\uses{inv:B1.1; \cite[Theorem 1.1 and Proposition 2.3]{VL26}; relative duality \cite{RRV71}.}

\begin{proof}
\par\smallskip\noindent\textbf{Step 1.}\ The cited proposition supplies exactly the free cyclic quotients used
in inv:B1.1; its Theorem 1.1 supplies the stated sphere topology for its
admissible choices. For the canonical bundle we check the natural evaluation
$f^*f_*\omega_{X/B}\to\omega_{X/B}$.
The invariant relative determinant has weights $a_1=2k$ and $a_2=l$
under the clockwise generators of orders $3k$ and $4l$. The pullback
of a canonical form is $F(s)\,ds\wedge e$, where the character of
$ds\wedge e$ is $\zeta_m^{1+a}$. Therefore its invariant coefficients
are $s^{m-1-a}$ times an invariant function and a nowhere vanishing
unit. Since the quotient of total spaces is etale, the evaluation
vanishes to order $m-1-a$ on the reduced fibre downstairs. These orders
are $k-1$ and $3l-1$. At the cusp the form
$d\log x_1\wedge d\log x_2\wedge dq$ is a nowhere vanishing canonical
frame in every smooth toric chart, so no cusp component occurs.

\par\smallskip\noindent\textbf{Step 2.}\ By inv:B1.1 and relative duality,
$f_*\omega_{X/B}=(R^2f_*\mathcal O_X)^*=\mathcal O(1)$.
The evaluation divisor just computed therefore gives
$\omega_{X/B}=f^*\mathcal O(1)\otimes
\mathcal O((k-1)S_1+(3l-1)S_2)$. Tensoring by $f^*K_B=f^*\mathcal O(-2)$
proves the formula.
\end{proof}

\begin{announcement}[inv:B1.2; A; All admissible twists have the same Hodge numbers]\label{inv:B1.2}
For every family in inv:B1.1 whose cusp is the one-component hexagonal filling,
the direct images of the absolute cotangent sheaves are
\[
\begin{array}{c|ccc}
 &R^0f'_*&R^1f'_*&R^2f'_*\\\hline
\Omega^1_{X'}&\mathcal O(-2)&\mathcal O\oplus\mathcal O(-1)^2&\mathcal O\oplus\mathcal O(-1)\\
\Omega^2_{X'}&\mathcal O(-2)\oplus\mathcal O(-1)&\mathcal O(-2)\oplus\mathcal O(-1)^2&\mathcal O
\end{array}
\]
and $R^qf'_*\mathcal O_{X'}=(\mathcal O,\mathcal O\oplus\mathcal O(-1),\mathcal O(-1))$.
In particular every admissible $(3,4)$ member with
$|12l_0-4l_1-3l_2|=1$ has the Hodge diamond of Section~\ref{sec:hodge}.
\end{announcement}
\uses{inv:B1.1; the repository residue calculation and relative duality.}

The repository compares the invariant extensions of the logarithmic cohomology bundles under the enlarged cyclic roots. The absolute residue maps then remove the same two finite skyscraper terms, and relative duality supplies the second cotangent row.
\par\noindent\emph{proof in the repository} (\texttt{INV}, B1.2).

\begin{announcement}[inv:B2.1-hodge; A; Hodge numbers for all coprime multiplicities]\label{inv:B2.1-hodge}
Let $X=X_{3k,4l}(u)$ with admissible translation data as in
\cite[Theorem 1.1 and Proposition 2.3]{VL26},
$\gcd(3k,4l)=1$, and any nonzero $u$ in its construction domain.
Then $(h^{0,q})_{q=0}^3=(1,1,0,0)$, $h^{p,0}=0$ for $p>0$,
and $(h^{1,1},h^{1,2})=(2,1)$ for every such pair.
Thus these Hodge numbers do not distinguish the multiplicity families.
\end{announcement}
\uses{inv:B1.2; \cite[Theorem 1.1 and Proposition 2.3]{VL26}.}

The announced direct-image comparison gives the structure-sheaf and cotangent cohomology by Leray on the base. These are the Hodge clauses of the same source statement as the preceding canonical formula.
\par\noindent\emph{proof in the repository} (\texttt{INV}, B2.1).

\begin{lemma}[inv:B3.1; T; Connected cyclic covers and the cusp comparison]\label{inv:B3.1}
For an admissible $(3,4)$ member with
$n=|12l_0-4l_1-3l_2|>1$, let $\pi:Y\to X$ be the universal
cover and $g=f\pi$. It is a finite etale holomorphic map of degree $n$.
Every fibre of $g$ is connected; its finite reductions are bielliptic,
its multiplicities are $3,4$, and its cusp $W_n$ normalizes to a
disjoint union of $n$ copies of $\mathrm{dP}_6$.
The pullback maps induce isomorphisms
\[
R^qf_*\mathcal O_X\simeq R^qg_*\mathcal O_Y,\qquad
R^qf_*\mathcal E_{\rm cot,X}\simeq R^qg_*\mathcal E_{\rm cot,Y}.
\]
\end{lemma}
\uses{The admissible affine and cusp fillings; inv:B1.1; the elementary covering and normalization calculations in the proof; \cite[III.4.12]{BS76}.}

\begin{proof}
\par\smallskip\noindent\textbf{Step 1.}\ We first derive the fundamental-group presentation from the fillings.
The smooth family over the punctured base has zero section, so its
fundamental group is $\Lambda\rtimes\langle x,y\rangle$, where
$x,y$ are the clockwise meridians and act by $A_1,A_2$.
On a finite filling the lifted affine generator has $m_j$th power
translation by $v_j$: its $m_j$ translations add to $v_j$ because
$A_jv_j=v_j$. Thus filling the normal disc imposes
$x^3=t_{v_1}$ and $y^4=t_{v_2}$, with the signs fixed by the
clockwise convention on the first page.

\par\smallskip\noindent\textbf{Step 2.}\ The toric covering space of the cusp tube is simply connected.
Indeed its dense torus has cocharacter lattice $\mathbb Z^3$; a loop may
be perturbed into that torus, and the normal discs of the three rays
$(0,1),(e_1,1),(e_2,1)$ kill a basis of this lattice.
This argument applies also to the infinite fan, since a compact loop
and a homotopy meet only finitely many of its locally finite charts.
The same conclusion holds for the small tube: if a loop in
$|t_c|<\epsilon$ bounds a disc in the toric space, the torus action
$(x_1,x_2,t_c)\mapsto(x_1,x_2,a t_c)$ for a sufficiently small
positive $a$ moves that compact disc into the tube, and the scaling
homotopy of its boundary stays in the tube.
The cusp quotient therefore has fundamental group
$\Lambda/\Lambda_{\rm tor}$, and its toric meridian bounds the
normal disc at fixed $x_1,x_2$. A translated collar makes this
meridian $t_\mu y^{-1}x^{-1}$, where $\gamma(\mu)=l_0$.
Van Kampen consequently kills $\Lambda_{\rm tor}$ and imposes
$xy=t_\mu$.

\par\smallskip\noindent\textbf{Step 3.}\ The normal closure of $\Lambda_{\rm tor}$ in the fibre lattice is
$\ker\gamma$: it contains $\hat w,\hat\delta$ and
$A_1\hat w+\hat w=\hat u$, and all these relations lie in
$\ker\gamma$. The remaining generator $c$, the image of
$\hat\gamma$, is central because $\gamma A_j=\gamma$.
The presentation is thus
$\langle c,x,y\mid c\text{ central},\ xy=c^{l_0},\
x^3=c^{l_1},\ y^4=c^{l_2}\rangle$.
Eliminating $y$ makes it abelian, with relation matrix
\[
 \begin{pmatrix}3&-l_1\\4&l_2-4l_0\end{pmatrix},
 \qquad\det=-p.
\]
Its first Smith divisor is one because its entries include three
and four; hence the group is $\mathbb Z/|p|$.
After imposing $c=1$ the quotient has $x^3=x^4=1$, so it is trivial.
Thus the image of the regular-fibre lattice, containing $c$, generates
the entire group $\mathbb Z/n$. The same is true of each reduced special
fibre. At a finite point the cyclic model is a smooth torus family
over a disc, with a free affine finite-group action. An equivariant
smooth trivialization, obtained by averaging a horizontal connection,
contracts the disc to its centre and descends to a deformation
retraction of the local total space onto $S_j$. Thus the image of
the regular-fibre lattice in the local fundamental group is contained
in the image of $\pi_1(S_j)$; the same holds after mapping to
$\pi_1(X)$. At the cusp $\hat\gamma$ survives in
$\Lambda/\Lambda_{\rm tor}$. The following source normal-crossing calculation contracts the tube onto $W$.

\par\smallskip\noindent\textbf{Step 4.}\ Let $N$ be the small cusp tube. The following normal-crossing contraction is the one supplied in inv:B3.3. In each normal-crossing chart
$t_c=z_1\cdots z_r$, $1\le r\le3$, choose the real vector field
whose $z_i$ component is
$-z_i|z_i|^{-2}/\sum_j|z_j|^{-2}$ off the divisor.
It satisfies $\dot t_c=-t_c$ and has norm at most
$C|t_c|^{1/r}$. A finite partition of unity preserves the first
identity and gives a global field with norm at most
$C'|t_c|^{1/3}$ after shrinking the tube. Extend it by zero on $W$.
Its flow exists for all positive time, and its tail length is
bounded by $3C'\epsilon^{1/3}e^{-T/3}$ after time $T$.
Thus the flow converges uniformly to $W$ and gives a continuous
strong deformation retraction fixing $W$. The same estimate proves
continuity at the divisor and at the endpoint of the retraction.

\par\smallskip\noindent\textbf{Step 5.}\ It identifies their fundamental groups, so the same generator is in the image of $\pi_1(W)$.
The covering-space criterion therefore makes the inverse image of
every fibre connected. Admissibility gives $3\nmid l_1$ and odd $l_2$,
so $\gcd(n,12)=1$. The etale cover preserves fibre multiplicities.
On a finite reduction it cannot kill the nontrivial order-$3$ or
order-$4$ linear holonomy, since the covering degree is prime to it;
the resulting finite torus quotient is again bielliptic.

\par\smallskip\noindent\textbf{Step 6.}\ Normalization commutes with the finite etale base change $Y\to X$.
The pullback of the simply connected normalization $\mathrm{dP}_6$
is its trivial degree-$n$ cover. The double curves and triple points
also pull back etale, and their incidence maps are the lifts of the
original incidence maps. Consequently the dual complex is an actual
degree-$n$ covering of the original real two-torus, rather than merely
a cell complex with the same cell counts. It is connected because
$W_n$ is connected. The normalization resolution for $\mathcal O_{W_n}$
has acyclic components, double curves and point strata. Its global
cochain complex is consequently the cellular cochain complex of this
dual torus, with $n,3n,2n$ cells. Thus its coherent cohomology is
$(1,2,1)$, and pullback from $W$ is an isomorphism over $\mathbb C$:
on a finite cover of real tori it is induced by a full-rank map of
their integral first cohomology lattices.

\par\smallskip\noindent\textbf{Step 7.}\ Modulo $d\log q$, the two toric forms $d\log x_i$ are invariant
under all cusp gluings as computed in inv:B1.1. Hence
$\mathcal E_{\rm cot,X}|_W=\mathcal O_W^2$ and $\mathcal E_{\rm cot,Y}|_{W_n}=\mathcal O_{W_n}^2$,
and the same pullback argument applies to their cohomology.
Over a regular fibre, an isogeny of complex tori induces isomorphisms
on $H^q(\mathcal O)$ and $H^q(\Omega^1)$ over $\mathbb C$.
At a finite fibre, make the cyclic base change. The map is an isogeny
of the smooth torus families and induces isomorphisms of their
cohomology bundles. Taking the same cyclic invariant eigenspaces
descends these isomorphisms, including all nilpotent layers of the
scheme fibre. Grauert base change at the cusp, and this explicit
calculation at the finite points, prove the asserted bundle maps are
isomorphisms on all of $B$.
\end{proof}

\begin{theorem}[inv:B3.2; T; Cotangent direct images and Hodge numbers of the covers]\label{inv:B3.2}
Write $T_n=\mathbb C_{p_0}^{\,n-1}$. Then
\[
\begin{array}{c|ccc}
 &R^0g_*&R^1g_*&R^2g_*\\\hline
\Omega^1_Y&\mathcal O(-2)&\mathcal O\oplus\mathcal O(-1)^2\oplus T_n&\mathcal O\oplus\mathcal O(-1)\\
\Omega^2_Y&\mathcal O(-2)\oplus\mathcal O(-1)&\mathcal O(-2)\oplus\mathcal O(-1)^2&\mathcal O\oplus T_n
\end{array}
\]
The splittings involving $T_n$ are not canonical.
The only nonzero Hodge numbers of $Y$ are
\[
h^{0,0}=h^{0,1}=h^{3,2}=h^{3,3}=1,\quad
h^{1,2}=h^{2,1}=1,\quad h^{1,1}=h^{2,2}=n+1.
\]
The original finite-fundamental-group threefold $X$ itself has the
same Hodge diamond as the simply connected $(3,4)$ member.
\end{theorem}
\uses{inv:B3.1; inv:B1.2 (A, source dependency retained); inv:A1.3; derived relative duality \cite{RRV71}.}

\begin{proof}
\par\smallskip\noindent\textbf{Step 1.}\ inv:B3.1 identifies the logarithmic cotangent sequence on cohomology,
and therefore its connecting maps and direct images, with those of
the original family. The normalized residue sheaf differs only at
the cusp: its degree-zero cohomology there is now $\mathbb C^n$.
The residue of a base logarithmic form is the diagonal vector
$(1,\ldots,1)$, since each cusp component has multiplicity one.
Its cokernel is $T_n$. The degree-one residues at the finite fibres
cancel the same two skyscrapers as in inv:B1.2. The long exact sequence
therefore gives
\[
0\longrightarrow T_n\longrightarrow R^1g_*\Omega^1_Y
\longrightarrow\mathcal O\oplus\mathcal O(-1)^2\longrightarrow0
\]
and the remaining entries of the first row. The extension splits
because $\operatorname{Ext}^1(V,T_n)=H^1(V^*\otimes T_n)=0$
for a vector bundle $V$ on $B$.

\par\smallskip\noindent\textbf{Step 2.}\ For clarity relative duality must now be used in its derived form.
For a sheaf $F$ locally free on the total space, it gives
\[
Rg_*(F^*\otimes\omega_{Y/B})
\simeq R\mathcal Hom_B(Rg_*F,\mathcal O_B)[-2].
\]
For $F=\Omega^1_Y$, the torsion in $R^1g_*F$ contributes
$\mathcal Ext^1_B(T_n,\mathcal O_B)$ in degree two of the right side.
This is a length-$(n-1)$ skyscraper at $p_0$. After tensoring by
$\mathcal O(-2)$ it still has that length. The locally free parts dualize
in complementary degrees, giving exactly the second row; the
extension by the locally free quotient splits for the same reason
as above. Coherent Leray computes the displayed Hodge numbers.
For $X$ with one-component cusp, inv:B1.2 already applies independently
of its fundamental group.
\end{proof}

\begin{theorem}[inv:B3.3; T; Canonical rings and rational topology of the covers]\label{inv:B3.3}
For $X$ and $Y$ as above,
\[
K_Y=\pi^*K_X=\mathcal O_Y(-2S'_2),\quad
\bigoplus_{m\ge0}H^0(Y,-mK_Y)=\mathbb C[U,V],\quad
\deg(U,V)=(1,2).
\]
The same ring holds for $X$.
Over $\mathbb C$, the Betti numbers of $Y$ are
$(1,0,n-1,0,n-1,0,1)$; its Fr\"olicher sequence degenerates at $E_2$.
For every nontrivial character of the deck group, the character
part of Dolbeault cohomology consists of one class in $(1,1)$ and
one in $(2,2)$ and no others.
\end{theorem}
\uses{inv:B3.1, inv:B3.2; the source valuation and monomial argument of inv:B2.2, printed below for $(k,l)=(1,1)$; the marked lattice and affine data.}

\begin{proof}
\par\smallskip\noindent\textbf{Step 1.}\ The canonical formula pulls back through an etale map, and
$4S'_2=g^*p_2$. inv:B3.1 gives $g_*\mathcal O_Y=\mathcal O_B$.
The valuation argument in inv:B2.2 therefore applies unchanged and proves
the ring statement with all multiplication maps.

\par\smallskip\noindent\textbf{Step 2.}\ For the invoked source argument, take $m_1=3$, $m_2=4$, $k=l=1$, $d=6$, and integers $r,s$ with $4r+3s=1$, and put $\mathcal L_{\rm root}=\mathcal O(rS_1+sS_2)$. In the following valuation calculation, $f$ denotes either $f:X\to B$ or $g:Y\to B$, and the $S_j$ denote its reduced finite fibres. The weighted root coordinates are $x$ of weight $3$ and $y$ of weight $4$.

\par\smallskip\noindent\textbf{Step 3.}\ The divisor of $f^*(t/(t-1))$ is $m_1S_1-m_2S_2$.
It follows, retaining this rational function as the comparison, that
$\mathcal L_{\rm root}^{m_2}\simeq\mathcal O(S_1)$,
$\mathcal L_{\rm root}^{m_1}\simeq\mathcal O(S_2)$ and
$\mathcal L_{\rm root}^{m_1m_2}\simeq f^*\mathcal O(1)$.
The canonical formula then gives $-K_X=\mathcal L_{\rm root}^d$.

\par\smallskip\noindent\textbf{Step 4.}\ We first establish the required section algebra for every integer
power of $\mathcal L_{\rm root}$. If a meromorphic function represents a section
of $\mathcal O(aS_1+bS_2)$, it is holomorphic on every regular compact fibre
and is constant there. It descends to the base on the smooth locus.
On a cyclic cover $q=s^m$, the finite pole bound gives a Laurent
expansion bounded below, and invariance leaves only powers divisible
by $m$. Thus descent extends meromorphically across the finite point,
with the exact valuation inequality $m\operatorname{ord}_q(h)+a\ge0$.
The identical argument at the reduced cusp has $m=1$.
Conversely each base function satisfying these bounds pulls back to
an allowed section. Consequently, as subsheaves of meromorphic functions,
\[
f_*\mathcal O(aS_1+bS_2)=
\mathcal O_B\left(\left\lfloor a/m_1\right\rfloor p_1+
\left\lfloor b/m_2\right\rfloor p_2\right).
\]
The cohomology used in this valuation calculation is $h^0(\mathbb P^1,\mathcal O(e))=\max\{e+1,0\}$.

\par\smallskip\noindent\textbf{Step 5.}\ For the algebra, let $y,x$ be the canonical sections of
$\mathcal O(S_1),\mathcal O(S_2)$ under the chosen identifications with powers of
$\mathcal L_{\rm root}$. Their powers $y^{m_1},x^{m_2}$ pull back two independent
sections of $\mathcal O_B(1)$, with zeros at the two distinct points.
For any integer $j$, choose the unique residues
$0\le a_0<m_1$, $0\le b_0<m_2$ for which
$j-m_2a_0-m_1b_0=e m_1m_2$ is integral.
The valuation formula identifies $H^0(\mathcal L_{\rm root}^j)$ with
$y^{a_0}x^{b_0}H^0(\mathcal O_B(e))$. When $e\ge0$ a basis is
\[
y^{a_0+m_1i}x^{b_0+m_2(e-i)},\qquad0\le i\le e;
\]
when $e<0$ there are no sections. These are exactly the monomials of
weighted degree $j$. Their multiplication is ordinary polynomial
multiplication, proving the algebra assertion in every degree.
When $k=l=1$, the weight-$6n$ monomials are generated by $x^2$ and
$y^3$, of regraded degrees $1,2$.

\par\smallskip\noindent\textbf{Step 6.}\ The comparisons of inv:B3.1 are pullback maps and identify the logarithmic
cohomology bundles with the invariant part; the deck group acts
trivially on those bundles. At the cusp its action on the $n$
normalization components is the regular permutation representation.
The quotient by the diagonal has each nontrivial character once.
Thus the first-row additional skyscraper in inv:B3.2 has precisely these
characters, and duality gives the conjugate characters in its
second-row additional skyscraper. All other cohomology is invariant
and equal to that of $X$. Coherent Leray now proves the assertion
about the Dolbeault character spaces.

\par\smallskip\noindent\textbf{Step 7.}\ The de Rham complex and its Fr\"olicher sequence split under the
idempotents of the finite deck group. A nontrivial character has just
the two stated Dolbeault entries. No differential can connect them
or connect either to another nonzero entry, so both survive.
The invariant summand is the de Rham complex of $X$ by finite
etale descent and averaging. We compute its rational topology here
from the actual attachment maps; the integral homology assertion
of AC26 is not used as an input. In the remainder of this proof
all homology and cohomology coefficients are $\mathbb Q$.

\par\smallskip\noindent\textbf{Step 8.}\ Let $N$ be a small closed cusp tube, $M=\partial N$, and
$U=X\setminus\operatorname{int}N$. The space $U$ is covered by
two pieces retracting onto $S_1,S_2$, whose intersection retracts
onto a regular fibre $F$. The retractions are obtained from the
equivariant smooth torus trivializations used in inv:B3.1. The maps
$F\to S_j$ are the central finite covering maps; translation of
their origins induces the identity on their rational cohomology.
Finite-cover transfer identifies their pullbacks with the invariant
subspaces for $A_j$. Put
\[
 \psi_1=2u_{\Lambda}+w_{\Lambda}+3\delta,\qquad \psi_2=u_{\Lambda}+w_{\Lambda}+2\delta,
 \qquad \eta=u_{\Lambda}\wedge w_{\Lambda}+6\gamma\wedge\delta.
\]
Direct substitution in $A_1,A_2$ shows that the invariant
degree-one spaces are $\langle\gamma,\psi_j\rangle$ and the
invariant degree-two spaces are
$\langle\gamma\wedge\psi_j,\eta\rangle$.
Their dimensions are two: the nontrivial eigenvalues of $A_j$ on
the complementary plane are the inverse primitive third or fourth
roots, so the exterior-square fixed space has exactly the two
displayed generators.
In the transfer-normalized dual homology bases, the two
Mayer--Vietoris maps are therefore the following actual matrices:
\[
 \alpha_1=\begin{pmatrix}
 1&0&0&0\\0&2&1&3\\-1&0&0&0\\0&-1&-1&-2
 \end{pmatrix},\qquad
 \alpha_2=\begin{pmatrix}
 2&1&3&0&0&0\\0&0&6&1&0&0\\
 -1&-1&-2&0&0&0\\0&0&-6&-1&0&0
 \end{pmatrix}.
\]
The second matrix uses the ordered wedge basis
$(\hat\gamma\hat u,\hat\gamma\hat w,\hat\gamma\hat\delta,
\hat u\hat w,\hat u\hat\delta,\hat w\hat\delta)$.
Both matrices have rank three, and
$\ker\alpha_1=\mathbb Q\nu_0$, with
$\nu_0=-\hat u-\hat w+\hat\delta$.
Consequently
\[
 H_1(U)=\mathbb Q,\qquad
 0\longrightarrow\operatorname{coker}\alpha_2=\mathbb Q
 \longrightarrow H_2(U)
 \xrightarrow{\partial}\mathbb Q\nu_0\longrightarrow0.       \tag{B.1}
\]
The generator $\Theta\in H^1(U)$ is normalized by
$\Theta|_F=\gamma$. The finite affine generators have
$3\bar g_1=l_1c$, $4\bar g_2=l_2c$ in rational homology, where
$c$ is the image of $\hat\gamma$ and $\Theta(c)=1$.
For the toric meridian $a_0=t_\mu y^{-1}x^{-1}$ this gives the
explicit collar value
\[
 \Theta(j_*a_0)=l_0-l_1/3-l_2/4=p/12.                 \tag{B.2}
\]

\par\smallskip\noindent\textbf{Step 9.}\ The contraction of the cusp tube used here was established in the proof of inv:B3.1 by the source normal-crossing vector-field calculation.

\par\smallskip\noindent\textbf{Step 10.}\ The constant-sheaf normalization resolution of $W$ has strata
$\mathrm{dP}_6$, three copies of $\mathbb P^1$, and two points.
Its row of degree-zero cohomology is
\[
 \mathbb Q\xrightarrow{0}\mathbb Q^3
 \xrightarrow{(a,b,c)\mapsto(a-b+c,a-b+c)}\mathbb Q^2.
\]
The strata have no degree-one cohomology. This resolution therefore
gives $H_1(N)=H_1(W)=\mathbb Q^2$, $H_4(N)=\mathbb Q[W]$ and $H_5(N)=0$;
it also gives $e(W)=6-3\cdot2+2=2$.
Poincare--Lefschetz duality for the compact smooth manifold with
boundary $(N,M)$ now shows that $H_1(M)\to H_1(N)$ and
$H_2(M)\to H_2(N)$ are surjective. For degree one the next
relative group is $H_1(N,M)=H^5(N)=0$.
For degree two the map to $H_2(N,M)=H^4(N)$ is the intersection
pairing with $[W]$, which is zero because $W$ is the divisor of
the global function $t_c$ on $N$. The Wang sequence for the
boundary mapping torus, with $\rank(M_0-I)=2$, gives
$\dim H_1(M)=3$. Its kernel to $H_1(N)$ is exactly
$\mathbb Q a_0$: this meridian is nonzero under projection to the base
circle and bounds its normal disc in $N$.

\par\smallskip\noindent\textbf{Step 11.}\ We exhibit two classes in
$K_2=\ker(H_2(M)\to H_2(N))$.
For $\nu=\hat w,\hat\delta$, let $\xi_\nu$ be the two-torus
obtained by translating the toric meridian around the angular
circle $S^1_\nu$. In the chart $t_c=z_0z_1z_2$ corresponding to
the rays $(0,1),(e_1,1),(e_2,1)$, these are the boundaries of
$D^2_{z_0}\times S^1_{z_i}$ with the other coordinate fixed
and the radii chosen so that $|t_c|=\epsilon$ on the boundary.
They therefore bound in $N$ with the literal toric meridian;
no choice of an unrecorded collar lift is being made.
Cutting the meridian into the two arcs of the cover defining
(B.1), the two cylinders have end circles differing by
$(I-A_1^{-1})\nu$. With this orientation of the Mayer--Vietoris
boundary, matrix multiplication gives
\[
 \partial j_*\xi_{\hat w}=-\nu_0,\qquad
 \partial j_*\xi_{\hat\delta}=0.                       \tag{B.3}
\]
The affine translation of either endpoint has no effect on its
homology class, which explains why this boundary formula is the
same for every translation twist.

\par\smallskip\noindent\textbf{Step 12.}\ Choose $\mathcal A\in H^2(U)$ restricting to $\eta$ on $F$;
it exists by Mayer--Vietoris because the two transfer-normalized
classes $\eta$ have equal restrictions. Fibrewise translation by
$S^1_{\hat\delta}$ defines a global smooth action on $U$:
$A_j\hat\delta=\hat\delta$, and all affine translations commute
with this circle. Slant product of its action map with the circle
defines $\iota:H^2(U)\to H^1(U)$. Naturality on $F$ gives
$\iota\mathcal A|_F=\iota_{\hat\delta}\eta=-6\gamma$.
The restriction $H^1(U)\to H^1(F)$ is injective by (B.1), so
$\iota\mathcal A=-6\Theta$. By the definition of slant product,
\[
 \langle\mathcal A,j_*\xi_{\hat\delta}\rangle
 =\langle\iota\mathcal A,j_*a_0\rangle=-p/2.           \tag{B.4}
\]
We may add the class pulled back from $(\mathbb Q\nu_0)^*$ under
the boundary map in (B.1) so that
$\langle\mathcal A,j_*\xi_{\hat w}\rangle=0$;
this does not change its restriction to $F$ or (B.4).
In the coordinates on $H_2(U)$ given by $\partial/\nu_0$ and
this normalized pairing, the two actual attachment cycles have
matrix
\[
 \bigl[j_*\xi_{\hat w}\ \ j_*\xi_{\hat\delta}\bigr]
 =\begin{pmatrix}-1&0\\0&-p/2\end{pmatrix}.             \tag{B.5}
\]
It is invertible over $\mathbb Q$ because admissibility implies
$p\ne0$. Thus the image of $K_2$ already fills $H_2(U)$.
Likewise (B.2) says the image of $\ker(H_1(M)\to H_1(N))$
is nonzero in $H_1(U)$. Using the surjectivity to $N$ just proved,
Mayer--Vietoris for $X=U\cup N$ now gives $H_1(X)=H_2(X)=0$.
Every other fibre has Euler number zero, so additivity over the
three special fibres and the smooth torus bundle gives $e(X)=2$.
Poincare duality then gives $b_4=b_5=0$ and $b_3=0$.
This proves the asserted rational homology of $X$ for all admissible
twists; the only changing determinant in (B.2)--(B.5) is a nonzero
constant multiple of $p$. The producing run reports that the integer matrices and all exterior-power identities were checked exactly in \texttt{code/b\_topology.py},
with output \texttt{outputs/b\_topology.txt}.

\par\smallskip\noindent\textbf{Step 13.}\ The argument of inv:C1.1 using this rational topology and the same Hodge
table now gives degeneration of the invariant summand at $E_2$.
Summing it with the nontrivial character spaces proves the Betti
numbers and the degeneration statement for $Y$. In particular
$e(Y)=2n$, as also follows from the covering degree.
\end{proof}

\subsection{Metric consequences}\label{sec:metric-announcements}

\begin{announcement}[inv:C3.5; A; Metric obstructions that do not require $h$]\label{inv:C3.5}
The threefold $X$ admits no balanced, strongly Gauduchon, K\"ahler,
locally conformally K\"ahler, or Hermitian-symplectic metric.
\end{announcement}
\uses{The fibration and sphere topology; Stokes' theorem.}

The source uses the nonzero semipositive exact form $f^*\omega_{FS}$ and the positive four-form supplied by a balanced or strongly Gauduchon metric. The remaining exclusions use $b_1=b_2=0$ and the same integral obstruction.
\par\noindent\emph{proof in the repository} (\texttt{INV}, C3.5).

\begin{announcement}[inv:C3.7; A; The fibre obstruction to SKT and astheno-K\"ahler metrics]\label{inv:C3.7}
The threefold $X$ admits no pluriclosed (SKT) metric and no
astheno-K\"ahler metric.  The geometric input is the value $h=1$
and the global Dolbeault class $q$ in inv:A3.2, whose restriction to a
regular fibre is $\eta=u_{\Lambda}\wedge w_{\Lambda}+6\gamma\wedge\delta$.
More generally, it is enough that the real image of
$H_A^{1,1}(X;\mathbb R)\to H_A^{1,1}(F;\mathbb R)$ be contained in
$\mathbb R\eta$ for one regular fibre $F$.
\end{announcement}
\uses{inv:A3.1; the global Dolbeault lift in repository inv:A3.2; inv:C2.1.}

The repository restricts the prospective metric class to a regular torus fibre. Its real Aeppli image lies in $\mathbb R\eta$, while the complex-oriented square of $\eta$ is $-12$; translation averaging would give a positive class with positive square.
\par\noindent\emph{proof in the repository} (\texttt{INV}, C3.7).

\subsection*{Correction concerning relative differentials}\label{sec:relative-correction}
The locators are \cite[Remark 1.3, printed p.~5; Remark 9.21, printed p.~81]{AC26}.
\emph{Paraphrase of the printed assertions:} the former attributes torsion
along the singular fibres to the relative differential sheaves without
separating their degrees; the latter includes the reduced cusp $W$ in the
torsion support of $\Omega^1_{X/B}$. This is a paraphrase, not a quotation:
the requirement not to reproduce the prose of \cite{AC26} prevents the
separately requested literal printed sentence from being supplied.
The correction is local and does not alter the preceding credits.


\begin{lemma}[inv:A4.1; T; All local modules needed to distinguish degrees one and two]\label{inv:A4.1}
Let $R=\mathbb C\{x,y,z\}$.  For $t=xy$,
\[
\Omega_{R/\mathbb C\{t\}}^1=R^3/R(y,x,0)
\simeq(x,y)\oplus R,
\]
and
\[
\Omega_{R/\mathbb C\{t\}}^2\simeq
(x,y)\oplus R/(x,y),\qquad
\Omega_{R/\mathbb C\{t\}}^3\simeq R/(x,y).
\]
For $t=xyz$,
$\Omega^1=R^3/R(yz,xz,xy)$ is torsion-free of rank two and is
not free at the triple point.  If $I=(xy,xz,yz)$, there is an exact
sequence
\[
0\to T\to\Omega^2\xrightarrow{\wedge dt} I\to0,
\qquad
T=\frac{R u\oplus R v}{(xu,\ y(u-v),\ zv)},
\quad IT=0,
\]
and $\Omega^3=R/I$.  Finally, for $t=s^m$ in
$R=\mathbb C\{s,y_1,y_2\}$,
\[
\begin{aligned}
\Omega^1&=R^2\oplus R/(s^{m-1})\,ds,\\
\Omega^2&=R\,dy_1\wedge dy_2
 \oplus\bigl(R/(s^{m-1})\bigr)^2,\\
\Omega^3&=R/(s^{m-1})\,ds\wedge dy_1\wedge dy_2.
\end{aligned}
\]
\end{lemma}
\uses{The presentations of ordinary relative Kähler differentials for $t=xy$, $t=xyz$, and $t=s^m$.}

\begin{proof}
\par\smallskip\noindent\textbf{Step 1.}\ The degree-one presentation is the quotient by $dt$; higher degrees
are its exterior powers.  At a double point the transverse quotient
$R^2/R(y,x)$ maps isomorphically to $(x,y)$ by
$(a,b)\mapsto xa-yb$ after an immaterial change of sign of one
basis vector.  More explicitly one may use
$(a,b)\mapsto xa-yb$ with relation $(y,x)$; its kernel is precisely
that relation, since $x,y$ are relatively prime.  In degree two,
with basis $A=dx\wedge dy,B=dx\wedge dz,C=dy\wedge dz$, the
relations are $xA=yA=0$, $yB+xC=0$.  This gives the displayed
torsion summand and the ideal $(x,y)$; wedging once more gives
$R/(x,y)$.

\par\smallskip\noindent\textbf{Step 2.}\ At the triple point the map from ordinary to logarithmic relative
one-forms sends $(a,b,c)$ to $(xa-zc,yb-zc)$.
Its kernel is precisely $R(yz,xz,xy)$: the two equations imply
$c=xyh$, $a=yzh$, $b=xzh$ by relative primality of the coordinates.
Thus it injects into $R^2$, proving torsion-freeness.  Its fibre at
the maximal ideal has dimension three although its rank is two,
so it is not free.  The same minimal-generator argument works at
the double locus.

\par\smallskip\noindent\textbf{Step 3.}\ In degree two, wedging with $dt$ is the row
$(xy,-xz,yz)$ in the basis $(A,B,C)$.
Its syzygies are generated freely by
$u=(z,y,0)$ and $v=(0,y,x)$.
To check generation, an equation $xyA-xzB+yzC=0$ implies
$x\mid C$, say $C=xC_1$; then $yA=z(B-yC_1)$ implies
$A=zA_1$, $B=y(A_1+C_1)$.  Thus the syzygy is $A_1u+C_1v$.
The three relations from $dt\wedge(dx,dy,dz)$ are
$-xu,y(u-v),zv$, proving the stated presentation of the kernel.
The ideal $I$ kills it by those relations.  Conversely the quotient
$I$ is torsion-free, so this kernel is exactly the torsion of
$\Omega^2$.  The presentation immediately gives $\Omega^3=R/I$.
All matrix identities have also been checked over $\mathbb Z[x,y,z]$ by
the supplied exact program.  For $t=s^m$ the sole degree-one
relation is $ms^{m-1}ds=0$; $m$ is a unit and taking exterior
powers gives the three formulas.
\end{proof}

\begin{remark}[inv:A4.2; T; The claimed degree-one torsion along the cusp]\label{inv:A4.2}
The assertion in AC26 Remark 9.21, PDF page 81, that
$\Omega_{X/B}^1$ has torsion along $W\cup S_1\cup S_2$ is false.
Its torsion is supported precisely on $S_1\cup S_2$; its degree-one
cusp modules are torsion-free but not locally free on the double
locus.  The plural statement in Remark 1.3, PDF page 5, must be
qualified by degree: $\Omega^2_{X/B}$ does have torsion on the cusp
double locus.  The workbench's degree-one correction is correct;
it must not be extended to degree two.
\end{remark}
\uses{inv:A4.1; \cite[Remarks 1.3 and 9.21]{AC26}.}

\begin{proof}
The producing run reports that both disputed sentences were checked on rendered pages of the original PDF. The present source check uses the supplied transcription at the exact locators above.  The
double and triple charts in inv:A4.1 exhaust the singular strata of
$W$.  They prove absence of degree-one torsion at the cusp and
exhibit nonzero degree-two torsion there.  At a multiple fibre the
finite etale comparison
$q^*\Omega^1_{N_j/B}=\Omega^1_{J'_j/B}$ is an isomorphism,
and its last summand in inv:A4.1 is nonzero for $m=3,4$.
Faithful flatness descends both its existence and its support.
There are no other critical fibres and the ordinary relative
cotangent is locally free on their complement.  Thus the stated
support is exact.  The erroneous sentence is at transcription
line 4330 (the final paragraph of Remark 9.21); the workbench's
correct degree-one replacement occupies lines 980--1006 of
\texttt{hodge\_frolicher.tex}.  No numerical Hodge value in that
workbench is refuted by this calculation.
\end{proof}

\section{Tangent cohomology and local deformations}\label{sec:tangent}
\subsection{The standard family: the known Kuranishi line}\label{sec:tangent-known}
The following statement is credited to \cite[Theorem 1.3 and Theorem 6.12]{DGKXv1}.
The first source theorem gives the Kuranishi and miniversality conclusions;
the tangent-cohomology theorem gives the displayed dimensions.
The supplied re-derivation is in Appendix~\ref{app:tangent-rederivation}.


\begin{cited}[fam:A.5.2; K; Tangent cohomology and the Kuranishi line]\label{fam:A.5.2}
For every admissible $c$, the tangent cohomology of the actual compact
member is
\[
 \bigl(h^0(X_c,T_{X_c}),h^1(X_c,T_{X_c}),h^2(X_c,T_{X_c}),
 h^3(X_c,T_{X_c})\bigr)=(1,1,1,0).
\]
Its Kuranishi germ is a smooth reduced disc. The family of fam:A.5.1 is
miniversal at every parameter: its Kodaira--Spencer map is an isomorphism,
and every sufficiently small deformation is induced from this family.
\end{cited}
\uses{The compact family fam:A.5.1 and its nonzero Kodaira--Spencer class fam:A.4.2; \cite[Theorems 1.3 and 6.12]{DGKXv1}.}

\subsection{Every integral Picard twist}\label{sec:tangent-twists}

\begin{definition}[cdp:notation-integral-twists; text; The integral Picard normalization]\label{cdp:notation-integral-twists}
In this subsection $X=X_c$ is the period construction of \cite{AC26}, compared with the quotient construction of \cite{E26}, with admissible parameter $c$, standard clockwise affine vectors
$v_1=(1,2,-4,0)^t$, $v_2=(-1,-3,3,0)^t$, and untwisted cusp collar. Thus $p=12l_0-4l_1-3l_2=-1$. We write $\mathfrak B=\mathbb P(3,4)$, $\varpi:\mathfrak B\to B$ for its coarse map, and $\bar f:X\to\mathfrak B$ for the root factorization. Let $\mathcal T=T_X(-\log W)$ and let $\nu:\mathrm dP_6\to X$ normalize $W$. The root generator is fixed by
\[
 M_n=\bar f^*\mathcal O_{\mathfrak B}(n)
     =\mathcal O_X(n(S_1-S_2)),\quad n\in\mathbb Z.
\]
The exponential isomorphism $H^1(X,\mathcal O_X)\to\Pic(X)$ identifies $\Pic(X)$ with $\mathbb C$. We scale its coordinate so that $M_z=\exp(z\gamma^{0,1})$ and $M_1=\mathcal O_X(S_1-S_2)$. The sign of this equality is fixed by the clockwise transgression $12\gamma\mapsto-p=1$. This is checked by evaluating the two finite lifts on $v_1,v_2$: their integral periods are $4l_1,3l_2$, while the cusp contributes $-12l_0$, whose sum is $-p$. The exponential sequence therefore gives $\exp(12\gamma^{0,1})=f^*\mathcal O_B(1)$. Its injectivity and the root relation give the stated $M_1$ equality.
\end{definition}
\uses{The standard clockwise affine data and exponential sequence; the root factorization.}

\begin{theorem}[cdp:C5b; T; Integral twists and the normalization-map exception]\label{cdp:C5b}
For $n\in\mathbb Z$ define
\[
 r_3(n)\equiv n\pmod3,\quad 0\le r_3(n)<3,
 \qquad r_4(n)\equiv3n\pmod4,\quad0\le r_4(n)<4,
\]
\begin{gather*}
 d(n)=\frac{n-4r_3(n)-3r_4(n)}{12},\qquad
 \epsilon_n=\begin{cases}0&3\mid n,\\1&3\nmid n.\end{cases}
\\
 u(n)=\max\{d(n)+1,0\},\qquad
 v(n)=\max\{-d(n)-1,0\}.
\end{gather*}
The coarse direct images of $T_X\otimes M_n$ are
\begin{align*}
 f_*(T_X\otimes M_n)&=\mathcal O(d(n))\oplus\mathcal O(d(n-1)),\\
 R^1f_*(T_X\otimes M_n)&=\mathcal O(d(n))\oplus\mathcal O(d(n-1))\oplus\mathcal O(d(n+2)),\\
 R^2f_*(T_X\otimes M_n)&=\mathcal O(d(n+6))\oplus\mathbb C_{p_1}^{\epsilon_n}.\tag{C.5}
\end{align*}
Their global cohomology dimensions are
\begin{align*}
 h^0&=u(n)+u(n-1),\\
 h^1&=v(n)+v(n-1)+u(n)+u(n-1)+u(n+2),\\
 h^2&=v(n)+v(n-1)+v(n+2)+u(n+6)+\epsilon_n,\\
 h^3&=v(n+6).                                             \tag{C.6}
\end{align*}
The natural map $H^2(X,T_X\otimes M_n)\to H^2(\mathrm dP_6,K_{\mathrm dP_6})=\mathbb C$ has rank one exactly for
\[
 n\in\{-6,-3,-2\}\cup\mathbb Z_{\ge0};                         \tag{C.7}
\]
it is zero for all other integers. In particular $M_{-1}$ has tangent cohomology $(0,0,1,0)$ and a zero map to the normalization target.
\end{theorem}
\uses{The root and normalization calculations below; inv:D1.3 for the source-expanded branch-normal and conormal computation; analytic base change \cite[III.4.10, III.4.12]{BS76} and relative duality \cite{RRV71}.}

\begin{proof}
\par\smallskip\noindent\textbf{Step 1.}\ We give the root-chart and cusp calculations used in this proof. The root charts at $p_1,p_2$ are the finite etale torus-family covers used in the construction; finite-group invariants over $\mathbb C$ are exact. At the cusp the local equations are $xy=q$ and $xyz=q$, and the vertical logarithmic fields are generated by $x\partial_x-y\partial_y$ and the remaining relative direction. Their affine transition derivatives depend only on the base. The invariant vertical field gives a trivial subline of the vertical bundle $\mathcal E_{\rm tan}$. The quotient has twelfth power $\mathcal O(-p_0)$: its automorphy is that of the reciprocal elliptic discriminant, whose only divisor is a simple pole at the cusp. As $\Pic(\mathfrak B)=\mathbb Z$, the quotient is $\mathcal O(-1)$; $H^1(\mathfrak B,\mathcal O(1))=0$ splits the extension. Thus $\mathcal E_{\rm tan}=\mathcal O\oplus\mathcal O(-1)$.

\par\smallskip\noindent\textbf{Step 2.}\ The normalization complex of $W$ has global scalar complex
\[
 \mathbb C\xrightarrow{0}\mathbb C^3
 \xrightarrow{(a,b,c)\mapsto(a-b+c,a-b+c)}\mathbb C^2.
\]
The rational normalization and the three rational double curves have no positive-degree structure-sheaf cohomology. The displayed complex therefore gives $(1,2,1)$ for the fibre cohomology and its exterior cup product; the two oriented triangular faces give the nondegenerate degree-one pairing. The cyclic-cover fibres have the same dimensions. Proper flat base change on these charts gives locally free $G_{\mathfrak B}=R^1\bar f_*\mathcal O_X$, with $R^2\bar f_*\mathcal O_X=\det G_{\mathfrak B}$. Relative duality and the determinant of the logarithmic sequence give $\det G_{\mathfrak B}=\mathcal O(-1)$. The integral class $\gamma$ maps to a nowhere-zero section of $G_{\mathfrak B}$: on a marked smooth torus the two identity columns of the period matrix prevent it from belonging to the span of the two holomorphic period rows, and at the cusp it is a nonzero cycle in the displayed complex. The same splitting argument yields
\[
\mathcal E_{\rm tan}=G_{\mathfrak B}=\mathcal O\oplus\mathcal O(-1),\quad R^2\bar f_*\mathcal O_X=\mathcal O(-1),
\quad0\to\bar f^*\mathcal E_{\rm tan}\to\mathcal T\to\bar f^*\mathcal O(-5)\to0. \tag{C.8}
\]
These are the actual root-chart and cusp comparisons used here; no coarse multiple fibre was replaced by its reduction in this argument.

\par\smallskip\noindent\textbf{Step 3.}\ Choose root sections $x$ of weight $3$, vanishing at $p_2$, and $y$ of weight $4$, vanishing at $p_1$. The Kodaira--Spencer morphism $\mathcal O(-5)\to \mathcal E_{\rm tan}\otimes G_{\mathfrak B}$ has, by its weights, the form
\[
 \kappa=\begin{pmatrix}0&a y\\b y&c_0x\end{pmatrix}.
\]
Its first coefficient vanishes because $H^0(\mathcal O(5))=0$. The logarithmic period derivative at the cusp is invertible, giving $ab\ne0$. At $p_1$, the elliptic coordinate $\tau-\rho$ has a simple zero on the root chart, so its period derivative is nonzero, giving $c_0\ne0$. Rescaling the four line frames makes the three nonzero coefficients one. Cup product with this extension class gives
\[
\delta_0=(0,y,y,x)^t,\qquad
\delta_1=\begin{pmatrix}-y&0\\-x&y\end{pmatrix}:
 \mathcal O(-5)\oplus\mathcal O(-6)\to\mathcal O(-1)\oplus\mathcal O(-2).
\]
The first is a subbundle inclusion since $x,y$ have no common zero, and its cokernel is $\mathcal O\oplus\mathcal O(-1)\oplus\mathcal O(2)$. The second is injective with determinant $-y^2$; its cokernel is
$Q=\mathcal O(-1)|_{2p_1}$, a length-two module on the order-three chart. The direct-image sequence of (C.8) yields
\[
 R^0\bar f_*\mathcal T=\mathcal E_{\rm tan},\quad
 R^1\bar f_*\mathcal T=\mathcal O\oplus\mathcal O(-1)\oplus\mathcal O(2),\quad
 R^2\bar f_*\mathcal T=\mathcal O(-6)\oplus Q.
\]
The last splitting follows from $H^1(Q(6))=0$; it is not designated canonical.

\par\smallskip\noindent\textbf{Step 4.}\ The branch-normal sequence is
\[
0\longrightarrow\mathcal T\longrightarrow T_X
\longrightarrow\nu_*K_{\mathrm dP_6}\longrightarrow0.
\]
To verify the quotient, in a chart $W=(z_1\cdots z_r=0)$ the map takes
a vector field to its normal coefficient on each of the $r$ separate
branches. The coefficients lift independently; its kernel is precisely
the vector fields tangent to every branch. Thus the quotient is the
pushforward of the normal line of $\nu$, not $\mathcal O_W(W)$.
The form $d\log x_1\wedge d\log x_2\wedge dq$ becomes a unit times
$dz_1\wedge dz_2\wedge dz_3$ in every unimodular chart and is invariant
under every cusp deck transformation. It trivializes $\nu^*K_X$.
Adjunction for the immersion gives its normal line $K_{\mathrm dP_6}$, proving
the normal-coefficient sequence. Since $\mathrm dP_6$ is rational, its $K_{\mathrm dP_6}$ cohomology is $(0,0,1)$ by
Serre duality. Finite pushforward preserves that cohomology.

\par\smallskip\noindent\textbf{Step 5.}\ It remains to identify the extension at $p_0$ in degree two. We give
a fibre-sensitive calculation. The kernel of
\[
 \Omega_X^1|_W\longrightarrow\nu_*\Omega_{\mathrm dP_6}^1
\]
is $\nu_*\mathcal O_{\mathrm dP_6}\otimes N_W^*$. At $xy=0$ its sections are
$a(y,z)y\,dx+b(x,z)x\,dy$; at $xyz=0$ they are
$a(y,z)yz\,dx+b(x,z)xz\,dy+c(x,y)xy\,dz$.
These identifications glue by the common conormal $dq$.
Since $N_W$ and $K_X|_W$ are trivial and $H^0(P,\Omega_{\mathrm dP_6}^1)=0$, the
Serre-dual space of $H^2(W,T_X|_W)$ is one-dimensional.
Top-degree base change identifies this group with the residue fibre
of $(R^2\bar f_*T_X)_{p_0}$. Its generic rank is one by the displayed logarithmic direct images.
A finite module over the discrete valuation ring $\mathbb C\{q\}$ of rank
one and residue dimension one is free of rank one. Thus the extension
by $\mathbb C_{p_0}$ increases the degree of its line part by $p_0$, without
adding torsion. 
Every $\nu^*M_n$ is trivial because $H^2(X,\mathbb Z)=0$ and
$\operatorname{Pic}^0(\mathrm dP_6)=0$. Thus the same quotient
calculation applies after each integral twist. The elementary extension
of $\mathcal O(-6)$ by $\mathbb C_{p_0}$ is $\mathcal O(6)$,
since $\mathcal O(p_0)=\mathcal O(12)$ on the root curve.
We have
\[
 R^0\bar f_*T_X=\mathcal E_{\rm tan},\quad
 R^1\bar f_*T_X=\mathcal O\oplus\mathcal O(-1)\oplus\mathcal O(2),\quad
 R^2\bar f_*T_X=\mathcal O(6)\oplus Q.                       \tag{C.9}
\]

\par\smallskip\noindent\textbf{Step 6.}\ For a root line, removing its two stabilizer characters gives
$\varpi_*\mathcal O(n)=\mathcal O_B(d(n))$. Averaging the finite stabilizers is exact, so there are no higher direct images of $\rho$. Its sections are equivalently monomials $x^ay^b$ with $3a+4b=n$, and Serre duality gives $h^1\mathcal O(n)=h^0\mathcal O(-n-7)$. At $p_1$, the invariant exponent of $Q(n)$ is $n-1$ modulo three. Its two possible coefficients have orders zero and one, so its coarse length is $\epsilon_n$ and it is killed by the coarse parameter. Apply projection formula to (C.9) and then $\varpi_*$. This proves (C.5), including the support and length of every torsion summand. Leray over $B$ has columns only zero and one, proving (C.6).

\par\smallskip\noindent\textbf{Step 7.}\ The normalization map on $R^2$ is the quotient in
\[
0\longrightarrow\mathcal O(n-6)\longrightarrow\mathcal O(n+6)
 \longrightarrow\mathbb C_{p_0}\longrightarrow0.
\]
Its global-section map is nonzero precisely when a section of $\mathcal O(n+6)$ is nonzero at $p_0$. Both $x,y$ are nonzero at that point, so this occurs precisely when $n+6$ lies in the additive semigroup generated by $3,4$. That semigroup is $\{0,3,4,6,7,8,\ldots\}$, giving (C.7). The $H^1(R^1)$ part maps to zero in the cohomology of a sheaf supported at one point, by the natural Leray filtration. At $n=-1$, (C.5) has only one nonzero cohomological contribution, the reduced $\mathbb C_{p_1}$ summand in degree two, while $H^3(X,\mathcal T\otimes M_{-1})=H^1(\mathcal O(-7))=\mathbb C$. The boundary from the normalization target is an isomorphism. This proves the stated counterexample to a surjection for every twist.
\end{proof}

\subsection{Coprime higher multiplicities}\label{sec:tangent-higher}
The tangent vector occurs in two producing runs, \texttt{PUSH} E.2.1
and \texttt{CDP} E11: \emph{two runs, no second reader}.
The statement below retains the full hypotheses and conclusions of
\texttt{PUSH} E.2.1; the comparison with \texttt{CDP} is limited to
their common higher-multiplicity tangent-cohomology conclusion.


\begin{theorem}[push:E.2.1; T; Tangent cohomology and local miniversality for coprime multiplicities]\label{push:E.2.1}
For every sphere member with multiplicities $(3k,4l)$ and the
admissible logarithmic data of push:E.1 (in particular, $\gcd(3k,4l)=1$), put $z=kl$ and $d=3k+4l$.
On $\mathfrak B_{k,l}=\mathbb P(3k,4l)$ the root lift
$\bar f:X\to\mathfrak B_{k,l}$ satisfies
\[
 \begin{split}
 \mathcal E_{\rm tan}&=\mathcal O\oplus\mathcal O(-z),\\
 R^q\bar f_*\mathcal O_X&=
 \mathcal O,\quad\mathcal O\oplus\mathcal O(-z),\quad\mathcal O(-z),\\
 K_X&=\bar f^*\mathcal O(z-d).
 \end{split}
\]
The tangent cohomology is $(h^0,h^1,h^2,h^3)(X,T_X)=(1,1,1,0)$.
With the groupoid derivative convention
$\kappa(\Phi)=(\partial_c\Phi)\Phi^{-1}$, the class of
$\partial/\partial c$ is $e\otimes[\gamma]$.  The class of
$\partial/\partial u$ is
$-(2\pi i u)^{-1}e\otimes[\gamma]$.  Both are nonzero.
The Kuranishi germ is a smooth reduced line, and the fixed-data
period family is locally miniversal.
\end{theorem}
\uses{The admissible higher-multiplicity data push:E.1; the displayed root-chart and residue computations; \cite[III.4.12]{BS76} and \cite{RRV71}. Two runs, no second reader.}

\begin{proof}
\par\smallskip\noindent\textbf{Step 1.}\ The root object is defined by the actual divisor identities
$f^*(0)=mS_1$ and $f^*(1)=nS_2$, with $m=3k$, $n=4l$.
It gives a holomorphic map to the analytic root stack.  Its base
changes to the two root discs are the actual smooth torus
families; the quotient actions are free.  At the cusp its charts
are $q=z_0z_1z_2$, so the source is smooth and the map is flat,
of relative dimension two, with Gorenstein fibres.  Properness is
checked on the root atlas, where it is properness of the compact
torus families and the original compact cusp filling.

\par\smallskip\noindent\textbf{Step 2.}\ Use weighted coordinates $x,y$ of weights $m,n$, with coarse
coordinate $t=y^m/(y^m-x^n)$.
The two root divisors have bundles $\mathcal O(n)$ and
$\mathcal O(m)$, and the cusp has bundle $\mathcal O(mn)$.
The root-forgetting morphism
\[
 r:\mathfrak B_{k,l}\longrightarrow\mathfrak B_{1,1},
 \qquad (x,y)\longmapsto(x^l,y^k),
\]
uses the group homomorphism $\lambda\mapsto\lambda^{kl}$ on
the weighted quotient.  It is a holomorphic morphism of analytic
stacks over $B$, with local maps $s\mapsto s^k$ and
$s\mapsto s^l$.  In particular $r^*\mathcal O(1)=\mathcal O(z)$.
For completeness, the analytic Picard group of this root curve has
the presentation
$\langle H,R_1,R_2\mid mR_1=H,\ nR_2=H\rangle$:
after subtracting the two stabilizer characters, a line bundle
descends to the coarse projective line.  Its only relations are
the two defining root identities.  The Smith form and
$\gcd(m,n)=1$ identify this group with $\mathbb Z$, with
$H=mn\mathcal O(1)$, $R_1=n\mathcal O(1)$, and
$R_2=m\mathcal O(1)$.  This proves the analytic torsion-free
Picard assertion used below.
Translations have identity vertical derivative and act trivially
on structure-sheaf cohomology.  Thus the vertical bundle computed
from the linear period factors is $r^*(\mathcal O\oplus
\mathcal O(-1))$.  The original determinant calculation can also
be repeated: its quotient line $L$ has
$L^{12}=\mathcal O(-p_0)=\mathcal O(-12z)$, so
$L=\mathcal O(-z)$ because the Picard group is torsion-free.
The fixed vector gives a subline $\mathcal O$; its extension
splits because $H^1(\mathcal O(z))=0$.

\par\smallskip\noindent\textbf{Step 3.}\ The cusp normalization resolution has terms
$\mathcal O_{\mathrm{dP}_6}$, three copies of
$\mathcal O_{\mathbb P^1}$, and two point sheaves.
Its global differential is $(a,b,c)\mapsto
(a-b+c,a-b+c)$, giving cohomology $(1,2,1)$; its alternating cup
pairing in degree one is nondegenerate.  The degree-one cocycle
$\gamma$ extends across the cusp because it vanishes on the toric
period lattice.  Its edge values $(0,-1,-1)$ have nonzero cup
product with $(1,0,-1)$.  On finite root charts its class is the
nonzero constant lattice cocycle $\gamma$, invariant under every
deck transformation.  These are explicit comparisons of the
same class, so it gives a nowhere-zero subline in
$G_{\mathfrak B}=R^1\bar f_*\mathcal O_X$.
The base-change and relative-duality argument applies on these charts, using \cite[III.4.12]{BS76} and \cite[pp.~264--266]{RRV71}.  It gives $R^2\bar f_*\mathcal O_X=L$ and
$\det G_{\mathfrak B}=L$.  The extension of $L$ by the specified
$[\gamma]$-subline splits by the same $H^1(\mathcal O(z))=0$.
This proves all direct-image formulas.

\par\smallskip\noindent\textbf{Step 4.}\ The logarithmic tangent sequence is
\[
 0\longrightarrow\mathcal V=\bar f^*\mathcal E_{\rm tan}
 \longrightarrow\mathcal T=T_X(-\log W)
 \longrightarrow\mathcal H=\bar f^*\mathcal O(d-12z)
 \longrightarrow0.
\]
Its last map is the relative differential on smooth root charts
and sends $\sum a_i z_i\partial_{z_i}$ to
$(\sum a_i)q\partial_q$ at the cusp.  It is surjective there.
Since $K_{\mathfrak B}=\mathcal O(-d)$, determinants give
$K_X=\bar f^*\mathcal O(z-d)$.

\par\smallskip\noindent\textbf{Step 5.}\ We give the weighted cohomology calculation explicitly; simply
replacing $1$ by $kl$ in an acyclicity assertion would be incorrect.
The two quotient charts $x\ne0$, $y\ne0$ are finite quotients
of affine lines, and their overlap is a finite quotient of a
punctured line.  Finite averaging and Stein acyclicity compute
coherent cohomology by this two-chart complex.  Its homogeneous
Laurent monomials give
\[
 \begin{split}
 h^0\mathcal O(a)&=\#\{(i,j)\in\mathbb Z_{\ge0}^2:mi+nj=a\},\\
 h^1\mathcal O(a)&=h^0\mathcal O(-a-d).
 \end{split}
\]
The second equality identifies the remaining monomials
$x^{-i-1}y^{-j-1}$; it does not require a positivity theorem.
The bundles $\mathcal O(-z),\mathcal O(-2z)$ are acyclic.
Indeed a possible degree-one monomial would give
$m(i+1)+n(j+1)=az$, with $a=1$ or $2$.
Coprimality forces $i+1=lA$, $j+1=kB$, hence
$3A+4B=a$ with $A,B\ge1$, which is impossible.
Positive degree has no degree-one monomials, so
$H^1\mathcal O(z)=0$ as used above.

\par\smallskip\noindent\textbf{Step 6.}\ Put $b=d-12z$.  Both $b$ and $b-z$ are negative.
Degree-one monomials for $\mathcal O(b)$ would satisfy
$3A+4B=12$, with $A,B\ge1$, and none exist.
For $\mathcal O(b-z)$ the equation is $3A+4B=13$;
its only positive solution is $(A,B)=(3,1)$, producing
$x^{3l-2}y^{k-2}$ in the dual description.  It is allowed
exactly when $k>1$; write $e_k=1$ in that case and $e_k=0$
otherwise.  Projection formula and the two-column Leray sequence
now give
\[
 h^\bullet(\mathcal V)=(1,1,0,0),\qquad
 h^\bullet(\mathcal H)=(0,0,e_k,e_k),\qquad
 h^\bullet(\mathcal T)=(1,1,e_k,e_k).
\]
The producing run reports that \texttt{e\_push\_exact.py} independently checks these
counts for all $1965$ coprime pairs with $k,l\le80$; the monomial
argument proves them for all pairs.

\par\smallskip\noindent\textbf{Step 7.}\ On the normalization $\nu:P_6=\mathrm{dP}_6\to W\subset X$,
the branchwise normal line is $K_{P_6}$: the invariant cusp form
$d\log x_1\wedge d\log x_2\wedge dq$ trivializes $K_X$
on the entire cusp neighbourhood.  Local normal coefficients
therefore give the exact sequence
\[
 0\longrightarrow\mathcal T\longrightarrow T_X
 \longrightarrow\nu_*K_{P_6}\longrightarrow0,
 \qquad h^\bullet(K_{P_6})=(0,0,1).
\]
We must determine its connecting map when $k>1$.
The dual logarithmic sequence, tensored with $K_X$, has horizontal
subline $\bar f^*\mathcal O(13z-2d)$ and vertical quotient
$\bar f^*(\mathcal O(z-d)\oplus\mathcal O(2z-d))$.
The latter has no global sections, by the same equations
$3A+4B=1,2$ above.  Consequently
\[
 H^0(\Omega_X^1(\log W)\otimes K_X)
 =H^0(\mathfrak B,\mathcal O(13z-2d)).
\]
For $k>1$ this is generated by the monomial
$x^{3l-2}y^{k-2}$ multiplying the horizontal logarithmic
differential.  Both $x$ and $y$ are nonzero over $p_0$, so its
residue on $P_6$ is nonzero.  The ordinary logarithmic residue
sequence, tensored with $K_X$, identifies this residue map with
the Serre dual of
$H^2(P_6,K_{P_6})\to H^3(X,\mathcal T)$.
One can check this identification directly in each chart
$q=z_0z_1z_2$: the residues of $dq/q$ are one on each branch,
and contraction with the invariant three-form is the normal
coefficient pairing.  Thus that connecting map has rank $e_k$.
The long exact sequence proves $h^2(T_X)=1$ and $h^3(T_X)=0$;
degrees zero and one are unchanged.

\par\smallskip\noindent\textbf{Step 8.}\ Finally $\partial_c\Pi=e\otimes\gamma$ gives the constant
groupoid cocycle $\lambda\mapsto\gamma(\lambda)e$ on every
regular torus.  On the enlarged finite roots the derivative of the
affine generator is $(r_j/m_j)e$, and the derivative of its
punctured connector is $(\log s_j)/(2\pi i)\,r_je$.
The identity $m_j(r_j/m_j)=\gamma(w_j)$ verifies the power
relation, while $\gamma A_j=\gamma$ verifies conjugation with
every lattice arrow.  Changing $\log s_j$ by $2\pi i N$
adds the derivative of the integral arrow $Nw_j$; the connector
formula is used only on the punctured overlap.  Its central
root-chart representative is the constant $(r_j/m_j)e$,
which is holomorphic.  At the cusp $\gamma$ kills the toric periods,
so the constant deck cocycle is exactly its extending cocycle
above.  Its image under the two isomorphisms
$H^1(\mathcal V)\simeq H^1(\mathcal T)\simeq H^1(T_X)$
is $e\otimes[\gamma]\ne0$.  The chain rule gives the asserted
$u$-class.  The Kuranishi local algebra is
$\mathbb C\{t\}/I$ since $h^1=1$.  The existing family induces
$t\mapsto g(c)$ with $g'(c_0)\ne0$.  The analytic inverse
function theorem makes this substitution injective; every
element of $I$ maps to zero, so $I=0$.  Its classifying map is
then an isomorphism of germs, proving local miniversality.
\end{proof}

\subsection{The algebraic-dimension corollary}\label{sec:algebraic-dimension-nearby}

\begin{corollary}[fam:A.5.2-corollary; T; Algebraic dimension near the standard family]\label{fam:A.5.2-corollary}
Every sufficiently small deformation of $X_c$ has algebraic dimension $1$.
\end{corollary}
\uses{fam:A.5.2; \cite[Theorem 9.1(i)]{AC26}.}

\begin{proof}
By fam:A.5.2, every sufficiently small deformation is induced by the compact period family. Each member has algebraic dimension one by \cite[Theorem 9.1(i)]{AC26}. This proves the stated local corollary.
\end{proof}


% ===== geometry_ends source-preserving fragment =====

% Source-faithful fragments for §§7--11. Generated by build_geometry.py.
\section{Automorphisms and real structures}
\label{sec:real}
The full holomorphic automorphism group is the theorem of
\cite[theorem \texttt{p2:thm:full-automorphisms}]{DGKXv1}; its identity
component and translation action are due to \cite[\S9.7, Proposition~9.23
and Corollary~9.24]{AC26}. Our lattice calculation is a re-derivation.
\begin{cited}[FAM B.4.1; K]\label{fam:B.4.1}
For each member of the retained $(3,4)$ family, $\operatorname{Aut}(X_c)$, including all its components, is $\mathbb C/\mathbb Z\simeq\mathbb C^*$. Its action on the ordinary fibres is $\zeta\mapsto\zeta+be$, and the exponential group isomorphism is $b\mapsto e^{2\pi ib}$.
\end{cited}
\noindent\emph{Uses:} FAM A.1.1, A.1.2, A.2.1, A.3.1 and the cusp valuation comparison in FAM B.2.\par
\noindent The re-derivation is in Appendix~B.\par

\begin{proposition}[FAM B.4.2; T]\label{fam:B.4.2}
Use the normalization $(\beta^{\rm part}+\tau)(p_0)=0$. There is an antiholomorphic diffeomorphism $X_c\to X_{c'}$ exactly when
\[
 c'-(1-\bar c)\in2\mathbb Z,
 \qquad\text{equivalently }w'=-\bar w.
\]
Consequently $\overline{X_c}\simeq X_{1-\bar c}$. The members admitting antiholomorphic self-maps are exactly those with $\operatorname{Re}c\in\tfrac12+\mathbb Z$, equivalently $w=-\bar w$. On each such member the group of holomorphic and antiholomorphic diffeomorphisms is
\[
 \operatorname{Aut}^{\pm}(X_c)\simeq\mathbb C^*\rtimes\mathbb Z/2,
 \qquad z\longmapsto1/\bar z
\]
for the action of the second factor. Its antiholomorphic maps form the nonidentity coset. For all other members the full group is the holomorphic group of \ref{fam:B.4.1}.
\end{proposition}
\noindent\emph{Uses:} FAM A.3.1, A.5.1, B.2.1 and B.4.1; the cusp-normalized period laws.\par

\begin{proof}
\par\noindent\textit{Step 1.} We first construct the comparison before returning its affine data to the canonical choice. Lift conjugation of the real orbifold base to the antiholomorphic involution $\iota$ of its universal cover whose cusp branch has $\tau(\iota z)=-\overline{\tau(z)}$. A lift exists by the universal covering property; it extends at the finite orbifold points, whose local equations are $s_j^{m_j}=t_j$. Choose the usual paths from a real point $t>1$, so reflection fixes the second centre and sends the first chosen lift to $g_2^{-1}z_1$. The induced involution of the deck group is
\[
 \phi(g_1)=g_2^{-1}g_1^{-1}g_2,\qquad
 \phi(g_2)=g_2^{-1},\qquad \phi(g_0)=g_0^{-1}.
\]
These identities also follow by reflecting the lifted peripheral paths; their product and their squares agree with $g_1g_2g_0=1$. The cusp branch fixes the lift uniquely.

\par\noindent\textit{Step 2.} All reflection identities for the periods can be checked before invoking uniqueness. Denote their affine transformation formulas by $G_1,G_2$ on the triple $(\tau,\mu,\beta)$ and put $K(\tau,\mu,\beta)=(-\bar\tau,1-\bar\mu,-\bar\beta)$. Direct rational substitution gives
\[
 KG_1K=G_2^{-1}G_1^{-1}G_2,\qquad KG_2K=G_2^{-1}.
\]
For example $G_2^{-1}(\tau,\mu,\beta)=(-1/\tau,(1-\mu)/\tau,\beta+3-6(1-\mu)^2/\tau)$ and
$G_1^{-1}(\tau,\mu,\beta)=(1/(1-\tau),1-\mu/(1-\tau),\beta-2+6\mu^2/(1-\tau))$. These formulas verify every constant in the two equalities. The reflected $\mu$ is cusp bounded and has the same affine laws, so their difference is a global section of $\mathcal O_B(-p_0)$ and vanishes. After this equality, the difference between the reflected and original $\beta^{\rm part}$ is a holomorphic invariant base function, hence constant; its cusp value is zero under the stated normalization. Therefore
\[
 \mu(\iota z)=1-\bar\mu(z),\qquad
 \beta^{\rm part}(\iota z)=-\overline{\beta^{\rm part}(z)}.
\]
The original $\tau$ identity follows from its $j$ equation and its specified cusp branch. In particular $D_{-\bar c}(\iota z)=D_c(z)$, so this comparison stays in the exact admissible domain.

\par\noindent\textit{Step 3.} Let
\[
 H=\begin{pmatrix}1&0&0&0\\0&1&0&0\\-6&0&-1&0\\0&-1&0&-1\end{pmatrix}.
\]
Then $H^2=I$, $HA_2H=A_2^{-1}$, $HA_1H=A_2^{-1}A_1^{-1}A_2$, and
\[
 \Pi_{-\bar c}(\iota z)H=-\overline{\Pi_c(z)}.
\]
Thus $(z,\zeta)\mapsto(\iota z,-\bar\zeta)$ is a well-defined antiholomorphic torus map. It also descends under the middle deck action. Indeed covariance and the displayed lattice conjugacies imply $R_{\phi(g)}(\iota z)=\overline{R_g(z)}$ by multiplying the full real-invertible period matrices. This proves exact compatibility, not merely an equality of fibre lattices.

\par\noindent\textit{Step 4.} At the first finite centre compose this map with the target deck identification $g_2$ to return the centre $g_2^{-1}z_1$ to $z_1$. Its lattice map becomes $A_2H$ and it conjugates the clockwise generator to the inverse generator. At the second centre its lattice map is $H$ with the same reversal. Consequently the target affine vectors are exactly
\[
 w'_1=-A_2H\varepsilon=-\varepsilon+d,\qquad
 w'_2=-H(-\varepsilon')=\varepsilon'-3d.
\]
The minus signs arise because the inverse of a generator with invariant translation $w'_j/m_j$ has translation $-w'_j/m_j$. This calculation verifies the full affine actions on the root covers. For the overlaps choose the root coordinates of the real local base coordinates $t$ and $t-1$ compatible with these local reflections. This choice has an explicit comparison with any original root coordinate: their ratio is a holomorphic unit, so changing it changes the collar by $(2\pi i)^{-1}\log h\,\Pi w_j$, a translation extending on the entire root family. Its inverse gives a holomorphic isomorphism of the two glued descriptions. Thus the extension property and final classification transfer to the original coordinates. Then $\log s'_j=\overline{\log s_j}$ on paired branches. Under $\zeta\mapsto-\bar\zeta$, conjugation of $1/(2\pi i)$ supplies the compensating sign, and the identity of the two logarithmic collar translations follows exactly. Two roots of the same fixed real coordinate differ by an $m_j$th root of unity; this changes the comparison only by a constant period-coordinate translation, and the additive correction of \ref{fam:A.3.1} glues those constants. The preceding holomorphic-unit comparison handles a different original local base coordinate. Its inverse has the same construction.

\par\noindent\textit{Step 5.} At the cusp choose $s=(2\pi i)^{-1}\log t_c$, so $s\mapsto-\bar s$. The full holomorphic unit matrix satisfies
\[
 C_{-\bar c}(\bar t_c)=-\overline{C_c(t_c)}+
       \begin{pmatrix}6&0\\0&1\end{pmatrix}.
\]
The integral extra matrix has trivial exponential. Hence $(x,t_c)\mapsto(\bar x,\bar t_c)$ intertwines the full cusp deck groups, with the same integral index $\lambda$. The toric fan is defined over $\mathbb Z$, so this map extends antiholomorphically on every toric chart and stratum. It preserves the untwisted cusp collar. We have therefore constructed an antiholomorphic map to the actual target datum with parameter $-\bar c$, $r'_1=-1,k'_1=1,r'_2=1,k'_2=-3,l'_0=0$, whose $p'$ is $+1$.

\par\noindent\textit{Step 6.} Apply \ref{fam:B.2.1} to normalize this target back to the canonical $p=-1$ datum. Take $a=-1$ and the odd shift $n=1$. This gives a holomorphic map to $X_{1-\bar c}$. Every other antiholomorphic map, composed with the inverse of this one, is a holomorphic comparison between canonical members. \ref{fam:A.5.1} gives exactly the stated congruence. Exponentiating it gives $w'=-\bar w$.

\par\noindent\textit{Step 7.} It remains to verify the group extension rather than infer its square from $H^2$. For a real member, choose any resulting antiholomorphic self-map $J$. Its derivative on ordinary fibre coordinates is $+\text{complex conjugation}$: the initial reflection has derivative $-\text{complex conjugation}$, while the normalization just used has linear part $-I_2$. Hence it conjugates translation by $b\in\mathbb C/\mathbb Z$ to translation by $\bar b$. Its square is a holomorphic translation $d$. Commuting $J$ with $J^2$ shows $d=\bar d$ in $\mathbb C/\mathbb Z$. This equality forces the imaginary part of a representative of $d$ to vanish, so take $d\in\mathbb R$. Replacing $J$ by the composition of translation by $-d/2$ with $J$ changes its square to $d-d/2-d/2=0$. We have an antiholomorphic involution. Under $b\mapsto e^{2\pi ib}$, its conjugation action is $z\mapsto1/\bar z$. This proves the split group and specifies its antiholomorphic coset.
\end{proof}

\begin{theorem}[ARI C.2.2; T]\label{ari:C.2.2}
An antiholomorphic equivalence $X_c\to X_{c'}$ exists exactly when
$c'-(1-\bar c)\in2\mathbb Z$. Consequently $X_c$ admits an
antiholomorphic involution exactly when $q_2\in\mathbb R\setminus\{0\}$.
\end{theorem}
\noindent\emph{Uses:} FAM B.4.1, B.4.2 and A.5.1; ARI C.2.1.\par

\begin{proof}
\par\noindent\textit{Step 1.} We check the additional reflection argument, rather than infer a real
structure from the reality of one elliptic invariant. On the period
cover, reflection acts by
\[
 (\tau,\mu,\beta)\longmapsto(-\bar\tau,1-\bar\mu,-\bar\beta).
\]
The theta product and the cusp normalization prove this equality of
continued period functions. For the integral monodromy matrices
$A_1,A_2$ of the retained construction, the lattice comparison is
\[
 H=\begin{pmatrix}1&0&0&0\\0&1&0&0\\-6&0&-1&0\\0&-1&0&-1\end{pmatrix},
 \quad H^2=I,\quad HA_2H=A_2^{-1},\quad
 HA_1H=A_2^{-1}A_1^{-1}A_2,
\]
\[
 \Pi_{-\bar c}(\iota z)H=-\overline{\Pi_c(z)}.
\]
Thus $(z,\zeta)\mapsto(\iota z,-\bar\zeta)$ is an actual torus-family
map. Write $d=(0,0,0,1)^t$, $\varepsilon=(1,2,-4,0)^t$ and
$\varepsilon'=(1,3,-3,0)^t$. Reversing the finite deck generators
changes the target affine vectors to
\[
 v'_1=-\varepsilon+d,\qquad v'_2=\varepsilon'-3d.
\]
The second equality includes the sign of the inverse affine action.
The file \path{Cspecial_reflection.py} checks these matrix identities
over the integers.

\par\noindent\textit{Step 2.} Normalize this reflected datum to the retained vectors
$(\varepsilon,-\varepsilon')$. The linear fibre map is $-I_2$, the
parameter shift is $1$, and the lattice map is $-I+E$,
$E=d(1,0,0,0)$. The differences between target and transported affine
vectors are $2d$ and $-4d$. They lie in the images of the respective
norm operators. Explicit integral correcting cocycles are
$\ell_1=(0,0,-2,0)^t$ and $\ell_2=(0,2,0,0)^t$; the norm equations
are $N_1\ell_1=-2d$ and $N_2\ell_2=4d$. Rational solutions of
$(I-A_j)r_j=(v_j-(-I+E)v'_j)/m_j+\ell_j$ can be chosen with first
coordinate zero: $r_1=(0,-2/3,-2/3,0)^t$ and
$r_2=(0,1,1,0)^t$. These give translations holomorphic across both
root discs. Their cusp cocycle is
$\ell_0=-A_0(\ell_1+A_1\ell_2)=(0,0,4,-2)^t$; its coordinates are
integral, so the corresponding shear and inversion preserve the
periodic triangular cusp fan. They therefore extend across its toric
charts. The differences of the local corrections form a cocycle in
the rank-two vector bundle whose invariant line and quotient are
$\mathcal O$ and $\mathcal O(-1)$ on the base sphere. Both have zero
first cohomology, so the local maps glue. Equivalently one splits the
Laurent series of the cocycle into its nonnegative and negative
parts in these two line-bundle frames. This checks the precise
gluing step of FAM B.4.2 (Statement~\ref{fam:B.4.2}), including the finite affine translations
and cusp extension. The resulting map has target $X_{1-\bar c}$.

\par\noindent\textit{Step 3.} Every other antiholomorphic comparison composed with its inverse is
holomorphic, so Statement~\ref{fam:A.5.1} gives exactly the congruence in the statement.
A self-map exists iff $\operatorname{Re}c\in1/2+\mathbb Z$, iff
$w=-\bar w$, iff $q_2=iw/216$ is real. Its square is a holomorphic
translation $d_0\in\mathbb C/\mathbb Z$, since the automorphism group
is the $\mathbb C^*$ of \cite[theorem \texttt{p2:thm:full-automorphisms}]{DGKXv1}. The constructed reflection acts on
this additive translation coordinate by conjugation. Commutation
with its square gives $d_0=\bar d_0$ in $\mathbb C/\mathbb Z$, hence
$d_0$ has a real representative. Composing the reflection with
translation by $-d_0/2$ makes its square the identity. Necessity and
sufficiency for an involution follow.
\end{proof}

\subsection{Announced real fixed loci}
\begin{definition}[MAPS A.3.setup]\label{maps:A.3.setup}
Here $d=\widehat\delta$, and $\rho_{\rm circ}(b)$ is translation by $be$ for $b\in\mathbb R/\mathbb Z$.
We write $D_1,D_2,D_{12}$ for the three cusp double curves belonging to
edges in directions $e_1,e_2,e_1-e_2$, respectively. Thus $D_2$ is the
fixed sphere of the circle action $\rho_{\rm circ}$, whereas $D_1$ is the double
curve along which the literal character has a limit. This distinction
is essential in the real calculation as well.
\end{definition}
\noindent\emph{Uses:} FAM B.4.1 and the cusp fan.\par

\begin{announcement}[MAPS A.3.1; A]\label{maps:A.3.1}
For a real member the antiholomorphic involutions have exactly two
conjugacy classes under holomorphic automorphisms. For any fixed
involutive splitting $\sigma$ they are represented by $\sigma$ and
$\rho_{\rm circ}(1/2)\sigma$. Every nonempty fixed set is a closed connected
orientable totally real three-manifold with the mod-two homology of
$S^3$. The circle $\rho_{\rm circ}$ preserves each such fixed set and acts
semifreely on it.
\end{announcement}
\noindent\emph{Uses:} FAM B.4.1 and B.4.2; the semifree circle action.\par

The mechanism combines the sign of the coefficient of an involution in the split group $\mathbb C^*\rtimes\mathbb Z/2$ with the fixed-set calculation for the circle action. The statement concerns the smooth fixed manifolds and the two conjugacy classes; the embedded knot type is not determined.

\noindent\emph{Proof:} proof in the repository, MAPS A.3.1 (\url{https://github.com/KokunoYumeto/s6-complex-structure-workbench}).

\begin{definition}[MAPS A.3.representatives]\label{maps:A.3.representatives}
For the explicit representatives, put
\[
 L_F=(2\pi i)^{-1}\log F,\qquad L_G=(2\pi i)^{-1}\log G,
 \qquad F^3=t,\quad G^4=t-1.
 \tag{MAPS A.3.1}
\]
The roots are positive over $t>1$. We continue from that interval
through the upper half $t$-plane. Their logarithms are taken on the
orbifold uniformizing cover. Write $\iota$ for its reflection lift.
Every point of $-i\R_{>0}$ has a representative $c=1/2+iy$, and every
point of $i\R_{>0}$ has a representative $c=-1/2+iy$, with $y$ in the
admissible range. Set
\[
\begin{array}{c|c|c}
 &\alpha&S(t)\\\hline
 c=1/2+iy&1/3&2L_F-4L_G-1/6\\
 c=-1/2+iy&1/2&-2L_G+1/4
\end{array}
\qquad b_c(z)=\alpha\Pi_c(z)\widehat\gamma+S(t(z))e.
\tag{MAPS A.3.2}
\]
Here the fibre vector $e=\Pi_c d=(0,1)^{\mathsf t}$ makes explicit the fibre-coordinate interpretation of $d$ in the source translation formula.
\end{definition}
\noindent\emph{Uses:} FAM B.4.2 and the specified root branches.\par

\begin{announcement}[MAPS A.3.2; A]\label{maps:A.3.2}
On the ordinary period charts the formulas
\[
 \sigma_c[z,\zeta]=[\iota z,\bar\zeta+b_c(\iota z)],
 \qquad \sigma_c^- =\rho_{\rm circ}(1/2)\sigma_c
 \label{maps:eq:A.3.3}
\]
extend to antiholomorphic involutions of the compact threefold.
With these representatives their complete fixed-set data are as follows.
Here $2T^2$ means the disjoint union of two real two-tori; every displayed
circle in a double curve avoids its two triple-point endpoints.
\begin{center}\small
\begin{tabular}{cclccc}
\toprule
$w$&involution&image in $\mathbb{RP}^1$&on $S_1$&on $S_2$&on $W$\\
\midrule
$-ir$, $r>0$&$\sigma_c$&all $\mathbb{RP}^1$&$2T^2$&$2T^2$&$T^2\sqcup S^1_{D_2}\sqcup S^1_{D_{12}}$\\
$-ir$, $r>0$&$\sigma_c^-$&$[1,\infty]$&$\varnothing$&$T^2$&$S^1_{D_1}\sqcup S^1_{D_2}$\\
$ir$, $r>0$&$\sigma_c$&$\varnothing$&$\varnothing$&$\varnothing$&$\varnothing$\\
$ir$, $r>0$&$\sigma_c^-$&$[\infty,1]$ through $0$&$2T^2$&$T^2$&$S^1_{D_1}\sqcup S^1_{D_{12}}$\\
\bottomrule
\end{tabular}
\end{center}
The first interval is the closure of $t>1$; the second is the closure
of $t<1$. On every open real arc contained in the indicated image the
real regular fibre is $2T^2$. No triple point is real for these structures.
\end{announcement}
\noindent\emph{Uses:} MAPS A.3.1 and its explicit representative data; FAM B.4.2.\par

The mechanism is the exact affine fixed-point congruence on the three real base arcs, followed by the finite-root quotient calculation at $0$ and $1$ and the sign calculation for the toric units at the cusp. These are calculations for the stated representatives, including their half-circle translates.

\noindent\emph{Proof:} proof in the repository, MAPS A.3.2 (\url{https://github.com/KokunoYumeto/s6-complex-structure-workbench}).

\begin{announcement}[MAPS A.3.3; A]\label{maps:A.3.3}
Every nonempty fixed set in Theorem \ref{maps:A.3.2} is diffeomorphic to $S^3$.
Its fundamental group is trivial and its integral homology is
$(\Z,0,0,\Z)$. For negative imaginary $w$, the restricted circle action
is the standard semifree action on $S^3$ with fixed circle
$D_2^{\sigma_c}$. For positive imaginary $w$, the nonempty real form
has a free circle action isomorphic, up to circle and orientation
conventions, to the Hopf action.
\end{announcement}
\noindent\emph{Uses:} MAPS A.3.1 and A.3.2.\par

For negative imaginary $w$, the mechanism uses the surface orbit space of the semifree action. For positive imaginary $w$, it uses the two solid-torus caps and the determinant-one comparison of their primitive meridians. No ambient-isotopy classification or Haefliger knot invariant of the embedded three-spheres is asserted.

\noindent\emph{Proof:} proof in the repository, MAPS A.3.3 (\url{https://github.com/KokunoYumeto/s6-complex-structure-workbench}).

\section{The monodromy group and Klingler's question}
\label{sec:monodromy}
The indefinite alternating form and the question about proper discontinuity
are those of \cite[Remark~8.4]{E26}. We retain the actual integral
monodromy of the constructed family.
\begin{definition}[FAM F.setup]\label{fam:F.setup}
We use the clockwise root deck generators and the lattice basis specified
in our conventions.  In this section a group acts properly discontinuously when, for every
compact subset $K$ of its space, only finitely many group elements $g$ satisfy
$gK\cap K\ne\varnothing$.  This definition permits finite stabilizers.  We
write $d=\widehat\delta$, $E=d\otimes\gamma$, and
\[
 J_2=\begin{pmatrix}0&1\\-1&0\end{pmatrix},\qquad
 \Omega=\begin{pmatrix}
 0&0&0&6\\0&0&1&0\\0&-1&0&0\\-6&0&0&0
 \end{pmatrix}.
\]
Thus the alternating form is
$\psi=u_{\Lambda}\wedge w_{\Lambda}+6\gamma\wedge\delta$, an integral alternating form of
type $(1,6)$ on $\Lambda$.  We use $\Omega$ for its matrix and retain $J$ for
a real complex structure.  The domain $\mathcal D$ consists of the real
endomorphisms $J$ of $\Lambda_{\mathbb R}$ satisfying
$J^2=-I$, $J^{\mathsf t}\Omega J=\Omega$, and
$\operatorname{sign}(\Omega J)=(2,2)$.  It is the stated
$\operatorname{Sp}_4(\mathbb R)/U(1,1)$, with its complex structure as an
open subset of the Lagrangian Grassmannian.  An integral symplectic matrix
acts by $J\mapsto gJg^{-1}$.
\end{definition}
\noindent\emph{Uses:} The retained integral marking and monodromy matrices.\par

\subsection{An exact arithmetic description}
\begin{theorem}[FAM F.1.1; T]\label{fam:F.1.1}
Let $P_{d,\mathbb Z}$ be the integral symplectic group fixing $d$ as a
vector, and let $P_{\mathbb Z}$ be the integral symplectic group preserving
$\mathbb Zd\subset\ker\gamma\subset\Lambda$.  Every element of
$P_{d,\mathbb Z}$ has a unique expression
\begin{equation}\label{fam:eq:F-g}
 g(B,x,t)=
 \begin{pmatrix}
 1&0&0\\6x&B&0\\t&-x^{\mathsf t}J_2B&1
 \end{pmatrix},\qquad
 B\in\operatorname{SL}_2(\mathbb Z),\quad x\in\mathbb Z^2,\quad t\in\mathbb Z,
\end{equation}
where the block sizes are $1,2,1$.  Put
\[
 S=\begin{pmatrix}0&-1\\1&0\end{pmatrix},\quad
 T=\begin{pmatrix}1&1\\0&1\end{pmatrix},\quad
 a=\begin{pmatrix}0&1\\-1&-1\end{pmatrix}.
\]
There is a homomorphism
$\alpha:\operatorname{SL}_2(\mathbb Z)\to\mathbb Z/12$ determined by
$\alpha(S)=9$ and $\alpha(T)=1$; in particular $\alpha(a)=4$.
With $Q(x)=x_1+x_2+x_1x_2$, the actual monodromy group is
\begin{equation}\label{fam:eq:F-Gamma}
 \Gamma=\left\{g(B,x,t):t+6Q(x)-\alpha(B)=0\ \text{in }\mathbb Z/12\right\}.
\end{equation}
Consequently $[P_{d,\mathbb Z}:\Gamma]=12$ and
$P_{\mathbb Z}=\{\pm I\}\times P_{d,\mathbb Z}$, so
$[P_{\mathbb Z}:\Gamma]=24$.  Its image in the Levi quotient is exactly
$\operatorname{SL}_2(\mathbb Z)$, with scalar on $d$ equal to $1$.
\end{theorem}
\noindent\emph{Uses:} The matrices $A_1,A_2$ and the integral symplectic form.\par

\begin{proof}
\par\noindent\textit{Step 1.} The matrices in the fixed conventions, in the fixed basis, are
\[
 A_1=\begin{pmatrix}1&0&0&0\\6&0&1&0\\-6&-1&-1&0\\-2&1&0&1\end{pmatrix},
 \qquad
 A_2=\begin{pmatrix}1&0&0&0\\0&0&-1&0\\-6&1&0&0\\3&0&1&1\end{pmatrix}.
\]
Multiplication gives $A_1^3=A_2^4=I$ and
$A_j^{\mathsf t}\Omega A_j=\Omega$.  In particular
$A_1=g(a,(1,-1)^{\mathsf t},-2)$ and
$A_2=g(S,(0,-1)^{\mathsf t},3)$.  Their images $a,S$ generate
$\operatorname{SL}_2(\mathbb Z)$ because $Sa=T$ and $S,T$ perform the
Euclidean algorithm on primitive integral columns.

\par\noindent\textit{Step 2.} A matrix fixing $d$ and preserving $\Omega$ preserves
$d^{\perp}=\ker\gamma$.  Its first row is $(1,0,0,0)$, its last column is
$d$, its middle block $B$ is symplectic on
$d^{\perp}/\mathbb Rd$, and its remaining middle column, say $v$, satisfies
$6r=-v^{\mathsf t}J_2B$.  Integrality of $r$ and unimodularity of $B$
give $v=6x$ with $x\in\mathbb Z^2$.  There is no further condition on
the lower left entry $t$.  This proves \eqref{fam:eq:F-g}.  Direct multiplication
gives
\begin{equation}\label{fam:eq:F-law}
 g(B,x,t)g(C,y,s)
 =g(BC,x+By,t+s-6x^{\mathsf t}J_2By).
\end{equation}

\par\noindent\textit{Step 3.} For completeness, the presentation of the Levi group needed here follows
from an elementary action on the real projective line.  The images of $a$
and $S$ in $\operatorname{PSL}_2(\mathbb Z)$ have orders $3$ and $2$.
On positive real numbers, $a(z)=-1/(z+1)$ has image $(-1,0)$ and
$a^2(z)=-(z+1)/z$ has image $(-\infty,-1)$; $S$ carries negative real
numbers to positive real numbers.  A cyclically reduced alternating word
starting with a nonidentity power of $a$ and ending with $S$ therefore
carries $(-\infty,0)$ into one of these two proper subintervals.  It cannot
be the identity.  Every reduced word is conjugate to such a word or to a
nonidentity element of one factor.  Together with the Euclidean algorithm,
this proves the presentation
$\operatorname{PSL}_2(\mathbb Z)=\langle\bar a,\bar S\mid
\bar a^3=\bar S^2=1\rangle$.  It follows that
\begin{equation}\label{fam:eq:F-SLpresentation}
 \operatorname{SL}_2(\mathbb Z)=
 \langle a,S\mid a^3=S^4=1,\ [a,S^2]=1\rangle.
\end{equation}
Indeed the group on the right modulo its central element $S^2$ is the
just-proved free product, and $S^2$ maps to the nonidentity matrix $-I$.
Its kernel in the displayed matrix group is consequently trivial.
The values $4,9$ respect all these relations, proving the assertion about
$\alpha$; $T=Sa$ has value $1$.

\par\noindent\textit{Step 4.} The quadratic function $Q$ modulo $2$ is zero at $(0,0)$ and one at all
three other points of $\mathbb F_2^2$.  Hence $Q(Bx)=Q(x)\pmod2$ for every
$B\in\operatorname{SL}_2(\mathbb Z)$.  Also
\[
 Q(x+y)-Q(x)-Q(y)\equiv x^{\mathsf t}J_2y\pmod2.
\]
Equation \eqref{fam:eq:F-law} now shows that
$h(g(B,x,t))=t+6Q(x)\pmod{12}$ is a homomorphism.
The homomorphism $\chi=h-\alpha\circ B$ is surjective, since
$\chi(g(I,0,1))=1$.  Its values on $A_1,A_2$ are zero.

\par\noindent\textit{Step 5.} We verify equality with its kernel; containment alone would not determine
the group.  Let $N_{\mathbb Z}$ be the kernel of the map to the middle
block.  The kernel of
$\langle g_1,g_2\mid g_1^3=g_2^4=1\rangle\to
\operatorname{SL}_2(\mathbb Z)$ is normally generated by
$[g_1,g_2^2]$, by \eqref{fam:eq:F-SLpresentation}.  Its image in $\Gamma$ is
therefore normally generated by
\[
 K=[A_1,A_2^2]=g(I,(0,-1)^{\mathsf t},6).
\]
Put
\[
 U=A_2KA_2^{-1}=g(I,(1,0)^{\mathsf t},-6),\quad
 V=K^{-1}=g(I,(0,1)^{\mathsf t},-6),\quad Z=I+12E.
\]
Equation \eqref{fam:eq:F-law} gives $[U,V]=Z^{-1}$.  Thus $U,V,Z$ generate
\begin{equation}\label{fam:eq:F-H}
 H=\{g(I,x,t):t\in6\mathbb Z,\quad t/6\equiv Q(x)\pmod2\}.
\end{equation}
More explicitly,
$U^rV^sZ^k=g(I,(r,s)^{\mathsf t},-6(r+s+rs)+12k)$, so these expressions
give every member of the set.  Conjugation by $g(B,x,t)$ sends
$g(I,y,s)$ to $g(I,By,s-12x^{\mathsf t}J_2By)$; the invariance of $Q$
modulo $2$ proves that $H$ is normal in $P_{d,\mathbb Z}$.  The normal
closure of $K$ is contained in $H$ and contains $U,V,Z$, so it is exactly
$H$.  We have proved $\Gamma\cap N_{\mathbb Z}=H$.  Both $\Gamma$ and
$\ker\chi$ surject onto the same Levi group and have intersection $H$
with $N_{\mathbb Z}$.  Given $g\in\ker\chi$, choose $\gamma\in\Gamma$
with the same middle block.  Then $g\gamma^{-1}\in H$, and $g\in\Gamma$.
This proves \eqref{fam:eq:F-Gamma} and the index $12$.

\par\noindent\textit{Step 6.} Finally an integral automorphism preserving the primitive line
$\mathbb Zd$ sends $d$ to $d$ or $-d$.  Multiplication by the central
matrix $-I$ converts the latter case to the former.  The intersection of
$\{\pm I\}$ with $P_{d,\mathbb Z}$ is trivial.  This proves the direct
product and the index $24$.
\end{proof}

\begin{proposition}[FAM F.1.2; T]\label{fam:F.1.2}
The group in \eqref{fam:eq:F-Gamma} has center
$\langle Z\rangle\cong\mathbb Z$ and is the following explicit central
extension of $\mathbb Z^2\rtimes\operatorname{SL}_2(\mathbb Z)$ by
$\mathbb Z$.  Choose the representative $a(B)\in\{0,\ldots,11\}$ of
$\alpha(B)$.  Its underlying set consists of triples $(B,x,k)$, with
product
\[
 (B,x,k)(C,y,l)=(BC,x+By,k+l+\kappa((B,x),(C,y))),
\]
where
\begin{align*}
12\kappa((B,x),(C,y))={}&a(B)+a(C)-a(BC)-6Q(x)-6Q(y)\\
 &-6x^{\mathsf t}J_2By+6Q(x+By).
\end{align*}
Its real and rational Zariski closures are respectively the full real
and rational symplectic stabilizers of the vector $d$.  They have
dimension $6$, with three-dimensional Heisenberg unipotent radical and
Levi quotient $\operatorname{SL}_2$.
\end{proposition}
\noindent\emph{Uses:} FAM F.1.1 and its exact Heisenberg kernel.\par

\begin{proof}
\par\noindent\textit{Step 1.} The numerator defining $\kappa$ is divisible by $12$, by the two
congruences proved above.  The bijection
\[
 (B,x,k)\longmapsto g(B,x,a(B)-6Q(x)+12k)
\]
intertwines the displayed multiplication with matrix multiplication by
\eqref{fam:eq:F-law}.  This proves associativity as well as the claimed abstract
description, without leaving an unspecified extension class.  The kernel
of the homomorphism to $\mathbb Z^2\rtimes\operatorname{SL}_2(\mathbb Z)$
is $\langle Z\rangle$.  That affine group has trivial center: commuting
with all translations forces its linear part to be $I$, and commuting
with $S,T$ forces its translation part to vanish.  Thus the center of
$\Gamma$ is exactly the stated kernel.

\par\noindent\textit{Step 2.} The logarithms $U-I,V-I,Z-I$ are linearly independent elements of the
three-dimensional Heisenberg Lie algebra.  The square of each is zero.
A polynomial vanishing on all integral powers $I+n(U-I)$ vanishes on
the whole one-parameter line; the same holds for $V$ and $Z$.  Hence the
Zariski closure contains the full unipotent radical.  Its Levi image
contains $S,T$; the powers of $T$ and of $STS^{-1}$ similarly supply the
two opposite unipotent subgroups of $\operatorname{SL}_2$.  Their Lie
algebras and their bracket span $\mathfrak{sl}_2$.  They generate the
connected algebraic group $\operatorname{SL}_2$.  A subgroup containing
the entire unipotent kernel and surjecting onto this Levi quotient is
the entire vector stabilizer.  Its dimension is $3+3=6$.  This reasoning
uses equations over $\mathbb Q$, so proves the rational assertion as
well as the real assertion.
\end{proof}

\subsection{Properness and the maximal open set}
\begin{definition}[FAM F.2.setup]\label{fam:F.2.setup}
Define the real-analytic function
\[
 q_d:\mathcal D\longrightarrow\mathbb R,\qquad q_d(J)=\psi(d,Jd),
 \qquad\mathcal O=\{J:q_d(J)\ne0\}.
\]
The function $q_d$ is invariant under $P_{\mathbb Z}$: replacing $d$ by
$-d$ changes both of its arguments.  Write $\mathcal O_+$ and
$\mathcal O_-$ for its two sign regions.
\end{definition}
\noindent\emph{Uses:} FAM F.setup.\par

\begin{proposition}[FAM F.2.1; T]\label{fam:F.2.1}
The assertion that $\Gamma$ acts properly discontinuously on all of
$\mathcal D$ is false.  The following point of $\mathcal D$ has an
explicit infinite stabilizer:
\[
 J_0=\begin{pmatrix}
 0&0&-1&0\\0&0&0&-6\\1&0&0&0\\0&1/6&0&0
 \end{pmatrix},\qquad
 R=\begin{pmatrix}
 1&0&0&0\\0&1&72&0\\0&0&1&0\\-12&0&0&1
 \end{pmatrix}\in\Gamma.
\]
Every power of $R$ fixes $J_0$, and all these powers are distinct.
\end{proposition}
\noindent\emph{Uses:} FAM F.1.1 and the explicit matrices $J_0,R$.\par

\begin{proof}
Direct multiplication gives $J_0^2=-I$ and
$J_0^{\mathsf t}\Omega J_0=\Omega$.  Moreover
\[
 \Omega J_0=\begin{pmatrix}
 0&1&0&0\\1&0&0&0\\0&0&0&6\\0&0&6&0
 \end{pmatrix}
\]
has eigenvalues $1,-1,6,-6$, so $J_0\in\mathcal D$ and $q_d(J_0)=0$.
The matrix $R$ equals $g(T^{72},0,-12)$ and belongs to \eqref{fam:eq:F-Gamma},
because $\alpha(T^{72})=0\pmod{12}$.  Its commutator with $J_0$ is zero.
Writing $R=I+N$ gives $N\ne0$ and $N^2=0$, so $R^m=I+mN$ and the
powers are distinct.  The compact singleton $\{J_0\}$ therefore violates
the stated definition of proper discontinuity.  These identities are
checked both with SymPy and with a separate rational matrix
implementation in \texttt{code/f\_exact.py}.
\end{proof}

\begin{theorem}[FAM F.2.2; T]\label{fam:F.2.2}
The set $\mathcal O$ is the largest $\Gamma$-invariant open subset of
$\mathcal D$ on which the action is properly discontinuous.  It is also
the largest such open subset for $P_{\mathbb Z}$.  Each sign region is
a complex manifold in the analytic topology, and its quotient by either
discrete group is a Hausdorff complex analytic orbifold.
\end{theorem}
\noindent\emph{Uses:} FAM F.1.1, F.1.2 and F.2.1; the homogeneous-space stabilizer calculation.\par

\begin{proof}
\par\noindent\textit{Step 1.} Let $P_{\mathbb R}$ be the real symplectic group preserving the real line
$\mathbb Rd$.  We first prove properness on each sign region.  For a fixed
$J$, the form
$h_J(v,w)=\psi(v,Jw)+i\psi(v,w)$ is a nondegenerate Hermitian form of
signature $(1,1)$ on the complex vector space $(\Lambda_{\mathbb R},J)$,
with the convention conjugate linear in the first entry.  When
$q_d(J)\ne0$, the complex line $\mathbb Cd$ is nondegenerate and its
Hermitian orthogonal complement is a complex line of the opposite sign.
Normalize $d$ by a positive real scalar and choose a unit vector in this
complement.  Comparing these two Hermitian bases for two complex
structures with the same sign gives a real symplectic map preserving
$\mathbb Rd$ and conjugating one complex structure to the other.  Thus
$P_{\mathbb R}$ is transitive on each sign region.

\par\noindent\textit{Step 2.} An element of the stabilizer of $J$ inside $P_{\mathbb R}$ is complex
linear and sends $d$ to $r d$ for a real nonzero $r$.  Preservation of
$q_d(J)$ gives $r^2=1$.  On the complementary complex line it acts by an
element of $U(1)$.  The stabilizer is consequently the compact group
$\{\pm1\}\times U(1)$.  Each sign region is therefore a homogeneous
space of $P_{\mathbb R}$ with compact stabilizer $L$.  Local sections of
$P_{\mathbb R}\to P_{\mathbb R}/L$, and a finite cover of a compact
subset $K$, give a compact subset $C$ of $P_{\mathbb R}$ mapping onto
$K$.  Every element carrying a point of $K$ back into $K$ belongs to
$CLC^{-1}$, which is compact.  This proves properness.  Restricting a
proper action to a closed discrete subgroup gives
the required finite-intersection condition.  Both $\Gamma$ and
$P_{\mathbb Z}$ are closed discrete matrix groups.  The two sign regions
are disjoint invariant open-and-closed subsets of $\mathcal O$, so the
same finite-intersection condition holds on their union.  Proper actions
by discrete groups have finite stabilizers and Hausdorff quotients; their
holomorphic local finite-group quotient charts supply the asserted
analytic orbifolds.

\par\noindent\textit{Step 3.} It remains to prove maximality, which requires more than one bad point.
The real vector stabilizer $P_{d,\mathbb R}$ acts transitively on
$\mathcal D_0=\{J:q_d(J)=0\}$.  Here is a direct comparison map.  In the
Hermitian space just described, $d$ is a nonzero null vector.  Choose
$v$ with $h_J(d,v)=1$, and replace $v$ by
$v-\tfrac12h_J(v,v)d$.  This preserves $h_J(d,v)=1$ and makes $v$ null.
The ordered complex basis $(d,v)$ has Gram matrix
$\bigl(\begin{smallmatrix}0&1\\1&0\end{smallmatrix}\bigr)$.
The complex-linear map between two such bases fixes $d$ and preserves
$h$, hence its imaginary part $\psi$.  As a real map it belongs to
$P_{d,\mathbb R}$ and conjugates the two complex structures.

\par\noindent\textit{Step 4.} The subgroup $P_{d,\mathbb Q}$ is dense in $P_{d,\mathbb R}$.
For the unipotent coordinates in \eqref{fam:eq:F-g} this is density of
$\mathbb Q^3$ in $\mathbb R^3$.  For $\operatorname{SL}_2$ one may
approximate $a,b,c$ by rational numbers in the chart $a\ne0$ and set
$d_{\rm aux}=(1+bc)/a$; translates of that chart cover the group.  Hence
\[
 \{rJ_0r^{-1}:r\in P_{d,\mathbb Q}\}
\]
is dense in $\mathcal D_0$.
Put $R_0=g(T^6,0,-1)$, so $R_0^{12}=R$ and
$(R_0-I)^2=0$.  At $rJ_0r^{-1}$ the nonzero rational nilpotent
$L=r(R_0-I)r^{-1}$ commutes with the complex structure and satisfies
$L^2=0$.  Choose a positive integer $d_{\rm aux}$ clearing its denominators.
The matrix $I+dL$ and its inverse $I-dL$ are integral, symplectic, and
fix $d$; this follows either by conjugating the identity
$(I+s(R_0-I))^{\mathsf t}\Omega(I+s(R_0-I))=\Omega$ or by expanding it.
The twelfth power of $I+dL$ lies in $\Gamma$, because
$P_{d,\mathbb Z}/\Gamma\cong\mathbb Z/12$.  It has infinite order and
fixes $rJ_0r^{-1}$.  Thus points with infinite $\Gamma$-stabilizer are
dense in $\mathcal D_0$.

\par\noindent\textit{Step 5.} Every open neighborhood in $\mathcal D$ of any point of $\mathcal D_0$
therefore contains a point with infinite stabilizer.  An open set on
which the action is proper cannot meet $\mathcal D_0$.  This proves the
largest-open-set assertion for $\Gamma$.  The same points have infinite
$P_{\mathbb Z}$-stabilizer, so it proves that assertion for the full
parabolic as well.
\end{proof}

\begin{proposition}[FAM F.2.3; T]\label{fam:F.2.3}
The positive sign region admits the global period chart
\[
 \mathcal O_+=\left\{
 (\tau,\mu,\beta)\in\mathbb C^3:
 \operatorname{Im}\tau>0,\quad
 \operatorname{Im}\beta-\frac{6(\operatorname{Im}\mu)^2}
 {\operatorname{Im}\tau}<0\right\},
\]
represented by the matrix $\Pi$ in the question.  In this chart,
\begin{equation}\label{fam:eq:F-qD}
 q_d(J)=-\frac6D.
\end{equation}
Every admissible period $\Pi_c(z)$ belongs to $\mathcal O_+$, for every
admissible $c$.  Thus this single maximal proper-action open set contains
all the period images in question.
\end{proposition}
\noindent\emph{Uses:} The global Siegel chart and the signature condition.\par

\begin{proof}
\par\noindent\textit{Step 1.} Let $J\in\mathcal O_+$.  The two vectors $\widehat w,d$ are complex-linearly
independent in $(\Lambda_{\mathbb R},J)$.  Otherwise
$\widehat w=r d+sJd$ for real $r,s$.  Pairing with $d$ gives
$0=sq_d(J)$; hence $s=0$, contradicting the real linear independence
of $\widehat w,d$.  There is consequently a unique complex-linear
isomorphism $(\Lambda_{\mathbb R},J)\to\mathbb C^2$ sending these two
vectors to $(1,0)$ and $(0,1)$.  The corresponding period matrix has
right block $I$.  The equation
$\Pi\Omega^{-1}\Pi^{\mathsf t}=0$ says that its upper left entry is
six times its lower middle entry, giving exactly the displayed matrix.
This construction is the ordinary holomorphic Grassmannian chart,
so its coordinates and inverse are holomorphic.

\par\noindent\textit{Step 2.} In the real symplectic basis
$(\widehat u,\widehat\gamma/6,\widehat w,d)$ the normalized symmetric
period matrix is
$\bigl(\begin{smallmatrix}\tau&\mu\\\mu&\beta/6\end{smallmatrix}\bigr)$.
Its imaginary part has signature $(1,1)$, equivalently
$6(\operatorname{Im}\mu)^2-
\operatorname{Im}\tau\operatorname{Im}\beta>0$.
Solving $\Pi(Jd)=i\Pi d=(0,i)$ in real coefficients gives
\[
 \gamma(Jd)=
 \frac{\operatorname{Im}\tau}
 {\operatorname{Im}\tau\operatorname{Im}\beta-
 6(\operatorname{Im}\mu)^2}.
\]
Since $\psi(d,v)=-6\gamma(v)$, positivity of $q_d$ now gives
$\operatorname{Im}\tau>0$ and $D<0$, and proves \eqref{fam:eq:F-qD}.
Conversely these inequalities make the imaginary period matrix
invertible with the stated signature and give $q_d>0$.  The admissibility
inequality in the question is $D_c<0$ everywhere, which proves the last
assertion by this exact comparison map.
\end{proof}

\subsection{The curve in the proper quotient}
\begin{theorem}[FAM F.3.1; T]\label{fam:F.3.1}
For every admissible $c$, the marked period map on the covering of
$B^\circ$ is a holomorphic immersion into $\mathcal O_+$.  The induced
map
\[
 p_c:B^\circ\longrightarrow\Gamma\backslash\mathcal O_+
\]
is injective and a holomorphic immersion.  It extends over $t=0,1$ to a
proper holomorphic embedding
\[
 \bar p_c:\mathbb C=B\setminus\{\infty\}
 \lhook\joinrel\longrightarrow\Gamma\backslash\mathcal O_+.
\]
The closure of $p_c(B^\circ)$ in this quotient is exactly
$\bar p_c(\mathbb C)$.
\end{theorem}
\noindent\emph{Uses:} FAM F.2.2 and F.2.3; the finite-root and cusp period expansions.\par

\begin{proof}
\par\noindent\textit{Step 1.} For $B=\bigl(\begin{smallmatrix}a&b\\c&d_{\rm aux}\end{smallmatrix}\bigr)$, changing
the marking by $g(B,x,t)$ replaces $\Pi$ by $\Pi g(B,x,t)^{-1}$ and then
renormalizes its right block.  On the quotient by the complex line of
$d$, the two periods are the row $(\tau,1)$.  Right multiplication by
$B^{-1}$ sends it to $(d_{\rm aux}\tau-c,a-b\tau)$.  Hence
\begin{equation}\label{fam:eq:F-tauaction}
 \tau\longmapsto\frac{d_{\rm aux}\tau-c}{a-b\tau}.
\end{equation}
The denominator never vanishes in the upper half plane.  The function
$j(\tau)$ is therefore a holomorphic invariant of this action and
descends to a holomorphic function $j_\Gamma$ on the proper quotient.
Here $j$ is the elliptic invariant, normalized by
$j(E)=1728\,4A^3/(4A^3+27B^2)$ for $y^2=x^3+Ax+B$; changing a lattice
basis preserves the isomorphism of elliptic curves and hence this
function.

\par\noindent\textit{Step 2.} We derive the required value in the present family.  The elliptic curve
whose period is $\tau$ has the equation stipulated in the question,
\[
 y^2=x^3-3t^3(t-1)x+2t^4(t-1)^2.
\]
Its discriminant is $1728t^8(t-1)^3$ and direct cancellation gives
\begin{equation}\label{fam:eq:F-j}
 j(\tau)=1728t\qquad(t\in B^\circ).
\end{equation}
Equivalently this is the defining lift of the elliptic modular map used
for $\tau$ in the period construction.  The computation is reproduced
exactly in \texttt{f\_exact.py}.  In a local base coordinate $t$, the
derivative of the composite $j\circ\tau$ equals $1728$.  Thus the
derivative of the marked period map is nonzero.  The same identity after
passage to the quotient shows both injectivity and injectivity of the
tangent map of $p_c$: $j_\Gamma/1728$ is a holomorphic left inverse.

\par\noindent\textit{Step 3.} At the two elliptic points, the period functions are holomorphic in the
cyclic covering coordinates $s_1,s_2$, with $t=s_1^3$ and
$t-1=s_2^4$ after choosing local coordinates as in the question.  Their
values are nondegenerate because $D_c<0$ there.  The maps from these
discs to $\mathcal O_+$ are equivariant for $A_1,A_2$.  Composing with
the quotient makes them invariant, so they descend to holomorphic maps
from the base discs.  This is the finite-quotient comparison map: a
holomorphic invariant germ on a disc is a convergent series in $s_j^{m_j}$,
which proves descent on local function rings.  The extensions agree
with $p_c$ off their centers and therefore with each other.  Identity
\eqref{fam:eq:F-j} extends by holomorphy, so the resulting $\bar p_c$ still
has holomorphic left inverse $j_\Gamma/1728$.

\par\noindent\textit{Step 4.} The quotient is Hausdorff by \ref{fam:F.2.2}.  For a compact subset $K$ of it,
$\bar p_c^{-1}(K)$ is closed in $\mathbb C$ and lies in the compact set
$j_\Gamma(K)/1728$.  Thus $\bar p_c$ is proper.  A holomorphic section of
a holomorphic map is locally a closed immersion: in local embeddings
into polydiscs, its coordinate along the base is $t$ and the remaining
coordinates are its holomorphic component functions, whose graph
defines the image.  Applying this to the left inverse shows that
$\bar p_c$ is a holomorphic embedding; properness makes its image
closed.  Since $\mathbb C\setminus\{0,1\}$ is dense in $\mathbb C$, its
image has the asserted closure.
\end{proof}

\begin{proposition}[FAM F.3.2; T]\label{fam:F.3.2}
At the two finite punctures the period map has the holomorphic
finite-cover limits
\[
 \Pi_1=
 \begin{pmatrix}2(2-\rho)&\rho&1&0\\
 \beta_c(z_1)&(2-\rho)/3&0&1\end{pmatrix},\qquad
 \Pi_2=
 \begin{pmatrix}3(1-i)&i&1&0\\
 \beta_c(z_2)&(1-i)/2&0&1\end{pmatrix},
 \quad \rho=e^{\pi i/3}.
\]
Their monodromies are the given order-three and order-four matrices.
At the cusp the curve leaves every compact subset of
$\Gamma\backslash\mathcal O_+$.  On the prescribed cusp branch its
Hodge plane in the compact dual tends to
$\mathbb C\gamma+\mathbb Cu$.  The latter plane has complex dimension
two in common with its conjugate and therefore does not represent a
point of $\mathcal D$.
\end{proposition}
\noindent\emph{Uses:} The finite special values and the cusp-normalized period matrix.\par

\begin{proof}
The finite values of $\tau,\mu$ follow by evaluating their transformation
laws at the fixed points: $\tau(z_1)=\rho$,
$\mu(z_1)=(2-\rho)/3$, $\tau(z_2)=i$,
$\mu(z_2)=(1-i)/2$.  Substitution in $\Pi$ gives the two matrices.
The finite covering comparison and its descent were proved in \ref{fam:F.3.1}.
The unipotent cusp transformation is
$\tau\mapsto\tau-1$, $\mu\mapsto\mu$, $\beta\mapsto\beta+1$.
On the cusp branch, $\operatorname{Im}\tau\to+\infty$,
$\mu\to\mu_0$, and $\beta_c+\tau\to b_0+c$, from the holomorphic
cusp expansions.  Consequently
\[
 D_c=-\operatorname{Im}\tau+O(1),\qquad
 q_d(\Pi_c)\longrightarrow0^+,
 \qquad j_\Gamma(p_c(t))=1728t\longrightarrow\infty.
\]
The last equality alone excludes a convergent subnet in any compact
subset of the proper quotient; the first two identities give the
position relative to the maximal proper-action open set.
The rows of $\Pi_c$, regarded as covectors, are
$6\mu\gamma+\tau u+w$ and $\beta_c\gamma+\mu u+\delta$.
Dividing each by $\tau$ gives limits $u$ and $-\gamma$.
Their exterior product has a nonzero limit, so their span converges
in the Grassmannian to the stated plane.  It is spanned by real
covectors, hence equals its conjugate.  A complex structure has
$H^{1,0}\cap H^{0,1}=0$; here the intersection is the entire
two-dimensional plane.  This explicit comparison proves the assertion
about the compact-dual limit.
\end{proof}

\begin{proposition}[FAM F.3.3; T]\label{fam:F.3.3}
The determinant used in the workbench is generically nonzero, with
the stated zero at $z_1$:
\[
 \det\dot Z=6\dot\mu^2-\dot\tau\dot\beta,\qquad
 \det\dot Z(z_1)=0.
\]
The immersion established in \ref{fam:F.3.1} holds at every point of $B^\circ$
without requiring this determinant to be everywhere nonzero.
\end{proposition}
\noindent\emph{Uses:} The derivative transformation laws and the finite special values.\par

\begin{proof}
In the cusp coordinate $s$, the expansion
$Z=s\bigl(\begin{smallmatrix}0&1\\-1&0\end{smallmatrix}\bigr)+C(t_c)$,
$t_c=e^{2\pi is}$, gives
$\det(dZ/ds)=1+O(t_c)$, so the determinant is not the zero function.
At $z_1$ the fixed-point transformation law of $\mu$ gives
$\dot\mu=-\dot\tau/3$.  One may obtain this by differentiating
$\mu\circ g_1=(1-\mu)/\tau$, using
$g_1'(z_1)=-\rho$ and $\rho^2-\rho+1=0$.
The law
$\beta\circ g_1=\beta+2-6(1-\mu)^2/\tau$ then gives
$\dot\beta=2\dot\tau/3$.  Substitution yields
$6\dot\tau^2/9-2\dot\tau^2/3=0$.
The derivative of a map from a curve to the three-dimensional period
chart is a vector of three coordinate derivatives.  Identity
\eqref{fam:eq:F-j} proves that this vector is nonzero over $B^\circ$;
the determinant tests the rank of the associated two-by-two
Kodaira--Spencer matrix.  The maps compared here are explicitly the
coordinate derivative and its determinant, and the displayed computation
states exactly what each supplies.
\end{proof}

\section{Special fibres, curves and surfaces}
\label{sec:fibres}
The N\'eron--Severi equation below is already
\cite[\S9.1, Lemma~9.2]{AC26}; the accompanying independence and
indefiniteness calculations occur in its Lemmas~9.3--9.4. Its
Proposition~9.5 and Corollary~9.8 give the very-general conclusion
$\operatorname{NS}=\mathbb Z\eta$ and algebraic dimension zero.
Those parts are credited to that source, with a derivation retained here
under the binding T status. The explicit quotient loci, examples and
exhaustive descriptions below have their distinct source statements.
\subsection{Elliptic quotients and the torsion locus}
\begin{proposition}[FAM D.1.1; T]\label{fam:D.1.1}
The specified linear projection \(\operatorname{pr}_1:\mathbb C^2\to
\mathbb C\) induces a holomorphic group homomorphism
\(F_t\to E_t=\mathbb C/(\mathbb Z\tau+\mathbb Z)\) exactly on
\(6\mu=A\tau+B\), \(A,B\in\mathbb Z\). It is then surjective with
connected kernel
\(\mathbb C/(\mathbb Z+(A\mu-\beta)\mathbb Z)\).
On the smooth base the two values of \(t\) are precisely
\[
 t_\pm=\frac{64(47\pm27\sqrt3)}{11}.
\]
\end{proposition}
\noindent\emph{Uses:} The fixed period matrix and the exact division-polynomial computation.\par

\begin{proof}
\par\noindent\textit{Step 1.} The images of the four lattice generators under the linear projection
are \(6\mu,\tau,1,0\). This proves descent in both directions.
The integral basis change replacing \(\widehat\gamma\) with
\(\widehat\gamma-A\widehat u-B\widehat w\) changes its period to
\((0,\beta-A\mu)\); the remaining vertical generator is \((0,1)\).
The other two periods map onto the entire target lattice, so the kernel
is connected. Its second modulus has imaginary part
\(A\operatorname{Im}\mu-\operatorname{Im}\beta=-D>0\).

\par\noindent\textit{Step 2.} The elliptic model may be scaled by the explicit isomorphism
\((X,Y)\mapsto(2X,2\sqrt2Y)\) to
\[
 Y^2=X^3-\frac34t^3(t-1)X+\frac14t^4(t-1)^2,
 \qquad P=(0,t^2(t-1)/2).
\]
Its sixth division polynomial at \(P\) is
\[
 \psi_6(P)=\frac9{65536}t^{24}(t-1)^9
 (11t^2-6016t+8192).
\]
For a smooth elliptic fibre, \(\psi_n(P)=0\) is exactly the condition
\(nP=O\), since the multiplication rational map has these and only these
poles in its affine coordinate. The two roots displayed above are simple
and avoid \(0,1\). We computed the same equation independently by the
chord-and-tangent addition algorithm, obtaining
\[
 \begin{aligned}
 x(6P)&=\frac{(89t^2+4736t-4096)Q(t)}
 {1296(11t^2-6016t+8192)^2},\\
 Q(t)&=881t^4+4745344t^3-12079104t^2\\
 &\hspace{2em}+24641536t-16777216.
 \end{aligned}
\]
The division-polynomial recurrence and the addition algorithm are in
separate supplied programs; all polynomial operations use exact rational
arithmetic. Thus the denominator has no unverified numerical root or
cancellation. The decimal values \(1.36510950191\) and
\(545.543981407\) are only convenient approximations to the exact roots.
\end{proof}

\begin{proposition}[FAM D.1.2; T]\label{fam:D.1.2}
The locus on \(B^\circ\) where \(P(t)\) is torsion is countable and dense.
Every fibre over this locus is an abelian surface. The divisor on the
smooth base cut out by \(6kP=O\) has degree \(3k^2-1\), including local
intersection multiplicity, for every positive integer \(k\).
\end{proposition}
\noindent\emph{Uses:} FAM D.1.1; E26 Lemma 2.1 and its height calculation.\par

\begin{proof}
\par\noindent\textit{Step 1.} On a simply connected open subset write uniquely
\(\mu=a_0\tau+b_0\), \(a_0,b_0\) real. The section is torsion precisely
when \((a_0,b_0)\in\mathbb Q^2\). In a local complex parameter \(z\), direct
differentiation gives
\[
 \det\frac{\partial(a_0,b_0)}{\partial(\operatorname{Re}z,
 \operatorname{Im}z)}
 =-\frac{|\mu'(z)-a_0\tau'(z)|^2}{\operatorname{Im}\tau}.
\]
This determinant vanishes on no open set. On an open set of vanishing,
the two real differential equations imply \(da_0=db_0=0\), hence
\(\mu=a_0\tau+b_0\) with constant real coefficients. Analytic continuation
would give this identity on the cover. The cusp law
\((\tau,\mu)\circ g_0=(\tau-1,\mu)\) then forces \(a_0=0\), and the
\(g_2\)-law makes a constant \(\mu\) impossible. Thus every open subset
contains a point where this real map is a local diffeomorphism, and its
image contains rational pairs. This proves density. For each integer
triple \((n,A,B)\), \(n\ne0\), the function \(n\mu-A\tau-B\) is
holomorphic and not identically zero by the same argument. Its zeros are
discrete; a countable union proves countability.

\par\noindent\textit{Step 2.} For \(N\mu=A\tau+B\), \(A,B,N\in\mathbb Z\), \(N>0\), use the linear
coordinate change \(\zeta'_2=\zeta_2-(A/N)\zeta_1\).
The four integral vectors
\[
 N\widehat u-B\widehat\delta,\quad
 N\widehat w+A\widehat\delta,\quad
 N\widehat\gamma-6A\widehat u-6B\widehat w,\quad
 \widehat\delta
\]
span a sublattice of index \(N^3\). Its first two periods lie on
\(\zeta'_2=0\), and its last two lie on \(\zeta_1=0\).
Both intersections are genuine rank-two lattices because \(D<0\).
Their product quotient is therefore a product of elliptic curves, and
inclusion of this sublattice induces a finite surjective holomorphic
homomorphism onto \(F_t\). A positive integral form on the product can
be summed over its finite translation kernel and multiplied by the
index to descend an integral positive form to \(F_t\). The Riemann
polarization criterion proves projectivity.

\par\noindent\textit{Step 3.} For the degree statement we use the height computation actually read in
\cite[Lemma 2.1 and proof]{E26}: \(\langle P,P\rangle=1/6\),
\(\chi(\mathcal O_S)=1\), the component groups at the two reducible fibres
have orders three and two, and the local corrections for \(P\) are
\(4/3\) and \(1/2\). The application to \(6kP\) is as follows.
The corrected divisor class of a section is its class minus \(O\) and
a fibre multiple, plus the unique rational combination of nonidentity
fibre components that makes it perpendicular to those components.
On restriction to the generic fibre this assignment is additive because
the divisor class \([Q]-[O]\) is the group law in \(\operatorname{Pic}^0\).
The perpendicular correction makes its self-intersection the negative of
the height, so the height is quadratic. The element \(6kP\) is in the
identity component at both reducible fibres; its correction terms vanish.
Hence
\[
 6k^2=(6k)^2\langle P,P\rangle
 =2+2((6kP)\mathbin{\cdot}O),
\]
which gives the formula. There are no missing intersection points at the
three special fibres. The computed rational functions for \(6P\) show
that it has finite affine coordinates at \(0,1\), so it is a nonzero point
of each additive identity component; multiplication by \(k\) remains
nonzero in \(\mathbb C\). At infinity its multiplicative coordinate is
\((-a)^{6k}=a^{6k}\ne1\). Thus all intersections counted by the height
occur on \(B^\circ\). This derives the cited height input in precisely
the form needed here; the first five degrees and the correction matrices
are checked exactly in \texttt{cd\_torsion\_independent.py}.
\end{proof}

\noindent The following counterexample addresses only the unrestricted wording in FAM Lead~6. The first-projection assertion of Statement~\ref{fam:D.1.1} stands.\par
\begin{theorem}[FAM D.1.3; T]\label{fam:D.1.3}
The assertion in Lead 6 that \emph{existence of a surjection onto the
specified \(E_t\) with elliptic kernel} is equivalent to \(6P(t)=O\)
is false. It remains true for the specified projection in \ref{fam:D.1.1}.
\end{theorem}
\noindent\emph{Uses:} FAM D.1.1 and the source bound FAM C.2.1.\par

\begin{proof}
\par\noindent\textit{Step 1.} Set \(t=32/5\). The exact fourth division polynomial is
\[
 \psi_4(P)=-\frac1{32}t^{10}(t-1)^4(5t-32),
\]
whereas \(\psi_2(P)=t^2(t-1)\ne0\). Thus \(P\) has exact order four.
On the rectangular branch write \(\tau=iT\), \(T>1\), with
\(j(iT)=1728\cdot32/5\). The chosen zero has real part \(1/2\), negative
imaginary part of absolute value less than \(T/2\), and order four;
therefore \(\mu=(2-\tau)/4\). Choose the concrete parameter
\[
 c=\left(\frac98-\frac{101}{4}\right)\tau-\frac94
       -\beta^{\rm part}(t).
\]
It is globally admissible. Indeed the positive-real Jacobi product at this
\((\mu,\tau)\) gives
\(\operatorname{Im}\beta^{\rm part}\geq-T-(3/\pi)\log3\).
Consequently
\(\operatorname{Im}c\leq-185T/8+(3/\pi)\log3<-22<-M\), by C.2.1.
This inequality proves admissibility without using the numerical estimate
\(c=-2.25-34.02109724218209472308\,i\).

\par\noindent\textit{Step 2.} The holomorphic linear map \(L:\mathbb C^2\to\mathbb C\),
\(L(\zeta)=3\zeta_1+4\zeta_2\), takes the four periods to
\[
 -101\tau,\qquad2\tau+2,\qquad3,\qquad4.
\]
These generate exactly \(\mathbb Z\tau+\mathbb Z\): the last two generate
\(1\), and the first two then generate \(\tau\), since
\(\gcd(101,2)=1\). Thus \(L\) descends to a surjective holomorphic group
homomorphism \(F_{32/5}\to E_{32/5}\) with connected one-dimensional
kernel. The kernel is compact as a closed subset of a compact torus,
hence an elliptic curve. Nevertheless \(6P=2P\ne O\).
Every equality, the division-polynomial noncancellation, and the four
lattice images are checked by exact arithmetic in the supplied code.
\end{proof}

\begin{definition}[FAM D.1.Hermitian]\label{fam:D.1.Hermitian}
The real vector-space isomorphism
$\Pi:\mathbb R^4\to\mathbb C^2$ transfers $E$ to the actual alternating
form $\widetilde E(v,w)=E(\Pi^{-1}v,\Pi^{-1}w)$ on the tangent space.
Equation \eqref{fam:eq:cd-NS} gives $J$-invariance. The explicit comparison is
\[
 H_E(v,w)=\widetilde E(v,Jw)-i\widetilde E(v,w).
\] We use $T=\operatorname{Im}\tau$, $U=\operatorname{Im}\mu$ and $D=\operatorname{Im}\beta-6U^2/T$ in the matrix calculation below.
\end{definition}
\noindent\emph{Uses:} FAM D.1.6, Hermitian-form convention.\par

\begin{proposition}[FAM D.1.4; T]\label{fam:D.1.4}
For every smooth fibre and every admissible \(c\), its Neron--Severi lattice
is the following explicit subgroup of \(\mathbb Z^6\):
\begin{equation}\label{fam:eq:cd-NS}
 \left\{E:\ E_{12}-E_{13}\tau+(6E_{23}-E_{14})\mu+
 E_{24}\beta+E_{34}(6\mu^2-\tau\beta)=0\right\}.
\end{equation}
The complex orientation has intersection form
\(E^2=-2\operatorname{Pf}(E)\), where
\(\operatorname{Pf}(E)=E_{12}E_{34}-E_{13}E_{24}+E_{14}E_{23}\).
In particular the invariant class \(\eta=6\gamma\wedge\delta+u_{\Lambda}\wedge w_{\Lambda}\)
has \(\int_{F_t}\eta^2=-12\). Its Hermitian form is indefinite.
The fibres with lattice exactly \(\mathbb Z\eta\) have no compact curves
and algebraic dimension zero.
\end{proposition}
\noindent\emph{Uses:} The period matrix; the Hermitian-form convention in FAM D.1.6; AC26 Lemmas 9.2 and 9.4.\par

\begin{proof}
\par\noindent\textit{Step 1.} On the universal cover the holomorphic two-form is the wedge of the two
rows of \(\Pi\). Its six coefficients in the integral coframe are
\(6\mu^2-\tau\beta,-\beta,6\mu,-\mu,\tau,1\).
An integral real two-form is of type \((1,1)\) exactly when its wedge
with this form vanishes: its \((0,2)\) part is then zero, and reality
also annihilates its \((2,0)\) part. Expanding the wedge gives
\eqref{fam:eq:cd-NS}. The exponential sequence identifies these integral
\((1,1)\) classes with Chern classes of holomorphic line bundles, proving
the lattice assertion.

\par\noindent\textit{Step 2.} The real period matrix in C.3.1 has determinant \(TD<0\). Thus the
integral orientation \((\widehat\gamma,\widehat u,\widehat w,
\widehat\delta)\) is opposite to the complex orientation. Algebraically
\(E\wedge E=2\operatorname{Pf}(E)\gamma\wedge u_{\Lambda}\wedge w_{\Lambda}\wedge\delta\),
which proves the sign in the displayed intersection formula. In
particular the expression ``\(\eta^2=12\)'' in \cite[Lemma~9.4]{AC26} is correct
as its displayed coefficient of \(\gamma\wedge u_{\Lambda}\wedge w_{\Lambda}\wedge\delta\),
and has the opposite sign as a complex-oriented numerical
self-intersection. We retain the coefficient identity and record the
orientation comparison explicitly.

\par\noindent\textit{Step 3.} For the precise Hermitian convention $H_E$ defined in \ref{fam:D.1.6}, direct
substitution of the four periods gives $H_\eta=-2G^{-1}$, where
\[
 G=\begin{pmatrix}2T&2U\\2U&\operatorname{Im}\beta/3\end{pmatrix},
 \qquad \det G=2TD/3<0.
\]
Thus $H_\eta$ has one positive and one negative eigenvalue, and every
nonzero multiple remains indefinite. The integration
current of an effective divisor, averaged under translations of the
compact torus, is a nonzero semipositive constant \((1,1)\)-form in its
Chern class. It cannot be an indefinite multiple of \(\eta\).
Thus no curve exists when \(\operatorname{NS}=\mathbb Z\eta\).
A nonconstant meromorphic function would have a nonempty polar divisor
(after subtracting a constant if necessary), so the algebraic dimension
is zero.
\end{proof}

\begin{proposition}[FAM D.1.5; T]\label{fam:D.1.5}
For a fixed \(c\), all but countably many smooth fibres have
\(\operatorname{NS}=\mathbb Z\eta\). Conversely, for every fixed smooth
\(t\), there are admissible parameters tending to the cusp of the
parameter disc for which \(F_t\) is an abelian surface, irrespective of
the order of \(P(t)\).
\end{proposition}
\noindent\emph{Uses:} FAM D.1.4; the five-function independence of AC26 Lemma 9.3.\par

\begin{proof}
\par\noindent\textit{Step 1.} We verify the required independence of five holomorphic functions.
Suppose \(A+B\tau+C\mu+D\beta+E(6\mu^2-\tau\beta)=0\).
Apply \(g_0\), using \((\tau,\mu,\beta)\circ g_0=(\tau-1,\mu,\beta+1)\),
and subtract. The result is \(-B+D+E(\beta-\tau+1)=0\).
A second cusp difference gives \(2E=0\), hence \(E=0,D=B\).
Applying \(g_2\) to the remaining relation and multiplying the difference
by \(-\tau\) gives
\[
 6B\mu^2+C(\tau-1)\mu+B(\tau^2+3\tau+1)-C\tau=0.
\]
Its cusp difference gives \(C\mu+2B(\tau+1)-C=0\), whose next cusp
difference forces \(B=0\). Then \(C\mu=C\); the \(g_2\)-law forces
\(C=0\), and finally \(A=D=0\). Thus an integral class outside
\(\mathbb Z\eta\) has a nonidentically-zero holomorphic left side in
\eqref{fam:eq:cd-NS}; its zeros are discrete. Countability of \(\mathbb Z^6\)
proves the first assertion, with explicit verification of its reduction.

\par\noindent\textit{Step 2.} For the second assertion put
\(\beta_n=-3n^2\tau+12n\mu\),
\(c_n=\beta_n-\beta^{\rm part}(t)\). Its imaginary part tends to
minus infinity, so the already proved finite bound on \(M\) makes every
sufficiently large \(n\) admissible. The integral alternating matrix
with nonzero upper entries
\[
 E_{13}=-3n^2,\quad E_{14}=6n,\quad E_{23}=-n,\quad E_{24}=1
\]
satisfies \eqref{fam:eq:cd-NS} exactly at \(\beta_n\), and
\(\operatorname{Pf}(E)=-3n^2<0\). Its complex-oriented square is positive.
A nondegenerate Hermitian form in complex dimension two with positive
determinant is definite. Choosing its sign gives a positive integral
Riemann form, so \(F_t\) is projective. This supplies an explicit family
of projective fibres even at any chosen non-torsion value of the section.
\end{proof}

\begin{theorem}[FAM D.1.6; T]\label{fam:D.1.6}
Fix any admissible $c$. For each integral alternating matrix $E$, let
$Z_E(c)$ be the zero set on the marking cover of the explicit holomorphic
function in \eqref{fam:eq:cd-NS}, with $\beta=\beta^{\rm part}+c$ given by
\eqref{fam:eq:cd-beta}. Let $\pi$ be the cover map to $B^\circ$, and define
\[
 \mathcal E_c=\pi\!\left(\bigcup_{E\ne0,\ \operatorname{Pf}(E)=0}Z_E(c)\right),
 \qquad
 \mathcal A_c=\pi\!\left(\bigcup_{\operatorname{Pf}(E)<0}Z_E(c)\right).
\]
These give the exact loci of all elliptic quotients and all abelian
surfaces, respectively. The algebraic dimension of every smooth fibre is
\[
 a(F_t)=
 \begin{cases}
 2,&t\in\mathcal A_c,\\
 1,&t\in\mathcal E_c\setminus\mathcal A_c,\\
 0,&t\notin\mathcal E_c\cup\mathcal A_c.
 \end{cases}
\]
Together with \eqref{fam:eq:cd-NS}, these are an exact classification for every
$t$, by countable unions of explicitly given analytic zero sets; they do
not assert that the loci have a finite algebraic description. They include
all quotient maps, without requiring preservation of the first coordinate.
\end{theorem}
\noindent\emph{Uses:} FAM D.1.4; the Appell--Humbert and Riemann criteria.\par

\begin{proof}
\par\noindent\textit{Step 1.} We apply the Appell--Humbert correspondence and the Riemann projectivity
criterion in \cite[Theorems 4.1 and 4.4]{Schnell19}; the reduction to the
present explicit periods is as follows. The real vector-space isomorphism
$\Pi:\mathbb R^4\to\mathbb C^2$ transfers $E$ to the actual alternating
form $\widetilde E(v,w)=E(\Pi^{-1}v,\Pi^{-1}w)$ on the tangent space.
Equation \eqref{fam:eq:cd-NS} gives $J$-invariance. The explicit comparison is
\[
 H_E(v,w)=\widetilde E(v,Jw)-i\widetilde E(v,w).
\]
Direct substitution gives $H_E(Jv,w)=iH_E(v,w)$ and
$H_E(w,v)=\overline{H_E(v,w)}$. Its associated real Chern form is
$-\operatorname{Im}H_E=\widetilde E$, since
$(i/2)(H_E(v,w)-H_E(w,v))=\widetilde E(v,w)$.
Thus the Appell--Humbert integrality condition is exactly integrality
of the given $E$, with its sign fully specified. The complex-oriented square computed in \ref{fam:D.1.4} shows that
$\operatorname{Pf}(E)<0$ is equivalent to definiteness of this Hermitian
form. Its sign can be chosen positive. The cited projectivity criterion
then gives precisely the asserted abelian locus, in both directions.

\par\noindent\textit{Step 2.} For nonzero $E$ with Pfaffian zero, its real rank is two. Its nullspace
$K_{\mathbb R}\subset\mathbb R^4$ is a rational plane because it is the
kernel of an integral matrix. $J$-invariance makes $\Pi K_{\mathbb R}$ a
complex line, and $K=K_{\mathbb R}\cap\mathbb Z^4$ is saturated. The maps
\[
 (\Pi K_{\mathbb R})/\Pi K\ \hookrightarrow F_t
 \ \longrightarrow\
 \mathbb C^2/(\Pi K_{\mathbb R}+\Pi\mathbb Z^4)
\]
are holomorphic group morphisms; the first is a closed embedded elliptic
curve and the second is its quotient onto an elliptic curve, with
connected kernel. A basis of the saturated $K$ extends to an integral basis of $\mathbb Z^4$.
The other two periods project to real-linearly independent vectors in the
one-dimensional complex quotient, because $\Pi$ is a real isomorphism.
They therefore form a discrete lattice there. Compactness follows by
taking the image of its compact real fundamental parallelogram. Conversely any elliptic quotient
pulls an integral positive area form back to a nonzero rank-two integral
$(1,1)$ form. This proves the characterization of $\mathcal E_c$.
Equivalently all its kernels can be found by saturated rank-two lattices
with a basis $v,w$ satisfying
$\det_{\mathbb C}(\Pi v,\Pi w)=0$ and real independence. This is a second
exact integer test, with explicit inclusion and quotient morphisms.

\par\noindent\textit{Step 3.} For the algebraic-dimension assertion, an effective divisor on a torus
has a semipositive Hermitian form: averaging its integration current under
translations gives that constant form. A nonzero semipositive form has
rank one or two. Rank two gives the abelian case. Rank one gives the
elliptic quotient just constructed. In the absence of a definite form,
there cannot be two different such kernel subtori. Indeed addition from
their product onto $F_t$ would be a surjective holomorphic homomorphism
with finite kernel, and the product polarization would descend after
multiplication by an integer, giving a definite form. Thus there is a
single kernel $K_t$ in this case. Every effective divisor is a sum of
its translates: its intersection number with $K_t$ is zero, whereas a
curve mapping nonconstantly to the elliptic quotient has strictly positive
intersection with a general quotient fibre. Every meromorphic function
therefore factors through this quotient. One can check the last statement
on a fibre avoiding its polar divisor, where it is a holomorphic function
on a compact elliptic curve and hence constant; analytic continuation
then gives the factorization. The quotient supplies nonconstant
meromorphic functions, so the algebraic dimension is exactly one.
If neither locus occurs, a nonconstant meromorphic function would have
a polar divisor of one of the two forbidden ranks. This proves dimension
zero. Finally monodromy acts by integral determinant-one changes of basis,
preserving \eqref{fam:eq:cd-NS} and its Pfaffian. Thus the stated unions descend
without any marking ambiguity. Each individual matrix and every locus
is specified; no expected identification is left in these formulas.
\end{proof}

\begin{proposition}[FAM D.1.7; T]\label{fam:D.1.7}
There is an explicit nonempty locus of admissible members with ordinary
fibres of algebraic dimension one. On a sufficiently small punctured
cusp neighbourhood set
\[
 c(z)=-6\tau(z)-12\mu(z)-\beta^{\rm part}(z).
\]
Outside a countable set of $z$, the fibre of $X_{c(z)}$ over $\pi(z)$ has
Neron--Severi lattice $\mathbb Z\eta+\mathbb ZE$, where the nonzero upper
entries of $E$ are
$E_{13}=-6$, $E_{14}=-6$, $E_{23}=1$, $E_{24}=1$.
Its algebraic dimension is one. The section $P(\pi(z))$ can simultaneously
be required to be non-torsion.
\end{proposition}
\noindent\emph{Uses:} FAM D.1.6 and the cusp expansions.\par

\begin{proof}
At the specified parameter $\beta=-6\tau-12\mu$. Both $E$ and $\eta$
satisfy \eqref{fam:eq:cd-NS}, and exact multiplication gives
$\operatorname{Pf}(a\eta+bE)=6a^2$.
For any further integral class $F$, substitution in its period equation
gives a linear relation among $1,\tau,\mu,(\mu+\tau)^2$.
These four functions are linearly independent: applying $g_0$ to a
relation $A+B\tau+C\mu+D(\mu+\tau)^2=0$ and subtracting yields
$-B+D(1-2(\mu+\tau))=0$. A further cusp difference forces $D=0$,
and then $B=0$; nonconstancy of $\mu$ forces $A=C=0$.
Comparison of coefficients now gives
$F_{34}=F_{12}=0$, $F_{13}=-6F_{24}$, and
$6F_{23}-F_{14}=12F_{24}$, exactly
$F\in\mathbb Z\eta+\mathbb ZE$. Every other integral class therefore
has a discrete zero set along this period locus. Excluding their
countable union gives the asserted exact Neron--Severi lattice.
It contains the nonzero isotropic $E$ and no matrix of negative
Pfaffian; \ref{fam:D.1.6} gives algebraic dimension one.
Finally
$c=-5\tau-12\mu-(\beta^{\rm part}+\tau)$ has imaginary part tending
to minus infinity at the cusp. A sufficiently small cusp neighbourhood
is therefore admissible by C.2.1. The torsion locus of $P$ is countable
by \ref{fam:D.1.2}, so its exclusion still leaves such points in every open subset
of that neighbourhood. This supplies actual members and fibres, with
all their parameter functions specified.
\end{proof}

\subsection{Compact curves and surfaces}
The cusp normalization and its boundary identifications used here are due to \cite[\S4.6, Proposition~4.6]{AC26}.
\begin{proposition}[FAM D.2.1; T]\label{fam:D.2.1}
For a standard-twist member \(X_c\) with underlying manifold \(S^6\),
every compact irreducible complex curve lies in one fibre of \(f\), and
every compact irreducible complex surface is one of the reduced fibres.
\end{proposition}
\noindent\emph{Uses:} $X_c\simeq S^6$; the K\"ahler geometry of the reduced fibres.\par

\begin{proof}
The given diffeomorphism with \(S^6\) implies \(H^2(X_c;\mathbb Z)=0\).
For a compact irreducible curve \(C\), let
\(\nu:\widetilde C\to C\hookrightarrow X_c\) be its normalization map.
The holomorphic map \(f\nu\) has
\(\deg (f\nu)^*\mathcal O_{\mathbb P^1}(1)=0\), since its Chern class is
pulled back from the zero class on \(X_c\). A nonconstant holomorphic map
of compact Riemann surfaces has positive degree: a regular value has a
positive number of preimages, all with positive local multiplicity.
Therefore \(f\nu\), and hence \(f|_C\), is constant.

An irreducible surface \(Z\) in the smooth threefold is an effective
Cartier divisor. If its image were all of \(B\), a general smooth fibre
would meet it in a nonempty effective divisor \(D\) whose line bundle
is \(\mathcal O_{X_c}(Z)|_{F_t}\). Its integral Chern class is zero.
Choose a flat K\"ahler form \(\omega\) on the torus \(F_t\). The equality
\(\int_{F_t}c_1(D)\wedge\omega=\int_D\omega\) has left side zero and right
side positive, a contradiction. Thus \(Z\) lies in a single fibre.
Each reduced fibre is irreducible of dimension two, so containment forces
equality. This uses the full map into \(X_c\), not a transfer of a
vanishing claim from its normalization.
\end{proof}

\begin{proposition}[FAM D.2.2; T]\label{fam:D.2.2}
There is no rational curve in \(X_c\setminus W_c\). There are elliptic
curves outside all three special fibres, in particular in each of the
two smooth fibres \(F_{t_\pm}\). All curves in \(W_c\) are precisely the
images of irreducible algebraic curves in \(\mathrm{dP}_6\) under its
finite normalization map. Distinct curves have the same image only for
the three paired boundary lines. All curves in a bielliptic fibre
\(S_j\) are images of irreducible curves in its finite product cover.
\end{proposition}
\noindent\emph{Uses:} FAM D.1.1 and D.2.1; the finite normalization and boundary identifications of AC26 Proposition 4.6.\par

\begin{proof}
\par\noindent\textit{Step 1.} By \ref{fam:D.2.1} every curve is vertical. A holomorphic map from
\(\mathbb P^1\) into a torus lifts to its vector-space universal cover,
since \(\mathbb P^1\) is simply connected, and its coordinate functions
are constant by compactness. Thus a smooth torus fibre contains no
rational curve. A map from \(\mathbb P^1\) into a bielliptic surface
lifts through its finite unramified torus cover by the same covering
criterion; the lifted map is constant. These arguments prove the first
assertion, including possibly singular rational images.
\ref{fam:D.1.1} gives an embedded connected elliptic kernel in each \(F_{t_\pm}\),
proving the second.

\par\noindent\textit{Step 2.} For the third assertion let \(\eta:\mathrm{dP}_6\to W_c\) be the finite
normalization morphism. Each component of the inverse image of a compact
curve is compact and one-dimensional and maps onto that curve, because
finite maps have zero-dimensional fibres. It is algebraic in the
projective surface \(\mathrm{dP}_6\). Conversely its image is a compact
analytic curve by properness. Away from the conductor, \(\eta\) is a
biholomorphism, so two distinct curves can have the same image only when
both lie in the conductor; its identifications are exactly the three
opposite pairs computed in FAM C.4.1; see also \cite[\S4.6, Proposition~4.6]{AC26}. This is also an explicit algebraic
parametrization: the curves of \(\mathrm{dP}_6\) are the three exceptional
lines and strict transforms of irreducible plane curves under the blow-up
of the three coordinate points. Their homogeneous equations and their
three multiplicities specify the strict transforms. The finite-product
cover statement for \(S_j\) follows from the identical proper-image and
finite-preimage argument applied to its unramified morphism
\(E_j\times F_j\to S_j\); curves in the product are algebraic.
\end{proof}

\begin{theorem}[FAM D.2.3; T]\label{fam:D.2.3}
The following descriptions give all compact irreducible curves of $X_c$.
In $W_c$ and $S_j$ they are exactly the images described in \ref{fam:D.2.2}.
In a smooth fibre of algebraic dimension zero there are none. In a
smooth fibre of algebraic dimension one they are exactly the translates
of its unique elliptic kernel $K_t$. In an abelian fibre they are exactly
the translates of elliptic subtori, together with the irreducible zero
divisors of the positive Appell--Humbert theta sections described below.
These descriptions, \ref{fam:D.1.6}, and \ref{fam:D.2.1} constitute a classification including
every smooth exceptional fibre and every compact irreducible surface.
\end{theorem}
\noindent\emph{Uses:} FAM D.1.6 and D.2.1; the Appell--Humbert theorem and its theta equation.\par

\begin{proof}
\par\noindent\textit{Step 1.} All curves are vertical by \ref{fam:D.2.1}, so the list of fibre types is exhaustive.
The dimension-zero and dimension-one assertions, including uniqueness,
were proved in \ref{fam:D.1.6}. Let now $F_t=\mathbb C^2/\Pi\Lambda$ be abelian.
For every positive definite Hermitian form $H$ whose imaginary part takes
integer values on $\Pi\Lambda$, and every function
$\alpha:\Pi\Lambda\to U(1)$ satisfying
\[
 \alpha(\lambda+\nu)=\alpha(\lambda)\alpha(\nu)
                 e^{\pi i\operatorname{Im}H(\lambda,\nu)},
\]
take the nonzero entire solutions $\theta:\mathbb C^2\to\mathbb C$ of
\begin{equation}\label{fam:eq:cd-theta-sections}
 \theta(\zeta+\lambda)=\alpha(\lambda)
 e^{\pi H(\zeta,\lambda)+\pi H(\lambda,\lambda)/2}\theta(\zeta).
\end{equation}
The cocycle identity follows by multiplying the displayed factors; its
integral imaginary cross-term has exponential one. It defines a
holomorphic line bundle by quotienting $\mathbb C^2\times\mathbb C$.
The invariant zero set of $\theta$ descends to a compact divisor on
$F_t$. Retain precisely those zero divisors that are reduced and
irreducible. This is an exact parametrization, with all periods and
multiplier equations given; multiplying $\theta$ by a nonzero scalar
does not change its divisor.

\par\noindent\textit{Step 2.} The correspondence \cite[Theorem 4.1 and equation (4.3)]{Schnell19} proves
that every holomorphic line bundle and every section on this actual
quotient has this form. For the cohomology matrix the exact comparison is
$H_E=\widetilde E(\,\cdot\,,J\,\cdot\,)-i\widetilde E$ from \ref{fam:D.1.6};
we choose the sign of $E$ making $H_E$ positive. Its imaginary part is
$-\widetilde E$, which gives the integral semicharacter equation displayed
above. A divisor class with positive Hermitian rank two gives exactly
these positive theta sections. A divisor of Hermitian rank one is pulled
back from its elliptic quotient and, when irreducible and reduced, is a
single fibre. Thus it is a translate of an elliptic subtorus. The null
lattice is computed by the explicit integer kernel in \ref{fam:D.1.6}. No other
rank is possible for an effective nonempty divisor, by positivity of its
averaged current. This proves exhaustiveness on the abelian fibre.

\par\noindent\textit{Step 3.} The same equations give the curves in the product covers of $S_j$:
compute them on its two-elliptic-curve lattice and take images under the
specified finite morphism. Conversely each curve downstairs is the image
of a component of its finite inverse image, so that this procedure omits
none. In $W_c$ the analogous procedure uses the homogeneous polynomial
parametrization of the normalization in \ref{fam:D.2.2}, with the three boundary
pairings already identified. The maps used throughout are finite or
closed embeddings, so they preserve compactness and give one-dimensional
images; no assertion is transferred from a normalization in an invalid
direction.
\end{proof}

\section{The outer boundary and the untwisted filling conjecture}
\label{sec:ends}
\subsection{The intrinsic obstruction at the outer circle}
Throughout this subsection the twists are the standard ones.
Since $q_2=iw/216$, its boundary circle is $|q_2|=1$.
\begin{definition}[PUSH D.1.setup]\label{push:D.1.setup}
We use the standard twists throughout this section. Put
\(H_*=(3/\pi)\log6\), fix a boundary parameter
\(c_*=x+iH_*\), and set \(a=x+3/2\).
\end{definition}
\noindent\emph{Uses:} The exact radius 216.\par

\begin{proposition}[PUSH D.1.1; T]\label{push:D.1.1}
For \(0<|w|<216\), the reduced multiplicity-four surface has
\[
 j(\operatorname{Alb}S_2(w))=J(iw/216),\qquad
 J(q)=q^{-1}+744+196884q+\cdots.
\]
This holomorphic function has no meromorphic continuation through any
point of \(|w|=216\).
\end{proposition}
\noindent\emph{Uses:} FAM C.1.1 and A.4.1; ARI C.2.1; modular invariance of j.\par

\begin{proof}
The fixed-coordinate periods on the product cover, including the
isogeny kernel and the affine deck translation, generate
\(\frac12\mathbb Z+\frac{\kappa_2}{4}\mathbb Z\).
Multiplication by two therefore identifies the Albanese lattice with
\(\mathbb Z+(\kappa_2/2)\mathbb Z\). This verifies the precise step of
FAM A.4.1 used here; taking only the invariant eigenspace would omit the
half-period. The values in FAM C.1.1 and \(\mu(z_2)=(1-i)/2\) give
\[
 \kappa_2=c+\frac32-iH_*,\qquad
 e^{-\pi i\kappa_2}=e^{-3\pi i/2}e^{-\pi H_*}e^{-\pi ic}
 =\frac{iw}{216}.
\]
The positive-imaginary representative is \(-\kappa_2/2\), proving the
identity with the stated Fourier normalization.

For completeness we reproduce the obstruction in FAM C.3.4, rather than regard a numerical Fourier plot as evidence. At a
reduced rational \(r=b/d\) choose a matrix in
\(\operatorname{SL}_2(\mathbb Z)\) with bottom row \((d,-b)\).
For \(z=r+iy\), its transformed imaginary part is \(1/(d^2y)\).
Modular invariance and the convergent Fourier series imply
\(\log|j(r+iy)|=2\pi/(d^2y)+o(1)\).
A meromorphic function near \(r\) grows at most polynomially in
\(1/y\), so no continuation exists there. Every real interval contains
such a rational, excluding continuation at irrational points as well.
A local logarithm of \(q=iw/216\) transfers this result to the circle.
This is a derivation of the earlier natural-boundary result, with the
normalizing constant checked exactly.
\end{proof}

\begin{lemma}[PUSH D.1.2; T]\label{push:D.1.2}
Let \(V\) be a connected Riemann surface and
\(\pi:\mathcal X\to V\) a flat proper holomorphic map of reduced
complex spaces. Suppose that over a nonempty open subset \(V_0\), the
fibres are smooth compact threefolds whose unique anticanonical divisor
is \(2S_v\), where \(S_v\) is a smooth connected surface with
\(b_1(S_v)=2\). Then there is an open subset \(V^{\rm good}\subset V\), with at
most countable complement, on which the reduced support of the
relative anticanonical divisor is a smooth proper surface family. Its
Albanese elliptic invariant is holomorphic.
\end{lemma}
\noindent\emph{Uses:} Grauert direct-image and proper-mapping theorems with the 1963 correction; Ehresmann; Buchdahl Theorem 11.\par

\begin{proof}
\par\noindent\textit{Step 1.} Flatness and the smoothness of the fibres over \(V_0\) imply that
\(\pi\) is a submersion there: the local analytic smoothness criterion
applies because the base is smooth. The relative nonsmooth locus is a
closed analytic subset, by the Jacobian criterion. Its proper image is
analytic by \cite[\S7, Satz~1]{Grauert1960D}, the proper mapping theorem.
It misses \(V_0\), hence is discrete. We remove it and work over the
connected remaining curve \(V'\). Removing a discrete subset from a
connected Riemann surface does not disconnect it. Thus
\(\mathcal X'\to V'\) is a smooth proper map.

\par\noindent\textit{Step 2.} Write \(L=K_{\mathcal X'/V'}^{-1}\). Grauert's direct-image theorem \cite[\S6, Hauptsatz~I]{Grauert1960D} makes \(F=\pi_*L\) and \(R^1\pi_*L\)
coherent. We prove the needed base change directly over the smooth
curve; we do not use the originally stated \S7, Satz~2, whose claimed
closed exceptional locus was corrected in Grauert's 1963 erratum \cite[item~2, p.~36]{Grauert1963ErratumD}.
For a local base parameter \(t\) vanishing at \(v\), smoothness and
invertibility of \(L\) give an exact sequence
\(0\to L\xrightarrow{t}L\to i_*L_v\to0\), where
\(i:X_v\hookrightarrow\mathcal X'\) is the closed embedding.
Its direct-image sequence gives, at \(v\),
\[
 0\longrightarrow F_v/tF_v\longrightarrow H^0(X_v,L_v)
 \longrightarrow (R^1\pi_*L)_v[t]\longrightarrow0.
\]
The last term denotes elements annihilated by \(t\). The support of
the torsion subsheaf of the coherent module \(R^1\pi_*L\) is closed
discrete, by the structure theorem for finitely generated modules
over the local discrete valuation rings of the smooth base curve.
Outside that support the canonical base-change map is an isomorphism.
Also \(F\) is torsion-free: multiplication by a nonzero base function
is injective on sections of \(L\). It is therefore locally free.
Its rank is one by choosing an interior point outside the discrete
torsion support. These statements concern the canonical map, not
merely equal dimensions.
Locally a frame gives a relative section \(\sigma\) of \(L\),
nowhere identically zero on a fibre. Its divisor is independent of the
frame. Let \(D\) be its reduced support. On \(V_0\), the divisor is
\(2S_v\) and \(D\) is exactly \(S_v\).

\par\noindent\textit{Step 3.} The map \(D\to V'\) is proper, since \(D\) is closed in a proper
space. There is no vertical divisor component: a vertical divisor in
the smooth total space would contain a whole irreducible fibre,
contrary to the nonzero restriction of \(\sigma\). Every horizontal
component has proper image all of the connected base. Since the
interior fibre is irreducible, there is exactly one horizontal
component, and its multiplicity in \(\operatorname{div}\sigma\) is
two. Thus \(\operatorname{div}\sigma=2D\) on this base. In the smooth
ambient space \(D\) is Cartier. It is flat over the curve because
its local rings are torsion-free over the base discrete valuation
ring. Restriction of the Cartier identity to an interior fibre gives
\(2D_v=2S_v\), whence \(D_v=S_v\) with its reduced structure.
Flatness and the smoothness of these fibres make \(D\to V'\) smooth
over \(V_0\). Its analytic nonsmooth locus therefore misses \(V_0\);
the proper mapping theorem once more makes its image discrete.
After deleting that image, \(D\to V^{\rm good}\) is a smooth proper
family of connected surfaces. Each deletion is closed discrete in the
then-current open curve. Their finite union is at most countable,
and the final good locus is open; we do not assert that successive
exceptional sets cannot accumulate at an earlier removed point.

\par\noindent\textit{Step 4.} Ehresmann's theorem \cite[\S1, Proposition, pp.~154--155]{Ehresmann52} identifies the smooth fibres along every path;
in particular their first Betti number is two. \cite[Theorem~11]{Buchdahl1999} then gives a K\"ahler metric on every fibre. Hence
\(h^{1,0}=1\), every holomorphic one-form is closed, and the integral
period map has a rank-two lattice image. Constant \(h^{1,0}\),
properness, and analytic base change supply a local relative
holomorphic one-form. Locally on a simply connected
base disc, a relative one-form and a basis of the free part of
\(H_1\) give two holomorphic periods. Their ratio lies in one
half-plane after orienting the basis. To see holomorphicity directly,
choose local primitives of the closed relative form on simply
connected fibre charts; their differences are holomorphic functions
of the base, and the finite sum around a transported cycle is its
period. The modular invariant of the ratio is independent of all
choices and is holomorphic. This proves the last assertion without
requiring a global choice of section of the Albanese family.
\end{proof}

\begin{theorem}[PUSH D.1.3; T]\label{push:D.1.3}
For every \(w_0\) with \(|w_0|=216\) and every connected open
neighbourhood \(U\) of \(w_0\), there is no flat proper holomorphic
map \(\mathcal X\to U\) from a reduced complex space whose fibres
over \(U\cap\{0<|w|<216\}\) are biholomorphic, fibre by fibre, to
\(X_w\). In particular an extension isomorphic to the given family
over that intersection does not exist.
\end{theorem}
\noindent\emph{Uses:} FAM A.1.1; PUSH D.1.1 and D.1.2.\par

\begin{proof}
We may restrict a putative extension to a small disc meeting the
circle, with closure away from zero. This is a holomorphic base
restriction; properness and flatness are preserved. The calculation
of the canonical evaluation in the earlier return, \ref{push:A.1.1}, gives
\(K_{X_w}=\mathcal O(-2S_2(w))\) and
\(h^0(K_{X_w}^{-1})=1\). Its sole divisor is \(2S_2(w)\), so it is
intrinsic under every fibrewise biholomorphism. The surface has
\(b_1=2\): on its finite product cover the rotation kills the moving
factor in invariant first cohomology and fixes the elliptic factor.
These are the exact FAM A.1.1 and FAM A.4.1 steps required in Lemma~\ref{push:D.1.2}.

The at most countable exceptional set in that lemma cannot contain
an arc of the circle. Choose a boundary point outside it and then a smaller disc
on which the reduced divisor is a smooth proper surface family. Its
holomorphic Albanese invariant agrees, pointwise on the inner part,
with \(J(iw/216)\). It is therefore an analytic continuation through
that point, contrary to Proposition~\ref{push:D.1.1}. No choice of the
fibrewise biholomorphisms was used. Thus the weaker hypothesis already
leads to the contradiction.
\end{proof}

\subsection{The specified pinching and its virtual Euler number}
\begin{definition}[PUSH A.1.setup]\label{push:A.1.setup}
We write $M_S=\mathcal O_S(P-O)$, $T=t_{6P}$,
$J_{\rm lin}=M_S\otimes T^*M_S^{-1}=\pi^*\mathcal O_B(1)$, and
$R=\mathbb P_S(\mathcal O\oplus M_S)$. The convention for this projective
bundle is that its zero and infinity sections have normal bundles
$M_S$ and $M_S^{-1}$. The identification defining $Q$ is
$s_0(x)\sim s_\infty(Tx)$. The equality defining $J_{\rm lin}$ is the equality
in \cite[Lemma~2.1]{E26}, with this direction of $T$; its intersection
calculation uses $P\cdot O=0$, $P^2=O^2=-1$,
$6P\cdot O=6P\cdot P=2$. Thus the degree on $O$ is
$(O-P)\cdot6P+(P-O)\cdot O=1$. We use the section
$\phi\in H^0(S,J_{\rm lin})$ whose divisor is $S_\infty$, as in \cite[Corollary~2.2]{E26}. 
\end{definition}
\noindent\emph{Uses:} E26 Lemma 2.1 and Corollary 2.2.\par

\begin{refuted}[E26 Conjecture 8.2; false as worded]\label{e26:Conjecture8.2}
There is a flat, proper filling $\mathcal X\to \Delta$ and the central fiber $X(0)$ is non-normal. The normalization
is $\mathbb P_S(\mathcal O\oplus M)$ and the normalization map glues the zero and infinity sections via $t_{6P}$.
\end{refuted}
\noindent\emph{Uses:} the untwisted family and the specified pinching of \cite[Conjecture~8.2]{E26}. In this quotation $M$ is Engel's bundle $M_S$.\par
\begin{lemma}[PUSH A.3.1; T]\label{push:A.3.1}
Let $Z$ be a compact reduced local-complete-intersection complex
space of pure dimension $n$. In any proper flat holomorphic
family over a disc with central fibre $Z$ and smooth nearby
fibres, their Euler number is
\[
 e_{\mathrm{vir}}(Z)=\langle c_n(T_Z^{\mathrm{vir}}),[Z]\rangle.
\]
This statement needs no projectivity or K\"ahler hypothesis.
\end{lemma}
\noindent\emph{Uses:} The local-complete-intersection deformation calculation; Levy Lemmas 3.1--3.2; \L{}ojasiewicz \S3 Theorems 1--2.\par

\begin{proof}
\par\noindent\textit{Step 1.} Flat deformation preserves the local-complete-intersection
property relative to the disc near $Z$. Locally this follows by
lifting an embedding and its regular sequence of equations;
Nakayama's lemma and flatness show that their lifts generate
the ideal and remain a regular sequence. Properness allows one
to shrink the disc so that these neighbourhoods cover the total
space. The relative cotangent complex is consequently perfect
of amplitude $[-1,0]$. Its dual defines a topological complex
$K$-class $V$ on that total space; local finite vector-bundle
complexes define the class by descent, or by stabilization and
partitions of unity; the analytic-to-topological comparison is
given in \cite[Lemmas 3.1--3.2 and Section 3.1]{Levy08}.
Its restriction to a smooth fibre is its
tangent bundle, and its restriction to $Z$ is its intrinsic
virtual tangent class.

\par\noindent\textit{Step 2.} The compact fundamental cycles of the central and a nearby
fibre are homologous in the total space. More explicitly,
the inverse image $K$ of a short closed radial segment
$[0,\epsilon]$ is compact and semianalytic. Away from
$\operatorname{Sing}Z$ it is an oriented real
$(2n+1)$-manifold with boundary the nearby fibre and
$-Z_{\mathrm{reg}}$, because flatness with smooth fibre
implies smoothness of the morphism at those points.
The exceptional set $\operatorname{Sing}Z$ has real
dimension at most $2n-2$. Triangulate compatibly with
this set and the two endpoints, using
\cite[Section 3, Theorems 1--2]{Loj64}. The sum of the
oriented top simplices has boundary exactly
$[X_\epsilon]-[Z]$: no real $2n$-simplex can be supported
on that exceptional set. Reducedness and flatness give
multiplicity one on every central component. Pairing the globally defined
$c_n(V)$ with these homologous cycles gives the equality of
Chern numbers. Gauss--Bonnet on the smooth compact complex
fibre gives its Euler number. This proves the analytic assertion
directly. The algebraic counterpart is \cite[Theorem~5.8 and Remark~5.9]{FG10}; those scheme statements alone would
not justify omitting the analytic comparison above.
\end{proof}

\begin{theorem}[PUSH A.3.2; T]\label{push:A.3.2}
For the specified pinching $Q$ one has $e_{\mathrm{vir}}(Q)=0$.
Every smooth proper-flat smoothing of $Q$ therefore has Euler
number zero. In particular there is no proper flat holomorphic
filling of Engel's untwisted family with central fibre exactly
$Q$, even if no map to $B$ is required. Thus Conjecture 8.2
of \cite[Conjecture~8.2]{E26} is false as worded.
\end{theorem}
\noindent\emph{Uses:} PUSH A.3.1; the specified pinching; the Euler number two of the chosen untwisted family.\par

\begin{proof}
\par\noindent\textit{Step 1.} Let $D_0,D_\infty\subset R$ be the conductor sections and
$\nu:R\to Q$ the normalization. On either branch of the
hypersurface $xy=0$, the derivative comparison with its
normalization gives in $K^0(R)$, where pullback of the
perfect tangent class means derived pullback,
\[
 \nu^*T_Q^{\mathrm{vir}}
 =T_R-i_{0*}(J_{\rm lin}|_{D_0})-i_{\infty*}(J_{\rm lin}|_{D_\infty}).
\]
Indeed in a smooth local ambient space the normal line $A$ of
the branch and the normal line of the whole hypersurface
differ by $D$: the latter is $A(D)$. Their difference in
$K$-theory is $-[i_*(A(D)|_D)]$, and $A(D)|_D$ is exactly
the product of the two branch normals, namely $\mathcal T^1_Q$.
These local comparisons are intrinsic and agree on overlaps.
On the infinity branch $J_{\rm lin}$ means its $T^{-1}$-pullback
from the double locus, canonically the same line bundle
because it comes from $B$. Thus $J_{\rm lin}$ extends from both sections by the pullback from $B$,
so $[i_*J_{\rm lin}]=[J_{\rm lin}]-[J_{\rm lin}(-D)]$.

\par\noindent\textit{Step 2.} Put $j=c_1(J_{\rm lin})=\pi_R^*F$, $d_i=[D_i]$. Since the sections are
disjoint, $d_0d_\infty=0$, and since $F$ is a fibre, $j^2=0$.
Consequently
\[
 c(\nu^*T_Q^{\mathrm{vir}})
 =c(T_R)\frac{(1+j-d_0)(1+j-d_\infty)}{(1+j)^2}
 =c(T_R)\{1-d_0-d_\infty+j(d_0+d_\infty)\}.
\]
The degree-three part has integral
\[
 e(R)-\sum_i\int_{D_i}c_2(T_R|_{D_i})
       +\sum_i\int_{D_i}j\,c_1(T_R|_{D_i}).
\]
We have $e(R)=2e(S)=24$. On a section with normal class
$n_i$, $c_2(T_R|_{D_i})=c_2(S)+c_1(S)n_i$ and
$c_1(T_R|_{D_i})=c_1(S)+n_i$. The two normal classes on
the unglued normalization are $m,-m$, where $m=c_1(M_S)$.
Their sum is zero; $c_1(S)=F$ and $F^2=0$. The displayed
number is therefore $24-24+0=0$. Because
$\nu_*[R]=[Q]$, this is the intrinsic virtual number on $Q$.
The projective-bundle calculation is independently checked
in \texttt{ab\_push\_local\_checks.py}, using
$H(H+m)=0$, $d_0=H+m$, $d_\infty=H$.

\par\noindent\textit{Step 3.} The untwisted quotient away from its cusp has Euler number zero:
scalar multiplication supplies a free circle action, including
the two additive places before logarithmic transformation. Indeed
a circle scalar identified with an iterate of $u\phi$ must be
the identity, since every positive iterate strictly contracts
the Hermitian norm and every negative iterate expands it.
Compact fibres with a free circle action have Euler number zero.
The chosen Mumford fibre has two zero-dimensional torus orbits
and all other torus strata have Euler number zero, so
$e(X(u))=2$. The same number follows independently from the
semistable nearby-cycle count in PUSH A.4.3 in the repository. Lemma~\ref{push:A.3.1}
therefore excludes such a smoothing of $Q$. No requirement
on its leading smoothing section, and no extension of $f$,
entered this argument.

\par\noindent\textit{Step 4.} There is also a direct check that zero is the correct smooth
number. Replace $\phi$ in the global construction of \ref{push:A.1.2}
by a section of $J_{\rm lin}$ vanishing on a smooth elliptic fibre of
$S$. The compact collar construction still applies. At its
zero divisor the fixed nonzero-parameter equation is $xy=uz$
with $dz\ne0$, so every nearby threefold is smooth. The
double-locus nearby Euler function is one on that elliptic
curve and zero off it; its integral is zero. This is an actual
proper-flat smoothing of $Q$, but is not Engel's family.
\end{proof}

\begin{announcement}[PUSH A.3.2, final sentence; A]\label{push:A.3.2.replacement}
The proper nodal
filling of \ref{push:A.1.2} and the two natural semistable fillings of
\ref{push:A.2.2} are precise replacements, with different central fibres.
\end{announcement}
\noindent\emph{Uses:} PUSH A.1.2 and A.2.2.\par
The mechanism changes the central fibre through the nodal construction or one of the two toroidal resolutions described below.

\noindent\emph{Proof:} proof in the repository, PUSH A.1.2 and A.2.2 (\url{https://github.com/KokunoYumeto/s6-complex-structure-workbench}).
\paragraph{Counterexample checklist.}
The target is the exact pinching $Q$ in the quotation. Its normalization
and gluing are fixed in Definition~\ref{push:A.1.setup}; its reduced local
complete-intersection structure is the local equation $xy=0$ used in
the proof. Statement~\ref{push:A.3.1} applies to every proper flat smooth
smoothing of that $Q$, without a projectivity or K\"ahler assumption.
Statement~\ref{push:A.3.2} gives the conflicting counts $0$ and $2$.
The announced replacements change the central fibre, and the nodal
replacement also changes the nonzero fibres by the specified contraction.
\subsection{Announced corrected fillings}
\begin{definition}[PUSH A.1.fan]\label{push:A.1.fan}
At fixed $h_{t_c},h_u\geq0$, the cone $P_{k,l}$ cuts out the rectangle
\[
 kh_{t_c}\leq y_1\leq(k+1)h_{t_c},\qquad
 l(h_{t_c}+h_u)\leq y_2\leq(l+1)(h_{t_c}+h_u).
\]
\end{definition}
\noindent\emph{Uses:} The cusp cocharacter lattice.\par

\begin{announcement}[PUSH A.1.1; A]\label{push:A.1.1}
The cones $P_{k,l}$ in Definition~\ref{push:A.1.fan} and their faces form a fan. Its two
period automorphisms on the cocharacter lattice are
\[
 G_1(y_1,y_2,h_{t_c},h_u)=(y_1+h_{t_c},y_2,h_{t_c},h_u),\qquad
 G_2(y_1,y_2,h_{t_c},h_u)=(y_1,y_2+h_{t_c}+h_u,h_{t_c},h_u).
\]
The affine chart of $P_{0,0}$ is
\[
 U_{0,0}=\operatorname{Spec}\mathbb C[x,y,z,v,u]/(xy-uzv),
 \qquad t_c=zv.
\]
In particular the unsubdivided chart is not a smooth simultaneous
model of the cusp.
\end{announcement}
\noindent\emph{Uses:} PUSH A.1.fan.\par

The mechanism is the rectangle tiling in the cocharacter lattice and the dual-monoid equation $xy=ut_c$.

\noindent\emph{Proof:} proof in the repository, PUSH A.1.1 (\url{https://github.com/KokunoYumeto/s6-complex-structure-workbench}).

\begin{announcement}[PUSH A.1.2; A]\label{push:A.1.2}
There is a reduced complex space $\mathcal X^\flat$ and a proper flat
holomorphic map
\[
 (f,\varpi):\mathcal X^\flat\longrightarrow B\times\Delta_u
\]
whose central threefold is $Q$. For $u\ne0$ its fibre is the
contraction of Engel's chosen $X(u)$ at one rational curve of the
Mumford fibre. This curve is the double curve corresponding to the
diagonal of the square in the chosen $A_2$ triangulation; in the
return's coordinates its direction is $e_1-e_2$. Its normal bundle
in $X(u)$ is $\mathcal O_{\mathbb P^1}(-1)^{\oplus2}$. Its image is
the unique ordinary double point of $\mathcal X^\flat_u$.
\end{announcement}
\noindent\emph{Uses:} PUSH A.1.1; E26 Lemma 2.1 and Corollary 2.2.\par

The mechanism glues a compact annular core to the global node bundle $a\otimes b=u\phi_x$, using the two specified collars. At the nodal elliptic fibre, the chart $xy=uzv$ produces the stated ordinary double point and its prescribed contraction.

\noindent\emph{Proof:} proof in the repository, PUSH A.1.2 (\url{https://github.com/KokunoYumeto/s6-complex-structure-workbench}).

\begin{definition}[PUSH A.2.rays]\label{push:A.2.rays}
Write the six rays of $P_{0,0}$ as
$A=(0,0,1,0)$, $B=(1,0,1,0)$, $C=(0,1,1,0)$,
$D=(1,1,1,0)$, $E=(0,0,0,1)$, $F=(0,1,0,1)$.
The return's $A_2$ diagonal is $BC$.
\end{definition}
\noindent\emph{Uses:} PUSH A.1.1.\par

\begin{announcement}[PUSH A.2.1; A]\label{push:A.2.1}
Among subdivisions with no new rays, the fan in \ref{push:A.1.1} has exactly
two smooth period-invariant subdivisions restricting to the chosen
$A_2$ triangulation at $h_u=0$. These are also all subdivisions
which are flat over the original $(t_c,u)$ bidisc, have reduced
$u$-central divisor, and retain that boundary triangulation.
Without the joint-flatness restriction there are infinitely many
smooth invariant subdivisions retaining that triangulation and
having reduced $u$-central divisor.
\end{announcement}
\noindent\emph{Uses:} PUSH A.1.1 and the retained bottom triangulation.\par

The mechanism enumerates the unimodular triangulations of the six-ray prism and imposes the period identifications. Joint flatness excludes a ray over the origin of the two-dimensional base; without that restriction, the stated period-invariant point stars give arbitrarily many refinements.

\noindent\emph{Proof:} proof in the repository, PUSH A.2.1 (\url{https://github.com/KokunoYumeto/s6-complex-structure-workbench}).

\begin{announcement}[PUSH A.2.2; A]\label{push:A.2.2}
The two minimal subdivisions of \ref{push:A.2.1} give proper flat holomorphic
families $\mathcal X^+,\mathcal X^-\to\Delta_u$ with smooth total
spaces. Their joint maps to $B\times\Delta_u$ are flat. Over
$\Delta_u^*$ they are isomorphic as families to the specified
untwisted family $X(u)$. Each central fibre $Q^\pm$ is reduced,
irreducible, and has normal crossings, with smooth double locus
$S$ and no triple points.

Its normalization $R^\pm$ is obtained by blowing up $R$ along
$S_\infty$ in, respectively, one of its two sections and then
taking the small resolution at the unique ordinary double point
of that blow-up prescribed by the fixed bottom diagonal.
The two proper transforms of the sections are still copies of
$S$, glued by $T$. Their normal-bundle product is trivial;
in particular $\mathcal T^1_{Q^\pm}\simeq\mathcal O_S$.
\end{announcement}
\noindent\emph{Uses:} PUSH A.1.2 and A.2.1.\par

The mechanism applies either of the two period-invariant subdivisions to the proper nodal construction and leaves the other charts unchanged. The central heights determine the reduced normal crossings, and the conductor normal-bundle calculation gives the trivial product.

\noindent\emph{Proof:} proof in the repository, PUSH A.2.2 (\url{https://github.com/KokunoYumeto/s6-complex-structure-workbench}).

\begin{announcement}[PUSH A.2.3; A]\label{push:A.2.3}
The unrestricted request that every normalized component
be a blow-up of $R$ is false. The first maximal-cone
star already contributes a $\mathbb P^3$ component,
which admits no proper bimeromorphic map to $R$, since
such a map would inject a rank-eleven rational $H^2$
into a rank-one group. \ref{push:A.2.4} gives every normalization
by explicit star charts; A.2.5 supplies the terminating
blow-up and blow-down algorithm with centres, order,
exceptional bundles and gluing maps. Thus the literal
shape is refuted and its general algorithmic
replacement is proved.
\end{announcement}
\noindent\emph{Uses:} PUSH A.2.4 and A.2.5.\par

The mechanism is the new complete toric component in a maximal-cone star. Its second Betti number obstructs the proposed bimeromorphic map to the ruled threefold; the finite star-chart description is retained in its place.

\noindent\emph{Proof:} proof in the repository, PUSH A.2.3 (\url{https://github.com/KokunoYumeto/s6-complex-structure-workbench}).

\begin{definition}[PUSH A.2.periods]\label{push:A.2.periods}
Write the period shears as
\begin{align*}
 G_1(a,b,k,\ell)&=(a+k,b,k,\ell),\\
 G_2(a,b,k,\ell)&=(a,b+k+n\ell,k,\ell).
\end{align*}

In a unimodular chart its
coordinates $z_i$ satisfy
\[
 r=\text{unit}\cdot\prod_{i=1}^4 z_i^{\ell_i},
 \qquad t_c=\text{unit}\cdot\prod_{i=1}^4 z_i^{k_i}.
 \tag{A.24d}
\]
\end{definition}
\noindent\emph{Uses:} PUSH A.2.4, finite fan coordinates.\par

\begin{announcement}[PUSH A.2.4; A]\label{push:A.2.4}
The following finite-data procedure parametrizes all smooth
period-invariant subdivisions of the cusp fan after
$u=r^n$, retaining the prescribed bottom triangulation and
having finitely many cones modulo the periods. Here $n=1$
gives PUSH Part A, and $n=2,12$ give the two cases of PUSH B.2.
Each output defines a proper flat family over $\Delta_r$
with smooth total space, whose punctured restriction is the
given untwisted family pulled back by $u=r^n$.

For a bound $H\geq1$, retain the four bottom rays
$A,B,C,D$ and consider all primitive vectors
\[
 \begin{gathered}
 v=(a,b,k,\ell),\qquad
 \ell\geq1,\quad k\geq0,\quad k+\ell\leq H,\\
 0\leq a\leq k,\qquad0\leq b\leq k+n\ell.
 \end{gathered}
 \tag{A.24a}
\]
Take any collection of four-element subsets of this finite
set satisfying the following exact tests:
\begin{enumerate}
\item Each four-column determinant has absolute value one.
\item After dividing each ray by $k+\ell$, the resulting
tetrahedra intersect pairwise in common faces, and their
volumes sum to $(n+2)/6$.
\item Their bottom boundary is the two triangles $ABC,BCD$.
Their boundary subdivisions on $a=0$ and $a=k$ agree under
$G_1$, and those on $b=0$ and $b=k+n\ell$ agree under $G_2$.
\end{enumerate}
All these are finite rational-arithmetic tests. Extend a
passing collection by the two periods, and discard duplicates
as $H$ increases. This is exhaustive. For a reduced central
divisor, retain only $\ell=1$ in (A.24a).

The central prime components are indexed by period orbits of
the selected rays with $\ell>0$, and the multiplicity of the
component indexed by $v$ is exactly $\ell$. The normalization
charts of that component are the star fan of $v$ in
$N/\mathbb Zv$, with the inherited bundle units and the
stabilizer action. For $k>0$ the stabilizer is trivial and this
normalization is a smooth complete toric threefold. All
multiple loci, including possible triple and quadruple loci,
are specified by the positive-$\ell$ rays belonging to each
cone. In the reduced case
$\mathcal T^1\simeq\mathcal O_{\operatorname{Sing}D}$,
so every output is d-semistable.
\end{announcement}
\noindent\emph{Uses:} PUSH A.1.1 and the explicit base-changed period shears.\par

The mechanism enumerates a finite height-bounded collection of primitive rays, tests unimodularity and rational face intersections, and imposes the two period identifications. Increasing the bound exhausts the finite invariant subdivisions, while the star fans record the central normalizations and their multiplicities.

\noindent\emph{Proof:} proof in the repository, PUSH A.2.4 (\url{https://github.com/KokunoYumeto/s6-complex-structure-workbench}).

\begin{announcement}[PUSH A.2.5; A]\label{push:A.2.5}
The requirement that every normalized component of every
smooth refinement be a blow-up of
$R=\mathbb P_S(\mathcal O\oplus M_S)$ is false, already
for a reduced refinement with $n=1$.
The following is a uniform replacement for any finite
fan certificate $\Sigma$ in \ref{push:A.2.4}.

The star charts and integer transition matrices of \ref{push:A.2.4}
give every normalized component, its gluing maps, and
all multiple loci directly. There is also a terminating
finite-data algorithm that expresses its comparison
with a reference model as an explicit sequence of
blow-ups and blow-downs along specified smooth strata.
At each blow-up the algorithm gives the normalized
strict transforms, the exceptional projective bundle,
its multiplicity, and every new intersection.
This algorithm allows blow-downs; no toric
strong-factorization assertion is used.

For a component of multiplicity $\ell$, the local
nonreduced central equation is the monomial in A.24d.
More generally, in a unimodular chart with central
equation $F=\prod_i z_i^{\ell_i}$, the cotangent-complex
first-order deformation sheaf is
\[
 \mathcal T^1_D
  =\mathcal O_D\Big/
    \left(\ell_i z_i^{\ell_i-1}
          \prod_{j\ne i}z_j^{\ell_j}:\ell_i>0\right).
 \tag{A.25a}
\]
Here $\mathcal T^1_D$ means
$\mathcal Ext^1(\mathbb L_D,\mathcal O_D)$; in the
reduced case it also equals
$\mathcal Ext^1(\Omega_D,\mathcal O_D)$.
The reduced case is the trivial line on the singular
scheme, as in \ref{push:A.2.4}. For a nonreduced centre the displayed
Jacobian quotient is the appropriate deformation sheaf;
the reduced normal-crossings notion of d-semistability
is not asserted for it.
\end{announcement}
\noindent\emph{Uses:} PUSH A.2.4 and its marked star charts.\par

The mechanism compares finite marked cone complexes through ordinary star blow-ups and their inverse moves. The local monomial Jacobian quotient records the deformation sheaf with its stated nonreduced scope.

\noindent\emph{Proof:} proof in the repository, PUSH A.2.5 (\url{https://github.com/KokunoYumeto/s6-complex-structure-workbench}).

\section{Announcements at the puncture and in the space of structures}
\label{sec:announcements}
\subsection{The fixed-family noncompletion statement}
\begin{definition}[PUSH B.4.33.scope]\label{push:B.4.33.scope}
The retained sphere family has multiplicities $(3,4)$ and the fixed data
$\epsilon=(1,2,-4,0)^t$, $\epsilon'=(1,3,-3,0)^t$,
$v_1=\epsilon$, $v_2=-\epsilon'$, $l_0=0$, and $p=12l_0-4l_1-3l_2=-1$.
The untwisted family is
$X(u)=M_S^{\times,\circ}/\langle u\phi\rangle$, completed at infinity
by its Mumford fibre; the sphere family is $X_w$.
\end{definition}
\noindent\emph{Uses:} The first-run retained data.\par

\begin{announcement}[PUSH B.4.33; A]\label{push:B.4.33}
The actual punctured six-sphere family has no proper
Hausdorff complex-analytic filling after any finite
cover of its degeneration parameter. In particular,
there is no positive integer $N$ for which the cover
$w=r^N$ has a proper filling with smooth total space
and reduced normal-crossings central fibre.

The assertion permits an arbitrary singular central
space, arbitrary additional components, and arbitrary
central multiplicities. It does not impose a
K\"ahler, projective, or logarithmic hypothesis on
the putative filling. Its proof derives the specific
logarithmic torsor needed for the contradiction from
differences of sections extending into that filling.
\end{announcement}
\noindent\emph{Uses:} PUSH B.4.31 and B.4.32; the actual fixed punctured sphere family.\par

The mechanism passes from a proposed filling to a marked proper map and controls the radial-loop class through differences of extending sections. The resulting specific logarithmic torsor contradicts the nonzero degree; an arbitrary affine origin is not asserted to extend.

\noindent\emph{Proof:} proof in the repository, PUSH B.4.33 (\url{https://github.com/KokunoYumeto/s6-complex-structure-workbench}).

The assertion has the fixed-family scope just stated. Whether the same
noncompletion holds for every admissible twist remains a question.
The proposed reading of a monodromy integer as $J=P\mathbin{.}O$ is a
reader's remark recorded in the intake; it is unverified and is not a
theorem or a premise of this paper.
\subsection{Smooth monodromy: two runs and two methods}
\begin{definition}[PUSH C.orientation]\label{push:C.orientation}
The six-sphere has its complex orientation and a positive $w$-circle
is counterclockwise. We use
$M_\Phi=(S^6\times[0,1])/((x,1)\sim(\Phi(x),0))$ with orientation
$dt\wedge\operatorname{or}(S^6)$, and
$\Sigma_\Phi=D^7\cup_\Phi(-D^7)$ for the identification with
$\Theta_7$; its positive generator has Eells--Kuiper invariant $1/28$.
\end{definition}
\noindent\emph{Uses:} The complex orientation and the positive parameter circle.\par

\begin{definition}[PUSH C.4.1.scope]\label{push:C.4.1.scope}
The twist comparison used in PUSH C.3.24 has fixed multiplicities $(3,4)$ and the admissible sphere data $|p|=1$ of FAM B.2.1. Its monodromy convention is the positive $w$-circle with the complex orientation on every fibre.
\end{definition}
\noindent\emph{Uses:} FAM B.2.1; PUSH C.4.1.\par

\begin{announcement}[PUSH C.3.24; A]\label{push:C.3.24}
For the actual retained $S^6$ family and a positive generator
of its punctured $w$-parameter disc, the oriented smooth
monodromy satisfies
\[
 [\Phi]=0\quad\hbox{in }\pi_0\operatorname{Diff}^+(S^6)
                       \simeq\Theta_7\simeq\Z/28.
\]
Its mapping torus is consequently oriented-diffeomorphic to
$S^1\times S^6$. Every finite cyclic base change has trivial
smooth monodromy as well. By C.4.1's marked family comparison,
the same conclusion applies to the admissible sphere twist
data covered by that comparison.
\end{announcement}
\noindent\emph{Uses:} PUSH C.3.18, C.3.21--C.3.23 and C.4.1.\par

The mechanism is an integer calculation before reduction modulo $28$: the source gives $p_{\rm rel}=-24$, $s_\delta=-2$ and $J_\psi=\epsilon_{\rm cap}(-24-12(-2))/24=0$. Additivity compares the six-loop representative to six times the one-loop integer.

\noindent\emph{Proof:} proof in the repository, PUSH C.3.24 (\url{https://github.com/KokunoYumeto/s6-complex-structure-workbench}).

\begin{announcement}[PUSH C.3.27; A]\label{push:C.3.27}
The actual positive-loop smooth monodromy satisfies
$[\Phi]=0\in\Theta_7$. This determination uses the explicit
eight-dimensional coboundary of C.3.25 and the direct index
additivity of C.3.26. It does not use the three-knot integer,
Wall's formula for it, the Davis generator comparison, or the
conclusion of \ref{push:C.3.24}.
\end{announcement}
\noindent\emph{Uses:} PUSH C.3.25 and C.3.26; the shared characteristic calculations C.3.18--C.3.22.\par

The mechanism uses the explicit eight-dimensional spin$^c$ coboundary and the integer index $\lambda=(p_\psi-4k_\psi)/24$, with $p_\psi=-24\epsilon_{\rm cap}$ and $k_\psi=-6\epsilon_{\rm cap}$. The boundary-invariant and one-loop arguments differ from the orbit-knot method; the actual annular representative and characteristic inputs are shared.

\noindent\emph{Proof:} proof in the repository, PUSH C.3.27 (\url{https://github.com/KokunoYumeto/s6-complex-structure-workbench}).

\begin{announcement}[MAPS B.6.21; A]\label{maps:B.6.21}
For the original $(3,4)$ family, with its complex orientation,
the positive $w$ loop and the generator of $\Theta_7$ fixed by
$\mu=1/28$, the monodromy satisfies
\[
                         [\Phi]=0\in\mathbb Z/28.
 \tag{B.6.38}
\]
It is therefore smoothly isotopic to the identity.  Its sixfold
class is zero, and among the two resulting K7 residues $0,14$
the original class is $0$.  This statement does not yet assert
the same value for the other multiplicity pairs.
\end{announcement}
\noindent\emph{Uses:} MAPS B.6.15--B.6.20 and the specified mapping-torus comparison.\par

The mechanism uses the free-circle replacement with its actual peripheral identifications, the values $k=-1$ and $\ell=-4$, and the parametrized spin-surgery trace. Its normalized correction is zero, so the specified Eells--Kuiper invariant selects the original residue $0$.

\noindent\emph{Proof:} proof in the repository, MAPS B.6.21 (\url{https://github.com/KokunoYumeto/s6-complex-structure-workbench}).

\begin{definition}[PUSH E.1.1.scope]\label{push:E.1.1.scope}
For the higher-family announcements, $m=3k$, $n=4l$,
$\gcd(m,n)=1$, and $k,l>0$. A fixed sphere twist datum in the K1 scope
of MAPS means fixed integers $\rho_0,N_0,\rho_1,N_1$ with
\[
 w_1=N_0\varepsilon-\rho_0\widehat\delta,\qquad
 w_2=N_1\varepsilon'-(\rho_1+N_1)\widehat\delta,\qquad l_0=0,
\]
and $nN_0+mN_1=\pm1$, in the admissible period domain. The vectors are
the first-run clockwise-root vectors; the original K1 data use the
transported counterclockwise-root marking. The parameter period is
$2kl\mathbb Z$.
\end{definition}
\noindent\emph{Uses:} PUSH E.1.1--E.1.3; FAM B.3.1.\par

\begin{announcement}[MAPS B.6.24; A]\label{maps:B.6.24}
For every admissible coprime sphere pair $(3k,4l)$ and every
fixed twist datum of K1, the monodromy for the positive
generator $2kl$ of its parameter-period lattice satisfies
\[
                    [\Phi_{3k,4l}]=0\in\Theta_7.
 \tag{B.6.43}
\]
The identifications and orientation are those specified in
the question.  In particular this completes the numerical
part of B.6 for both the original and higher families.
\end{announcement}
\noindent\emph{Uses:} The fixed K1 twist datum; MAPS B.6.19 and B.6.22--B.6.23.\par

The mechanism is the higher-family free-circle quotient calculation $k_{\rm cup}=-kl$, $\ell=-4kl$, followed by the signed spin-surgery trace and the mapping-torus comparison. This is the source statement for every fixed K1 datum in the stated scope.

\noindent\emph{Proof:} proof in the repository, MAPS B.6.24 (\url{https://github.com/KokunoYumeto/s6-complex-structure-workbench}).

The last announcement has a broader scope than MAPS B.6.21, whose
last sentence is retained. The supplied open-problem list also calls
``the monodromy of the other families'' open. Those wordings are not
silently identified: MAPS B.6.24 announces zero for the explicitly
specified coprime sphere pairs and fixed K1 data, while an open question
beyond that scope is not answered here.
\subsection{Families and components of the space of structures}
\begin{definition}[PUSH E.2.5.scope]\label{push:E.2.5.scope}
The smooth-completion hypotheses used below are those of PUSH E.2.5:
$\pi:\mathcal X\to\Delta$ is a proper holomorphic submersion,
all fibres are diffeomorphic to $S^6$, and every noncentral fibre belongs
to one fixed coprime multiplicity family $(3k,4l)$ of PUSH E.2.1.
The integers in that statement are $A=12kl$ and $D=kl-3k-4l\ne0$.
For fixed discrete logarithmic data, the isomorphism parameter of
PUSH E.2.7 is denoted here by
\[
 w_{k,l}=\exp\!\left(-\frac{\pi i c}{kl}\right).
\]
The symbol $z$ in those source statements denotes this parameter,
not the uniformizing coordinate of the first-run conventions.
\end{definition}
\noindent\emph{Uses:} PUSH E.2.5 and E.2.7.\par

\begin{announcement}[PUSH E.2.21; A]\label{push:E.2.21}
Let $\pi:\mathcal X\to\Delta$ satisfy the hypotheses of
Definition~\ref{push:E.2.5.scope}, with fixed coprime multiplicities $(3k,4l)$ on
every noncentral fibre.  Then $X_0$ is a member of that
same constructed family.  Its isomorphism parameter is
the interior value $w_{k,l}(0)$ of PUSH E.2.7.

After choosing a holomorphic lift $c(s)$ of that parameter
and shrinking the disc, the actual family is isomorphic
to the prescribed family at $c(s)$.  The ambiguity in
this comparison is the connected translation automorphism
group already determined in E.1.
\end{announcement}
\noindent\emph{Uses:} PUSH E.2.5, E.2.7 and E.2.10--E.2.20.\par

The mechanism extends the intrinsic pencil and its nonzero interior period, identifies the ordinary marked Jacobian, and glues the finite-root and cusp comparisons using the connected translation group.

\noindent\emph{Proof:} proof in the repository, PUSH E.2.21 (\url{https://github.com/KokunoYumeto/s6-complex-structure-workbench}).

\begin{announcement}[PUSH E.2.22; A]\label{push:E.2.22}
Let $\pi:\mathcal X\to\Delta$ be a proper holomorphic
submersion with fibres diffeomorphic to $S^6$.  Suppose
there is a nonempty open subset $U\subset\Delta$ such
that every $X_s$, $s\in U$, is a member of the fixed
coprime multiplicity family $(3k,4l)$ of E.2.1.
Then every fibre over $\Delta$ is a member of that
same family.
\end{announcement}
\noindent\emph{Uses:} PUSH E.2.1 and E.2.21.\par

The mechanism propagates the geometric data outside successive discrete analytic exceptional sets and then fills those sets in reverse order by the smooth-completion comparison.

\noindent\emph{Proof:} proof in the repository, PUSH E.2.22 (\url{https://github.com/KokunoYumeto/s6-complex-structure-workbench}).

\begin{announcement}[PUSH E.2.23; A]\label{push:E.2.23}
In a proper smooth holomorphic family of compact complex
manifolds over a connected reduced complex analytic base,
if one fibre is a sphere member of E.2.1 with coprime
multiplicities $(3k,4l)$, then every fibre is a constructed
member with that same pair.

Consequently two of the constructed sphere members are
deformation equivalent, allowing chains of proper smooth
families over connected complex analytic bases, if and
only if their ordered multiplicity pairs agree.  In
particular the distinct coprime pairs give distinct
deformation classes.
\end{announcement}
\noindent\emph{Uses:} PUSH E.2.21 and E.2.22; local miniversality.\par

The mechanism combines propagation on Riemann surfaces with the local analytic-base argument. It uses the constructed fixed-pair locus and does not assume that a punctured analytic disc can be selected inside an arbitrary open set.

\noindent\emph{Proof:} proof in the repository, PUSH E.2.23 (\url{https://github.com/KokunoYumeto/s6-complex-structure-workbench}).

\begin{announcement}[PUSH E.2.24; A]\label{push:E.2.24}
Let $\mathcal I(S^6)$ be the space of smooth integrable
almost-complex structures on the fixed smooth sphere,
with its $C^\infty$ topology.  For a fixed coprime pair
$(3k,4l)$, let $\mathcal I_{k,l}$ consist of the structures
whose complex manifolds are biholomorphic to constructed
members with that pair.  Then $\mathcal I_{k,l}$ is open
and closed in $\mathcal I(S^6)$.

Consequently a continuous family of such structures
over any connected topological parameter space cannot
contain constructed members with two different pairs.
The same separation holds in the quotient spaces by
$\operatorname{Diff}_0(S^6)$ or by all orientation-preserving
diffeomorphisms, with their quotient topologies.
This statement does not identify the number of marked
connected components belonging to a single pair.
\end{announcement}
\noindent\emph{Uses:} PUSH E.2.23 and local completeness of the deformation charts.\par

The mechanism transfers the analytic-family separation to the local deformation charts in the smooth topology of \cite[\S2.1 and \S3.2, Theorem~3.4 and Remark~3.6]{Meersseman2019E}, and then to the indicated quotient spaces. It supplies no count of marked connected components within a single pair.

\noindent\emph{Proof:} proof in the repository, PUSH E.2.24 (\url{https://github.com/KokunoYumeto/s6-complex-structure-workbench}).


\section{Open questions}\label{sec:open}
The supplied status files leave the following questions open. We retain their scope rather than adding a conjectural answer.
\begin{enumerate}[leftmargin=*]
\item Does $S^6$ admit a complex structure of algebraic dimension $0$? The present period families have algebraic dimension $1$.
\item What is the smooth knot type of the real $S^3$ in $S^6$ in the real-form announcements of \S7?
\item What geometric object should replace a proper analytic limit at $q_2=0$? Does the asserted non-filling extend from the sphere family to every admissible free twist, including $|p|>1$, and to the higher-multiplicity families?
\item What is the smooth monodromy for the other constructions, including the two-parameter, comb and Liu--Rollenske families? MAPS B.6.24 already announces zero for its admissible coprime $(3k,4l)$ scope; this question concerns the families beyond that stated scope. The reading $J=\epsilon(P\cdot O)$ on a page surface is a reader's unverified remark, not a theorem of this paper.
\end{enumerate}
\noindent\textit{Sources:} \texttt{INTAKE\_FINAL.md}, \S6; \texttt{RESPONSE\_INTAKE.md}, \S\S1--2; \texttt{RESPONSE\_EXPERT.md}, \S4. The open questions are not strengthened versions of the announcements.

\section*{Limits and open points}\label{sec:limits}
\begin{itemize}[leftmargin=*]
\item \textbf{Unverified in this assembly: certificate files.} FAM C.2.3 and PUSH E.4 are retained with their binding T status. The supplied reviews report the interval re-execution and independent-library certification. The attached bundles omit the certificate data and verifier programs, so this assembly cannot supply their bytes, hashes or a fresh execution. The exact recorded counts and arithmetic assumptions are printed in \S4. No new numerical grid is substituted for a certificate.
\item \textbf{Weak point: scope of a known theorem.} FAM A.1.1 includes every admissible affine translation and cusp collar. DGKX Theorem 7.1 supplies the standard-data intrinsic pencil; its statement does not have that extra scope. We retain FAM's exact statement and proof and distinguish the two scopes in the credit table.
\item \textbf{Not proved here: announcements.} Every A-item has its exact source statement and a repository pointer. The higher-multiplicity dictionary, metric restrictions, square-zero product, corrected fillings, zero-end non-filling, monodromy-class and component assertions keep this status. The common producer is not an independent reader.
\item \textbf{Conditional dependency retained: covers.} The source proof of INV B3.2 invokes the higher-multiplicity Hodge calculation INV B1.2, which the binding table assigns status A. We retain both assignments and this dependency; no replacement proof is supplied. The unread local direct-image inputs INV A1.1--A1.2 and the ladder inputs ARI C.8.1--C.8.3 remain limitations of the historical review record.
\item \textbf{Source conflict: known content and proof placement.} The specific binding T rows and full-proof requirements include results already present in AC26 or standard torus theory, including parts of INV A1, FAM B.1.3 and FAM D.1--D.2. We credit those sources and retain the requested proofs in their prescribed sections. This differs from the general format rule placing every new proof of a known statement in an appendix; no claim that the known statement is new is made.
\item \textbf{Source conflict: quotations.} The request for the literal sentences of AC26 Remarks 1.3 and 9.21 conflicts with the prohibition on reproducing AC26 text. Their locators, the precise corrected assertions and the reason are printed, with a labelled paraphrase rather than a quotation. Engel's relevant wording is quoted from the supplied source.
\item \textbf{Source references and order.} The fixed section order leaves explicit references to supporting re-derivations in Appendix B and to the maximum in \S4. The dependency table records these rather than silently changing the source arguments. The all-maps-to-$\mathbb P^1$ companion is DGKX Theorem 7.2, not a separate assertion in FAM A.1.1; its supplied proof is credited accordingly. The standard base-change and duality citations replace the two unsuitable DGKX lemma citations in FAM A.5.2.
\item \textbf{Unverified references and interpretation.} Gray's article was not obtained; it is not used as an input. The repository URL is supplied by the prompt, but public retrieval failed during assembly; announcement pointers specify only the supplied run and statement number. The reader's $J=\epsilon(P\cdot O)$ interpretation remains unverified. The complete discrepancy and source-search record accompanies the paper.
\end{itemize}

\section*{Verification record}
The long record is in Appendix~\ref{app:verification}, one row per check. Historical rows are transcribed in scope from the six supplied status files. They distinguish a full read, a derivation, a numerical or exact computation, and a source opening; a statement read without its proof is not labelled a proof check. Current S-rows record the citations opened and matched for this assembly. The source-to-print ledger and the PDF audit are separate reproducibility records. No human has checked a statement or proof, and no proof has been formalized.

\section*{Provenance}
The user's writing prompt, \texttt{PROMPT(20261116-143923).md}, selected the statements, statuses, proof placement and title. The six returns were AI-produced by GPT-6 Astra (ChatGPT Pro mode); the supplied workbench proof was AI-produced by GPT-5.6 Sol and GPT-6 Astra (Codex reasoning effort ultra). GPT-6 Astra (ChatGPT Pro mode) assembled the text, applied the prescribed notation, moved the K-proofs, removed the excluded material, opened the cited sources and recorded every statement comparison. The source table is the numbering concordance. Input hashes, extraction scripts, comparison records, bibliography checks and build logs accompany this paper. The supplied source bundles remain the inputs; they are not reproduced in the deliverable. Repository: \url{https://github.com/KokunoYumeto/s6-complex-structure-workbench}.

\section*{Declarations}
\textbf{AI use.} This paper is AI-produced by GPT-6 Astra (OpenAI, ChatGPT Pro mode). The checking roles of Claude Fable 5.2 (effort high), Claude Fable 5.2 (effort max) and Claude Opus 5.5 (effort max) are exactly those in the verification record. The workbench contribution of GPT-5.6 Sol and GPT-6 Astra (Codex reasoning effort ultra) is identified in Appendix A. No human authorship or checking is attributed. \textbf{Human understanding: none. Formalization: none.} The text produced here is dedicated to the public domain under CC0 1.0; source works retain their own terms. The original construction and all known results are credited by statement and section. This is the assembled version described by the accompanying source and discrepancy records.


\appendix
\section{The workbench period-shift proof}\label{app:workbench}

The proof in this appendix is the supplied workbench fragment, produced by GPT-5.6 Sol (Codex, reasoning effort ultra) and GPT-6 Astra (Codex, reasoning effort ultra). Its mathematical content is retained. Reference labels for the canonical pencil, finite normalization and period sheaf are mapped to their supplied FAM counterparts; the concordance records these conversions.
\subsection{Exact global identifications of the retained period parameter}
\label{pshift:sec:classification}

Keep the original base $B$, coordinate $t$, values
$t(p_1)=0,t(p_2)=1,t(p_0)=\infty$, and the original functions
$\tau,\mu,\beta^{\mathrm{part}}$. Write
\[
 \beta_c=\beta^{\mathrm{part}}+c,\qquad
 \mathcal P=\{c\in\mathbb C:\operatorname{Im}c<-M\},\qquad
 \Pi_c(z)=\begin{pmatrix}6\mu(z)&\tau(z)&1&0\\
                      \beta_c(z)&\mu(z)&0&1\end{pmatrix}.
 \tag{PS1}\label{pshift:eq:periods}
\]
All assertions below concern the analytic pieces and their actual gluing.
No integral cohomology, sphere recognition, or CDP assertion is an input.
In particular, the global parameter classification is not inferred from
an equality between two fibre lattices.

Put $e=(0,1)^{\mathsf t}\in\mathbb C^2$ and retain the ordered lattice
\[
 \Lambda=\mathbb Z\langle\widehat\gamma,\widehat u,
                        \widehat w,\widehat\delta\rangle,
 \qquad E=\widehat\delta\otimes\gamma.
 \tag{PS2}\label{pshift:eq:rank-one-map}
\]
Thus $E$ is the matrix with its only nonzero entry, equal to one, in
row four and column one; $E^2=0$ and $\Pi_cE=e\gamma$.
The finite data, with every sign retained, are
\[
\begin{gathered}
 A_1=\begin{pmatrix}1&0&0&0\\6&0&1&0\\-6&-1&-1&0\\-2&1&0&1\end{pmatrix},
 \qquad
 A_2=\begin{pmatrix}1&0&0&0\\0&0&-1&0\\-6&1&0&0\\3&0&1&1\end{pmatrix},\\
 v_1=(1,2,-4,0)^{\mathsf t},\quad
 v_2=(-1,-3,3,0)^{\mathsf t},\quad
 (m_1,m_2)=(3,4),\quad
 (\epsilon_1,\epsilon_2)=(1,-1),\quad \epsilon_j=\gamma(v_j).
\end{gathered}
\tag{PS3}\label{pshift:eq:finite-data}
\]
The covariance is $R_g\Pi_c=\Pi_c(gz)A_g$, with
\[
 R_{g_1}=\begin{pmatrix}-1/\tau&0\\(1-\mu)/\tau&1\end{pmatrix},\qquad
 R_{g_2}=\begin{pmatrix}1/\tau&0\\-\mu/\tau&1\end{pmatrix},\qquad
 R_{g_0}=I_2.
 \tag{PS4}\label{pshift:eq:covariance}
\]
In particular $R_ge=e$ for every $g$.

\begin{cited}[WORKBENCH; K: The exact identification subgroup is $2\mathbb Z$]\label{pshift:thm:classification}
For $c,c'\in\mathcal P$, the retained analytic threefolds are
biholomorphic if and only if
\[
                         c'-c\in2\mathbb Z.
 \tag{PS5}\label{pshift:eq:classification}
\]
Every such biholomorphism preserves the original map to $B$. For every
even integer $n$, the proof below constructs an isomorphism
$X_c\longrightarrow X_{c+n}$ whose complex-linear part on the original
universal fibre coordinates is $I_2$. Its map on the marked fibre lattice
is $I-nE$.

The equality of the middle lattices for all integer shifts loses a
nonzero order-two finite-filling invariant: the shift $n$ changes that
invariant by $n\bmod2$. The cusp contributes no additional obstruction
to the displayed even-shift construction.
\end{cited}
\uses{The period data (PS1)--(PS4); the numbered steps below and the retained construction.}


\subsubsection{The middle family and its complete rational centralizer}

For an integer $n$, put $S_n=I+nE$. Since $S_n^{-1}=S_{-n}$,
\[
 \Pi_{c+n}=\Pi_c S_n,\qquad
 \Pi_{c+n}\Lambda=\Pi_c\Lambda.
 \tag{PS6}\label{pshift:eq:integer-shear}
\]
Both original matrices fix $\widehat\delta$ and preserve $\gamma$.
Consequently $A_jE=EA_j=E$, and the same commutation holds for every
monodromy word. The identity of $\mathbb C^2$ therefore identifies
the two full middle torus families with their actual $R_g$-action.
Its source-to-target lattice map is $S_{-n}$, since
$\Pi_{c+n}S_{-n}=\Pi_c$.

\begin{cited}[WORKBENCH; K: All middle-family linear parts]\label{pshift:lem:centralizer}
The simultaneous rational centralizer of $A_1,A_2$ is exactly
\[
 \{aI+bE:a,b\in\mathbb Q\}.
 \tag{PS7}\label{pshift:eq:centralizer}
\]
A biholomorphism between the middle families for $c,c'$ over the
identity of $B^\circ$ forces $c'-c\in\mathbb Z$. Its complex-linear
part on every fibre is either $I_2$ or $-I_2$.
\end{cited}
\uses{The period data (PS1)--(PS4); the numbered steps below and the retained construction.}


\begin{proof}
\par\smallskip\noindent\textbf{Step 1.} Solving the two fixed-vector systems from (PS3) gives
\[
 (\Lambda\otimes\mathbb Q)^{A_1}
   =\mathbb Q(1,2,-4,0)^{\mathsf t}+\mathbb Q\widehat\delta,
 \quad
 (\Lambda\otimes\mathbb Q)^{A_2}
   =\mathbb Q(1,3,-3,0)^{\mathsf t}+\mathbb Q\widehat\delta.
\]
Their intersection is $\mathbb Q\widehat\delta$. The fixed covectors
are respectively
\[
 \mathbb Q\gamma+\mathbb Q(2u_{\Lambda}+w_{\Lambda}+3\delta),\qquad
 \mathbb Q\gamma+\mathbb Q(u_{\Lambda}+w_{\Lambda}+2\delta).
\]
Their intersection is $\mathbb Q\gamma$: comparing the $u_{\Lambda}$ and $w_{\Lambda}$
coefficients of a common multiple gives $2x=y$ and $x=y$, hence
$x=y=0$. Every commuting endomorphism therefore preserves
$\mathbb Q\widehat\delta$ and $\ker\gamma$.

\par\smallskip\noindent\textbf{Step 2.} On $\ker\gamma/\mathbb Q\widehat\delta$, in the classes of
$(\widehat u,\widehat w)$, the two matrices are
\[
 Q_1=\begin{pmatrix}0&1\\-1&-1\end{pmatrix},\qquad
 Q_2=\begin{pmatrix}0&-1\\1&0\end{pmatrix}.
\]
A matrix commuting with $Q_2$ has the form
$\left(\begin{smallmatrix}b&d\\-d&b\end{smallmatrix}\right)$.
Commutation with $Q_1$ gives $d=0$. Thus a commuting endomorphism $H$
has columns
\[
 H\widehat\gamma=(a,r,s,k)^{\mathsf t},\quad
 H\widehat u=b\widehat u+u_0\widehat\delta,\quad
 H\widehat w=b\widehat w+w_0\widehat\delta,\quad
 H\widehat\delta=d_0\widehat\delta.
\]
Use the actual columns $A_1\widehat u=-\widehat w+\widehat\delta$,
$A_1\widehat w=\widehat u-\widehat w$, and
$A_2\widehat u=\widehat w$. Commutation on these columns gives
\[
 d_0-w_0=b+u_0,\qquad u_0=2w_0,\qquad w_0=u_0.
\]
Hence $u_0=w_0=0$ and $d_0=b$. Commutation on
$\widehat\gamma$ with $A_1$ gives
\[
 r=2(a-b),\qquad s=-4(a-b).
\]
Commutation on that column with $A_2$ also gives $r=-s$.
Therefore $a=b$, $r=s=0$, and $H=aI+kE$. Conversely these matrices
commute, by the fixed-vector and fixed-covector identities preceding
(PS6). This proves (PS7).

\par\smallskip\noindent\textbf{Step 3.} For the geometric assertion, a holomorphic map between two compact
complex tori lifts to their universal covers. The derivative of the
lift is periodic, so its entries descend to holomorphic functions on
the compact source torus and are constant. The lift is consequently
affine, with a complex-linear part $L$. In a holomorphic family these
linear parts vary holomorphically. Their lattice maps are integral
and locally constant, hence give an isomorphism of the two marked
local systems. Because the map is over the identity of the base,
that integral map $H$ commutes with both $A_j$. Thus
$H=aI+bE$, with $a,b\in\mathbb Z$. Its determinant is $a^4$;
invertibility over $\mathbb Z$ gives $a=1$ or $a=-1$.

\par\smallskip\noindent\textbf{Step 4.} The complete period equation is
\[
                   L(z)\Pi_c(z)=\Pi_{c'}(z)(aI+bE).
 \tag{PS8}\label{pshift:eq:period-morphism}
\]
Its third and fourth columns give $L(z)=aI_2$. Its first column
then gives $a(c'-c)+b=0$. Therefore $c'-c=-b/a\in\mathbb Z$,
and the asserted two linear parts exhaust the possibilities.
\end{proof}

\subsubsection{The exact finite obstruction, including normalization}

The two cyclic norm matrices are, by multiplication of (PS3),
\[
 N_1=I+A_1+A_1^2=
 \begin{pmatrix}3&0&0&0\\6&0&0&0\\-12&0&0&0\\0&2&1&3\end{pmatrix},
\]
\[
 N_2=I+A_2+A_2^2+A_2^3=
 \begin{pmatrix}4&0&0&0\\12&0&0&0\\-12&0&0&0\\0&2&2&4\end{pmatrix}.
 \tag{PS9}\label{pshift:eq:norms}
\]
For example these formulas give
$N_1\widehat w=\widehat\delta$,
$N_2\widehat u=N_2\widehat w=2\widehat\delta$,
$N_1\widehat\delta=3\widehat\delta$, and
$N_2\widehat\delta=4\widehat\delta$.

We first justify why a local isomorphism must satisfy the exact affine
equations on the original finite cover. The actual ramified base change
$t_j=s_{\rm new}^{m_j}$ has local equation
\[
 s_j^{m_j}-s_{\rm new}^{m_j}
   =\prod_{a=0}^{m_j-1}(s_j-e^{2\pi ia/m_j}s_{\rm new})=0.
\]
Its normalization separates all these smooth branches, each with
coordinates $(z_1,z_2,s_{\rm new})$. The original map
$y\mapsto(q_j(y),s_j(y))$ identifies the original torus family $J'_j$
with this normalization, as proved by the explicit branches in
\cref{fam:A.1.2}. A biholomorphism over the identity of
$D_j$ base-changes over the identity of the $s_{\rm new}$-disc and lifts uniquely
to the normalization: on local rings it carries integral elements in
the total fraction ring to integral elements, and its inverse does
the same. This lift commutes with the original deck action because
both composites lift the same map $(x,s_{\rm new})\mapsto(F(x),\zeta_j s_{\rm new})$.
In particular one may not replace the generator by an arbitrary
coprime power.

Let $c'=c+n$, where $n$ is an integer, and suppose the linear part of
the lifted map is $aI_2$, $a\in\{1,-1\}$. The lifted map on the torus
family is $[z,\zeta]\mapsto[z,a\zeta+B_j(z)]$. Its translating
section has a holomorphic lift $B_j:\Delta_j\to\mathbb C^2$:
$\Delta_j\times\mathbb C^2\to J'_j$ is a covering map and the disc is
simply connected. Local inverses of this covering show that the lift
is holomorphic. The deck-equivariance identity holds modulo an
actual lattice vector, so, after choosing this lift, it is
\[
 B_j(g_jz)-R_{g_j}(z)B_j(z)
  =\Pi_c(g_jz)\left(
       \frac{(1-a)v_j+n\epsilon_j\widehat\delta}{m_j}
          +\ell_j\right),\qquad \ell_j\in\Lambda.
 \tag{PS10}\label{pshift:eq:finite-equation}
\]
The integer column $\ell_j$ is constant: it is locally constant by
discreteness of the marked lattice and the disc is connected.
Iterating $m_j$ times, with $A_jv_j=v_j$ and
$A_j\widehat\delta=\widehat\delta$, yields the necessary equality
\[
 N_j\ell_j=-(1-a)v_j-n\epsilon_j\widehat\delta.
 \tag{PS11}\label{pshift:eq:finite-norm-obstruction}
\]
Here the left side of the iterated equation is zero because the
base point and the original vector coordinate return after $m_j$
steps; real invertibility of $\Pi_c$ makes the lattice equality exact.

If $a=-1$, applying $\gamma$ to (PS11) gives
\[
                   m_j\gamma(\ell_j)=-2\epsilon_j.
\]
This is impossible already at $j=1$, since $3$ does not divide $2$.
Thus a global isomorphism cannot have negative linear part.

For $a=1$, the first row of (PS11) gives $\gamma(\ell_j)=0$.
Equation (PS9) then proves that the order-three filling imposes no
restriction on $n$, whereas the order-four filling imposes precisely
$n\in2\mathbb Z$. More intrinsically, the finite obstruction is the
class of $-n\epsilon_j\widehat\delta$ in
\[
 \Lambda^{A_j}/N_j\Lambda.
 \tag{PS12}\label{pshift:eq:finite-class}
\]
In the original fixed bases this quotient is
$\mathbb Z/3$ at $j=1$, with $\widehat\delta$ zero, and
$\mathbb Z/4\oplus\mathbb Z/2$ at $j=2$, with
$\widehat\delta$ generating the second summand. Indeed (PS9) gives
\[
 N_1\Lambda=3\mathbb Z v_1\oplus\mathbb Z\widehat\delta,
 \quad
 N_2\Lambda=4\mathbb Z\varepsilon'\oplus
                      2\mathbb Z\widehat\delta,
 \qquad \varepsilon'=(1,3,-3,0)^{\mathsf t}.
\]
This proves the exact parity obstruction rather than merely failure
of the untranslated identity.

\subsubsection{Local maps for every even shift}

Let $n\in2\mathbb Z$. Define the following original rational columns,
integral cocycles, and holomorphic translation vectors:
\[
\begin{aligned}
 r_1&=-\frac n3(\widehat u+\widehat w),&
 \ell_1&=-n\widehat w,&
 b_1(z)&=\Pi_c(z)r_1=-\frac n3(\tau(z)+1,\mu(z))^{\mathsf t},\\
 r_2&= \frac n4(\widehat u+\widehat w),&
 \ell_2&=\frac n2\widehat u,&
 b_2(z)&=\Pi_c(z)r_2=\frac n4(\tau(z)+1,\mu(z))^{\mathsf t}.
\end{aligned}
\tag{PS13}\label{pshift:eq:local-corrections}
\]
Direct multiplication of the complete matrices in (PS3) gives
\[
 (I-A_j)r_j=\frac{n\epsilon_j}{m_j}\widehat\delta+\ell_j.
 \tag{PS14}\label{pshift:eq:local-matrix-identity}
\]
For $j=1$ the four entries are $(0,0,-n,n/3)$, and for $j=2$
they are $(0,n/2,0,-n/4)$. Pairing with the full $\Pi_c(g_jz)$
shows that $b_j$ satisfies (PS10) for $a=1$.
Translation by $b_j$ is therefore already an isomorphism of the
two finite local fillings. The remaining issue is its exact overlap
with the fixed middle identification.

Keep the original linearizing coordinate, radius, and phase
$s_j(g_jz)=e^{-2\pi i/m_j}s_j(z)$. On a logarithm branch put
\[
 \ell^{\log}_j(z)=\frac{\log s_j(z)}{2\pi i},\qquad
 a_j(z)=b_j(z)+n\epsilon_j \ell^{\log}_j(z)e.
 \tag{PS15}\label{pshift:eq:finite-middle-sections}
\]
Changing the logarithm adds $n\epsilon_j\widehat\delta$ under the
period map, so $a_j$ is branch independent as a torus section.
For the specified negative generator, $\ell^{\log}_j(g_jz)=\ell^{\log}_j(z)-1/m_j$.
Equations (PS4) and (PS14) consequently give the exact equality
\[
 a_j(g_jz)=R_{g_j}(z)a_j(z)+\Pi_c(g_jz)\ell_j.
 \tag{PS16}\label{pshift:eq:local-affine-cocycle}
\]
This is why the logarithmic term in (PS15) cannot be omitted.
In fact the original finite overlap maps are translations by
$\sigma_{j,c}=\ell^{\log}_j\Pi_c v_j$, and
\[
                  \sigma_{j,c+n}-\sigma_{j,c}
                     =n\epsilon_j\ell^{\log}_je.
 \tag{PS17}\label{pshift:eq:actual-overlap-difference}
\]
Thus $a_j$ is exactly the middle translation induced by the local
translation $b_j$, with the original gluing maps retained.

At the cusp retain $s=\tau-h$, $e^{2\pi is}=t_c$, $s(g_0z)=s(z)-1$,
and the complete period matrix
\[
 Z_c=sB_0+C_c(t_c),\quad
 B_0=\begin{pmatrix}0&1\\-1&0\end{pmatrix},\quad
 C_c=\begin{pmatrix}6\mu&h\\ b^{\mathrm{part}}-h+c&\mu\end{pmatrix}.
 \tag{PS18}\label{pshift:eq:cusp-periods}
\]
For $\lambda=(\lambda_1,\lambda_2)^{\mathsf t}\in\mathbb Z^2$,
$C_{c+n}\lambda-C_c\lambda=(0,n\lambda_1)^{\mathsf t}$. Therefore
\[
 \exp(2\pi iC_{c+n}\lambda)t_c^{B_0\lambda}
       =\exp(2\pi iC_c\lambda)t_c^{B_0\lambda}.
 \tag{PS19}\label{pshift:eq:cusp-actions-identical}
\]
The entire cusp actions, including their nonconstant holomorphic units,
are identical for an integer shift.

The affine cocycle convention in (PS16) has multiplication law
\[
                    \ell_{gh}=\ell_g+A_g\ell_h.
 \tag{PS20}\label{pshift:eq:cocycle-law}
\]
Since $g_0=(g_1g_2)^{-1}$, it gives
\[
\begin{aligned}
 \ell_1+A_1\ell_2
   &=-\frac{3n}{2}\widehat w+\frac n2\widehat\delta,\\
 \ell_0=-A_{g_0}(\ell_1+A_1\ell_2)
   &=\frac{3n}{2}\widehat w-\frac n2\widehat\delta.
\end{aligned}
\tag{PS21}\label{pshift:eq:cusp-cocycle}
\]
Here $A_{g_0}$ fixes both toric lattice vectors, as is also immediate
from $A_{g_0}=\left(\begin{smallmatrix}I&0\\B_0&I\end{smallmatrix}\right)$.
Put
\[
 q_n=(-3n/2,n/2)^{\mathsf t}\in\mathbb Z^2,
                  \qquad a_0(z)=s(z)q_n.
 \tag{PS22}\label{pshift:eq:cusp-translation}
\]
Then $a_0(g_0z)-a_0(z)=-q_n=\Pi_c(g_0z)\ell_0$, exactly as
required by (PS16). Exponentiating gives the map
$(x,t_c)\mapsto(t_c^{q_n}x,t_c)$ on the dense cusp torus.
It extends holomorphically and invertibly to the original infinite fan:
the integral map
\[
 (y,y_3)\longmapsto(y+y_3q_n,y_3)
\]
translates each height-one triangle by the integer vector $q_n$ and
therefore permutes the actual cones, with inverse given by $-q_n$.
It commutes with the cusp quotient action in (PS19). This proves
extension across every cusp stratum, not only across its dense torus.

\subsubsection{The global additive correction and all overlap identities}

We give the global correction explicitly as an additive gluing
construction and prove its solvability. Let $\widetilde B^\circ$ be
the universal cover of the three-punctured base. Its fundamental group
is free on the original two negative peripheral generators. Pull back
the original $\Pi_c,R$ to this cover. Assign $\ell_1,\ell_2$ from
(PS13) to its free generators and extend by (PS20). Freeness makes
this an actual integral cocycle; its third peripheral value is
precisely (PS21). The maps
\[
 (z,y)\longmapsto
       (gz,R_g(z)y+\Pi_c(gz)\ell_g)
 \tag{PS23}\label{pshift:eq:affine-bundle-action}
\]
form an affine action. To check the action law, covariance gives
$R_g(hz)\Pi_c(hz)=\Pi_c(ghz)A_g$; substituting this equality
in the composition gives exactly (PS20).
The quotient is an affine holomorphic bundle over $B^\circ$, since
the deck action on the base is free and properly discontinuous and
one may use evenly covered base opens as affine charts.

On a finite punctured disc, (PS16) says that $a_j$ is a section of
this affine bundle. On the cusp punctured disc, (PS21)--(PS22) say
the same for $a_0$. Define a sheaf $\mathscr H_n$ on $B$ as follows.
Away from the special points it is the sheaf of sections of (PS23).
Near $p_j$, a section is required to have the form $a_j+\eta_j$,
where $\eta_j:\Delta_j\to\mathbb C^2$ is holomorphic and obeys
$\eta_j(g_jz)=R_{g_j}(z)\eta_j(z)$. Near the cusp it is required
to have the form $a_0+\eta_0$, with $\eta_0(t_c)$ holomorphic on
the full cusp disc. These are compatible definitions on the punctured
overlaps by the equal affine cocycles just proved. They also prove
local nonemptiness, including at all three special points.

The differences of these local sections form a vector sheaf
$\mathcal V$. Its ordinary sections have covariance $\eta(gz)=R_g\eta(z)$;
at the two finite points and the cusp its extensions are precisely
the holomorphic extensions specified above. Projection to the first
coordinate gives an exact sequence on the original compact base
\[
 0\longrightarrow\mathcal O_B e\longrightarrow\mathcal V
        \longrightarrow\mathcal L_{\rm per}\longrightarrow0,
                  \qquad \mathcal L_{\rm per}\simeq\mathcal O_B(-p_0).
 \tag{PS24}\label{pshift:eq:vector-sheaf}
\]
Here $\mathcal L_{\rm per}$ is exactly the homogeneous period sheaf from
\cref{fam:A.3.1}, not a newly chosen extension. Indeed
the first diagonal entries in (PS4) are its factors $-1/\tau$ and
$1/\tau$, and at the cusp both extensions require ordinary
holomorphicity. The original function
\[
 \Theta=\frac{(E_4\circ\tau)^2(E_6\circ\tau)^{1/2}}
                 {\Delta_{\mathrm{mod}}\circ\tau}
\]
has those exact factors, orders two and one in $s_1,s_2$, and a
simple pole in $t_c$. Dividing by $\Theta$ identifies the ordinary
and finite stalks with $\mathcal O_B$, and the cusp stalk with
$\mathfrak m_{p_0}$. This recovers the displayed $\mathcal L_{\rm per}$
with its original cusp condition. These transformation and order
calculations are the ones used to construct the retained $\mu$;
no topological conclusion about $X_c$ is involved.

For completeness, exactness on the finite stalks in (PS24) retains
the off-diagonal entry of $R_g$. Given an equivariant first component
$h_{\rm loc}(z)$, begin with the vector $(h_{\rm loc}(z),0)^{\mathsf t}$ and average it
under the finite group action on vector-valued holomorphic functions.
The resulting vector is equivariant, and its first component is
still $h_{\rm loc}$. Thus projection is surjective there. Its kernel is
exactly the invariant second-component functions, which are
holomorphic functions of $s_j^{m_j}=t_j$, multiplied by $e$.
The ordinary and cusp stalks are immediate. This also proves that
$\mathcal V$ is locally free of rank two: lift a local frame of
$\mathcal L_{\rm per}$ and join it to $e$.

Both $H^1(B,\mathcal O_B)$ and
$H^1(B,\mathcal O_B(-p_0))$ vanish. One can compute these two
groups without a degree approximation: on the two standard affine
charts of $\mathbb P^1$, split a holomorphic Laurent series on
their overlap into its nonnegative and negative powers. For
$\mathcal O_B$ these are the two coboundary contributions, and
for $\mathcal O_B(-p_0)$ its transition factor shifts the second
contribution by one, still leaving no missing exponent. The positive
part converges on the full first affine chart and the transformed
negative part on the full second affine chart. This proves the two
vanishings, equivalently the usual additive Cousin calculation on
these charts. The exact sequence (PS24) now gives
\[
                             H^1(B,\mathcal V)=0.
 \tag{PS25}\label{pshift:eq:no-analytic-obstruction}
\]

Here is how to choose the correction entering the maps. Choose
local sections $h_\alpha$ of $\mathscr H_n$, taking the displayed
$a_j,a_0$ on sufficiently small special discs and affine charts
elsewhere. Their differences $d_{\alpha\beta}=h_\beta-h_\alpha$
are a Cech $1$-cocycle with values in $\mathcal V$. Project it to
$\mathcal L_{\rm per}$ and solve that additive cocycle, using its just proved
$H^1$ vanishing. Lift the resulting local $\mathcal L_{\rm per}$ sections to
$\mathcal V$ by the local lifting construction in (PS24). Subtracting
their coboundary from $d$ leaves a scalar $\mathcal O_Be$ cocycle.
Solve that scalar additive cocycle by its $H^1$ vanishing. This
produces sections $v_\alpha$ with
$d_{\alpha\beta}=v_\beta-v_\alpha$. The local sections
$h_\alpha-v_\alpha$ agree and define a global section
$a\in\mathscr H_n(B)$. This is a two-stage construction of the
required correction, not an unproved additional gluing condition.
It gives holomorphic vectors $\eta_j,\eta_0$ with
\[
 a=a_j+\eta_j\quad(j=1,2),\qquad a=a_0+\eta_0
                              \quad\hbox{at the cusp}.
 \tag{PS26}\label{pshift:eq:global-correction}
\]
Any choices made in the additive solutions give maps with the same
asserted properties. Their differences as sections of this affine
bundle lie in $H^0(B,\mathcal V)=\mathbb C e$: this follows from
(PS24), $H^0(B,\mathcal L_{\rm per})=0$, and $H^0(B,\mathcal O_B)=\mathbb C$.

Projecting $a$ modulo $\Pi_c\Lambda$ gives a holomorphic section of
the middle torus family because every $\ell_g$ in (PS23) is integral.
The promised three kinds of local maps are therefore
\[
\begin{aligned}
 F_{\mathrm{mid}}[z,\zeta]&=[z,\zeta+a(z)],\\
 F_j[z,\zeta]&=[z,\zeta+b_j(z)+\eta_j(z)]\quad(j=1,2),\\
 F_0[x,t_c]&=[\exp(2\pi i\eta_0(t_c))t_c^{q_n}x,t_c].
\end{aligned}
\tag{PS27}\label{pshift:eq:full-isomorphism}
\]
The first descends by (PS23) and equality of the integer-shift
lattices. The finite maps descend by (PS10), (PS14), and the exact
covariance of $\eta_j$. The cusp map descends by (PS19); its
holomorphic-unit factor extends by torus multiplication on every
original chart, and its integer factor extends by the fan map in
(PS22). Each map is a biholomorphism, with inverse given by the
negative translation or the inverse torus multiplier and fan map.

On a finite overlap, write the original gluing as
$h_{j,c}(\zeta)=\zeta+\sigma_{j,c}$. Then (PS17), (PS26), and
(PS27) give the equality of actual torus maps
\[
                     h_{j,c+n}F_j=F_{\mathrm{mid}}h_{j,c}.
 \tag{PS28}\label{pshift:eq:finite-overlap-identity}
\]
Both sides add
$\sigma_{j,c}+a=\sigma_{j,c+n}+b_j+\eta_j$.
On the cusp overlap the original map is exponentiation of $\zeta$;
(PS22) and (PS26) give
$\exp(2\pi ia)=t_c^{q_n}\exp(2\pi i\eta_0)$, exactly the
third line of (PS27). Thus that overlap also agrees, including
all nonconstant period units. The maps and their inverses glue to
a global biholomorphism $X_c\to X_{c+n}$ over $B$.

The choices can be made coherently for every parameter and every even
shift. Both initial cocycle columns in (PS13) have $\gamma=0$.
Because every $A_g$ preserves $\gamma$, (PS20) shows that
$\gamma(\ell_g)=0$ for every word. Therefore the affine term
$\Pi_c(gz)\ell_g$ in (PS23) is independent of $c$. The displayed
$a_j,a_0$, and the vector sheaf $\mathcal V$ are also independent
of $c$. Perform the additive construction once for $n=2$, keeping
one resulting section $a^{(2)}$. For $n=2k$ choose
$a^{(2k)}=k a^{(2)}$ and multiply its local correcting vectors by
$k$. All their affine cocycles and local formulas are then exactly
those for $n=2k$. The resulting maps satisfy the full group law
\[
 F_{2h,c+2k}\circ F_{2k,c}=F_{2(h+k),c},\qquad
                 F_{0,c}=\mathrm{id}_{X_c}.
 \tag{PS28a}\label{pshift:eq:coherent-shift-action}
\]
This follows by adding the actual translation vectors in the first
two lines of (PS27), and by multiplying their units and adding the
integer exponents in its third line. The formulas are jointly
holomorphic in $c$, since their corrections are independent of it.

\subsubsection{Intrinsic necessity and the resulting parameter space}

The analytically computed complete linear system
$|\omega_{X_c}^{-2}|$ is the original map $f_c$, by
\cref{fam:A.1.1}. A biholomorphism preserves its
canonical bundle, the complete space of sections, and its evaluation
map. Consequently it induces $\alpha\in\operatorname{PGL}_2(\mathbb C)$
with $f_{c'}F=\alpha f_c$. The three special fibres have distinct
intrinsic types: the cusp is reduced and nonnormal, while the other
two have multiplicities three and four. All remaining fibres are
smooth reduced tori. Thus $\alpha$ fixes $p_0,p_1,p_2$ individually.
In the original coordinate, a fractional-linear map fixing
$\infty,0,1$ has zero lower-left and upper-right entries and equal
diagonal entries; hence $\alpha=\mathrm{id}_B$.

Now \cref{pshift:lem:centralizer} forces $c'-c=n\in\mathbb Z$
and linear part $\pm I_2$. Equations (PS10)--(PS11), obtained on
the original normalized finite covers, exclude $-I_2$ and then
exclude odd $n$ for $+I_2$. Conversely (PS27)--(PS28) construct
the global map for every even $n$. This proves
\cref{pshift:thm:classification} without converting a failed
particular map into a nonisomorphism assertion.

As a complex parameter space for the isomorphism classes represented
by this retained family, the quotient $\mathcal P/(2\mathbb Z)$ has
the explicit coordinate
\[
 w=\exp(-\pi ic),\qquad 0<|w|<\exp(-\pi M).
 \tag{PS29}\label{pshift:eq:parameter-disc}
\]
Indeed $|w|=\exp(\pi\operatorname{Im}c)$, and two complex numbers
have the same value of this exponential exactly when they differ
by an even integer. Every point in the displayed punctured disc
has a logarithm giving a $c\in\mathcal P$. The exponential has
nonzero derivative, so these local inverse logarithms give the
claimed quotient complex structure. The integer-period quotient
seen by the middle family instead has coordinate $w^2$; it is
covered two-to-one by (PS29). Its two lifts are distinguished by
the exact order-two summand in (PS12). This statement classifies
the parameters of the retained family; it does not assert that
the family exhausts all deformations or all complex threefolds
with the same singular-fibre types.


\section{Re-derivations of known results}\label{app:known}

\subsection{Intrinsic fibration and even shifts}\label{app:known-periods}
The results cited as K in Sections 1 and 2 are known with the credit given there. The claim made for the text below concerns its re-derivation from the retained period charts.
\subsubsection{Re-derivation of FAM A.1.1}
\begin{proof}
\par\smallskip\noindent\textbf{Step 1.} We compute the line and its evaluation map, including the multiple fibres. A holomorphic function restricts to a constant on every connected compact reduced fibre. To descend functions on a neighbourhood we first use regular fibres; equality then extends from their dense union by the identity theorem on the smooth total space. Away from the multiple fibres, evaluation on local sections descends it to a holomorphic base function. On a finite root cover this follows by evaluation on the zero section; invariance under $s\mapsto\zeta s$ leaves precisely the convergent series in $s^m$. At the cusp, a smooth point of the reduced fibre admits a local section, and constancy on each nearby compact fibre gives the same descent there by the identity theorem. Consequently $f_*\mathcal O_X=\mathcal O_B$.

\par\smallskip\noindent\textbf{Step 2.} Put $e_\wedge=d\zeta_1\wedge d\zeta_2$ in relative fibre coordinates. At the finite centre $j$ let $a_1=2,a_2=1$ and $\zeta_j=e^{-2\pi i/m_j}$. The canonical frame $ds_j\wedge e_\wedge$ has character $\chi_j(s)=\zeta_j\det R_j(s)$, with $\chi_j(0)=\zeta_j^{1+a_j}$. Take its holomorphic logarithm $\chi_j=\zeta_j^{1+a_j}\exp\psi_j$, where $\psi_j(0)=0$. The finite group identity says $\sum_{r=0}^{m_j-1}\psi_j(\zeta_j^r s)=0$. Writing $\psi_j=\sum c_ns^n$, we obtain $c_n=0$ for $m_j\mid n$. Therefore
\[
 \varphi_j(s)=\sum_{m_j\nmid n}\frac{c_n}{1-\zeta_j^n}s^n
\]
is convergent and satisfies $\varphi_j(s)-\varphi_j(\zeta_js)=\psi_j(s)$. The invariant coefficient module of the relative canonical frame
$e^{\varphi_j}(ds_j\wedge e_\wedge)\otimes(dt_j)^{-1}$ is exactly
$s_j^{m_j-1-a_j}\mathbb C\{s_j^{m_j}\}$. Notice that $(dt_j)^{-1}$ denotes a pulled-back line-bundle frame; no division by a vanishing differential is made. Since the total-space cover is etale, the evaluation divisor downstairs is $(m_j-1-a_j)S_j$, namely $0$ and $2S_2$.

\par\smallskip\noindent\textbf{Step 3.} At the cusp the ordinary form
\[
 \Omega=\frac{dx_1}{x_1}\wedge\frac{dx_2}{x_2}\wedge dt_c
\]
is nowhere zero. Indeed in a maximal unimodular cone chart, with coordinates $z_0,z_1,z_2$, the height-one condition gives $t_c=z_0z_1z_2$ and unimodularity gives $\Omega=\pm dz_0\wedge dz_1\wedge dz_2$. The same equality after inverting coordinates treats the other charts. Integral shears and holomorphic torus multipliers depending on $t_c$ preserve this form: all their extra differential terms contain $dt_c$ and vanish in the wedge. Thus the form descends to the cusp quotient and evaluation has no zero there. This calculation also shows that changing any collar by a fibre translation does not alter the canonical transition factors.

\par\smallskip\noindent\textbf{Step 4.} We next determine $\mathfrak K=f_*\omega_{X/B}$. On the simply connected period cover let $h_6^2=E_6(\tau)$, and put $\Theta=E_4(\tau)^2h_6/\Delta_{\rm mod}(\tau)$. The relation $j(\tau)=1728t$ gives local degrees $\tau-\rho\sim s_1$, $\tau-i\sim s_2^2$ and $q_\tau=e^{2\pi i\tau}\sim t_c/1728$. The identities $E_4^3=1728t\Delta_{\rm mod}$ and $E_6^2=1728(t-1)\Delta_{\rm mod}$, together with the nonvanishing product $\Delta_{\rm mod}=q_\tau\prod_{n\ge1}(1-q_\tau^n)^{24}$ on the upper half-plane, give orders $2,1,-1$ for $\Theta$. The same identities exclude every other zero or pole. The square root exists because the zeros of $E_6(\tau)$ have even order and the cover is simply connected. Its signs are determined locally: at $\rho$, $\rho^3=-1$ and $h_6\ne0$, whereas at $i$, $h_6$ has order one in $s_2$. The weight transformation laws consequently give
\[
 \Theta(g_1z)=-\Theta(z)/\tau(z),\qquad
 \Theta(g_2z)=\Theta(z)/\tau(z),\qquad \Theta(g_0z)=\Theta(z).
\]
The invariant relative form $\Theta^{-1}e_\wedge$ is a unit in the finite local generator just calculated and has a simple zero in the cusp generator $e_\wedge=(2\pi i)^{-2}\Omega\otimes(dt_c)^{-1}$. It defines an isomorphism $\mathfrak K\simeq\mathcal O_B(p_0)$. The evaluation divisor now yields
\[
 \omega_X\simeq f^*\mathcal O_B(p_0)\otimes f^*\omega_B\otimes\mathcal O_X(2S_2).
\]
In the fixed coordinate $t$, the function $(t-1)^{-1}$ has divisor $p_0-p_2$, while $dt/(t-1)^2$ has divisor $-2p_2$. These explicit maps identify the last display with $\mathcal O_X(-4S_2+2S_2)$.

\par\smallskip\noindent\textbf{Step 5.} A meromorphic representative of a section of $\mathcal O_X(kS_2)$ is holomorphic off $S_2$, so the constancy argument makes it a pullback there. On the root cover its value along the zero section is meromorphic, with order at least $-k$, and is invariant under $s\mapsto-is$. Its Laurent series has only powers $s^{4r}$. The necessary and sufficient inequality is $4\operatorname{ord}_{p_2}b+k\ge0$, which proves the direct-image formula, including negative $k$.

\par\smallskip\noindent\textbf{Step 6.} Choose the canonical divisor section $s_{\rm can}$ of $\mathcal O_X(2S_2)$ and identify this line with $\omega_X^{-1}$ using the preceding maps. Then $s_{\rm can}^2=f^*\sigma_{p_2}$. Choose $\ell\in H^0(B,\mathcal O_B(p_2))$ with $\ell(p_2)\ne0$, and let $t_{\rm can}=f^*\ell$ in degree two. In degree $m=2r$ the monomials $s_{\rm can}^{2r-2j}t_{\rm can}^j$, $0\le j\le r$, are the pullbacks of the usual basis $\sigma_{p_2}^{r-j}\ell^j$. In degree $2r+1$, multiplication by $s_{\rm can}$ gives the same basis because the allowed base pole order remains $r$. These are all sections by the direct-image formula. Thus $U\mapsto s_{\rm can},V\mapsto t_{\rm can}$ is a graded algebra isomorphism, not only a Hilbert-series equality. The two degree-two sections have no common zero and give $f$ in the projective coordinate $[\sigma_{p_2}:\ell]$.

\par\smallskip\noindent\textbf{Step 7.} A biholomorphism pulls back the canonical line, its sections and their evaluations; it therefore induces $\alpha\in\operatorname{PGL}_2(\mathbb C)$ with $f'F=\alpha f$. The cusp fibre is reduced with nonnormal points. The two finite fibres have smooth reductions and nilpotence orders exactly three and four in their local equations $s^3=0,s^4=0$. These are intrinsic local-ring properties of the full fibres and distinguish all three. Every other fibre is a smooth reduced torus. Consequently $\alpha$ fixes $\infty,0,1$ individually, and the equation of a fractional linear transformation gives $\alpha=1$.
\end{proof}
\begin{proposition}[FAM A.1.2; T]\label{fam:A.1.2}
The normalized base change $N_j\times_{D_j}\Delta_j\to\Delta_j$, $t_j=s_j^{m_j}$, is the original smooth torus family $J'_j$. The maps $E_1\times F_1\to S_1$ and $E_2\times F_2\to S_2$ described below are finite etale covers of degrees $9$ and $8$.
\end{proposition}
\uses{Definition~\ref{def:conventions}; FAM A.1.1 and preceding appendix calculations.}
\begin{proof}
\par\smallskip\noindent\textbf{Step 1.} On $J'_j$, the deck generator is
$(s,\zeta)\mapsto(\zeta_js,R_j(s)\zeta+\Pi_c(g_js)v_j/m_j)$. It has exact order $m_j$. Its $r$th power changes the invariant circle coordinate $\gamma$ by $r\gamma(v_j)/m_j$; this is nonintegral for $0<r<m_j$, so the action is free. The quotient map is therefore a local biholomorphism. Its local base function is $s^{m_j}$ and its fibre multiplicity is exactly $m_j$. The map
\[
 J'_j\longrightarrow N_j\times_{D_j}\Delta_j,\qquad y\longmapsto(q_j(y),s_j(y))
\]
is finite, is an isomorphism over $s\ne0$, and has normal source. Locally at the central fibre the base-change ring has equation $z^{m_j}-s^{m_j}=0$. Its normalization is the disjoint union of branches $z=\zeta_j^rs$. The displayed map gives those branches in the covering charts. It is thus the normalization everywhere, with its specified root coordinate.

\par\smallskip\noindent\textbf{Step 2.} Put $\rho=e^{\pi i/3}$, $\mu_1=(2-\rho)/3$, $\mu_2=(1-i)/2$, $\kappa_1=\beta_c(z_1)+2\mu_1$, and $\kappa_2=\beta_c(z_2)+3\mu_2$. The fixed-point transformation laws for $\mu$ give the two displayed values directly. We have $\operatorname{Im}\kappa_j=D_c(z_j)<0$. Define
\[
 E_j=\mathbb C/(\mathbb Z+\kappa_j\mathbb Z),\quad
 F_1=\mathbb C/(\mathbb Z(\rho-2)-3\mathbb Z),\quad
 F_2=\mathbb C/(\mathbb Z(i-1)-2\mathbb Z).
\]
Let $e_1=(0,1,-2,0)^t,f_1=(0,0,-3,1)^t,e_2=(0,1,-1,0)^t,f_2=(0,0,-2,1)^t$. Direct solution of the fixed equations and inspection of the columns give
\[
 \Lambda^{A_1}=\mathbb Z\varepsilon+\mathbb Zd,\quad
 \Lambda^{A_2}=\mathbb Z\varepsilon'+\mathbb Zd,\quad
 (A_j-I)\Lambda=\mathbb Ze_j+\mathbb Zf_j.
\]
The image lattices are saturated: they are the integer solutions of $\gamma=0,2u_{\Lambda}+w_{\Lambda}+3\delta=0$ and $\gamma=0,u_{\Lambda}+w_{\Lambda}+2\delta=0$, respectively. The determinants of the fixed-plus-moving basis matrices are $-3,-2$.

\par\smallskip\noindent\textbf{Step 3.} The complex-linear maps $L_1(\xi,\eta)=(\eta,\xi-\eta/3)$ and $L_2(\xi,\eta)=(\eta,\xi-\eta/2)$ carry the product periods to those fixed and moving sublattices. They induce isogenies to the original central tori, with kernel generated respectively by $(2/3,\rho)$ and $(1/2,i)$, of exact orders three and two. For example $\widehat u=2d/3+e_1-2f_1/3=d/2+e_2-f_2/2$ gives both kernel representatives.

\par\smallskip\noindent\textbf{Step 4.} The affine deck maps lift to
\[
 (\xi,\eta)\longmapsto(\xi+\kappa_1/3,(\rho-1)\eta),\qquad
 (\xi,\eta)\longmapsto(\xi-\kappa_2/4,-i\eta).
\]
Substitution in $L_j$ gives exactly $R_jL_j+\Pi_cv_j/m_j$. They commute with their kernel translations: the moving differences are periods, since $(\rho-2)\rho=-1-\rho=-(\rho-2)-3$ can equally be checked in the lattice basis, and $(-i-1)i=1-i=-(i-1)$. More directly these assertions follow by the two-dimensional period matrices of the indicated rotations. A nonidentity element of either resulting product group translates the $E_j$ coordinate by $2a/3+b\kappa_1/3$ or $a/2-b\kappa_2/4$, respectively. Real independence of $1,\kappa_j$ proves that this translation vanishes only for the identity element. Hence the groups $C_3\times C_3$ and $C_2\times C_4$ act freely and have quotient precisely $S_j$.
\end{proof}
\begin{cited}[FAM A.2.1; K]\label{fam:A.2.1}
The rational centralizer of $(A_1,A_2)$ is $\mathbb QI+\mathbb QE$. Every biholomorphism $X_c\to X_{c'}$ has $c'-c=n\in2\mathbb Z$, complex-linear fibre part $I_2$, and source-to-target lattice map $I-nE$. In the full fixed lattices,
\[
 \Lambda^{A_1}/N_1\Lambda\simeq\mathbb Z/3,\qquad
 \Lambda^{A_2}/N_2\Lambda\simeq\mathbb Z/4\oplus\mathbb Z/2,
\]
with $d$ generating the second factor of the latter group and vanishing in the former.
\end{cited}
\uses{Definition~\ref{def:conventions}; FAM A.1.1 and preceding appendix calculations.}
\begin{proof}
\par\smallskip\noindent\textbf{Step 1.} The common invariant vector space is $\mathbb Qd$, and the common invariant covector space is $\mathbb Q\gamma$; this follows by intersecting the fixed bases in A.1.2 and the rows $\langle\gamma,2u_{\Lambda}+w_{\Lambda}+3\delta\rangle$, $\langle\gamma,u_{\Lambda}+w_{\Lambda}+2\delta\rangle$. A commuting endomorphism preserves $\mathbb Qd\subset\ker\gamma$. Its induced matrix on $\ker\gamma/\mathbb Qd$ commutes with
$Q_1=\left(\begin{smallmatrix}0&1\\-1&-1\end{smallmatrix}\right)$ and $Q_2=\left(\begin{smallmatrix}0&-1\\1&0\end{smallmatrix}\right)$. Commutation with $Q_2$ makes it $\left(\begin{smallmatrix}b&r\\-r&b\end{smallmatrix}\right)$, and the $Q_1$ equation sets $r=0$. Write the remaining columns as
$H\widehat\gamma=(a,x,y,k)^t$, $H\widehat u=b\widehat u+u_0d$, $H\widehat w=b\widehat w+w_0d$, $Hd=d_0d$. The equations on $\widehat u,\widehat w$ give $d_0-w_0=b+u_0$, $u_0=2w_0$ and $u_0=w_0$, hence $u_0=w_0=0,d_0=b$. The first-column equations give $x=2(a-b),y=-4(a-b),x=-y$. Thus $a=b,x=y=0$, proving the assertion. Conversely $A_jE=EA_j=E$.

\par\smallskip\noindent\textbf{Step 2.} A fibre biholomorphism lifts to an affine map on $\mathbb C^2$: the derivative of a lifted map is lattice periodic, descends to a holomorphic function on a compact torus, and is constant. Its lattice map is locally constant and commutes with both monodromies. Write it as $H=aI+bE$; integrality and $\det H=a^4=1$ give $a=\pm1,b\in\mathbb Z$. The equation $L\Pi_c=\Pi_{c'}H$, in columns three and four, gives $L=aI_2$; column one gives $a(c'-c)+b=0$. Thus $n=c'-c\in\mathbb Z$.

\par\smallskip\noindent\textbf{Step 3.} The normalized base-change construction in A.1.2 is functorial for biholomorphisms over the base: the base-changed isomorphism transports integral elements of the total quotient ring to integral elements, and its inverse does the same. Its lift is over the identity of the root disc. It commutes with the prescribed deck generator because both compositions lift the same map $(x,s)\mapsto(F(x),\zeta_js)$ on the unnormalized base change. This proves exact commutation, rather than commutation after an unchosen power of the generator.

\par\smallskip\noindent\textbf{Step 4.} Lift the translation of the resulting affine torus map holomorphically to a vector $B_j(s)$ on the simply connected root disc. Its deck equation, expressed in the source periods, is
\[
 B_j(g_jz)-R_j(z)B_j(z)=\Pi_c(g_jz)
 \left(\frac{(1-a)v_j+n\gamma(v_j)d}{m_j}+\ell_j\right),\qquad\ell_j\in\Lambda.
\]
The integral vector is constant by continuity and real invertibility of the period matrix. Iterating the equation $m_j$ times gives
\[
 N_j\ell_j=-(1-a)v_j-n\gamma(v_j)d.
\]
Exact multiplication yields
\[
 N_1=\begin{pmatrix}3&0&0&0\\6&0&0&0\\-12&0&0&0\\0&2&1&3\end{pmatrix},\quad
 N_2=\begin{pmatrix}4&0&0&0\\12&0&0&0\\-12&0&0&0\\0&2&2&4\end{pmatrix}.
\]
In fixed bases their images are $3\mathbb Z\varepsilon+\mathbb Zd$ and $4\mathbb Z\varepsilon'+2\mathbb Zd$. These explicit generators prove the asserted quotient groups. Applying $\gamma$ to the norm equation excludes $a=-1$ already at $m_1=3$, since $3\gamma(\ell_1)=-2$. For $a=1$, the second norm equation says exactly $2\mid n$; the first imposes no condition. Program \texttt{code/ab\_exact.py} verifies the matrices, Smith forms, centralizer and membership independently by Hermite-form comparisons. All its computations are over $\mathbb Z$ or $\mathbb Q$.
\end{proof}
\begin{cited}[FAM A.3.1; K]\label{fam:A.3.1}
For each $n\in2\mathbb Z$ there is a biholomorphism $F_{n,c}:X_c\to X_{c+n}$, holomorphic jointly in $c$, with $F_{2h,c+2k}F_{2k,c}=F_{2(h+k),c}$. The quotient family over $0<|w|<e^{-\pi M}$, $w=e^{-\pi ic}$, is proper and holomorphic, with pairwise non-biholomorphic fibres.
\end{cited}
\uses{Definition~\ref{def:conventions}; FAM A.1.1 and preceding appendix calculations.}
\begin{proof}
\par\smallskip\noindent\textbf{Step 1.} Integer shifts satisfy $\Pi_{c+n}=\Pi_c(I+nE)$, so the middle lattices agree. Set
\[
 r_1=-\tfrac n3(\widehat u+\widehat w),\quad
 r_2=\tfrac n4(\widehat u+\widehat w),\quad
 \ell_1=-n\widehat w,\quad \ell_2=\tfrac n2\widehat u,\quad b_j=\Pi_cr_j.
\]
The four-entry equations are $(I-A_1)r_1=(0,0,-n,n/3)^t$ and $(I-A_2)r_2=(0,n/2,0,-n/4)^t$. Consequently $B_j=b_j$ solves the preceding deck equation for $a=1$. The map $(z,\zeta)\mapsto(z,\zeta+b_j(z))$ extends on the entire root family and descends, together with its inverse, to each finite filling.

\par\smallskip\noindent\textbf{Step 2.} Its middle translation on the punctured overlap is
\[
 a_j=b_j+n\gamma(v_j)\frac{\log s_j}{2\pi i}e.
\]
This follows from the actual overlap $h_{j,c}(\zeta)=\zeta+\sigma_{j,c}$, where $\sigma_{j,c}=(\log s_j)/(2\pi i)\,\Pi_cv_j$: precisely $\sigma_{j,c+n}-\sigma_{j,c}=n\gamma(v_j)(\log s_j)/(2\pi i)e$. Changes of logarithm add periods. Since the chosen generator changes $(\log s_j)/(2\pi i)$ by $-1/m_j$, we have $a_j(g_jz)=R_ja_j(z)+\Pi_c(g_jz)\ell_j$.

\par\smallskip\noindent\textbf{Step 3.} On the universal cover of $B^\circ$, extend $\ell_1,\ell_2$ by the integral cocycle law $\ell_{gh}=\ell_g+A_g\ell_h$. Freeness of $\pi_1(B^\circ)$ makes this well defined. At $g_0=(g_1g_2)^{-1}$ it gives
\[
 \ell_0=-A_0(\ell_1+A_1\ell_2)=\tfrac{3n}2\widehat w-\tfrac n2d.
\]
Put $q_n=(-3n/2,n/2)^t$ and $a_0=sq_n$. Then $a_0(g_0z)-a_0(z)=\Pi_c\ell_0$. In cusp torus coordinates its map is $(x,t_c)\mapsto(t_c^{q_n}x,t_c)$. The integral linear map $(y,h)\mapsto(y+hq_n,h)$ of the fan lattice translates each height-one triangle by an integer vector. It therefore permutes every cone of the prescribed periodic triangulation. The map and its inverse extend across all strata, including both triple points and all three double curves after quotient. Moreover the cusp deck multipliers are identical at $c$ and $c+n$, since $C_{c+n}\lambda-C_c\lambda=(0,n\lambda_1)^t$ has integral exponential. Thus the extended map descends with all its nonconstant units retained.

\par\smallskip\noindent\textbf{Step 4.} We supply the global analytic correction. Quotient the trivial affine vector bundle on the universal cover by
$(z,y)\mapsto(gz,R_g(z)y+\Pi_c(gz)\ell_g)$. Its action law is exactly the cocycle law by period covariance. Specify at the finite points sections of the form $a_j+\eta_j$, with $\eta_j$ holomorphic equivariant on the full root disc; at the cusp specify $a_0+\eta_0(t_c)$ with $\eta_0$ holomorphic on the full disc. These prescriptions define a locally nonempty torsor under a rank-two locally free sheaf $\mathcal V$ on $B$.

\par\smallskip\noindent\textbf{Step 5.} Projection to the first coordinate gives
\[
 0\longrightarrow\mathcal O_Be\longrightarrow\mathcal V
 \longrightarrow\mathcal O_B(-p_0)\longrightarrow0.
\]
Indeed its factors are $-1/\tau,1/\tau,1$. Multiplication by $\Theta$ identifies this first-coordinate sheaf with holomorphic functions at ordinary and finite points and functions vanishing at $p_0$: the finite vanishing orders are $2,1$, and the cusp pole is one. At a finite stalk, any equivariant first component has an equivariant vector lift by averaging a preliminary lift $(h_{\rm loc},0)$ over the finite group. This proves stalkwise surjectivity and local freeness, including the off-diagonal entries of $R_j$.

\par\smallskip\noindent\textbf{Step 6.} Both $H^1(\mathbb P^1,\mathcal O)$ and $H^1(\mathbb P^1,\mathcal O(-1))$ vanish. To make the analytic gluing explicit, split a smooth cocycle with a partition of unity, solve $\bar\partial$ locally by the Cauchy integral, and compare the two resulting solutions on an annulus. A Laurent series $\sum a_nt^n$ is the difference of its nonnegative and negative parts for $\mathcal O$; for $\mathcal O(-1)$ use the second frame $t^{-1}$, so the negative terms become $t^{-1}(-\sum_{n<0}a_nt^{n+1})$. Both series converge on their respective discs and leave no exponent unsolved. This kills the cocycle for any refinement. The exact sequence gives $H^1(B,\mathcal V)=0$. Project the torsor cocycle to $\mathcal O(-1)$, solve it, lift the solutions, and solve the remaining scalar cocycle. We obtain one global section $a$, with $a=a_j+\eta_j$ at finite points and $a=sq_n+\eta_0$ at the cusp.

\par\smallskip\noindent\textbf{Step 7.} The three actual maps are now
\[
 \zeta\mapsto\zeta+a,\qquad
 \zeta\mapsto\zeta+b_j+\eta_j,\qquad
 (x,t_c)\mapsto(e^{2\pi i\eta_0(t_c)}t_c^{q_n}x,t_c).
\]
The finite overlap identity is the equality of sums
$\sigma_{j,c}+a=\sigma_{j,c+n}+b_j+\eta_j$. The cusp overlap identity is $e^{2\pi ia}=t_c^{q_n}e^{2\pi i\eta_0}$. The holomorphic-unit factors act on every toric chart. The local inverses satisfy the same identities; hence the maps glue to a biholomorphism.

\par\smallskip\noindent\textbf{Step 8.} For coherence, perform the additive solution only for $n=2$ and multiply that solution by $k$ for $n=2k$. The cocycles have $\gamma(\ell_g)=0$, so their affine terms $\Pi_c\ell_g$ are independent of $c$. The reference translations, $\mathcal V$, and the correction can all be chosen independent of $c$. Adding the translations and multiplying the toric factors proves the stated cocycle identity exactly.

\par\smallskip\noindent\textbf{Step 9.} Finally the original gluing constructs a holomorphic family over any parameter disc compactly contained in $\mathcal P$. In ordinary and finite charts, real invertibility of $\Pi_c$ is uniform on compact parameter sets, so a fixed compact real parallelepiped proves properness and the lattice quotient is holomorphic. At the cusp write $Z=sB_0+C(t_c,c)$ with $B_0=\left(\begin{smallmatrix}0&1\\-1&0\end{smallmatrix}\right)$. On a compact parameter disc $C$ is uniformly bounded; shrink the cusp radius until
$\|(-2\pi\operatorname{Im}C)/\log|t_c|\|<1/2$. The effective logarithmic lattice is then within $1/2$ of the unimodular $B_0$, uniformly. Reduce its logarithmic valuation modulo the effective logarithmic lattice to a bounded parallelogram. Choose a height-one triangle containing that reduced valuation. Its barycentric coordinates are nonnegative and sum to one; the corresponding three unimodular cone characters therefore have modulus at most one. This puts the point in the closed unit polydisc of that cone. Only finitely many triangles meet the fixed bounded parallelogram, uniformly on the parameter compactum. Boundary orbits have the same representatives by translating their height-one components to these triangles. Thus every orbit has a representative in this finite union of cone polydiscs. Those closed polydiscs times a compact parameter set form a compact surjective representative set. This proves properness at the cusp and proper discontinuity by the same uniform lower bound on nonzero lattice translations. The quotient charts are open subsets of a fixed smooth toric variety times the parameter disc, so projection to the parameter is a submersion. The explicit overlap maps are jointly holomorphic. These local families agree after shrinking cusp radii, and glue over $\mathcal P$. The coherent $2\mathbb Z$ action is properly discontinuous because its action on $\mathcal P$ is. Local inverse logarithms for $w=e^{-\pi ic}$ identify the quotient family with these proper local families. A.2.1 gives pairwise non-biholomorphic fibres.
\end{proof}
\begin{proposition}[FAM A.4.1; T]\label{fam:A.4.1}
With $\kappa_j(c)=\kappa_j^{\rm part}+c$, the Albanese elliptic curves of the reduced finite fibres are
\[
 \operatorname{Alb}(S_1)=\mathbb C/(\mathbb Z+\kappa_1(c)\mathbb Z),\qquad
 \operatorname{Alb}(S_2)=\mathbb C/(\mathbb Z+\tfrac12\kappa_2(c)\mathbb Z).
\]
Writing $J(q_\tau)=q_\tau^{-1}+744+196884q_\tau+\cdots$, their invariants are
\[
 j_1(w)=J(e^{-2\pi i\kappa_1^{\rm part}}w^2),\qquad
 j_2(w)=J(e^{-\pi i\kappa_2^{\rm part}}w).
\]
These are convergent formulas because the displayed nome moduli are below one. Outside the explicitly defined countable exceptional set in the proof, equality of both Albanese invariants forces $c'-c\in2\mathbb Z$. Thus the full generic Albanese cross-check of Lead 3 holds independently of the centralizer calculation.
\end{proposition}
\uses{Definition~\ref{def:conventions}; FAM A.1.1 and preceding appendix calculations.}
\begin{proof}
\par\smallskip\noindent\textbf{Step 1.} In the product covers of A.1.2 the invariant holomorphic one-forms are exactly the multiples of $d\xi$: the moving factor is multiplied by a primitive third or fourth root. Integration thus identifies the Albanese with the quotient of the fixed line by the fixed-coordinate projections of every deck translation. These are, respectively,
\[
 \mathbb Z+\kappa_1\mathbb Z+\tfrac23\mathbb Z+\tfrac{\kappa_1}3\mathbb Z
   =\tfrac13\mathbb Z+\tfrac{\kappa_1}3\mathbb Z,
\quad
 \mathbb Z+\kappa_2\mathbb Z+\tfrac12\mathbb Z+\tfrac{\kappa_2}4\mathbb Z
   =\tfrac12\mathbb Z+\tfrac{\kappa_2}4\mathbb Z.
\]
Multiplication by three and two gives the stated maps of complex tori. In particular the Albanese computation includes the kernel of the product isogeny and the affine quotient, rather than just its fixed eigenspace. Using the positive-imaginary representatives $-\kappa_1$ and $-\kappa_2/2$ in the standard modular $j$-series gives the formulas. At sufficiently small nome $J$ is injective: $1/J(q_\tau)=q_\tau+O(q_\tau^2)$ has nonzero derivative at zero. Thus $j_2$ alone distinguishes parameters within a sufficiently small punctured subdisc when both parameters lie there. For the full generic assertion put $x=-\kappa_2(c)/2$ and $d_*=\kappa_1^{\rm part}-\kappa_2^{\rm part}$, so $-\kappa_1(c)=2x-d_*$. The independently certified interval in \texttt{cd\_special\_intervals.py} gives $-0.2<\operatorname{Im}d_*<-0.1$; in particular it is nonzero. Equality of the two elliptic invariants at $c,c'$ gives matrices $M,N\in\operatorname{PSL}_2(\mathbb Z)$ with $x'=M(x)$ and
\[
 N(2x-d_*)=2M(x)-d_*.
\]
Define the exceptional set to be the $c\in\mathcal P$ obtained from roots of these equations for all pairs for which the two rational functions are not identical. Each equation has at most two roots after clearing denominators, so this is a countable, explicitly specified set. For a pair giving an identity, let $C(x)=2x-d_*$. Then $N=CMC^{-1}$. The real matrix $M$ preserves both the real circle $\mathbb R\cup\{\infty\}$ and $C^{-1}(\mathbb R\cup\{\infty\})$, the parallel line with imaginary part $\operatorname{Im}d_*/2$ together with infinity. Their unique intersection is infinity, so $M$ fixes infinity. An element of $\operatorname{PSL}_2(\mathbb Z)$ fixing infinity is $x\mapsto x+b$, $b\in\mathbb Z$, and then $N(y)=y+2b$. It follows that $c'-c=-2b$. This proves the generic necessity with its actual exceptional set; A.2.1 proves the necessity also on that set.
\end{proof}
\begin{cited}[FAM A.4.2; K]\label{fam:A.4.2}
For every $c\in\mathcal P_{\rm KS}$, the Kodaira--Spencer class of $\partial/\partial c$ in $H^1(X_c,T_{X_c})$ is nonzero. On any fixed ordinary fibre, put $T=\operatorname{Im}\tau$, $m=\operatorname{Im}\mu$, $a=m/T$, $D=\operatorname{Im}\beta_c-6m^2/T$. In the convention $\mathrm{KS}(\partial_c)=[\bar\partial h]$ for a smooth $(1,0)$ lift $h$, its restricted torus-deformation representative is
\[
 -\frac{1}{2iD}\,\partial_{\zeta_2}\otimes
       (d\bar\zeta_2-a\,d\bar\zeta_1).
\]
\end{cited}
\uses{Definition~\ref{def:conventions}; FAM A.1.1 and preceding appendix calculations.}
\begin{proof}
\par\smallskip\noindent\textbf{Step 1.} Retain the same real marking $r\in\mathbb R^4$ and set $\zeta=\Pi_cr$. The two imaginary-part equations give
$r_1=(\operatorname{Im}\zeta_2-a\operatorname{Im}\zeta_1)/D$. Thus at $c+\delta$ the marked real-linear identification is
\[
 \zeta^{\rm new}=\zeta+\frac{\delta}{2iD}P_{\rm KS}(\zeta-\bar\zeta),\qquad
 P_{\rm KS}=\begin{pmatrix}0&0\\-a&1\end{pmatrix},\qquad P_{\rm KS}^2=P_{\rm KS}.
\]
It is exactly $\Pi_{c+\delta}\Pi_c^{-1}$ as a real map, so it sends every marked lattice vector to the corresponding vector. Its determinant is $(D+\operatorname{Im}\delta)/D>0$. The transported $(0,1)$ graph coefficient is $\nu_\delta=\delta P_{\rm KS}/(2iD+\delta)$. The derivative at zero is $P_{\rm KS}/(2iD)$. Equivalently the smooth lift of $\partial_c$ in moving coordinates has vertical part $r_1\partial_{\zeta^{\rm new}_2}$ and $\bar\partial r_1=-(d\bar\zeta_2-a d\bar\zeta_1)/(2iD)$, giving the stated sign.

\par\smallskip\noindent\textbf{Step 2.} A nonzero constant tangent-valued $(0,1)$ form on a compact torus cannot be $\bar\partial$ of a global vector field. Write that field in the invariant holomorphic frame and integrate each derivative of its periodic coefficients against invariant probability volume. Every such derivative has integral zero, whereas the coefficient of $\partial_{\zeta_2}\otimes d\bar\zeta_2$ above has nonzero integral. This proves nonvanishing for the fibre deformation.

\par\smallskip\noindent\textbf{Step 3.} We explain its transfer to the threefold. Over the dual numbers $A=\mathbb C[\epsilon]/(\epsilon^2)$ at $c$, the chart proof of A.1.1 gives the relative line isomorphism $\omega^{-2}_{X_A/A}=f_A^*\mathcal O_B(p_2)$. Extra parameter differentials disappear in relative forms and all determinant characters are independent of $c$. The two pulled-back sections form an $A$-basis of $H^0(X_A,\omega^{-2})$: reduce a section modulo $\epsilon$, subtract its linear combination of the two lifts, and use the exact sequence $0\to L\xrightarrow{\epsilon}L_A\to L\to0$ to write the remainder uniquely as $\epsilon$ times another such combination. The same argument proves independence.

\par\smallskip\noindent\textbf{Step 4.} A first-order trivialization would preserve this complete system and give a base automorphism reducing to the identity. It fixes all three special values even over $A$. On the finite etale charts the relative critical ideal is $(s_j^{m_j-1})$, and $t_j=s_j^{m_j}$ is zero on that scheme; its constant parameter ring $A$ injects. At a cusp triple chart, the critical ideal of $t_c=z_0z_1z_2$ is $(z_0z_1,z_0z_2,z_1z_2)$ and has the same property. There are no other critical values. Therefore an infinitesimal displacement of a special value would remain nonzero in its critical chart, a contradiction. Writing a base infinitesimal automorphism as $t+\epsilon(b_0+b_1t+b_2t^2)$, fixing $0,1,\infty$ gives $b_0=b_1=b_2=0$. The trivialization would consequently trivialize the displayed nonzero torus deformation. Finally a zero Kodaira--Spencer class produces such a first-order trivialization by subtracting from local product charts the vector fields whose coboundary is the differentiated transition cocycle. This proves the claim.
\end{proof}
\subsubsection{Re-derivation of FAM A.5.1}
\begin{proof}
\par\smallskip\noindent\textbf{Step 1.} A.1.1 makes the base intrinsic. A.2.1 gives the necessary parity and excludes the negative linear part. A.3.1 supplies the inverse maps, all collar identities, and descent of the jointly holomorphic family. A.4.2 proves its infinitesimal effectiveness. Finally $|e^{-\pi ic}|=e^{\pi\operatorname{Im}c}$ and equality of exponentials is exactly difference in $2\mathbb Z$, so the coordinate has precisely the stated image and fibres.
\end{proof}
\subsubsection{The universal map statement}
The following source proof is that of \cite[Theorem 7.2]{DGKXv1}, with its parameter renamed; it is reproduced with that credit rather than attributed to the workbench.
\begin{proof}
\par\smallskip\noindent\textbf{Step 1.} Fix $c$ and let $F$ be a smooth torus fibre of $f_c$. Since $X_c$ is
diffeomorphic to $S^6$, $H^2(X_c,\mathbb Z)=0$. Thus
$c_1((h|_F)^*\mathcal O_{\mathbb P^1}(1))=0$. If $h|_F$ were nonconstant, the
pullback of the Fubini--Study form would be semipositive and nonzero
on an open subset of the K\"ahler surface $F$. Its wedge with a
K\"ahler form on $F$ would have strictly positive integral,
contradicting the vanishing of that Chern class. Hence $h$ is
constant on every smooth fibre.

\par\smallskip\noindent\textbf{Step 2.} The proper holomorphic map $(f_c,h):X_c\to\mathbb P^1\times\mathbb P^1$ has,
by Remmert's theorem, an irreducible compact analytic image $\Gamma$
of dimension one. Over the complement of the three special values,
$\Gamma$ is the graph of a holomorphic map. The first projection
$p:\Gamma\to\mathbb P^1$ is therefore generically one-to-one. It is also
finite: a positive-dimensional fibre of a map from an irreducible
compact analytic curve would be the whole curve, whereas $p$ is
surjective. Thus $p$ is finite and bimeromorphic. On a coordinate
disc of $\mathbb P^1$, the finite algebra $p_*\mathcal O_\Gamma$ sits in the
meromorphic function field of the disc and is integral over its
holomorphic functions. Normality of the disc forces equality, so
$p$ is an isomorphism. The second projection gives $\varphi$ and
$h=\varphi\circ f_c$. Surjectivity of $f_c$ gives uniqueness.
\end{proof}

\subsection{Tangent cohomology and the Kuranishi line}\label{app:tangent-rederivation}
\uses{fam:A.1.2, fam:A.3.1, fam:A.4.2; the cyclic and cusp charts; \cite[III.4.10, III.4.12]{BS76}; \cite[pp.~261, 264--266]{RRV71}; \cite[p.~91, Proposition 2]{RR70}.}
\begin{proof}[Re-derivation of the known result \ref{fam:A.5.2}]
\par\smallskip\noindent\textbf{Step 1.}\ This is the supplied re-derivation of the known statement
\cite[Theorems 1.3 and 6.12]{DGKXv1}; the body retains its K status.
The analytic base-change and duality applications below are supported by
the opened standard sources \cite[III.4.10, III.4.12]{BS76},
\cite[pp.~261, 264--266]{RRV71}, and \cite[p.~91, Proposition 2]{RR70}.

\par\smallskip\noindent\textbf{Step 2.}\ \emph{The base and its line bundles.}
Write $X=X_c$ and form the analytic root curve
\[
 \mathfrak B=\sqrt[3]{(B,p_1)}\mathbin{\times_B}\sqrt[4]{(B,p_2)},
 \qquad \varpi:\mathfrak B\longrightarrow B.
\]
A root object is a line bundle with section whose indicated power is the
pullback of the section cutting out $p_j$. The actual equalities
$f^*p_1=3S_1$ and $f^*p_2=4S_2$ supply these objects as
$(\mathcal O_X(S_j),1_{S_j})$, and hence supply a morphism
$\bar f:X\to\mathfrak B$ with $\varpi\bar f=f$.
The explicit weighted quotient
\[
 \mathfrak B\simeq[(\mathbb C^2\setminus\{0\})/\mathbb C^*],
 \qquad \lambda(x,y)=(\lambda^3x,\lambda^4y),
 \qquad t=\frac{y^3}{y^3-x^4},
\]
identifies the roots with $(\mathcal O(4),y)$ and
$(\mathcal O(3),x)$ and the ordinary point $p_0$ with
$y^3-x^4=0$. In particular $\varpi^*\mathcal O_B(1)=\mathcal O(12)$.
This isomorphism follows on each root chart by the displayed root
sections; away from their zeros both sides have the same ordinary coarse
coordinate, so the chart isomorphisms glue.

\par\smallskip\noindent\textbf{Step 3.}\ For completeness the analytic line-bundle computation can be made on
the coarse curve. Removing the two stabilizer characters from any line
bundle makes it descend to $B$: on each root chart, averaging a local
frame whose central value is nonzero gives an invariant frame still
nonzero at the centre, and its invariant coefficients are holomorphic
in $s_j^{m_j}$. These frames give the descended line and its transition
functions. Consequently its Picard group has
generators $H,R_1,R_2$ and relations $3R_1=H=4R_2$.
Thus it is freely generated by $R_1-R_2$, denoted $\mathcal O(1)$;
$H=12\mathcal O(1)$, $R_1=4\mathcal O(1)$ and
$R_2=3\mathcal O(1)$. The chart descent just used is exact because
invariants of a finite group over $\mathbb C$ are an exact functor.
With the clockwise root generators, a local invariant section of
$\mathcal O(k)$ has vanishing exponents
\[
 r_3(k)\in\{0,1,2\},\quad r_3(k)\equiv k\pmod3,
 \qquad
 r_4(k)\in\{0,1,2,3\},\quad r_4(k)\equiv-k\pmod4.
\]
Dividing by those local powers gives the precise coarse direct image
\begin{equation}\label{fam:eq:ab-root-degree}
 \varpi_*\mathcal O(k)=\mathcal O_B\bigl(d(k)\bigr),\qquad
 d(k)=\frac{k}{12}-\frac{r_3(k)}3-\frac{r_4(k)}4.
\end{equation}
Indeed $k-4r_3-3r_4=12d(k)$ identifies the line bundle with
$\varpi^*\mathcal O_B(d(k))\otimes R_1^{r_3}\otimes R_2^{r_4}$.
The canonical root sections generate the invariant direct image of
the last two factors; division by their local powers is the displayed
line-bundle isomorphism on the coarse charts. Finite invariants also give
$R^q\varpi_*=0$ for $q>0$, as seen on each Stein quotient chart.
For $k=0$ the degree is zero; for
$k=1,-1,-2,-5,-6$ it is $-1$. It follows directly from the
cohomology of $\mathbb P^1$ that these five line bundles are acyclic,
whereas $H^0(\mathcal O)=\mathbb C$ and $H^1(\mathcal O)=0$.
The local equation $dt_j=m_js_j^{m_j-1}ds_j$ gives
\[
 K_{\mathfrak B}=\mathcal O(-24+2\cdot4+3\cdot3)
 =\mathcal O(-7),\qquad
 T_{\mathfrak B}(-p_0)=\mathcal O(-5).
\]

\par\smallskip\noindent\textbf{Step 4.}\ \emph{The actual base changes and cusp cohomology.}
On a finite root disc $\Delta_{s_j}$ the local presentations are
$[J'_j/\mu_{m_j}]\to[\Delta_{s_j}/\mu_{m_j}]$.
The fibre-product identity
\[
 [J'_j/\mu_{m_j}]
 \mathbin{\times_{[\Delta_{s_j}/\mu_{m_j}]}}\Delta_{s_j}
 \simeq J'_j
\]
is induced by the quotient torsor and its equivariant map to the disc;
its inverse sends a torsor with a chosen root-chart trivialization to
the corresponding point of $J'_j$. The action on $J'_j$ is free by
fam:A.1.2. Thus this is an etale atlas base change and is exactly the
smooth torus family, with its given deck generator. Over the cusp the
base change is the original cusp family. Locally it is
$t_c=z_0z_1z_2$; the homomorphism
$\mathbb C\{t_c\}\to\mathbb C\{z_0,z_1,z_2\}$ is flat because its
target is torsion-free over the discrete valuation ring.
These charts prove that $\bar f$ is representable, proper and flat,
with connected Gorenstein fibres of dimension two. Properness on the
cusp is the compact representative construction of fam:A.3.1, and on the
finite charts it is the compact torus quotient. Also
$\bar f^*p_0=W$ as reduced Cartier divisors.

\par\smallskip\noindent\textbf{Step 5.}\ Let $\nu:P_6\to W$ be the normalization, $P_6=\mathrm{dP}_6$, let
$C_h,C_v,C_d$ be the three double-curve normalizations, and let
$P_+,P_-$ be the two triple points. Each $C_e$ is $\mathbb P^1$.
Write $\rho_e^0,\rho_e^1:C_e\to P_6$ for its two paired boundary maps
and $j_e=\nu\rho_e^0=\nu\rho_e^1$. The sheaf sequence is
\begin{equation}\label{fam:eq:ab-cusp-resolution}
 0\longrightarrow\mathcal O_W\longrightarrow\nu_*\mathcal O_{P_6}
 \xrightarrow{\delta^0}\bigoplus_{e=h,v,d}j_{e*}\mathcal O_{C_e}
 \xrightarrow{\delta^1}\mathbb C_{P_+}\oplus\mathbb C_{P_-}
 \longrightarrow0,
\end{equation}
where $\delta^0(s)_e=(\rho_e^1)^*s-(\rho_e^0)^*s$ and
$\delta^1$ evaluates $t_h-t_v+t_d$ at the two ordered triple points.
Its exactness is local: for $xyz=0$ it is the alternating restriction
complex of the three coordinate planes, their pairwise lines and their
common point. Decomposing convergent power series by the variables
occurring in each monomial gives the augmented simplex complex, which
is exact. The one- and two-branch cases are the same calculation with
fewer planes. This verifies exactness for the actual boundary maps,
including the torus multipliers in their gluing.

\par\smallskip\noindent\textbf{Step 6.}\ All positive-degree structure-sheaf cohomology of $P_6,C_e$ vanishes.
Consequently \eqref{fam:eq:ab-cusp-resolution} computes cohomology by
\[
 \mathbb C\xrightarrow{0}\mathbb C^3
 \xrightarrow{\left(\begin{smallmatrix}1&-1&1\\1&-1&1\end{smallmatrix}\right)}
 \mathbb C^2.
\]
It gives $(h^0,h^1,h^2)(\mathcal O_W)=(1,2,1)$.
We also need its multiplication. The branch-intersection resolution
is the augmented semisimplicial sheaf algebra with its restriction face
maps; its local exactness makes the augmentation a multiplicative
quasi-isomorphism. Its Alexander--Whitney product therefore computes
the sheaf-cohomology cup product. In the ordered triangles
$[0,e_1,e_2]$ and $[e_1,e_2,e_1+e_2]$ it is
\[
 \alpha\smile\beta=(\alpha_h\beta_d,\alpha_d\beta_h),
 \qquad [\alpha]\smile[\beta]
 =\alpha_h\beta_d-\alpha_d\beta_h
 \quad\hbox{in }\mathbb C^2/\mathbb C(1,1).
\]
The cocycles $(1,1,0)$ and $(0,1,1)$ pair to one. Thus
$\bigwedge^2H^1(\mathcal O_W)\to H^2(\mathcal O_W)$ is an
isomorphism, including its product rather than just its dimensions.

\par\smallskip\noindent\textbf{Step 7.}\ \emph{Vertical fields and the direct images.}
Set $\mathcal T=T_X(-\log W)$ and
$\mathcal V=\ker(d\bar f^{\log})$.
The vertical transitions on the core are precisely $R_g(z)$.
Finite collar translations have vertical differential the identity,
and at the cusp the frame $\partial_{\zeta_i}$ corresponds to
$2\pi i x_i\partial_{x_i}$. The cusp transformations
$x\mapsto u(t_c)t_c^{B_0\lambda}x$ fix the coefficients of these
relative logarithmic fields. These exact identifications construct
a bundle $\mathcal E_{\rm tan}$ on $\mathfrak B$ with
$\mathcal V=\bar f^*\mathcal E_{\rm tan}$.
The fixed vector $e=(0,1)^t$ is nonzero in every logarithmic frame,
so there is an exact sequence
\[
 0\longrightarrow\mathcal O\xrightarrow{e}\mathcal E_{\rm tan}
 \longrightarrow L\longrightarrow0.
\]
The factors of $L$ are $-1/\tau$ and $1/\tau$. Therefore the usual
weight-twelve modular discriminant makes $1/\Delta_{\rm mod}(\tau)$ a
meromorphic section of $L^{12}$. It has no finite zero or pole,
and its cusp pole is simple: $\Delta_{\rm mod}(\tau)=t_c\,u(t_c)$ with
$u(0)\ne0$. Hence $L^{12}=\mathcal O(-p_0)=\mathcal O(-12)$.
The Picard calculation gives $L=\mathcal O(-1)$, and the extension
splits since $H^1(\mathcal O(1))=0$. Thus
\[
 \mathcal E_{\rm tan}\simeq\mathcal O\oplus\mathcal O(-1).
\]
There is also the exact sequence
\begin{equation}\label{fam:eq:ab-log-tangent}
 0\longrightarrow\bar f^*\mathcal E_{\rm tan}\longrightarrow\mathcal T
 \longrightarrow\bar f^*\mathcal O(-5)\longrightarrow0.
\end{equation}
On the smooth root charts it is the relative tangent sequence. At the
cusp its last map sends $\sum a_i z_i\partial_{z_i}$ to
$(\sum a_i)t_c\partial_{t_c}$, so it is onto and its kernel is the
two height-zero lattice directions already used above. Taking its
determinant and using $W=\bar f^*p_0$ gives the actual relative
dualizing line
\begin{equation}\label{fam:eq:ab-relative-dualizing}
 \omega_{X/\mathfrak B}=\bar f^*(\det\mathcal E_{\rm tan})^{-1}.
\end{equation}
In particular this computation gives
$K_X=\bar f^*\mathcal O(-6)=\mathcal O_X(-2S_2)$, with the last
identification supplied by $\bar f^*R_2=\mathcal O_X(S_2)$;
it agrees through this actual pullback with fam:A.1.1.

\par\smallskip\noindent\textbf{Step 8.}\ The fibres of $\bar f$ all have structure-sheaf cohomology dimensions
$(1,2,1)$: on ordinary and finite root charts they are tori, and at
the cusp we have just computed them. Analytic proper flat base change
on these complex-space charts gives locally free $R^q\bar f_*\mathcal O_X$
with these fibres. Its canonical comparison maps agree on the atlas
double overlaps, so they descend to $\mathfrak B$. The unit gives
$R^0\bar f_*\mathcal O_X=\mathcal O$.
The source charts are smooth and the fibres are Cartier divisors,
so relative analytic duality applies in its Gorenstein form. The
relative trace pairing, compatible on the same overlaps, and
\eqref{fam:eq:ab-relative-dualizing} give
\[
 (R^2\bar f_*\mathcal O_X)^\vee
 =\bar f_*\omega_{X/\mathfrak B}
 =(\det\mathcal E_{\rm tan})^{-1},
 \qquad R^2\bar f_*\mathcal O_X=\mathcal O(-1).
\]
Here local freeness removes the higher sheaf-Ext terms in the duality
spectral sequence. The base-change application is \cite[III.4.10 and III.4.12]{BS76}; the relative trace is supplied by \cite[pp.~261, 264--266]{RRV71}, with the Gorenstein dualizing line as in \cite[p.~91, Proposition 2]{RR70}. These replace the producing run's incorrect DGKX lemma citations.

\par\smallskip\noindent\textbf{Step 9.}\ Put $G_{\mathfrak B}=R^1\bar f_*\mathcal O_X$. Cup product and the verified fibre
comparisons show $\det G_{\mathfrak B}=R^2\bar f_*\mathcal O_X=\mathcal O(-1)$;
at $W$ this uses the nonzero pairing above.
We construct the required subline of $G_{\mathfrak B}$. On an ordinary fibre,
constant lattice cocycles identify
\[
 H^1(F_t,\mathcal O_{F_t})
 =\Lambda_{\mathbb C}^*/\operatorname{rowspan}_{\mathbb C}\Pi_c.
\]
Indeed a coboundary among constant cocycles is given by an entire
function with constant period differences. Its derivatives are
periodic holomorphic functions on the compact torus, hence constant;
the function is affine linear. This gives exactly the displayed
kernel, and both sides then have dimension two. The class of
$\gamma=(1,0,0,0)$ is nonzero, since evaluation on columns
$\widehat w,\widehat\delta$ would force both coefficients of a
putative period-row combination to vanish.
The matrices satisfy $\gamma A_j=\gamma$. Translations act trivially
on $H^1(\mathcal O)$, so this holomorphic class extends on both finite
root charts and respects their gluing.

\par\smallskip\noindent\textbf{Step 10.}\ At the cusp $\gamma$ kills
$\Lambda_{\mathrm{tor}}=\mathbb Z\langle\widehat w,\widehat\delta\rangle$.
For local sections of the free toric quotient with changes
$s_i=\Psi_{\lambda_{ij}}s_j$, the constants
$\gamma(\lambda_{ij})$ therefore form a holomorphic Cech cocycle
on the entire neighbourhood. With
$B_0=\left(\begin{smallmatrix}0&1\\-1&0\end{smallmatrix}\right)$,
the edges $e_1,e_2,e_2-e_1$ have deck vectors
$(0,1),(-1,0),(-1,-1)$. The central cocycle values are consequently
$(0,-1,-1)$, which pair to one with $(1,0,-1)$.
Its value on $W$ is nonzero. On the punctured overlap the exponential
map kills precisely $\Lambda_{\mathrm{tor}}$, so its full-lattice
cocycle is exactly $\lambda\mapsto\gamma(\lambda)$, proving equality
with the preceding class. Over the Stein cusp disc, the Leray edge
map $H^1(\bar f^{-1}D_0,\mathcal O_X)\to H^0(D_0,G_{\mathfrak B})$ is an
isomorphism, since $H^1(D_0,\mathcal O)=H^2(D_0,\mathcal O)=0$.
Thus this is a holomorphic, nonvanishing extension of $[\gamma]$.
It supplies a global subline and exact sequence
\[
 0\longrightarrow\mathcal O\xrightarrow{[\gamma]}G_{\mathfrak B}
 \longrightarrow\mathcal O(-1)\longrightarrow0.
\]
Again $H^1(\mathcal O(1))=0$ splits it. We have proved
\begin{equation}\label{fam:eq:ab-direct-images}
 R^q\bar f_*\mathcal O_X=
 \begin{cases}
 \mathcal O,&q=0,\\
 \mathcal O\oplus\mathcal O(-1),&q=1,\\
 \mathcal O(-1),&q=2.
 \end{cases}
\end{equation}

\par\smallskip\noindent\textbf{Step 11.}\ \emph{From logarithmic to ordinary tangent cohomology.}
The projection formula gives, for $q=0,1,2$ respectively,
\[
 R^q\bar f_*\mathcal V:
 \quad\mathcal O\oplus\mathcal O(-1),\quad
 \mathcal O\oplus2\mathcal O(-1)\oplus\mathcal O(-2),\quad
 \mathcal O(-1)\oplus\mathcal O(-2).
\]
For $\bar f^*\mathcal O(-5)$ the three expressions are
$\mathcal O(-5)$, $\mathcal O(-5)\oplus\mathcal O(-6)$ and
$\mathcal O(-6)$. Formula \eqref{fam:eq:ab-root-degree} makes every
nontrivial summand in this list acyclic. Leray, whose base has
cohomological dimension one, therefore gives
$h^\bullet(\mathcal V)=(1,1,0,0)$ and zero cohomology for the
horizontal term. Sequence \eqref{fam:eq:ab-log-tangent} yields
$h^\bullet(\mathcal T)=(1,1,0,0)$.

\par\smallskip\noindent\textbf{Step 12.}\ The normalization map $\nu:P_6\to X$ is an immersion on each local
branch. Let $N_\nu=\operatorname{coker}(T_{P_6}\to\nu^*T_X)$.
The exact normal-derivative map is
\begin{equation}\label{fam:eq:ab-normal-quotient}
 0\longrightarrow\mathcal T\longrightarrow T_X
 \longrightarrow\nu_*N_\nu\longrightarrow0.
\end{equation}
For a local equation $z_0\cdots z_{r-1}=0$, it sends
$\sum b_i\partial_{z_i}$ to the $r$ independent normal coefficients
$b_i|_{z_i=0}$. These can be lifted independently, proving
surjectivity; the kernel is exactly the logarithmic tangent module.
This also specifies the global map under changes of branch charts.
The cusp three-form
\[
 \Omega=d\log x_1\wedge d\log x_2\wedge dt_c
\]
is invariant under $x\mapsto u(t_c)t_c^{B_0\lambda}x$ because the
extra logarithmic differentials are multiples of $dt_c$. In any
unimodular maximal cone the height-one equation is $t_c=z_0z_1z_2$,
and the determinant-one character change gives
$\Omega=\pm dz_0\wedge dz_1\wedge dz_2$.
It therefore trivializes $K_X$ on the entire cusp neighbourhood,
including every branch intersection. Taking determinants of the
immersion tangent sequence now gives the explicit line identification
$N_\nu=K_{P_6}\otimes(\nu^*K_X)^{-1}\simeq K_{P_6}$.
Since $P_6$ is the rational projective surface $\mathrm{dP}_6$,
Serre duality gives $h^\bullet(P_6,K_{P_6})=(0,0,1)$.
The cohomology sequence of \eqref{fam:eq:ab-normal-quotient} proves
$h^\bullet(X,T_X)=(1,1,1,0)$, as asserted.

\par\smallskip\noindent\textbf{Step 13.}\ \emph{The reduced deformation germ.}
fam:A.4.2 gives a nonzero Kodaira--Spencer vector at this same $c$;
the just computed $h^1=1$ makes its map an isomorphism.
Embed the actual Kuranishi germ in its one-dimensional tangent space
and write its local analytic algebra as $\mathbb C\{t\}/I$.
The existing proper family has a local classifying map, whose pullback
sends $t$ to $g(s)$ with $g'(0)\ne0$. The holomorphic inverse-function
theorem makes $\mathbb C\{t\}\to\mathbb C\{s\}$, $t\mapsto g(s)$,
an isomorphism. Every element of $I$ maps to zero, so $I=0$.
This proves smoothness and reducedness even though $h^2(T_X)=1$.
The classifying map is consequently a local isomorphism of germs;
completeness of the Kuranishi deformation then gives the claimed
miniversality and exhaustion of sufficiently small deformations.
\end{proof}

% Re-derivation of the credited full automorphism theorem.
\subsection{Re-derivation of FAM B.4.1}
\begin{proof}
\par\noindent\textit{Step 1.} \ref{fam:A.1.1} and \ref{fam:A.2.1} give base identity, $n=0$ and fibre-linear part $I$. Let $\ell_1,\ell_2$ be the integral defects of a lift of its translating section to the universal cover of $B^\circ$. Finite extension gives $N_j\ell_j=0$. The saturated image calculations in \ref{fam:A.1.2} give the integral equality $\ker N_j=(I-A_j)\Lambda$. The cusp valuation argument (B2.3) gives $\ell_0\in\Lambda_{\rm tor}$. Write
\[
 \ell_1=ae_1+bf_1,\qquad \ell_2=ce_2+df_2.
\]
The $\widehat u$ coordinate of $\ell_1+A_1\ell_2$ is $a-c-2d$, so the cusp condition is exactly $a=c+2d$. Changing the global lift by a period $\lambda$ changes this pair by the simultaneous coboundary $((I-A_1)\lambda,(I-A_2)\lambda)$. In the four displayed coefficient coordinates its matrix is
\[
 C=\begin{pmatrix}-6&1&-1&0\\2&-1&0&0\\0&1&1&0\\-3&0&-1&0\end{pmatrix}.
\]
The lattice of all allowed pairs has basis matrix
$K=\left(\begin{smallmatrix}0&1&2\\1&0&0\\0&1&0\\0&0&1\end{smallmatrix}\right)$. The first three columns of $C$ expressed in this basis have determinant one. Hence every allowed pair is an integral coboundary; the computation is exact, including its saturation, and is independently checked by Hermite form in \texttt{ab\_exact.py}. After changing the lift, all its defects vanish.

\par\noindent\textit{Step 2.} The corrected lift extends as an equivariant holomorphic vector on the finite discs. A holomorphic local lift exists on the whole root disc by lifting the translating section through its covering vector bundle. Its defect is in $\ker N_j=(I-A_j)\Lambda$; subtracting an integral period kills that defect while preserving holomorphy. The corrected global lift and this local defect-zero lift differ by an invariant period, which is holomorphic through the centre. At the cusp, the local toric logarithm in (B2.3) is $sq+\eta(t_c)$. The chosen global lift differs from it by $\Pi_c(s)\lambda=sB_0\bar\lambda+C(t_c)\bar\lambda+\lambda_{\rm tor}$. Defect zero gives exactly $q+B_0\bar\lambda=0$, so the logarithmic terms cancel and the remaining vector is holomorphic. This proves the extension at the cusp as well. We therefore obtain a global section of $\mathcal V$. Its exact sequence in \ref{fam:A.3.1} gives $H^0(B,\mathcal V)=\mathbb Ce$. Conversely each such section defines the displayed translations on the finite covers and torus multiplication at the cusp, hence a global automorphism. Its kernel consists of global invariant periods. The common fixed lattice is $\mathbb Zd$, and $\Pi_cd=e$, giving kernel exactly $\mathbb Z$. This proves the full group assertion and its connectedness.
\end{proof}


\section{Sources, status, credit and notation}\label{app:audits}

\subsection{Per-result source and status concordance}\label{app:source-status}
The complete statement text is in the referenced result and, verbatim as TeX, in \path{S6_07_PAPER_source_table.csv}. The table below gives one row per printed statement or explicitly recorded setup. T, K and A are the binding statuses; \emph{text} denotes definitions or setup. The credit checks and precise wording changes are retained per row in that CSV, with the full source-to-print diffs in the accompanying records. Source filenames are relative to the supplied bundles.
\begingroup\small\setlength{\tabcolsep}{3pt}
\begin{longtable}{@{}L{.18\linewidth}L{.48\linewidth}L{.29\linewidth}@{}}
\toprule Result / status & Statement locator and source & Kind / proof\\\midrule\endhead
FAM \nolinkurl{conventions}\newline text & Statement~\ref{def:conventions}. Base, lattice, periods, clockwise deck generators and transported reference.\newline \path{returns/S6_02_FAMILY/sections/conventions.tex} & Definition\newline §1\\
FAM \nolinkurl{A.1.1}\newline K & Statement~\ref{fam:A.1.1}. Intrinsic pencil, canonical ring and base preservation.\newline \path{returns/S6_02_FAMILY/sections/AB.tex} & New proof\newline Appendix B\\
DGKXv1 Theorem 7.2; thm:\allowbreak{}all-maps-factor\newline K & Statement~\ref{dgkx:all-maps-factor}. Every holomorphic map to P¹ factors uniquely through f.\newline \path{papers/arXiv_2610.00250v1_Du_Guo_Kusakabe_Xie.tex} & New proof\newline Appendix B (source re-derivation)\\
FAM \nolinkurl{B.1.1}\newline T & Statement~\ref{fam:B.1.1}. Elliptic-surface quotient and Mordell--Weil generator.\newline \path{returns/S6_02_FAMILY/sections/AB.tex} & Theorem\newline §1\\
FAM \nolinkurl{B.1.2}\newline T & Statement~\ref{fam:B.1.2}. The exact admissible domain for the reference quotient.\newline \path{returns/S6_02_FAMILY/sections/AB.tex} & Theorem\newline §1\\
FAM \nolinkurl{B.1.3}\newline T & Statement~\ref{fam:B.1.3}. The distinguished elliptic logarithm is a zero of the Weierstrass function.\newline \path{returns/S6_02_FAMILY/sections/AB.tex} & New proof\newline §1\\
FAM \nolinkurl{B.1.4}\newline T & Statement~\ref{fam:B.1.4}. The transported reference and the square relation between parameters.\newline \path{returns/S6_02_FAMILY/sections/AB.tex} & Theorem\newline §1\\
FAM \nolinkurl{B.2.3}\newline T & Statement~\ref{fam:B.2.3}. Correction of the global descent claim in Engel Remark 8.1.\newline \path{returns/S6_02_FAMILY/sections/AB.tex} & Correction\newline §1\\
FAM \nolinkurl{A.5.1}\newline K & Statement~\ref{fam:A.5.1}. The known exact identification for the standard data.\newline \path{returns/S6_02_FAMILY/sections/AB.tex} & New proof\newline Appendices A and B\\
FAM \nolinkurl{B.2.1}\newline T & Statement~\ref{fam:B.2.1}. Complete invariants for all admissible discrete data.\newline \path{returns/S6_02_FAMILY/sections/AB.tex} & Theorem\newline §2\\
FAM \nolinkurl{B.2.2}\newline T & Statement~\ref{fam:B.2.2}. Two classes at a fixed Engel parameter.\newline \path{returns/S6_02_FAMILY/sections/AB.tex} & Theorem\newline §2\\
FAM \nolinkurl{B.3.1}\newline T & Statement~\ref{fam:B.3.1}. Identification subgroup for the stated higher-multiplicity model.\newline \path{returns/S6_02_FAMILY/sections/AB.tex} & Theorem\newline §2\\
PUSH \nolinkurl{E.1.1}\newline A & Statement~\ref{push:E.1.1}. The coordinate and primitive-circle dictionary.\newline \path{returns/S6_02_PUSH/sections/E_push.tex} & Theorem\newline Repository:\allowbreak{} S6\_\allowbreak{}02\_\allowbreak{}PUSH E.1.1\\
PUSH \nolinkurl{E.1.2}\newline A & Statement~\ref{push:E.1.2}. The higher-multiplicity parameter dictionary.\newline \path{returns/S6_02_PUSH/sections/E_push.tex} & Theorem\newline Repository:\allowbreak{} S6\_\allowbreak{}02\_\allowbreak{}PUSH E.1.2\\
FAM \nolinkurl{A.1.2}\newline T & Statement~\ref{fam:A.1.2}. Normalized finite base changes and product covers.\newline \path{returns/S6_02_FAMILY/sections/AB.tex} & Theorem\newline Appendix B\\
FAM \nolinkurl{A.2.1}\newline K & Statement~\ref{fam:A.2.1}. Centralizer, finite norms and the parity obstruction.\newline \path{returns/S6_02_FAMILY/sections/AB.tex} & New proof\newline Appendix B\\
FAM \nolinkurl{A.3.1}\newline K & Statement~\ref{fam:A.3.1}. Coherent even shifts and the proper quotient family.\newline \path{returns/S6_02_FAMILY/sections/AB.tex} & New proof\newline Appendix B\\
FAM \nolinkurl{A.4.1}\newline T & Statement~\ref{fam:A.4.1}. Albanese lattices and the generic two-invariant comparison.\newline \path{returns/S6_02_FAMILY/sections/AB.tex} & Theorem\newline Appendix B\\
FAM \nolinkurl{A.4.2}\newline K & Statement~\ref{fam:A.4.2}. Nonvanishing of the Kodaira--Spencer class.\newline \path{returns/S6_02_FAMILY/sections/AB.tex} & New proof\newline Appendix B\\
WORKBENCH \nolinkurl{pshift:lem:centralizer}\newline K & Statement~\ref{pshift:lem:centralizer}. Complete rational centralizer and middle-family linear parts.\newline \path{workbench/period_shift_classification.tex} & New proof\newline Appendix A\\
WORKBENCH \nolinkurl{pshift:thm:classification}\newline K & Statement~\ref{pshift:thm:classification}. Exact global identification of the retained period parameter.\newline \path{workbench/period_shift_classification.tex} & New proof\newline Appendix A\\
FAM \nolinkurl{C.1.1}\newline T & Statement~\ref{fam:C.1.1}. Theta formula, distinguished Weierstrass zero and special values.\newline \path{returns/S6_02_FAMILY/sections/CD.tex} & Theorem\newline §3\\
FAM \nolinkurl{C.1.2}\newline T & Statement~\ref{fam:C.1.2}. Special-value evaluations and the period formula.\newline \path{returns/S6_02_FAMILY/sections/CD.tex} & Computation\newline §3\\
ARI \nolinkurl{C.1.1}\newline T & Statement~\ref{ari:C.1.1}. distribution under isogenies\newline \path{returns/S6_05_ARITHMETIC/sections/C_special.tex} & Theorem\newline §3\\
ARI \nolinkurl{C.1.2}\newline T & Statement~\ref{ari:C.1.2}. the phases and the two periods\newline \path{returns/S6_05_ARITHMETIC/sections/C_special.tex} & Theorem\newline §3\\
ARI \nolinkurl{C.2.1}\newline T & Statement~\ref{ari:C.2.1}. exact nomes and continuation radius\newline \path{returns/S6_05_ARITHMETIC/sections/C_special.tex} & Theorem\newline §3\\
ARI \nolinkurl{C.4.2}\newline T & Statement~\ref{ari:C.4.2}. the full division identity\newline \path{returns/S6_05_ARITHMETIC/sections/C_special.tex} & Theorem\newline §3\\
ARI \nolinkurl{C.4.3}\newline T & Statement~\ref{ari:C.4.3}. algebraicity and explicit fields for every order\newline \path{returns/S6_05_ARITHMETIC/sections/C_special.tex} & Theorem\newline §3\\
ARI \nolinkurl{C.4.4}\newline T & Statement~\ref{ari:C.4.4}. the five values and all orders through eight\newline \path{returns/S6_05_ARITHMETIC/sections/C_special.tex} & Theorem\newline §3\\
ARI \nolinkurl{C.8.6}\newline T & Statement~\ref{ari:C.8.6}. the Mahler-measure identities\newline \path{returns/S6_05_ARITHMETIC/sections/C_integral.tex} & Theorem\newline §3\\
ARI \nolinkurl{C.8.7}\newline T & Statement~\ref{ari:C.8.7}. the two complex beta integrals\newline \path{returns/S6_05_ARITHMETIC/sections/C_integral.tex} & Theorem\newline §3\\
ARI \nolinkurl{C.8.8}\newline A & Statement~\ref{ari:C.8.8}. the signed Gaussian identity\newline \path{returns/S6_05_ARITHMETIC/sections/C_integral.tex} & Theorem\newline §3 (mechanism only; repository proof)\\
ARI \nolinkurl{C.8.9}\newline A & Statement~\ref{ari:C.8.9}. C.8.9; Corrected sign and measure normalization\newline \path{returns/S6_05_ARITHMETIC/sections/C_measure.tex} & Theorem\newline §3 (mechanism only; repository proof)\\
ARI \nolinkurl{C.8.10}\newline A & Statement~\ref{ari:C.8.10}. the correctly normalized equidistribution\newline \path{returns/S6_05_ARITHMETIC/sections/C_measure.tex} & Theorem\newline §3 (mechanism only; repository proof)\\
ARI \nolinkurl{C.8.11}\newline A & Statement~\ref{ari:C.8.11}. the two first-derivative identities\newline \path{returns/S6_05_ARITHMETIC/sections/C_measure.tex} & Theorem\newline §3 (mechanism only; repository proof)\\
FAM \nolinkurl{C.2.3}\newline T & Statement~\ref{fam:C.2.3}. The exact maximum and the radius of the parameter disc.\newline \path{returns/S6_02_FAMILY/sections/CD.tex} & Computation\newline §4\\
PUSH \nolinkurl{E.4.1}\newline T & Statement~\ref{push:E.4.1}. independent certification of every box\newline \path{returns/S6_02_PUSH/sections/E4_push.tex} & Computation\newline §4\\
ARI \nolinkurl{C.3.1}\newline T & Statement~\ref{ari:C.3.1}. fourth order at the maximum\newline \path{returns/S6_05_ARITHMETIC/sections/C_special.tex} & Theorem\newline §4\\
ARI \nolinkurl{C.8.1}\newline T & Statement~\ref{ari:C.8.1}. three uniform bounds on the elliptic norm\newline \path{returns/S6_05_ARITHMETIC/sections/C_analytic.tex} & Theorem\newline §4\\
ARI \nolinkurl{C.8.2}\newline T & Statement~\ref{ari:C.8.2}. an explicit local maximum on a larger disc\newline \path{returns/S6_05_ARITHMETIC/sections/C_analytic.tex} & Theorem\newline §4\\
ARI \nolinkurl{C.8.3}\newline T & Statement~\ref{ari:C.8.3}. the comparison of the two local discs\newline \path{returns/S6_05_ARITHMETIC/sections/C_analytic.tex} & Theorem\newline §4\\
ARI \nolinkurl{C.8.4}\newline T & Statement~\ref{ari:C.8.4}. the first three normalized duplication steps\newline \path{returns/S6_05_ARITHMETIC/sections/C_analytic.tex} & Theorem\newline §4\\
ARI \nolinkurl{C.8.5}\newline T & Statement~\ref{ari:C.8.5}. Maximum by explicit analytic estimates\newline \path{returns/S6_05_ARITHMETIC/sections/C_analytic.tex} & Theorem\newline §4\\
ARI \nolinkurl{C.8.15}\newline A & Statement~\ref{ari:C.8.15}. every critical point and every level set\newline \path{returns/S6_05_ARITHMETIC/sections/C_critical.tex} & Theorem\newline §4 (mechanism only; repository proof)\\
FAM C.1 setup\newline text & \nolinkurl{periods:theta-definitions}. C.1 setup\newline \path{returns/S6_02_FAMILY/sections/CD.tex} & Definition/\allowbreak{}record\newline Setup\\
FAM C.1.2 artifact-availability record\newline text & \nolinkurl{periods:special-computation-availability}. C.1.2 artifact-availability record\newline \path{returns/S6_02_FAMILY/sections/CD.tex} & Artifact record\newline Setup\\
ARI C.1 setup; C.8 critical setup\newline text & \nolinkurl{periods:norm-definition}. C.1 setup; C.8 critical setup\newline \path{returns/S6_05_ARITHMETIC/sections/C_special.tex}\newline \path{returns/S6_05_ARITHMETIC/sections/C_critical.tex} & Definition/\allowbreak{}record\newline Setup\\
FAM C.2.1 proof, eq:\allowbreak{}cd-norm\newline text & \nolinkurl{fam:eq:cd-norm}. C.2.1 proof, eq:\allowbreak{}cd-norm\newline \path{returns/S6_02_FAMILY/sections/CD.tex} & Definition/\allowbreak{}record\newline Setup\\
ARI C.4 setup and local-height notation in C.4.1\newline text & \nolinkurl{periods:division-model}. C.4 setup and local-height notation in C.4.1\newline \path{returns/S6_05_ARITHMETIC/sections/C_special.tex} & Definition/\allowbreak{}record\newline Setup\\
ARI C.8 measure conventions\newline text & \nolinkurl{periods:measure-conventions}. C.8 measure conventions\newline \path{returns/S6_05_ARITHMETIC/sections/C_measure.tex} & Definition/\allowbreak{}record\newline Setup\\
FAM;PUSH FAM C.2.3; PUSH E.4.1--3; REVIEW\_\allowbreak{}S6\_\allowbreak{}02\_\allowbreak{}FAMILY §2\newline text & \nolinkurl{periods:certificate-record}. FAM C.2.3; PUSH E.4.1--3; REVIEW\_\allowbreak{}S6\_\allowbreak{}02\_\allowbreak{}FAMILY §2\newline \path{returns/S6_02_FAMILY/sections/CD.tex}\newline \path{returns/S6_02_PUSH/sections/E4_push.tex}\newline \path{status/REVIEW_S6_02_FAMILY.md} & Computation record\newline Setup\\
ARI C.3 setup\newline text & \nolinkurl{periods:local-expansion-constants}. C.3 setup\newline \path{returns/S6_05_ARITHMETIC/sections/C_special.tex} & Definition/\allowbreak{}record\newline Setup\\
ARI C.8 analytic setup\newline text & \nolinkurl{periods:ladder-modular-conventions}. C.8 analytic setup\newline \path{returns/S6_05_ARITHMETIC/sections/C_analytic.tex} & Definition/\allowbreak{}record\newline Setup\\
INV \nolinkurl{notation-cotangent}\newline text & Statement~\ref{inv:notation-cotangent}. Cotangent notation\newline \path{returns/S6_01_INVARIANTS/fragments/part_A.tex} & Definition\newline definition\\
INV \nolinkurl{A1.0}\newline T & Statement~\ref{inv:A1.0}. The structure-sheaf input in the precise form used below\newline \path{returns/S6_01_INVARIANTS/fragments/part_A.tex} & New proof\newline main text\\
INV \nolinkurl{A1.1}\newline T & Statement~\ref{inv:A1.1}. Local freeness and the determinant line\newline \path{returns/S6_01_INVARIANTS/fragments/part_A.tex} & Theorem\newline main text\\
INV \nolinkurl{A1.2}\newline T & Statement~\ref{inv:A1.2}. The logarithmic direct images by vertical twisting and duality\newline \path{returns/S6_01_INVARIANTS/fragments/part_A.tex} & Theorem\newline main text\\
INV \nolinkurl{notation-period-branch}\newline text & Statement~\ref{inv:notation-period-branch}. The order-four branch\newline \path{returns/S6_01_INVARIANTS/fragments/part_A.tex} & Definition\newline definition\\
INV \nolinkurl{A2.1}\newline T & Statement~\ref{inv:A2.1}. The complete raw period determinant\newline \path{returns/S6_01_INVARIANTS/fragments/part_A.tex} & Theorem\newline main text\\
INV \nolinkurl{A2.2}\newline T & Statement~\ref{inv:A2.2}. The evaluated scalar and the two local leading terms\newline \path{returns/S6_01_INVARIANTS/fragments/part_A.tex} & Theorem\newline main text\\
INV \nolinkurl{A2.4}\newline T & Statement~\ref{inv:A2.4}. The global rank and Lead 2\newline \path{returns/S6_01_INVARIANTS/fragments/part_A.tex} & Theorem\newline main text\\
INV \nolinkurl{A1.3}\newline T & Statement~\ref{inv:A1.3}. All absolute cotangent direct images\newline \path{returns/S6_01_INVARIANTS/fragments/part_A.tex} & Theorem\newline main text\\
INV \nolinkurl{A3.1}\newline T & Statement~\ref{inv:A3.1}. The complete Hodge table for every admissible parameter\newline \path{returns/S6_01_INVARIANTS/fragments/part_A.tex} & Theorem\newline main text\\
INV \nolinkurl{notation-hodge-input}\newline text & Statement~\ref{inv:notation-hodge-input}. The formal Hodge parameter\newline \path{returns/S6_01_INVARIANTS/fragments/part_C.tex} & Definition\newline definition\\
INV \nolinkurl{C1.1}\newline T & Statement~\ref{inv:C1.1}. The complete Frölicher differential in terms of h.\newline \path{returns/S6_01_INVARIANTS/fragments/part_C.tex} & Theorem\newline main text\\
INV \nolinkurl{C2.1}\newline T & Statement~\ref{inv:C2.1}. Bott--Chern and Aeppli tables as functions of h.\newline \path{returns/S6_01_INVARIANTS/fragments/part_C.tex} & Theorem\newline main text\\
INV \nolinkurl{C3.3}\newline A & Statement~\ref{inv:C3.3}. The complete Bott--Chern algebra, for every h.\newline \path{returns/S6_01_INVARIANTS/fragments/part_C.tex} & Theorem\newline repository:\allowbreak{} \path{returns/S6_01_INVARIANTS/fragments/part_C.tex}\\
INV \nolinkurl{B1.1}\newline T & Statement~\ref{inv:B1.1}. The relative cohomology comparison\newline \path{returns/S6_01_INVARIANTS/fragments/part_B.tex} & Theorem\newline main text\\
INV \nolinkurl{B2.1}\newline T & Statement~\ref{inv:B2.1-canonical}. Canonical bundles for all coprime multiplicities\newline \path{returns/S6_01_INVARIANTS/fragments/part_B.tex} & Theorem\newline main text\\
INV \nolinkurl{B1.2}\newline A & Statement~\ref{inv:B1.2}. All admissible twists have the same Hodge numbers\newline \path{returns/S6_01_INVARIANTS/fragments/part_B.tex} & Theorem\newline repository:\allowbreak{} \path{returns/S6_01_INVARIANTS/fragments/part_B.tex}\\
INV \nolinkurl{B2.1}\newline A & Statement~\ref{inv:B2.1-hodge}. Hodge numbers for all coprime multiplicities\newline \path{returns/S6_01_INVARIANTS/fragments/part_B.tex} & Theorem\newline repository:\allowbreak{} \path{returns/S6_01_INVARIANTS/fragments/part_B.tex}\\
INV \nolinkurl{B3.1}\newline T & Statement~\ref{inv:B3.1}. Connected cyclic covers and the cusp comparison\newline \path{returns/S6_01_INVARIANTS/fragments/part_B.tex} & Theorem\newline main text\\
INV \nolinkurl{B3.2}\newline T & Statement~\ref{inv:B3.2}. Cotangent direct images and Hodge numbers of the covers\newline \path{returns/S6_01_INVARIANTS/fragments/part_B.tex} & Theorem\newline main text\\
INV \nolinkurl{B3.3}\newline T & Statement~\ref{inv:B3.3}. Canonical rings and rational topology of the covers\newline \path{returns/S6_01_INVARIANTS/fragments/part_B.tex} & Theorem\newline main text\\
INV \nolinkurl{C3.5}\newline A & Statement~\ref{inv:C3.5}. Metric obstructions that do not require h.\newline \path{returns/S6_01_INVARIANTS/fragments/part_C.tex} & Theorem\newline repository:\allowbreak{} \path{returns/S6_01_INVARIANTS/fragments/part_C.tex}\\
INV \nolinkurl{C3.7}\newline A & Statement~\ref{inv:C3.7}. The fibre obstruction to SKT and astheno-Kähler metrics.\newline \path{returns/S6_01_INVARIANTS/fragments/part_C.tex} & Theorem\newline repository:\allowbreak{} \path{returns/S6_01_INVARIANTS/fragments/part_C.tex}\\
INV \nolinkurl{A4.1}\newline T & Statement~\ref{inv:A4.1}. All local modules needed to distinguish degrees one and two\newline \path{returns/S6_01_INVARIANTS/fragments/part_A.tex} & Theorem\newline main text\\
INV \nolinkurl{A4.2}\newline T & Statement~\ref{inv:A4.2}. The claimed degree-one torsion along the cusp\newline \path{returns/S6_01_INVARIANTS/fragments/part_A.tex} & Correction\newline main text\\
FAM \nolinkurl{A.5.2}\newline K & Statement~\ref{fam:A.5.2}. Tangent cohomology and the Kuranishi line\newline \path{returns/S6_02_FAMILY/sections/AB.tex} & New proof\newline Appendix B, tangent cohomology subsection\\
CDP \nolinkurl{notation-integral-twists}\newline text & Statement~\ref{cdp:notation-integral-twists}. The integral Picard normalization\newline \path{returns/S6_03_CDP/sections/part_C5.tex} & Definition\newline definition\\
CDP \nolinkurl{C5b}\newline T & Statement~\ref{cdp:C5b}. Integral twists and the normalization-map exception\newline \path{returns/S6_03_CDP/sections/part_C5.tex} & Theorem\newline main text\\
PUSH \nolinkurl{E.2.1}\newline T & Statement~\ref{push:E.2.1}. Tangent cohomology and local miniversality for coprime multiplicities\newline \path{returns/S6_02_PUSH/sections/E_push.tex} & Theorem\newline main text\\
FAM F4 consequence of A.5.2; binding row 6d\newline T & Statement~\ref{fam:A.5.2-corollary}. Algebraic dimension near the standard family\newline \path{status/REVIEW_S6_02_FAMILY.md} & Theorem\newline main text\\
FAM \nolinkurl{B.4.1}\newline K & Statement~\ref{fam:B.4.1}. The full automorphism group is C*.\newline \path{returns/S6_02_FAMILY/sections/AB.tex} & New proof\newline Appendix B, re-derivation of FAM B.4.1\\
FAM \nolinkurl{B.4.2}\newline T & Statement~\ref{fam:B.4.2}. Conjugation sends c to 1−conjugate(c), and the real members have split extended automorphism group.\newline \path{returns/S6_02_FAMILY/sections/AB.tex} & Theorem\newline §7\\
ARI \nolinkurl{C.2.2}\newline T & Statement~\ref{ari:C.2.2}. Real structures occur exactly at nonzero real Albanese nome q2.\newline \path{returns/S6_05_ARITHMETIC/sections/C_special.tex} & Theorem\newline §7\\
MAPS A.3 (opening conventions)\newline text & Statement~\ref{maps:A.3.setup}. Cusp double curves and the real circle action.\newline \path{returns/S6_04_MAPS_MANIFOLDS/sections/real_forms.tex}\newline \path{returns/S6_02_FAMILY/sections/AB.tex} & Definition\newline §7\\
MAPS \nolinkurl{A.3.1}\newline A & Statement~\ref{maps:A.3.1}. There are two holomorphic-conjugacy classes of antiholomorphic involutions.\newline \path{returns/S6_04_MAPS_MANIFOLDS/sections/real_forms.tex} & Theorem\newline §7; proof in repository:\allowbreak{} MAPS A.3.1\\
MAPS A.3, equations (A.3.1)--(A.3.2)\newline text & Statement~\ref{maps:A.3.representatives}. Explicit roots, branches and translating vectors for the two real representatives.\newline \path{returns/S6_04_MAPS_MANIFOLDS/sections/real_forms.tex} & Definition\newline §7\\
MAPS \nolinkurl{A.3.2}\newline A & Statement~\ref{maps:A.3.2}. Exact real fixed sets on the ordinary, multiple and cusp fibres.\newline \path{returns/S6_04_MAPS_MANIFOLDS/sections/real_forms.tex} & Theorem\newline §7; proof in repository:\allowbreak{} MAPS A.3.2\\
MAPS \nolinkurl{A.3.3}\newline A & Statement~\ref{maps:A.3.3}. Every nonempty real fixed manifold is S3, with the specified circle action.\newline \path{returns/S6_04_MAPS_MANIFOLDS/sections/real_forms.tex} & Theorem\newline §7; proof in repository:\allowbreak{} MAPS A.3.3\\
FAM F (opening conventions)\newline text & Statement~\ref{fam:F.setup}. The indefinite period domain and the convention for proper discontinuity.\newline \path{returns/S6_02_FAMILY/sections/F.tex} & Definition\newline §8\\
FAM \nolinkurl{F.1.1}\newline T & Statement~\ref{fam:F.1.1}. The monodromy group is an explicit index-12 character kernel.\newline \path{returns/S6_02_FAMILY/sections/F.tex} & Theorem\newline §8\\
FAM \nolinkurl{F.1.2}\newline T & Statement~\ref{fam:F.1.2}. The centre, extension cocycle and Zariski closures are explicit.\newline \path{returns/S6_02_FAMILY/sections/F.tex} & Theorem\newline §8\\
FAM F.2 (definition preceding F.2.1)\newline text & Statement~\ref{fam:F.2.setup}. The invariant function q\_\allowbreak{}d and its two nonzero sign regions.\newline \path{returns/S6_02_FAMILY/sections/F.tex} & Definition\newline §8\\
FAM \nolinkurl{F.2.1}\newline T & Statement~\ref{fam:F.2.1}. The action is not properly discontinuous on the entire domain.\newline \path{returns/S6_02_FAMILY/sections/F.tex} & Counterexample\newline §8\\
FAM \nolinkurl{F.2.2}\newline T & Statement~\ref{fam:F.2.2}. The complement of q\_\allowbreak{}d=0 is the largest invariant proper-discontinuity domain.\newline \path{returns/S6_02_FAMILY/sections/F.tex} & Theorem\newline §8\\
FAM \nolinkurl{F.2.3}\newline T & Statement~\ref{fam:F.2.3}. All admissible period images lie in a global chart of the positive region.\newline \path{returns/S6_02_FAMILY/sections/F.tex} & Theorem\newline §8\\
FAM \nolinkurl{F.3.1}\newline T & Statement~\ref{fam:F.3.1}. The period curve is an injective immersion and has a proper embedded extension.\newline \path{returns/S6_02_FAMILY/sections/F.tex} & Theorem\newline §8\\
FAM \nolinkurl{F.3.2}\newline T & Statement~\ref{fam:F.3.2}. The finite limits lie in the domain, while the cusp limit lies outside it.\newline \path{returns/S6_02_FAMILY/sections/F.tex} & Theorem\newline §8\\
FAM \nolinkurl{F.3.3}\newline T & Statement~\ref{fam:F.3.3}. The determinant of the derivative is generically nonzero and vanishes at z1.\newline \path{returns/S6_02_FAMILY/sections/F.tex} & Theorem\newline §8\\
FAM \nolinkurl{D.1.1}\newline T & Statement~\ref{fam:D.1.1}. The first projection descends exactly at the six-division points.\newline \path{returns/S6_02_FAMILY/sections/CD.tex} & Theorem\newline §9\\
FAM \nolinkurl{D.1.2}\newline T & Statement~\ref{fam:D.1.2}. The torsion locus is countable and dense, with exact intersection degree.\newline \path{returns/S6_02_FAMILY/sections/CD.tex} & Theorem\newline §9\\
FAM \nolinkurl{D.1.3}\newline T & Statement~\ref{fam:D.1.3}. An unrestricted elliptic-quotient wording is refuted at t=32/\allowbreak{}5.\newline \path{returns/S6_02_FAMILY/sections/CD.tex} & Counterexample\newline §9\\
FAM D.1.6 (proof, Hermitian-form convention)\newline text & Statement~\ref{fam:D.1.Hermitian}. The transferred alternating form and its explicit Hermitian convention.\newline \path{returns/S6_02_FAMILY/sections/CD.tex} & Definition\newline §9\\
FAM \nolinkurl{D.1.4}\newline T & Statement~\ref{fam:D.1.4}. The six-coefficient Néron--Severi equation and the indefinite generator.\newline \path{returns/S6_02_FAMILY/sections/CD.tex} & New proof\newline §9\\
FAM \nolinkurl{D.1.5}\newline T & Statement~\ref{fam:D.1.5}. Generic rank one and explicit abelian fibres arbitrarily far into the c half-plane.\newline \path{returns/S6_02_FAMILY/sections/CD.tex} & Theorem\newline §9\\
FAM \nolinkurl{D.1.6}\newline T & Statement~\ref{fam:D.1.6}. Exact elliptic-quotient, projective and algebraic-dimension loci.\newline \path{returns/S6_02_FAMILY/sections/CD.tex} & Theorem\newline §9\\
FAM \nolinkurl{D.1.7}\newline T & Statement~\ref{fam:D.1.7}. Actual algebraic-dimension-one fibres, including nontorsion P.\newline \path{returns/S6_02_FAMILY/sections/CD.tex} & Theorem\newline §9\\
FAM \nolinkurl{D.2.1}\newline T & Statement~\ref{fam:D.2.1}. Every compact curve is vertical and every compact surface is a reduced fibre.\newline \path{returns/S6_02_FAMILY/sections/CD.tex} & Theorem\newline §9\\
FAM \nolinkurl{D.2.2}\newline T & Statement~\ref{fam:D.2.2}. Rational and elliptic curves and the singular-fibre descriptions.\newline \path{returns/S6_02_FAMILY/sections/CD.tex} & Theorem\newline §9\\
FAM \nolinkurl{D.2.3}\newline T & Statement~\ref{fam:D.2.3}. All compact curves are classified by the algebraic dimension of their fibre.\newline \path{returns/S6_02_FAMILY/sections/CD.tex} & Theorem\newline §9\\
PUSH D (opening definitions)\newline text & Statement~\ref{push:D.1.setup}. The boundary parameter and its height.\newline \path{returns/S6_02_PUSH/sections/D_push.tex} & Definition\newline §10\\
PUSH \nolinkurl{D.1.1}\newline T & Statement~\ref{push:D.1.1}. The Albanese j-invariant has a natural boundary at |w|=216.\newline \path{returns/S6_02_PUSH/sections/D_push.tex} & Theorem\newline §10\\
PUSH \nolinkurl{D.1.2}\newline T & Statement~\ref{push:D.1.2}. An intrinsic reduced anticanonical surface family exists off a countable exceptional set.\newline \path{returns/S6_02_PUSH/sections/D_push.tex} & Theorem\newline §10\\
PUSH \nolinkurl{D.1.3}\newline T & Statement~\ref{push:D.1.3}. No flat proper analytic extension crosses any point of the outer circle, even fibrewise.\newline \path{returns/S6_02_PUSH/sections/D_push.tex} & Theorem\newline §10\\
PUSH A.1 (opening definitions)\newline text & Statement~\ref{push:A.1.setup}. The ruled normalization, specified pinching Q and linearization line.\newline \path{returns/S6_02_PUSH/sections/AB_push.tex} & Definition\newline §10\\
E26 Conjecture 8.2\newline false as worded & Statement~\ref{e26:Conjecture8.2}. The specified nonnormal pinching is conjectured to fill the untwisted family.\newline \path{papers/arXiv_2609.38442_Engel_src/S6.tex} & Counterexample\newline Refuted by PUSH A.3.2 in text\\
PUSH \nolinkurl{A.3.1}\newline T & Statement~\ref{push:A.3.1}. The intrinsic virtual Euler number equals that of every smooth proper-flat smoothing.\newline \path{returns/S6_02_PUSH/sections/AB_push.tex} & Theorem\newline §10\\
PUSH \nolinkurl{A.3.2}\newline T & Statement~\ref{push:A.3.2}. The specified pinching has virtual Euler number zero and cannot fill Engel's family.\newline \path{returns/S6_02_PUSH/sections/AB_push.tex} & Counterexample\newline §10\\
PUSH \nolinkurl{A.3.2}\newline A & Statement~\ref{push:A.3.2.replacement}. The nodal and two semistable constructions are precise replacements with different central fibres.\newline \path{returns/S6_02_PUSH/sections/AB_push.tex} & Counterexample\newline Repository:\allowbreak{} PUSH A.1.2 and A.2.2\\
PUSH A.1.1 (proof, cone definition)\newline text & Statement~\ref{push:A.1.fan}. The cones of the two-parameter cusp fan.\newline \path{returns/S6_02_PUSH/sections/AB_push.tex} & Definition\newline §10\\
PUSH \nolinkurl{A.1.1}\newline A & Statement~\ref{push:A.1.1}. The unsubdivided cusp chart is xy=uzv, with cusp coordinate t\_\allowbreak{}c=zv.\newline \path{returns/S6_02_PUSH/sections/AB_push.tex} & Theorem\newline §10; proof in repository:\allowbreak{} PUSH A.1.1\\
PUSH \nolinkurl{A.1.2}\newline A & Statement~\ref{push:A.1.2}. A proper nodal filling contracts the specified (-1,-1) curve on each nonzero fibre.\newline \path{returns/S6_02_PUSH/sections/AB_push.tex} & Theorem\newline §10; proof in repository:\allowbreak{} PUSH A.1.2\\
PUSH A.2.1 (proof, six rays and bottom diagonal)\newline text & Statement~\ref{push:A.2.rays}. The six cusp rays and retained bottom diagonal.\newline \path{returns/S6_02_PUSH/sections/AB_push.tex} & Definition\newline §10\\
PUSH \nolinkurl{A.2.1}\newline A & Statement~\ref{push:A.2.1}. Exactly two subdivisions satisfy the joint-flatness and reducedness restrictions.\newline \path{returns/S6_02_PUSH/sections/AB_push.tex} & Theorem\newline §10; proof in repository:\allowbreak{} PUSH A.2.1\\
PUSH \nolinkurl{A.2.2}\newline A & Statement~\ref{push:A.2.2}. The two natural semistable fillings have the stated normalized central fibres.\newline \path{returns/S6_02_PUSH/sections/AB_push.tex} & Theorem\newline §10; proof in repository:\allowbreak{} PUSH A.2.2\\
PUSH \nolinkurl{A.2.3}\newline A & Statement~\ref{push:A.2.3}. The unrestricted blow-up shape is false; star charts supply its replacement.\newline \path{returns/S6_02_PUSH/sections/AB_push.tex} & Counterexample\newline §10; proof in repository:\allowbreak{} PUSH A.2.3\\
PUSH A.2.4 (proof, base-changed period shears)\newline text & Statement~\ref{push:A.2.periods}. The period shears after the base change u=r\textasciicircum{}n.\newline \path{returns/S6_02_PUSH/sections/AB_push.tex} & Definition\newline §10\\
PUSH \nolinkurl{A.2.4}\newline A & Statement~\ref{push:A.2.4}. Finite fan data enumerate all invariant smooth refinements in the stated scope.\newline \path{returns/S6_02_PUSH/sections/AB_push.tex} & Theorem\newline §10; proof in repository:\allowbreak{} PUSH A.2.4\\
PUSH \nolinkurl{A.2.5}\newline A & Statement~\ref{push:A.2.5}. A terminating blow-up/\allowbreak{}blow-down algorithm describes normalization comparisons.\newline \path{returns/S6_02_PUSH/sections/ab_push_normalization_algorithm.tex} & Counterexample\newline §10; proof in repository:\allowbreak{} PUSH A.2.5\\
PUSH Conventions; B.4.33 (scope)\newline text & Statement~\ref{push:B.4.33.scope}. The fixed sphere family and its distinction from Engel's untwisted family.\newline \path{returns/S6_02_PUSH/sections/conventions_push.tex} & Definition\newline §11\\
PUSH \nolinkurl{B.4.33}\newline A & Statement~\ref{push:B.4.33}. The retained family has no proper Hausdorff analytic filling after a finite parameter cover.\newline \path{returns/S6_02_PUSH/sections/B_global_noncompletion.tex} & Theorem\newline §11; proof in repository:\allowbreak{} PUSH B.4.33\\
PUSH Conventions (mapping-torus orientation)\newline text & Statement~\ref{push:C.orientation}. Mapping-torus, sphere-gluing and Theta7 orientation conventions.\newline \path{returns/S6_02_PUSH/sections/conventions_push.tex} & Definition\newline §11\\
PUSH C.4.1 (scope)\newline text & Statement~\ref{push:C.4.1.scope}. The exact twist-comparison scope used in PUSH C.3.24.\newline \path{returns/S6_02_PUSH/sections/C_push.tex} & Definition\newline §11\\
PUSH \nolinkurl{C.3.24}\newline A & Statement~\ref{push:C.3.24}. The actual positive-loop mapping class is zero in Theta7, by the orbit-knot method.\newline \path{returns/S6_02_PUSH/sections/C_monodromy_value.tex} & Theorem\newline §11; proof in repository:\allowbreak{} PUSH C.3.24\\
PUSH \nolinkurl{C.3.27}\newline A & Statement~\ref{push:C.3.27}. A spin-c coboundary and direct index additivity give a second zero determination.\newline \path{returns/S6_02_PUSH/sections/C_independent_monodromy_value.tex} & Theorem\newline §11; proof in repository:\allowbreak{} PUSH C.3.27\\
MAPS \nolinkurl{B.6.21}\newline A & Statement~\ref{maps:B.6.21}. The MAPS calculation gives the original (3,4) positive-loop class zero.\newline \path{returns/S6_04_MAPS_MANIFOLDS/sections/maps_B.tex} & Theorem\newline §11; proof in repository:\allowbreak{} MAPS B.6.21\\
PUSH E.1.1 (sphere data); FAM B.3.1 (parameter lattice)\newline text & Statement~\ref{push:E.1.1.scope}. The fixed coprime sphere-pair and twist scope of the higher-family announcements.\newline \path{returns/S6_02_PUSH/sections/E_push.tex}\newline \path{returns/S6_02_FAMILY/sections/AB.tex} & Definition\newline §11\\
MAPS \nolinkurl{B.6.24}\newline A & Statement~\ref{maps:B.6.24}. All stated coprime sphere-pair monodromy classes vanish in Theta7.\newline \path{returns/S6_04_MAPS_MANIFOLDS/sections/maps_B.tex} & Theorem\newline §11; proof in repository:\allowbreak{} MAPS B.6.24\\
PUSH E.2.5 (hypotheses); E.2.7 (parameter)\newline text & Statement~\ref{push:E.2.5.scope}. The precise smooth-specialization hypotheses and fixed-pair parameter.\newline \path{returns/S6_02_PUSH/sections/E_smooth_specialization.tex}\newline \path{returns/S6_02_PUSH/sections/E_zero_period_obstruction.tex} & Definition\newline §11\\
PUSH \nolinkurl{E.2.21}\newline A & Statement~\ref{push:E.2.21}. A smooth completion is an actual member at an interior fixed-pair parameter.\newline \path{returns/S6_02_PUSH/sections/E_global_completion.tex} & Theorem\newline §11; proof in repository:\allowbreak{} PUSH E.2.21\\
PUSH \nolinkurl{E.2.22}\newline A & Statement~\ref{push:E.2.22}. A fixed constructed pair propagates from a nonempty open subset of a disc.\newline \path{returns/S6_02_PUSH/sections/E_disc_propagation.tex} & Theorem\newline §11; proof in repository:\allowbreak{} PUSH E.2.22\\
PUSH \nolinkurl{E.2.23}\newline A & Statement~\ref{push:E.2.23}. Proper smooth deformation chains preserve the ordered coprime multiplicity pair.\newline \path{returns/S6_02_PUSH/sections/E_arbitrary_deformation_bases.tex} & Theorem\newline §11; proof in repository:\allowbreak{} PUSH E.2.23\\
PUSH \nolinkurl{E.2.24}\newline A & Statement~\ref{push:E.2.24}. Each fixed-pair locus is open and closed in the space of integrable structures.\newline \path{returns/S6_02_PUSH/sections/E_continuous_deformations.tex} & Theorem\newline §11; proof in repository:\allowbreak{} PUSH E.2.24\\
\bottomrule\end{longtable}\endgroup
\subsection{Known and new, by section}\label{app:credits}
This table incorporates the source openings of AC26 that the supplied expert table had not performed. \emph{Selected run results} identifies the supplied proof source, not an assertion of absence from all literature.
\begingroup\small\setlength{\tabcolsep}{3pt}\begin{longtable}{@{}L{.05\linewidth}L{.34\linewidth}L{.56\linewidth}@{}}
\toprule § & Known source and locator & Scope and credit\\\midrule\endhead
1 & AC26 Main Theorem; Theorem 3.4; Remark 3.12; DGKXv1 Theorems 7.1--7.2; E26 Section 2. & Construction, canonical bundle, standard intrinsic pencil, factorization and elliptic-logarithm interpretation are known. FAM gives the explicit transported-reference dictionary and corrects E26 Remark 8.1. Its wp-zero geometric reading is credited to AC26.\\
2 & DGKXv1 Theorem 1.4 and bare-period corollary; AC26 Remark 6.4; Viaclovsky Proposition 4.7. & The standard-data 2Z statement is known. The all-discrete-data invariants and two classes at fixed u are the selected FAM results. The 2klZ theorem retains its stated model; the Viaclovsky coordinate dictionary is A.\\
3 & AC26 Theorem 3.4 and Remark 3.12; Eichler--Zagier p.399; Duke--Imamoglu Section 2; Shioda Theorem 8.6. & Finite period values and the elliptic-logarithm interpretation have their source credits. The normalized beta formula, exact nomes, division evaluations and specified integrals are the selected run results. The general division identity is classical local-height quasi-quadraticity, re-derived here; it is not credited as a new general theorem.\\
4 & DGKXv1 Corollary p2:\allowbreak{}cor:\allowbreak{}bare-period; E26 Proposition 2.4. & The abstract punctured disc and sufficiently small geometric parameter are known. The maximum, unique maximizer, exact radius 216 and normalized unit-disc nome are the selected new evaluations. Certificate scope and missing files are explicit.\\
5 & AC26 Proposition 9.13; Theorem 9.1; Remarks 1.3 and 9.21; DGKXv1 Proposition 6.10; Ugarte Theorem 3.2; McHugh Section 7.1; Stelzig Theorem A. & The structure-sheaf direct images, h0q, hp0 and h11=h12+1 are known. The selected cotangent calculation determines h12=1. INTAKE\_\allowbreak{}S6\_\allowbreak{}01 Section 4 also credits the numerical Hodge table to the workbench (GPT-5.6 Sol, Codex reasoning effort ultra), with a degree-count argument; INV supplies the closed scalar and determinant identity. Conditional background formulas for a hypothetical complex S6 are credited where applied; exact Frolicher ranks, evaluated Bott--Chern/\allowbreak{}Aeppli tables and specified cover calculations have their run proofs. Product and metric statements remain A.\\
6 & DGKXv1 Theorem 1.3 and tangent-cohomology theorem; standard proper base change and relative duality. & The tuple (1,1,1,0), Kuranishi line and miniversality are known and re-derived. Integral twists are CDP C5b and INV D4.1. Higher multiplicities carry the two-runs, no-second-reader qualifier. The nearby a=1 corollary is the specified consequence.\\
7 & AC26 Proposition 9.23 and Corollary 9.24; DGKXv1 full-automorphisms theorem. & Aut0 and its action are AC26; the full group is DGKXv1. The conjugate-parameter and real-nome criterion are the selected T results. Fixed loci and classes of real involutions are A.\\
8 & AC26 period monodromies; E26 Remark 8.4. & The construction matrices and Klingler question have source credit. The explicit index-12 subgroup, maximal proper-action open set and period embedding are the selected FAM results.\\
9 & AC26 Lemmas 9.2--9.4, Proposition 9.5, Corollary 9.8, Proposition 4.6; E26 Lemma 2.1; Shioda and Appell--Humbert theory. & The NS equation, independence, generic indefinite class and generic NS rank are already in AC26; the normalization and conductor geometry of W are also there. These results are credited as known content despite the binding T/\allowbreak{}proof placement. The selected torsion-point calculations, all-curve/\allowbreak{}surface statements and special-parameter examples retain their stated scope.\\
10 & E26 Conjecture 8.2; Buchdahl Theorem 11; Grauert base-change and analytic flatness results; analytic specialization sources. & The outer-boundary obstruction and virtual Euler count refute the stated target in their exact source scope. The candidate centre has virtual Euler number zero whereas the smooth sphere fibre has Euler number two. Corrected fillings are A.\\
11 & The supplied PUSH and MAPS returns, with the bounded intake checks. & Zero-end non-filling, trivial Theta7 monodromy and unions of components are A. Two runs use different methods for the monodromy value. The interpretation J=epsilon P.O is an unverified reader remark.\\
12 & INTAKE\_\allowbreak{}FINAL Section 6; RESPONSE\_\allowbreak{}INTAKE; RESPONSE\_\allowbreak{}EXPERT. & The four open question groups are retained as questions, with the MAPS B.6.24 announced scope distinguished from the other families.\\
\bottomrule\end{longtable}\endgroup
\subsection{Dependency concordance}\label{app:dependencies}
The fixed section order retains the source arguments. A pointer to Appendix B is a dependency on its printed re-derivation. Forward references within the selected source order are retained and disclosed in Limits and open points. Every A row points to its repository proof rather than contributing a printed proof.
\begingroup\small\begin{longtable}{@{}L{.21\linewidth}L{.73\linewidth}@{}}
\toprule Result & Uses / dependence\\\midrule\endhead
FAM \nolinkurl{conventions} (\ref{def:conventions}) & AC26 §§2--6; E26 §§2--5\\
FAM \nolinkurl{A.1.1} (\ref{fam:A.1.1}) & Definition~\ref{def:conventions}; \cite[Theorem 7.1]{DGKXv1}; \cite[Theorem 9.1]{AC26}.\\
DGKXv1 Theorem 7.2; thm:\allowbreak{}all-maps-factor (\ref{dgkx:all-maps-factor}) & H²(S⁶;Z)=0; proper image theorem\\
FAM \nolinkurl{B.1.1} (\ref{fam:B.1.1}) & Definition~\ref{def:conventions}; FAM A.3.1 and preceding dictionary identities; \cite[Remark 3.12]{AC26}; \cite[Lemma 2.1; Corollary 2.2; Definition 2.3]{E26}.\\
FAM \nolinkurl{B.1.2} (\ref{fam:B.1.2}) & Definition~\ref{def:conventions}; FAM A.3.1 and preceding dictionary identities; \cite[Remark 3.12]{AC26}; \cite[Lemma 2.1; Corollary 2.2; Definition 2.3]{E26}.\\
FAM \nolinkurl{B.1.3} (\ref{fam:B.1.3}) & Definition~\ref{def:conventions}; FAM A.3.1 and preceding dictionary identities; \cite[Remark 3.12]{AC26}; \cite[Lemma 2.1; Corollary 2.2; Definition 2.3]{E26}.\\
FAM \nolinkurl{B.1.4} (\ref{fam:B.1.4}) & Definition~\ref{def:conventions}; FAM A.3.1 and preceding dictionary identities; \cite[Remark 3.12]{AC26}; \cite[Lemma 2.1; Corollary 2.2; Definition 2.3]{E26}.\\
FAM \nolinkurl{B.2.3} (\ref{fam:B.2.3}) & FAM A.2.1, A.3.1, A.4.1 and B.2.1; \cite[Remark 8.1]{E26}.\\
FAM \nolinkurl{A.5.1} (\ref{fam:A.5.1}) & FAM A.1.1, A.2.1, A.3.1 and A.4.2; \cite[Theorems 1.3, 1.4 and 7.1; punctured-disc corollary]{DGKXv1}.\\
FAM \nolinkurl{B.2.1} (\ref{fam:B.2.1}) & Definition~\ref{def:conventions}; FAM A.1.1--A.3.1 and the preceding discrete-data calculation; \cite[Section 6.4]{AC26}; \cite[Sections 3 and 5]{E26}.\\
FAM \nolinkurl{B.2.2} (\ref{fam:B.2.2}) & Definition~\ref{def:conventions}; FAM A.1.1--A.3.1 and the preceding discrete-data calculation; \cite[Section 6.4]{AC26}; \cite[Sections 3 and 5]{E26}.\\
FAM \nolinkurl{B.3.1} (\ref{fam:B.3.1}) & FAM A.1.2, A.2.1, B.2.1; \cite[Proposition 4.7]{VL26}.\\
PUSH \nolinkurl{E.1.1} (\ref{push:E.1.1}) & FAM B.3.1; PUSH E.1.1--E.1.2; \cite[Section 2]{VL26}.\\
PUSH \nolinkurl{E.1.2} (\ref{push:E.1.2}) & FAM B.3.1; PUSH E.1.1--E.1.2; \cite[Section 2]{VL26}.\\
FAM \nolinkurl{A.1.2} (\ref{fam:A.1.2}) & Definition~\ref{def:conventions}; FAM A.1.1 and preceding appendix calculations.\\
FAM \nolinkurl{A.2.1} (\ref{fam:A.2.1}) & Definition~\ref{def:conventions}; FAM A.1.1 and preceding appendix calculations.\\
FAM \nolinkurl{A.3.1} (\ref{fam:A.3.1}) & Definition~\ref{def:conventions}; FAM A.1.1 and preceding appendix calculations.\\
FAM \nolinkurl{A.4.1} (\ref{fam:A.4.1}) & Definition~\ref{def:conventions}; FAM A.1.1 and preceding appendix calculations.\\
FAM \nolinkurl{A.4.2} (\ref{fam:A.4.2}) & Definition~\ref{def:conventions}; FAM A.1.1 and preceding appendix calculations.\\
WORKBENCH \nolinkurl{pshift:lem:centralizer} (\ref{pshift:lem:centralizer}) & Retained period construction; canonical pencil; finite normalization; explicit centralizer and gluing\\
WORKBENCH \nolinkurl{pshift:thm:classification} (\ref{pshift:thm:classification}) & Retained period construction; canonical pencil; finite normalization; explicit centralizer and gluing\\
FAM \nolinkurl{C.1.1} (\ref{fam:C.1.1}) & FAM B.1.3; the cusp normalization; the Jacobi transformation laws\\
FAM \nolinkurl{C.1.2} (\ref{fam:C.1.2}) & FAM C.1.1; the theta product and its displayed tail bound\\
ARI \nolinkurl{C.1.1} (\ref{ari:C.1.1}) & the theta product and Haar probability measure\\
ARI \nolinkurl{C.1.2} (\ref{ari:C.1.2}) & ARI C.1.1; FAM C.1.1; the fixed cusp branch\\
ARI \nolinkurl{C.2.1} (\ref{ari:C.2.1}) & FAM A.4.1 and A.5.1; ARI C.1.2; the maximum theorem stated in the Result\\
ARI \nolinkurl{C.4.2} (\ref{ari:C.4.2}) & ARI C.1.1; the displayed rational Weierstrass model and standard division-polynomial normalization\\
ARI \nolinkurl{C.4.3} (\ref{ari:C.4.3}) & ARI C.4.2; the nonreal cusp value of the chosen elliptic logarithm\\
ARI \nolinkurl{C.4.4} (\ref{ari:C.4.4}) & ARI C.1.2, C.4.2 and C.4.3; the division-polynomial recurrence\\
ARI \nolinkurl{C.8.6} (\ref{ari:C.8.6}) & ARI C.1.1; the maximum theorem stated in the Result\\
ARI \nolinkurl{C.8.7} (\ref{ari:C.8.7}) & ARI C.8.6; the gamma duplication and triplication formulas\\
ARI \nolinkurl{C.8.8} (\ref{ari:C.8.8}) & the theta series and Poisson summation\\
ARI \nolinkurl{C.8.9} (\ref{ari:C.8.9}) & the normalized theta norm; local root expansions and the cusp branch\\
ARI \nolinkurl{C.8.10} (\ref{ari:C.8.10}) & ARI C.8.9; the height metric and its equidistribution theorem\\
ARI \nolinkurl{C.8.11} (\ref{ari:C.8.11}) & the theta heat equation and the Weierstrass differential equation\\
FAM \nolinkurl{C.2.3} (\ref{fam:C.2.3}) & FAM C.1.1 and C.2.2; the interval cover described below\\
PUSH \nolinkurl{E.4.1} (\ref{push:E.4.1}) & the supplied rational rectangle endpoints; the theta norm; Arb interval arithmetic\\
ARI \nolinkurl{C.3.1} (\ref{ari:C.3.1}) & the theta heat equation; the half-period Weierstrass equation; the cusp-normalized branch\\
ARI \nolinkurl{C.8.1} (\ref{ari:C.8.1}) & the theta product, Poisson summation and Fourier majorants\\
ARI \nolinkurl{C.8.2} (\ref{ari:C.8.2}) & ARI C.8.1 and C.3.1; the theta heat equation and Cauchy estimates\\
ARI \nolinkurl{C.8.3} (\ref{ari:C.8.3}) & ARI C.8.2; the Ramanujan differential identities and Rouche theorem\\
ARI \nolinkurl{C.8.4} (\ref{ari:C.8.4}) & ARI C.4.2, C.4.4 and C.3.1\\
ARI \nolinkurl{C.8.5} (\ref{ari:C.8.5}) & ARI C.8.1--C.8.4; the norm values of ARI C.1.2\\
ARI \nolinkurl{C.8.15} (\ref{ari:C.8.15}) & ARI C.8.11--C.8.14 and the two local root expansions\\
FAM C.1 setup (\nolinkurl{periods:theta-definitions}) & named source setup\\
FAM C.1.2 artifact-availability record (\nolinkurl{periods:special-computation-availability}) & named source setup\\
ARI C.1 setup; C.8 critical setup (\nolinkurl{periods:norm-definition}) & named source setup\\
FAM C.2.1 proof, eq:\allowbreak{}cd-norm (\nolinkurl{fam:eq:cd-norm}) & named source setup\\
ARI C.4 setup and local-height notation in C.4.1 (\nolinkurl{periods:division-model}) & named source setup\\
ARI C.8 measure conventions (\nolinkurl{periods:measure-conventions}) & named source setup\\
FAM;PUSH FAM C.2.3; PUSH E.4.1--3; REVIEW\_\allowbreak{}S6\_\allowbreak{}02\_\allowbreak{}FAMILY §2 (\nolinkurl{periods:certificate-record}) & named source setup\\
ARI C.3 setup (\nolinkurl{periods:local-expansion-constants}) & named source setup\\
ARI C.8 analytic setup (\nolinkurl{periods:ladder-modular-conventions}) & named source setup\\
INV \nolinkurl{notation-cotangent} (\ref{inv:notation-cotangent}) & The family and marking fixed in Section~\ref{sec:family}.\\
INV \nolinkurl{A1.0} (\ref{inv:A1.0}) & The cyclic and cusp charts; analytic base change \cite[III.4.10, III.4.12]{BS76}; analytic duality \cite[pp.~261, 264--266]{RRV71} and \cite[p.~91, Proposition~2]{RR70}.\\
INV \nolinkurl{A1.1} (\ref{inv:A1.1}) & The local cyclic and normal-crossing charts; inv:A1.0.\\
INV \nolinkurl{A1.2} (\ref{inv:A1.2}) & inv:A1.0, inv:A1.1; analytic duality \cite[pp.~261, 264--266]{RRV71}.\\
INV \nolinkurl{notation-period-branch} (\ref{inv:notation-period-branch}) & The period functions and the first-run clockwise convention.\\
INV \nolinkurl{A2.1} (\ref{inv:A2.1}) & inv:A1.0; the period transformations \cite[Definition 3.1 and Theorem 3.4(i)]{AC26}; the modular identities stated in the proof.\\
INV \nolinkurl{A2.2} (\ref{inv:A2.2}) & inv:A2.1; the already identified elliptic section fam:B.1.3; the displayed Ramanujan identities.\\
INV \nolinkurl{A2.4} (\ref{inv:A2.4}) & inv:A1.2, inv:A2.1, inv:A2.2; for the explicitly counterfactual final clause, the residue calculation in inv:A1.3.\\
INV \nolinkurl{A1.3} (\ref{inv:A1.3}) & inv:A1.0--inv:A1.2, inv:A2.1, inv:A2.2, inv:A2.4; the normalized residue sequence; analytic relative duality \cite{RRV71}.\\
INV \nolinkurl{A3.1} (\ref{inv:A3.1}) & inv:A1.0, inv:A1.3; coherent Leray on $\mathbb P^1$ and Serre duality.\\
INV \nolinkurl{notation-hodge-input} (\ref{inv:notation-hodge-input}) & inv:A1.0, inv:A1.3, inv:A3.1; sphere topology.\\
INV \nolinkurl{C1.1} (\ref{inv:C1.1}) & The Hodge table in inv:notation-hodge-input; convergence of the Fr\"olicher spectral sequence to de Rham cohomology.\\
INV \nolinkurl{C2.1} (\ref{inv:C2.1}) & inv:C1.1; the bounded-double-complex decomposition \cite[Theorem 3 and Proposition 6(1)]{Stelzig21}.\\
INV \nolinkurl{C3.3} (\ref{inv:C3.3}) & inv:C1.1, inv:C2.1; repository inv:C1.2 and inv:C3.1; the graded Leibniz rule.\\
INV \nolinkurl{B1.1} (\ref{inv:B1.1}) & The fixed period and cusp charts; inv:A1.0--inv:A1.2; smooth proper base change \cite[III.4.12]{BS76}.\\
INV \nolinkurl{B2.1} (\ref{inv:B2.1-canonical}) & inv:B1.1; \cite[Theorem 1.1 and Proposition 2.3]{VL26}; relative duality \cite{RRV71}.\\
INV \nolinkurl{B1.2} (\ref{inv:B1.2}) & inv:B1.1; the repository residue calculation and relative duality.\\
INV \nolinkurl{B2.1} (\ref{inv:B2.1-hodge}) & inv:B1.2; \cite[Theorem 1.1 and Proposition 2.3]{VL26}.\\
INV \nolinkurl{B3.1} (\ref{inv:B3.1}) & The admissible affine and cusp fillings; inv:B1.1; the elementary covering and normalization calculations in the proof; \cite[III.4.12]{BS76}.\\
INV \nolinkurl{B3.2} (\ref{inv:B3.2}) & inv:B3.1; inv:B1.2 (A, source dependency retained); inv:A1.3; derived relative duality \cite{RRV71}.\\
INV \nolinkurl{B3.3} (\ref{inv:B3.3}) & inv:B3.1, inv:B3.2; the source valuation and monomial argument of inv:B2.2, printed below for $(k,l)=(1,1)$; the marked lattice and affine data.\\
INV \nolinkurl{C3.5} (\ref{inv:C3.5}) & The fibration and sphere topology; Stokes' theorem.\\
INV \nolinkurl{C3.7} (\ref{inv:C3.7}) & inv:A3.1; the global Dolbeault lift in repository inv:A3.2; inv:C2.1.\\
INV \nolinkurl{A4.1} (\ref{inv:A4.1}) & The presentations of ordinary relative Kähler differentials for $t=xy$, $t=xyz$, and $t=s^m$.\\
INV \nolinkurl{A4.2} (\ref{inv:A4.2}) & inv:A4.1; \cite[Remarks 1.3 and 9.21]{AC26}.\\
FAM \nolinkurl{A.5.2} (\ref{fam:A.5.2}) & The compact family fam:A.5.1 and its nonzero Kodaira--Spencer class fam:A.4.2; \cite[Theorems 1.3 and 6.12]{DGKXv1}.\\
CDP \nolinkurl{notation-integral-twists} (\ref{cdp:notation-integral-twists}) & The standard clockwise affine data and exponential sequence; the root factorization.\\
CDP \nolinkurl{C5b} (\ref{cdp:C5b}) & The root and normalization calculations below; inv:D1.3 for the source-expanded branch-normal and conormal computation; analytic base change \cite[III.4.10, III.4.12]{BS76} and relative duality \cite{RRV71}.\\
PUSH \nolinkurl{E.2.1} (\ref{push:E.2.1}) & The admissible higher-multiplicity data push:E.1; the displayed root-chart and residue computations; \cite[III.4.12]{BS76} and \cite{RRV71}. Two runs, no second reader.\\
FAM F4 consequence of A.5.2; binding row 6d (\ref{fam:A.5.2-corollary}) & fam:A.5.2; \cite[Theorem 9.1(i)]{AC26}.\\
FAM \nolinkurl{B.4.1} (\ref{fam:B.4.1}) & FAM A.1.1, A.1.2, A.2.1, A.3.1 and the cusp valuation comparison in FAM B.2\\
FAM \nolinkurl{B.4.2} (\ref{fam:B.4.2}) & FAM A.3.1, A.5.1, B.2.1 and B.4.1; the cusp-normalized period laws\\
ARI \nolinkurl{C.2.2} (\ref{ari:C.2.2}) & FAM B.4.1, B.4.2 and A.5.1; ARI C.2.1\\
MAPS A.3 (opening conventions) (\ref{maps:A.3.setup}) & FAM B.4.1 and the cusp fan\\
MAPS \nolinkurl{A.3.1} (\ref{maps:A.3.1}) & FAM B.4.1 and B.4.2; the semifree circle action\\
MAPS A.3, equations (A.3.1)--(A.3.2) (\ref{maps:A.3.representatives}) & FAM B.4.2 and the specified root branches\\
MAPS \nolinkurl{A.3.2} (\ref{maps:A.3.2}) & MAPS A.3.1 and its explicit representative data; FAM B.4.2\\
MAPS \nolinkurl{A.3.3} (\ref{maps:A.3.3}) & MAPS A.3.1 and A.3.2\\
FAM F (opening conventions) (\ref{fam:F.setup}) & The retained integral marking and monodromy matrices\\
FAM \nolinkurl{F.1.1} (\ref{fam:F.1.1}) & The matrices $A_1,A_2$ and the integral symplectic form\\
FAM \nolinkurl{F.1.2} (\ref{fam:F.1.2}) & FAM F.1.1 and its exact Heisenberg kernel\\
FAM F.2 (definition preceding F.2.1) (\ref{fam:F.2.setup}) & FAM F.setup\\
FAM \nolinkurl{F.2.1} (\ref{fam:F.2.1}) & FAM F.1.1 and the explicit matrices $J_0,R$\\
FAM \nolinkurl{F.2.2} (\ref{fam:F.2.2}) & FAM F.1.1, F.1.2 and F.2.1; the homogeneous-space stabilizer calculation\\
FAM \nolinkurl{F.2.3} (\ref{fam:F.2.3}) & The global Siegel chart and the signature condition\\
FAM \nolinkurl{F.3.1} (\ref{fam:F.3.1}) & FAM F.2.2 and F.2.3; the finite-root and cusp period expansions\\
FAM \nolinkurl{F.3.2} (\ref{fam:F.3.2}) & The finite special values and the cusp-normalized period matrix\\
FAM \nolinkurl{F.3.3} (\ref{fam:F.3.3}) & The derivative transformation laws and the finite special values\\
FAM \nolinkurl{D.1.1} (\ref{fam:D.1.1}) & The fixed period matrix and the exact division-polynomial computation\\
FAM \nolinkurl{D.1.2} (\ref{fam:D.1.2}) & FAM D.1.1; E26 Lemma 2.1 and its height calculation\\
FAM \nolinkurl{D.1.3} (\ref{fam:D.1.3}) & FAM D.1.1 and the source bound FAM C.2.1\\
FAM D.1.6 (proof, Hermitian-form convention) (\ref{fam:D.1.Hermitian}) & FAM D.1.6, Hermitian-form convention\\
FAM \nolinkurl{D.1.4} (\ref{fam:D.1.4}) & The period matrix; the Hermitian-form convention in FAM D.1.6; AC26 Lemmas 9.2 and 9.4\\
FAM \nolinkurl{D.1.5} (\ref{fam:D.1.5}) & FAM D.1.4; the five-function independence of AC26 Lemma 9.3\\
FAM \nolinkurl{D.1.6} (\ref{fam:D.1.6}) & FAM D.1.4; the Appell--Humbert and Riemann criteria\\
FAM \nolinkurl{D.1.7} (\ref{fam:D.1.7}) & FAM D.1.6 and the cusp expansions\\
FAM \nolinkurl{D.2.1} (\ref{fam:D.2.1}) & $X_c\simeq S^6$; the K\"ahler geometry of the reduced fibres\\
FAM \nolinkurl{D.2.2} (\ref{fam:D.2.2}) & FAM D.1.1 and D.2.1; the finite normalization and boundary identifications of AC26 Proposition 4.6\\
FAM \nolinkurl{D.2.3} (\ref{fam:D.2.3}) & FAM D.1.6 and D.2.1; the Appell--Humbert theorem and its theta equation\\
PUSH D (opening definitions) (\ref{push:D.1.setup}) & The exact radius 216\\
PUSH \nolinkurl{D.1.1} (\ref{push:D.1.1}) & FAM C.1.1 and A.4.1; ARI C.2.1; modular invariance of j\\
PUSH \nolinkurl{D.1.2} (\ref{push:D.1.2}) & Grauert direct-image and proper-mapping theorems with the 1963 correction; Ehresmann; Buchdahl Theorem 11\\
PUSH \nolinkurl{D.1.3} (\ref{push:D.1.3}) & FAM A.1.1; PUSH D.1.1 and D.1.2\\
PUSH A.1 (opening definitions) (\ref{push:A.1.setup}) & E26 Lemma 2.1 and Corollary 2.2\\
E26 Conjecture 8.2 (\ref{e26:Conjecture8.2}) & E26 §8 specified untwisted family\\
PUSH \nolinkurl{A.3.1} (\ref{push:A.3.1}) & The local-complete-intersection deformation calculation; Levy Lemmas 3.1--3.2; \L{}ojasiewicz \S3 Theorems 1--2\\
PUSH \nolinkurl{A.3.2} (\ref{push:A.3.2}) & PUSH A.3.1; the specified pinching; the Euler number two of the chosen untwisted family\\
PUSH \nolinkurl{A.3.2} (\ref{push:A.3.2.replacement}) & PUSH A.1.2 and A.2.2\\
PUSH A.1.1 (proof, cone definition) (\ref{push:A.1.fan}) & The cusp cocharacter lattice\\
PUSH \nolinkurl{A.1.1} (\ref{push:A.1.1}) & PUSH A.1.fan\\
PUSH \nolinkurl{A.1.2} (\ref{push:A.1.2}) & PUSH A.1.1; E26 Lemma 2.1 and Corollary 2.2\\
PUSH A.2.1 (proof, six rays and bottom diagonal) (\ref{push:A.2.rays}) & PUSH A.1.1\\
PUSH \nolinkurl{A.2.1} (\ref{push:A.2.1}) & PUSH A.1.1 and the retained bottom triangulation\\
PUSH \nolinkurl{A.2.2} (\ref{push:A.2.2}) & PUSH A.1.2 and A.2.1\\
PUSH \nolinkurl{A.2.3} (\ref{push:A.2.3}) & PUSH A.2.4 and A.2.5\\
PUSH A.2.4 (proof, base-changed period shears) (\ref{push:A.2.periods}) & PUSH A.2.4, finite fan coordinates\\
PUSH \nolinkurl{A.2.4} (\ref{push:A.2.4}) & PUSH A.1.1 and the explicit base-changed period shears\\
PUSH \nolinkurl{A.2.5} (\ref{push:A.2.5}) & PUSH A.2.4 and its marked star charts\\
PUSH Conventions; B.4.33 (scope) (\ref{push:B.4.33.scope}) & The first-run retained data\\
PUSH \nolinkurl{B.4.33} (\ref{push:B.4.33}) & PUSH B.4.31 and B.4.32; the actual fixed punctured sphere family\\
PUSH Conventions (mapping-torus orientation) (\ref{push:C.orientation}) & The complex orientation and the positive parameter circle\\
PUSH C.4.1 (scope) (\ref{push:C.4.1.scope}) & FAM B.2.1; PUSH C.4.1\\
PUSH \nolinkurl{C.3.24} (\ref{push:C.3.24}) & PUSH C.3.18, C.3.21--C.3.23 and C.4.1\\
PUSH \nolinkurl{C.3.27} (\ref{push:C.3.27}) & PUSH C.3.25 and C.3.26; the shared characteristic calculations C.3.18--C.3.22\\
MAPS \nolinkurl{B.6.21} (\ref{maps:B.6.21}) & MAPS B.6.15--B.6.20 and the specified mapping-torus comparison\\
PUSH E.1.1 (sphere data); FAM B.3.1 (parameter lattice) (\ref{push:E.1.1.scope}) & PUSH E.1.1--E.1.3; FAM B.3.1\\
MAPS \nolinkurl{B.6.24} (\ref{maps:B.6.24}) & The fixed K1 twist datum; MAPS B.6.19 and B.6.22--B.6.23\\
PUSH E.2.5 (hypotheses); E.2.7 (parameter) (\ref{push:E.2.5.scope}) & PUSH E.2.5 and E.2.7\\
PUSH \nolinkurl{E.2.21} (\ref{push:E.2.21}) & PUSH E.2.5, E.2.7 and E.2.10--E.2.20\\
PUSH \nolinkurl{E.2.22} (\ref{push:E.2.22}) & PUSH E.2.1 and E.2.21\\
PUSH \nolinkurl{E.2.23} (\ref{push:E.2.23}) & PUSH E.2.21 and E.2.22; local miniversality\\
PUSH \nolinkurl{E.2.24} (\ref{push:E.2.24}) & PUSH E.2.23 and local completeness of the deformation charts\\
\bottomrule\end{longtable}\endgroup
\subsection{Notation concordance}\label{app:notation}
The standing conventions are those of FAM. The companion notation sheet gives every retained label and the scoped symbol substitutions; the mathematical signs and normalizations are unchanged unless the source explicitly compares two conventions.
\paragraph{root}
FAM conventions are unchanged:\allowbreak{} rho=exp(pi i/\allowbreak{}3); clockwise root deck generators, positive base meridians, Pi\_\allowbreak{}c and w=exp(-pi i c). P\_\allowbreak{}E is renamed P\_\allowbreak{}mark in statements and proofs. Dual covectors u,w appearing in linear combinations and their coefficient comparisons are u\_\allowbreak{}Lambda,w\_\allowbreak{}Lambda; hatted homology vectors retain their names. The shorthand macro for hatted delta is dhat, avoiding the standard T1 command dh. The modular nome is q\_\allowbreak{}tau and the modular discriminant is Delta\_\allowbreak{}mod, agreeing with the periods section; finite discs Delta\_\allowbreak{}j, vector q\_\allowbreak{}n and the auxiliary division-model discriminant keep their distinct meanings.
Scoped renamings avoid collisions with the global parameters u=w\textasciicircum{}2 and w. In FAM A.1.1 the square root of E\_\allowbreak{}6 is h\_\allowbreak{}6, and the two canonical-ring sections u,v are s\_\allowbreak{}can,t\_\allowbreak{}can (the abstract polynomial generators U,V are unchanged). In FAM A.3.1 the preliminary first-coordinate lift is h\_\allowbreak{}loc. In FAM A.4.2 the moving torus coordinate w is zeta\textasciicircum{}new and its projector P is P\_\allowbreak{}KS. In FAM B.1.1, B.1.2 and B.1.4 the line-bundle fibre coordinate W is v, exactly as in the first-run conventions; it is distinct from the cusp fibre W. The reciprocal cusp coordinate v=1/\allowbreak{}t is r\_\allowbreak{}cusp, including the differential in B.1.3. The B.1.3 blow-up parameter u=t-1 is v\_\allowbreak{}t, agreeing with the duplication section. The holomorphic torus-valued unit u(t\_\allowbreak{}c) of B.2.1 is v\_\allowbreak{}unit(t\_\allowbreak{}c). These edits change symbols only, never hypotheses or algebraic identities.
Workbench mathscr V → mathcal V; its homogeneous line mathcal L → mathcal L\_\allowbreak{}per; its local logarithms w\_\allowbreak{}j → ell\_\allowbreak{}j\textasciicircum{}log; its finite base-change variable u → s\_\allowbreak{}new; its local equivariant first-component function u(z) → h\_\allowbreak{}loc(z). The covector coefficient comparison uses u\_\allowbreak{}Lambda,w\_\allowbreak{}Lambda consistently. Both u replacements are confined to the source paragraphs defining those local variables; the global Engel parameter is unchanged. No sign is changed. Its missing labels ffcert:\allowbreak{}thm:\allowbreak{}normalization, prop:\allowbreak{}period-input-audit and thm:\allowbreak{}canonical-ring-recovery are mapped respectively to FAM A.1.2, A.3.1 and A.1.1, whose displayed calculations have been matched.
Every retained FAM st:\allowbreak{}ID label becomes fam:\allowbreak{}ID. The workbench pshift labels are retained. The companion universal map-factorization statement comes from DGKXv1 thm:\allowbreak{}all-maps-factor; it is not part of the supplied FAM A.1.1 statement. DGKX a → c and cA → mathcal P. PUSH E.1.1 and E.1.2 are printed as announcements as mandated, with their statement bodies retained.
\paragraph{periods}
These are literal notation changes in statements and proofs. No change is made to a parameter range, quantifier, equality, inequality or equality locus. Source equation labels receive the namespace fam:\allowbreak{}, ari:\allowbreak{} or push:\allowbreak{}. The retained source statement numbers remain visible.
\begingroup\small\setlength{\tabcolsep}{3pt}\begin{longtable}{@{}L{0.24\linewidth}L{0.25\linewidth}L{0.45\linewidth}@{}}\toprule
Source notation & Paper notation & Meaning and reason\\\midrule\endhead
FAM a=5+2√6 & a\_\allowbreak{}cusp=5+2√6 & Cusp logarithmic constant, distinguished from affine twist variables.\\
FAM C.1 proof u=e\textasciicircum{}(2πiμ₀) & v\_\allowbreak{}cusp & Multiplicative cusp value; avoids Engel's global u=w².\\
FAM C.1 proof scalar f & f\_\allowbreak{}sc & The Weierstrass scaling factor already named f\_\allowbreak{}sc in FAM B.1.1; avoids the threefold projection f.\\
FAM and ARI modular nome q=e\textasciicircum{}(2πiτ) & q\_\allowbreak{}τ & Distinguishes the elliptic modulus from the Albanese nomes.\\
ARI q₁,q₂ & unchanged & q₁=-w²/\allowbreak{}48³, q₂=iw/\allowbreak{}216; q₁=(27/\allowbreak{}64)q₂².\\
ARI η & η\_\allowbreak{}D & Dedekind function; first-run notation, distinguished from the integral cohomology class.\\
ARI N & 𝒩 & The Gaussian-normalized theta norm; FAM C.2 uses 𝒩.\\
ARI L(t)=log N(P\_\allowbreak{}t) & 𝓛(t)=log 𝒩(P\_\allowbreak{}t) & Distinguished from a line bundle. U\_\allowbreak{}Λ retains the full torus Green function.\\
ARI elliptic lattice Λ,Λ′ & Λ\_\allowbreak{}ell,Λ′\_\allowbreak{}ell & Distinguished from FAM's marked rank-four period lattice.\\
ARI Δ₀ & Δ\_\allowbreak{}mod & The modular discriminant q\_\allowbreak{}τ∏(1−q\_\allowbreak{}τⁿ)²⁴.\\
ARI rational model E\_\allowbreak{}t, coordinates x,y, point P & E\_\allowbreak{}t\textasciicircum{}div, x\_\allowbreak{}div,y\_\allowbreak{}div, P\textasciicircum{}div & Auxiliary division-polynomial coordinates, not the FAM cubic coordinates.\\
ARI Δ(t)=27t⁸(t−1)³ & Δ\_\allowbreak{}div(t) & Discriminant of the auxiliary rational model. Generic Δ\_\allowbreak{}E and Δ\_\allowbreak{}Λell mean the discriminant of their named Weierstrass model/\allowbreak{}lattice.\\
ARI C.4.3 radicals W,V & W\_\allowbreak{}rad,V\_\allowbreak{}rad & Products of discriminants and division values, distinguished from the cusp fibre and the integral lattice.\\
ARI C.3.1 G=g₂(i) & G\_\allowbreak{}square & Square-lattice invariant, distinguished from the local function G(h,τ) in the maximum proof.\\
ARI C.3.1 Taylor coefficients A₂,A₄ & α₂,α₄ & Avoids the monodromy matrices A₁,A₂.\\
ARI C.3.1/\allowbreak{}C.8.2 scalar E=π²E₂/\allowbreak{}3 & 𝓔₂ & Scaled quasimodular scalar, distinguished from the rank-one lattice endomorphism E.\\
ARI C.8.2 Fourier sums A\_\allowbreak{}k & 𝒮\_\allowbreak{}k & Scalar theta-series bounds, distinguished from monodromy matrices and duplication polynomials.\\
ARI C.8.4 coefficient bound M & M\_\allowbreak{}coeff & Bound for five polynomial coefficients, distinguished from the global admissibility maximum M.\\
ARI C.8.4 local u=t−1 & v\_\allowbreak{}t & Avoids Engel's multiplicative family parameter.\\
ARI C.8.4 duplication polynomials A\_\allowbreak{}k & 𝒜\_\allowbreak{}k & Avoids the first-run monodromy matrices.\\
ARI C.8.6 integral coordinate u=1/\allowbreak{}℘(z) & v\_\allowbreak{}℘ & Avoids the Engel parameter and the ladder coordinate.\\
ARI C.8.9 local base coordinate w & z\_\allowbreak{}B & Avoids the global moduli parameter w=e\textasciicircum{}(−πic).\\
ARI C.8.11 W & 𝒲 & Logarithmic theta derivative; avoids the cusp fibre W.\\
ARI C.8.11 R=μ′−a & ℛ\_\allowbreak{}μ & Derivative defect; distinguished from radii and the theta quotient.\\
\bottomrule\end{longtable}\endgroup
\paragraph{Cubic normalization and section sign}
FAM B.1.1 uses y\_\allowbreak{}F²=x\_\allowbreak{}F³−3t³(t−1)x\_\allowbreak{}F+2t⁴(t−1)². FAM B.1.4 fixes the cusp branch with distinguished section (0,−√2 t²(t−1)). ARI C.4 uses y\_\allowbreak{}div²=x\_\allowbreak{}div³−(3/\allowbreak{}4)t³(t−1)x\_\allowbreak{}div+(1/\allowbreak{}4)t⁴(t−1)² and section (0,t²(t−1)/\allowbreak{}2). The coordinate substitution x\_\allowbreak{}F=2x\_\allowbreak{}div, y\_\allowbreak{}F=−2√2 y\_\allowbreak{}div identifies both equations and these distinguished sections. This is a coordinate dictionary, not a new geometric theorem. The norm is independent of the Weierstrass homothety and is even in the point, as in the ARI C.4 setup and C.1.1 proof. Generic sigma/\allowbreak{}division-polynomial arguments retain their explicitly stated normalization x=℘, y=℘′/\allowbreak{}2, ψ₁=1, ψ₂=2y.
\paragraph{Special source conventions}
ρ=(1+i√3)/\allowbreak{}2=e\textasciicircum{}(πi/\allowbreak{}3), not the other primitive cube root. The finite values are μ(z₁)=(2−ρ)/\allowbreak{}3, μ(z₂)=(1−i)/\allowbreak{}2; the cusp value is μ₀=1/\allowbreak{}2−i log(5+2√6)/\allowbreak{}(2π). All phases are continued from that branch. The period half-plane is Im c<−M with M=−(3/\allowbreak{}π)log6; the positive-imaginary Albanese representatives are −κ₁ and −κ₂/\allowbreak{}2. The interval cover uses 0≤Re τ≤1/\allowbreak{}2, |τ|≥1 and not a change of root-deck orientation.
The current convention is exactly the corrected ARI C.8.9 convention:\allowbreak{} ddᶜ=(i/\allowbreak{}π)∂\ensuremath{\bar\partial}, positive Betti measure db∧da, total mass 1/\allowbreak{}12, probability measure 12m\_\allowbreak{}P. No sign is changed in the mathematical statements.
\paragraph{cohomology}
The source statements retain the first-run convention:\allowbreak{} the coarse base is B=\textbackslash{}mathbb P\textasciicircum{}1\_\allowbreak{}t, its points are p\_\allowbreak{}1=0, p\_\allowbreak{}2=1, p\_\allowbreak{}0=\textbackslash{}infty, the standard finite multiplicities are 3 and 4, the root generators are clockwise, and \textbackslash{}beta=\textbackslash{}beta\_\allowbreak{}\{\textbackslash{}mathrm\{part\}\}+c. No additional special-point swap is required in these fragments. The alternate Engel marking Q in INV is not used; the paper retains the FAM marking P\_\allowbreak{}\{\textbackslash{}mathrm\{mark\}\}.
\paragraph{Symbol concordance}
\begingroup\small\setlength{\tabcolsep}{3pt}\begin{longtable}{@{}L{0.24\linewidth}L{0.25\linewidth}L{0.45\linewidth}@{}}\toprule
Source symbol & Printed symbol & Reason and scope\\\midrule\endhead
Dual lattice covectors u,w & u\_\allowbreak{}\{\textbackslash{}Lambda\},w\_\allowbreak{}\{\textbackslash{}Lambda\} & Distinguishes the period covectors from family parameters; the hatted lattice vectors stay unchanged. Local module generators named u are untouched.\\
INV logarithmic relative cotangent E or \textbackslash{}mathcal E & \textbackslash{}mathcal E\_\allowbreak{}\{\textbackslash{}rm cot\} & Keeps the relative cotangent bundle separate from the tangent bundle and elliptic curves.\\
FAM/\allowbreak{}PUSH/\allowbreak{}CDP root vertical tangent \textbackslash{}mathcal E & \textbackslash{}mathcal E\_\allowbreak{}\{\textbackslash{}rm tan\} & Distinguishes tangent and cotangent source conventions.\\
INV coarse direct-image bundle G & G\_\allowbreak{}B & Separates the coarse and root-curve extensions.\\
FAM/\allowbreak{}PUSH/\allowbreak{}CDP root direct-image bundle G & G\_\allowbreak{}\{\textbackslash{}mathfrak B\} & Same distinction.\\
INV root curve \textbackslash{}mathcal B & \textbackslash{}mathfrak B & Matches the first-run convention.\\
Root-to-coarse map \textbackslash{}rho & \textbackslash{}varpi & Reserves \textbackslash{}rho for the sixth root e\textasciicircum{}\{\textbackslash{}pi i/\allowbreak{}3\}.\\
INV cotangent determinant line N & N\_\allowbreak{}\{\textbackslash{}rm det\} & Separates this line from the Bott--Chern ideal and from a cusp tube.\\
INV square-zero ideal N & N\_\allowbreak{}\{BC\} & Only the ideal in inv:\allowbreak{}C3.3.\\
FAM/\allowbreak{}PUSH cusp normalization surface E & P\_\allowbreak{}6 & The degree-six del Pezzo surface; not an elliptic curve or a vector bundle.\\
Source B2.2 weighted variables x of weight 4, y of weight 3 & x of weight 3, y of weight 4 & Swapped consistently only in the valuation argument imported into inv:\allowbreak{}B3.3.\\
Source B2.2 auxiliary line \textbackslash{}mathcal L & \textbackslash{}mathcal L\_\allowbreak{}\{\textbackslash{}rm root\} & Separates the valuation argument from the integral Picard generator.\\
CDP residue indices r\_\allowbreak{}1(n),r\_\allowbreak{}2(n) & r\_\allowbreak{}3(n),r\_\allowbreak{}4(n) & The subscripts now record the actual root orders; the residue formulas are unchanged.\\
\bottomrule\end{longtable}\endgroup
Original statement IDs are retained with inv:\allowbreak{}, fam:\allowbreak{}, push:\allowbreak{}, or cdp:\allowbreak{} prefixes. Internal equation labels receive the same run prefix. Label and reference keys are protected before symbol substitutions; no mathematical replacement is allowed inside an identifier. The source's references to Part A now point to sec:\allowbreak{}cotangent-images, except the reference to its Hodge diamond, which points to the more precise sec:\allowbreak{}hodge.
\paragraph{Environments, packages, and macros}
The fragments have no preamble or bibliography. They use theorem, lemma, proposition, corollary, definition, remark, announcement, cited, and proof, with the supplied \textbackslash{}uses\{...\} macro after each formal statement. The central preamble provides these environments and \textbackslash{}uses. Standard requirements are amsmath, amssymb, amsthm, mathtools, mathrsfs, and the central table packages. Existing math operators such as \textbackslash{}Hom, \textbackslash{}Ext, \textbackslash{}rank, \textbackslash{}coker, \textbackslash{}im, \textbackslash{}Spec, \textbackslash{}Proj, \textbackslash{}Pic, \textbackslash{}End, \textbackslash{}Sym, \textbackslash{}Res, \textbackslash{}ord, \textbackslash{}tr, and \textbackslash{}lcm come from that preamble. The proof shorthands use the central \textbackslash{}dbar and \textbackslash{}dhat; the master normalizes any legacy \textbackslash{}dh spelling to \textbackslash{}dhat to avoid redefining the standard T1 text command. No standard LaTeX command is redefined here.
\paragraph{Content-preserving layout edits}
Long proofs are divided into numbered steps by placing step labels at existing source paragraph boundaries. Source formulas are not altered by that operation. References and notation are normalized before recording statement\_\allowbreak{}as\_\allowbreak{}printed; every stored printed statement is matched against the generated fragment. Exact original source statements are retained separately in source\_\allowbreak{}statement.
\paragraph{Final Appendix B cusp-coordinate correction}
The final visual check identified a remaining notation collision in the FAM A.5.2 re-derivation. In that appendix proof, the two occurrences of the weight-twelve modular discriminant \textbackslash{}Delta(\textbackslash{}tau) become \textbackslash{}Delta\_\allowbreak{}\{\textbackslash{}rm mod\}(\textbackslash{}tau). The finite root-disc notation \textbackslash{}Delta\_\allowbreak{}\{s\_\allowbreak{}j\} is unchanged.
The source expressly calls the local equation q=z\_\allowbreak{}0z\_\allowbreak{}1z\_\allowbreak{}2 the original cusp family. Only that coordinate and its subsequent occurrences in the same cusp calculations are renamed to t\_\allowbreak{}c:\allowbreak{} the convergent-series ring, the toric gluing u(q)q\textasciicircum{}\{B\_\allowbreak{}0\textbackslash{}lambda\}, the relation \textbackslash{}Delta(\textbackslash{}tau)=q\textbackslash{},u(q), the field q\textbackslash{}partial\_\allowbreak{}q, and the differentials dq. The last three now read \textbackslash{}Delta\_\allowbreak{}\{\textbackslash{}rm mod\}(\textbackslash{}tau)=t\_\allowbreak{}c\textbackslash{},u(t\_\allowbreak{}c), t\_\allowbreak{}c\textbackslash{}partial\_\allowbreak{}\{t\_\allowbreak{}c\}, and dt\_\allowbreak{}c. Cohomological indices q in higher direct images, degree conditions and cases are unchanged. No equality between t\_\allowbreak{}c and the modular nome q\_\allowbreak{}\textbackslash{}tau is introduced.
The exact scope and replacement counts are recorded in \path{appendix_notation_changes.json}; appendix\_\allowbreak{}notation\_\allowbreak{}fix.diff records the before/\allowbreak{}after appendix fragment. The generating script applies the same finite list of replacements with count assertions, preserving all statements, hypotheses and proof steps.
The same final scan found exactly one additional modular-discriminant token in the main INV A1.0 proof. With the assembling editor's explicit authorization, \textbackslash{}Delta(\textbackslash{}tau)\textasciicircum{}\{-1\} there is likewise renamed to \textbackslash{}Delta\_\allowbreak{}\{\textbackslash{}rm mod\}(\textbackslash{}tau)\textasciicircum{}\{-1\}. Its surrounding cusp coordinate was already t\_\allowbreak{}c; no other main-fragment notation is changed. This additional one-token correction is included in \path{appendix_notation_changes.json}.
\paragraph{geometry ends}
All mathematical conventions are those of FAM, \path{returns/S6_02_FAMILY/sections/conventions.tex}. The changes below are notation, reference or layout changes; the source statements, hypotheses and scopes remain separately recorded in \path{records.json}. Exact statement and proof diffs are in \path{statement_comparisons.json} and \path{proof_comparisons.json}.
\paragraph{Common conventions retained}
The base is B=P\textasciicircum{}1\_\allowbreak{}t, with p\_\allowbreak{}1=0, p\_\allowbreak{}2=1, p\_\allowbreak{}0=infinity; the retained multiplicities are (3,4). Homology columns have order (hat-gamma,hat-u,hat-w,hat-delta). The deck generators are clockwise root transformations, g\_\allowbreak{}1g\_\allowbreak{}2g\_\allowbreak{}0=1; positive geometric parameter loops are counterclockwise. The period matrix is Pi\_\allowbreak{}c=((6mu,tau,1,0),(beta\_\allowbreak{}part+c,mu,0,1)), with (beta\_\allowbreak{}part+tau)(p\_\allowbreak{}0)=0, and w=exp(-pi*i*c). The Albanese nome is q\_\allowbreak{}2=iw/\allowbreak{}216. The fibre vector is e=(0,1)\textasciicircum{}t=Pi\_\allowbreak{}c hat-delta; the lattice vector is d=hat-delta.
\paragraph{Symbol mappings}
\begingroup\small\setlength{\tabcolsep}{3pt}\begin{longtable}{@{}L{0.17\linewidth}L{0.21\linewidth}L{0.22\linewidth}L{0.34\linewidth}@{}}\toprule
Source and scope & Source notation & Printed notation & Reason /\allowbreak{} sign convention\\\midrule\endhead
FAM F, lattice & e=hat-delta & d=hat-delta & Distinguishes the primitive lattice vector from the fibre vector e already used in FAM AB. Every concatenated occurrence Je, C e, R e, Z e, stabilizer subscript and q\_\allowbreak{}e is changed consistently.\\
FAM F & E=e tensor gamma, q\_\allowbreak{}e, P\_\allowbreak{}\{e,Z/\allowbreak{}R/\allowbreak{}Q\} & E=d tensor gamma, q\_\allowbreak{}d, P\_\allowbreak{}\{d,Z/\allowbreak{}R/\allowbreak{}Q\} & Same operator, invariant function and stabilizer.\\
FAM F, auxiliary scalars & matrix entry/\allowbreak{}denominator d & d\_\allowbreak{}aux & Prevents the local scalar from colliding with the renamed lattice vector. Matrix formulas and denominator-clearing integers are otherwise unchanged.\\
FAM F.3.3 & cusp q=exp(2pi*i*s) & t\_\allowbreak{}c=exp(2pi*i*s) & Keeps the cusp coordinate distinct from the Albanese nome. Exponentials e\textasciicircum{}(pi*i/\allowbreak{}3) and e\textasciicircum{}(2pi*i*s) remain exponentials.\\
FAM F and D, dual coframe & u,w in wedge forms & u\_\allowbreak{}Lambda,w\_\allowbreak{}Lambda & Matches the root paper's dual-coframe notation; hatted homology vectors and parameter w are unchanged.\\
FAM D & beta & beta=beta\_\allowbreak{}part+c & Retains the actual period, not the particular solution alone.\\
FAM D.1.4 & T,U,D in G & T=Im tau, U=Im mu, D=Im beta-6U\textasciicircum{}2/\allowbreak{}T stated explicitly & T is defined in FAM C; U is not locally defined in \path{CD.tex}. This notational ambiguity is recorded in the discrepancy report. No entry of G, determinant or intersection sign is changed.\\
FAM D.1.4/\allowbreak{}D.1.6 & Hermitian form & H\_\allowbreak{}E(v,w)=Etilde(v,Jw)-i Etilde(v,w) & The source convention, linear in its first entry, is retained. Its Chern form is -Im H\_\allowbreak{}E; no sign is silently imported from AC26.\\
MAPS A.3 & homology g & hat-gamma in the translation formula & The short alias is removed.\\
MAPS A.3 representative formula & S(t)d in a fibre-coordinate translation, although d=hat-delta & S(t)e, with e=Pi\_\allowbreak{}c d=(0,1)\textasciicircum{}t explicitly stated & Exposes the source's lattice/\allowbreak{}fibre-vector shorthand. All coefficients, branches and real representatives are unchanged.\\
MAPS A.3 & circle action rho & rho\_\allowbreak{}circ(b) is translation by b e, b in R/\allowbreak{}Z & Explicitly identifies the source circle action with FAM B.4.1 and reserves the unadorned rho=exp(pi*i/\allowbreak{}3) for the period constant. Every rho(1/\allowbreak{}2) is renamed consistently.\\
MAPS A.3 & D\_\allowbreak{}1,D\_\allowbreak{}2,D\_\allowbreak{}12 & unchanged & Directions e\_\allowbreak{}1,e\_\allowbreak{}2,e\_\allowbreak{}1-e\_\allowbreak{}2; D\_\allowbreak{}2 is the circle-fixed curve. These are not interchanged.\\
PUSH A & Engel line bundle L & M\_\allowbreak{}S=O\_\allowbreak{}S(P-O) & Matches FAM conventions. The cotangent complex mathbb L\_\allowbreak{}D in A.2.5 remains mathbb L\_\allowbreak{}D; it is not the line bundle.\\
PUSH A & linearization line J & J\_\allowbreak{}lin & Avoids confusing the line bundle with a real complex structure or the monodromy integer. Its Chern class is still the lower-case j.\\
PUSH A fan & cusp q, height h\_\allowbreak{}q & t\_\allowbreak{}c, h\_\allowbreak{}\{t\_\allowbreak{}c\} & Same cusp base coordinate and same primitive ray heights. The degeneration parameter u is unchanged.\\
E26 Conjecture 8.2 quotation & M & M inside the verbatim quotation, identified with M\_\allowbreak{}S immediately afterward & Preserves the exact conjecture wording; outside that quotation M remains the period maximum.\\
PUSH E.2.7/\allowbreak{}E.2.21 & fixed-pair parameter z=exp(-pi*i*c/\allowbreak{}(kl)) & w\_\allowbreak{}\{k,l\} & Separates the parameter from the first-run uniformizing coordinate z.\\
MAPS B.6.24 & positive generator 2kl of the parameter-period lattice & unchanged & The source positive closing-translation convention is retained. The positive geometric w loop in the retained family lifts to c -> c-2; no nonzero sign-sensitive mapping class is inferred.\\
FAM/\allowbreak{}PUSH bibliography & DGKX26v1, DuGuoKusak\allowbreak{}abeXie2026\allowbreak{}v1 & DGKXv1 & Common bibliography key. The full-automorphism locator is the theorem labelled p2:\allowbreak{}thm:\allowbreak{}full-automorphisms, not an unqualified period-classification citation.\\
FAM period cross-reference & eq:\allowbreak{}cd-beta & fam:\allowbreak{}eq:\allowbreak{}cd-beta & Matches the period agent's sole equation label; no second label is inserted into its equation.\\
\bottomrule\end{longtable}\endgroup
\paragraph{Statement and equation labels}
Every retained numbered source statement has the label run:\allowbreak{}original-ID:\allowbreak{} fam:\allowbreak{}B.4.1, fam:\allowbreak{}B.4.2, ari:\allowbreak{}C.2.2, maps:\allowbreak{}A.3.1--maps:\allowbreak{}A.3.3, fam:\allowbreak{}F.1.1--fam:\allowbreak{}F.3.3, fam:\allowbreak{}D.1.1--fam:\allowbreak{}D.2.3, the retained push:\allowbreak{}A.*, push:\allowbreak{}D.1.*, and all specified push:\allowbreak{}B.4.33, push:\allowbreak{}C.3.24, push:\allowbreak{}C.3.27, maps:\allowbreak{}B.6.21, maps:\allowbreak{}B.6.24, push:\allowbreak{}E.2.21--push:\allowbreak{}E.2.24.
The exact E26 target is labelled e26:\allowbreak{}Conjecture8.2. The final replacement sentence of PUSH A.3.2 is labelled push:\allowbreak{}A.3.2.replacement and separately typed A; its source ID remains A.3.2. The rest of that source statement is the T-item push:\allowbreak{}A.3.2.
FAM displayed labels acquire the fam:\allowbreak{} prefix, including fam:\allowbreak{}eq:\allowbreak{}F-g, fam:\allowbreak{}eq:\allowbreak{}F-Gamma, fam:\allowbreak{}eq:\allowbreak{}F-law, fam:\allowbreak{}eq:\allowbreak{}F-SLpresentation, fam:\allowbreak{}eq:\allowbreak{}F-H, fam:\allowbreak{}eq:\allowbreak{}F-qD, fam:\allowbreak{}eq:\allowbreak{}F-tauaction, fam:\allowbreak{}eq:\allowbreak{}F-j, fam:\allowbreak{}eq:\allowbreak{}cd-NS and fam:\allowbreak{}eq:\allowbreak{}cd-theta-sections. Contextual definitions use labels ending .setup, .scope, .representatives, .fan, .rays, .periods, .Hermitian or .orientation; their metadata identifies the original unnumbered context or proof extract explicitly, rather than inventing an original theorem number.
\paragraph{Scoped reuse retained from the sources}
The symbols B (base /\allowbreak{} local 2-by-2 block), T (modular generator /\allowbreak{} imaginary period /\allowbreak{} translation t\_\allowbreak{}\{6P\}), and E (rank-one marking operator /\allowbreak{} alternating class) have explicitly separated local scopes, as in the first run. The local S(t) in the real representative table is a scalar translation function, whereas S in the filling section is Engel's rational elliptic surface. The record does not identify these objects.
\subsection{Verification record, one row per check}\label{app:verification}
R: statement read without following its proof. F: proof or explicitly delimited proof portion followed. D: steps re-derived. C: numerical or exact recomputation. S: source opened and matched. No L check is recorded. A producing-run computation is identified as such. The historical records do not constitute a fresh check of this assembly.
\begingroup\footnotesize\setlength{\tabcolsep}{3pt}\begin{longtable}{@{}L{.08\linewidth}L{.31\linewidth}L{.56\linewidth}@{}}
\toprule Check & By / date / type / scope & Verdict, findings, independence and report\\\midrule\endhead
V001 & Claude Fable 5.2, effort high\newline 2026-11-12; F\newline FAM A.1.1 & correct. Characters chi\_\allowbreak{}j(0) = zeta\_\allowbreak{}j\textasciicircum{}(1+a\_\allowbreak{}j), a = (2,1), from det R\_\allowbreak{}j at tau = rho, i; orders (2,1,-1) of Theta = E4\textasciicircum{}2 sqrt(E6)/\allowbreak{}Delta; divisor of Theta\textasciicircum{}-1 e dt = 2S2 - W, hence K = O(-2S2), equal to [AC26] Theorem (4) (T l. 156); dimension count floor(m/\allowbreak{}2)+1 for C[U,V]. K\_\allowbreak{}X formula is in [AC26]; ring and direct image are the workbench's (F17).\newline Given the return, bibliography and programs; one AI reader.\newline \path{status/REVIEW_S6_02_FAMILY.md}, Section 3, A.1.1\\
V002 & Claude Fable 5.2, effort high\newline 2026-11-12; F,C\newline FAM A.1.2 & correct. Fixed and moving lattices and their saturation by program; degrees 9, 8 from the isogeny kernels of orders 3, 2. Standard (workbench FF1-FF46).\newline Given the return, bibliography and programs; one AI reader.\newline \path{status/REVIEW_S6_02_FAMILY.md}, Section 3, A.1.2\\
V003 & Claude Fable 5.2, effort high\newline 2026-11-12; F,D,C\newline FAM A.2.1 & correct. Whole argument by hand; centralizer QI + QE, N1, N2, quotients Z/\allowbreak{}3 and Z/\allowbreak{}4 + Z/\allowbreak{}2 by my program from T1, T2 of [AC26]; the deck equation and its m\_\allowbreak{}j-fold iterate.\newline Given the return, bibliography and programs; one AI reader.\newline \path{status/REVIEW_S6_02_FAMILY.md}, Section 3, A.2.1\\
V004 & Claude Fable 5.2, effort high\newline 2026-11-12; F,C\newline FAM A.3.1 & correct. (I-A\_\allowbreak{}j)r\_\allowbreak{}j identities, l\_\allowbreak{}0 = (3n/\allowbreak{}2)w - (n/\allowbreak{}2)d, N\_\allowbreak{}j l\_\allowbreak{}j by program; first-coordinate sheaf = O(-p0) from the orders of Theta; H\textasciicircum{}1(V) = 0, H\textasciicircum{}0(V) = C e; invariance of the A2 fan under integer shears; equality of cusp multipliers at c and c+n. Properness at the cusp:\allowbreak{} read.\newline Given the return, bibliography and programs; one AI reader.\newline \path{status/REVIEW_S6_02_FAMILY.md}, Section 3, A.3.1\\
V005 & Claude Fable 5.2, effort high\newline 2026-11-12; F\newline FAM A.4.1 & correct. Albanese lattices; Im(kappa1\textasciicircum{}part - kappa2\textasciicircum{}part) = -0.13735807... by my own evaluation (inside the certified (-0.2,-0.1)).\newline Given the return, bibliography and programs; one AI reader.\newline \path{status/REVIEW_S6_02_FAMILY.md}, Section 3, A.4.1\\
V006 & Claude Fable 5.2, effort high\newline 2026-11-12; F\newline FAM A.4.2 & correct. P\textasciicircum{}2 = P, derivative P/\allowbreak{}(2iD); transfer to the threefold through the bi-anticanonical pencil:\allowbreak{} read. Consistent with A.5.2, where the generator of H\textasciicircum{}1(T) is e tensor [gamma].\newline Given the return, bibliography and programs; one AI reader.\newline \path{status/REVIEW_S6_02_FAMILY.md}, Section 3, A.4.2\\
V007 & Claude Fable 5.2, effort high\newline 2026-11-12; F\newline FAM A.5.1 & known. Assembly checked; the review identifies Theorem 1.4 of DGKXv1 and notes its absence in v2. No date comparison is carried into the paper.\newline Given the return, bibliography and programs; one AI reader.\newline \path{status/REVIEW_S6_02_FAMILY.md}, Section 3, A.5.1\\
V008 & Claude Fable 5.2, effort high\newline 2026-11-12; F,D\newline FAM A.5.2 & known. Every term recomputed:\allowbreak{} Pic of the root curve P(3,4), K = O(-7), L = O(-1), E = O + O(-1), R\textasciicircum{}q f\_\allowbreak{}* O = O, O + O(-1), O(-1); acyclicity of O(k), k = 1,-1,-2,-5,-6; h(T(-log W)) = (1,1,0,0); quotient nu\_\allowbreak{}* K\_\allowbreak{}dP6; result (1,1,1,0). Independent check (mine):\allowbreak{} Hirzebruch-Riemann-Roch gives chi(T\_\allowbreak{}X) = c3/\allowbreak{}2 = 1. Stated in arXiv:\allowbreak{}2610.00250v1 (Theorem 1.3 and introduction); absent from v2.\newline Given the return, bibliography and programs; one AI reader.\newline \path{status/REVIEW_S6_02_FAMILY.md}, Section 3, A.5.2\\
V009 & Claude Fable 5.2, effort high\newline 2026-11-12; F,C\newline FAM B.1.1 & correct. Line bundle O(P-O), factors (1, e\textasciicircum{}(2 pi i mu)); discriminant -108 t\textasciicircum{}8 (t-1)\textasciicircum{}3 and j = 1728t by program. MW = Z P with height 1/\allowbreak{}6 is the Oguiso-Shioda entry for IV*, III, I1 (used by Engel, Lemma 2.1).\newline Given the return, bibliography and programs; one AI reader.\newline \path{status/REVIEW_S6_02_FAMILY.md}, Section 3, B.1.1\\
V010 & Claude Fable 5.2, effort high\newline 2026-11-12; F,C\newline FAM B.1.2 & correct. det of the real period matrix = T D (numerical check).\newline Given the return, bibliography and programs; one AI reader.\newline \path{status/REVIEW_S6_02_FAMILY.md}, Section 3, B.1.2\\
V011 & Claude Fable 5.2, effort high\newline 2026-11-12; F,C\newline FAM B.1.3 & correct. Height 2 - 4/\allowbreak{}3 - 1/\allowbreak{}2 = 1/\allowbreak{}6; wp((1-i)/\allowbreak{}2, i) = wp((2-rho)/\allowbreak{}3, rho) = 0 numerically. The return's correction of the course label (P\_\allowbreak{}* meets E\_\allowbreak{}(2,-), not E\_\allowbreak{}(2,+)) is right:\allowbreak{} Y2 = y/\allowbreak{}t\textasciicircum{}2 = sqrt2 (t-1) -> -sqrt2; course l. 680 says E\_\allowbreak{}(2,+).\newline Given the return, bibliography and programs; one AI reader.\newline \path{status/REVIEW_S6_02_FAMILY.md}, Section 3, B.1.3\\
V012 & Claude Fable 5.2, effort high\newline 2026-11-12; F\newline FAM B.1.4 & correct. div of the transported section = F\_\allowbreak{}infinity:\allowbreak{} read. lambda0 = 1 holds because the reference linearization is defined to be the transported one; see 5.4.\newline Given the return, bibliography and programs; one AI reader.\newline \path{status/REVIEW_S6_02_FAMILY.md}, Section 3, B.1.4\\
V013 & Claude Fable 5.2, effort high\newline 2026-11-12; F,C\newline FAM B.2.1 & correct. Norm criterion against (B2.2) by brute force, 604,800 data, no mismatch; p = -r1 mod 3, p = r2 mod 4; cusp necessity (B2.3) and sufficiency:\allowbreak{} read.\newline Given the return, bibliography and programs; one AI reader.\newline \path{status/REVIEW_S6_02_FAMILY.md}, Section 3, B.2.1\\
V014 & Claude Fable 5.2, effort high\newline 2026-11-12; F\newline FAM B.2.2 & correct. Two square roots of u; signs from a\_\allowbreak{}(j-1) = -P\_\allowbreak{}E v\_\allowbreak{}j/\allowbreak{}m\_\allowbreak{}j and gamma(v1) = 1, gamma(v2) = -1. P\_\allowbreak{}E conjugates A\_\allowbreak{}j to Engel's T\_\allowbreak{}0 (Thm 5.1) and T\_\allowbreak{}1 ((6.1), display garbled in my text copy).\newline Given the return, bibliography and programs; one AI reader.\newline \path{status/REVIEW_S6_02_FAMILY.md}, Section 3, B.2.2\\
V015 & Claude Fable 5.2, effort high\newline 2026-11-12; F\newline FAM B.2.3 & correct. Engel's Remark 8.1 read:\allowbreak{} "a proper, holomorphic submersion Y -> Delta\textasciicircum{}* whose fiber over u is X(u, a0, a1)". By Construction 3.5 and Cor. 3.9, psi(a0) is in (1/\allowbreak{}3)Z minus Z, so the datum is not invariant under the u-monodromy. Second witness:\allowbreak{} j(Alb S2) has a w\textasciicircum{}-1 term.\newline Given the return, bibliography and programs; one AI reader.\newline \path{status/REVIEW_S6_02_FAMILY.md}, Section 3, B.2.3\\
V016 & Claude Fable 5.2, effort high\newline 2026-11-12; F,D\newline FAM B.3.1 & correct with a gap. Arithmetic (N\_\allowbreak{}3k = k N1, N\_\allowbreak{}4l = l N2, lcm(k,2l) = 2kl, CRT) checked by hand. Gap:\allowbreak{} the map from Viaclovsky's data (m\_\allowbreak{}i, a\_\allowbreak{}i), v\_\allowbreak{}i = r\_\allowbreak{}i delta + n\_\allowbreak{}i u\_\allowbreak{}i, to (r\_\allowbreak{}j, s\_\allowbreak{}j, l\_\allowbreak{}0) is not written; sufficiency is "as in B.2.1".\newline Given the return, bibliography and programs; one AI reader.\newline \path{status/REVIEW_S6_02_FAMILY.md}, Section 3, B.3.1\\
V017 & Claude Fable 5.2, effort high\newline 2026-11-12; F,C\newline FAM B.4.1 & known. Coboundary lattice has determinant 1 in the allowed pairs (program); H\textasciicircum{}0(V) = C e. "Aut(X\_\allowbreak{}a) = C\textasciicircum{}*" is in Theorem 1.4 of arXiv:\allowbreak{}2610.00250v1.\newline Given the return, bibliography and programs; one AI reader.\newline \path{status/REVIEW_S6_02_FAMILY.md}, Section 3, B.4.1\\
V018 & Claude Fable 5.2, effort high\newline 2026-11-12; F,C\newline FAM B.4.2 & correct. H\textasciicircum{}2 = I, H A2 H = A2\textasciicircum{}-1, H A1 H = A2\textasciicircum{}-1 A1\textasciicircum{}-1 A2, Pi\_\allowbreak{}(-cbar) H = -conj(Pi), twisted targets (-eps + d, eps' - 3d) by program. Independent check (mine):\allowbreak{} with Re kappa2\textasciicircum{}part = 3/\allowbreak{}2 and Re kappa1\textasciicircum{}part = 1/\allowbreak{}2, the Albanese moduli of the conjugate manifold match those of X\_\allowbreak{}(c') exactly for c' = 1 - cbar mod 2 and not for c' = -cbar.\newline Given the return, bibliography and programs; one AI reader.\newline \path{status/REVIEW_S6_02_FAMILY.md}, Section 3, B.4.2\\
V019 & Claude Fable 5.2, effort high\newline 2026-11-12; F,C\newline FAM C.1.1 & correct. The two Jacobi laws give (beta1) of [AC26] exactly (hand, and numerically with the explicit branch); cusp limit from sin\textasciicircum{}2(pi mu0) = 3; R(z2) = e\textasciicircum{}(pi/\allowbreak{}4) sqrt2 to 40 digits.\newline Given the return, bibliography and programs; one AI reader.\newline \path{status/REVIEW_S6_02_FAMILY.md}, Section 3, C.1.1\\
V020 & Claude Fable 5.2, effort high\newline 2026-11-12; F,C\newline FAM C.1.2 & correct. All five values reproduced to 40 digits by different code.\newline Given the return, bibliography and programs; one AI reader.\newline \path{status/REVIEW_S6_02_FAMILY.md}, Section 3, C.1.2\\
V021 & Claude Fable 5.2, effort high\newline 2026-11-12; F,D,C\newline FAM C.2.3 & correct. Section 2:\allowbreak{} full re-run of the interval proof. Formulas and tails of `wp\_\allowbreak{}iv`, `N2\_\allowbreak{}iv` checked against the text; half-period identity log N = Re H - pi (Im h)\textasciicircum{}2/\allowbreak{}T checked; cusp bound 1.0814 < sqrt2 by hand. The rigorous part rests on mpmath's interval arithmetic; no second interval library was used.\newline Given the return, bibliography and programs; one AI reader.\newline \path{status/REVIEW_S6_02_FAMILY.md}, Section 3, C.2.3\\
V022 & Claude Fable 5.2, effort high\newline 2026-11-12; F\newline FAM D.1.1 & correct. x(6P) has denominator 648 (11t\textasciicircum{}2 - 6016t + 8192)\textasciicircum{}2 (my own chord-tangent computation); roots 64(47 +- 27 sqrt3)/\allowbreak{}11.\newline Given the return, bibliography and programs; one AI reader.\newline \path{status/REVIEW_S6_02_FAMILY.md}, Section 3, D.1.1\\
V023 & Claude Fable 5.2, effort high\newline 2026-11-12; F\newline FAM D.1.2 & known. Abelian at torsion points:\allowbreak{} index-N\textasciicircum{}3 product sublattice checked; 3k\textasciicircum{}2 - 1 from Engel's Lemma 2.1 (read in the text).\newline Given the return, bibliography and programs; one AI reader.\newline \path{status/REVIEW_S6_02_FAMILY.md}, Section 3, D.1.2\\
V024 & Claude Fable 5.2, effort high\newline 2026-11-12; F,D,C\newline FAM D.1.3 & correct. At t = 32/\allowbreak{}5, y(2P) = 0 (exact order 4); L-images (-101 tau, 2 tau + 2, 3, 4); c = -2.25 - 34.0210972421\allowbreak{}8209472308 i reproduced to 20 digits; Im c < -22 by hand. Scope in 5.3.\newline Given the return, bibliography and programs; one AI reader.\newline \path{status/REVIEW_S6_02_FAMILY.md}, Section 3, D.1.3\\
V025 & Claude Fable 5.2, effort high\newline 2026-11-12; F\newline FAM D.1.4 & correct. NS equation from the wedge of the rows of Pi; det = T D < 0; eta\textasciicircum{}2 = -12 in the complex orientation.\newline Given the return, bibliography and programs; one AI reader.\newline \path{status/REVIEW_S6_02_FAMILY.md}, Section 3, D.1.4\\
V026 & Claude Fable 5.2, effort high\newline 2026-11-12; F\newline FAM D.1.5 & correct. Pf = -3n\textasciicircum{}2 for the displayed class; independence argument read.\newline Given the return, bibliography and programs; one AI reader.\newline \path{status/REVIEW_S6_02_FAMILY.md}, Section 3, D.1.5\\
V027 & Claude Fable 5.2, effort high\newline 2026-11-12; F\newline FAM D.1.6 & known. Standard theory of 2-tori (Appell-Humbert, Riemann conditions); the cited lecture notes not opened by me.\newline Given the return, bibliography and programs; one AI reader.\newline \path{status/REVIEW_S6_02_FAMILY.md}, Section 3, D.1.6\\
V028 & Claude Fable 5.2, effort high\newline 2026-11-12; F,D\newline FAM D.1.7 & correct. Pf(a eta + b E) = 6a\textasciicircum{}2 by hand.\newline Given the return, bibliography and programs; one AI reader.\newline \path{status/REVIEW_S6_02_FAMILY.md}, Section 3, D.1.7\\
V029 & Claude Fable 5.2, effort high\newline 2026-11-12; F\newline FAM D.2.1 & correct. Uses H\textasciicircum{}2(X;Z) = 0; argument complete.\newline Given the return, bibliography and programs; one AI reader.\newline \path{status/REVIEW_S6_02_FAMILY.md}, Section 3, D.2.1\\
V030 & Claude Fable 5.2, effort high\newline 2026-11-12; F\newline FAM D.2.2 & correct. Read.\newline Given the return, bibliography and programs; one AI reader.\newline \path{status/REVIEW_S6_02_FAMILY.md}, Section 3, D.2.2\\
V031 & Claude Fable 5.2, effort high\newline 2026-11-12; F\newline FAM D.2.3 & known. A parametrization (theta divisors) more than a list; standard; the cited notes not opened by me.\newline Given the return, bibliography and programs; one AI reader.\newline \path{status/REVIEW_S6_02_FAMILY.md}, Section 3, D.2.3\\
V032 & Claude Fable 5.2, effort high\newline 2026-11-12; F,D,C\newline FAM F.1.1 & correct. By program:\allowbreak{} Omega-invariance, A1 = g(a,(1,-1),-2), A2 = g(S,(0,-1),3), alpha(S,T,a) = (9,1,4), chi a homomorphism (300 random products), chi = 0 on 300 random words, K, U, V, [U,V] = Z\textasciicircum{}-1, chi(I+E) = 1. The index argument by hand.\newline Given the return, bibliography and programs; one AI reader.\newline \path{status/REVIEW_S6_02_FAMILY.md}, Section 3, F.1.1\\
V033 & Claude Fable 5.2, effort high\newline 2026-11-12; F\newline FAM F.1.2 & correct. Read; follows from the group law.\newline Given the return, bibliography and programs; one AI reader.\newline \path{status/REVIEW_S6_02_FAMILY.md}, Section 3, F.1.2\\
V034 & Claude Fable 5.2, effort high\newline 2026-11-12; F,C\newline FAM F.2.1 & correct. J0\textasciicircum{}2 = -I, J0\textasciicircum{}t Omega J0 = Omega, eigenvalues (1,-1,6,-6), R = g(T\textasciicircum{}72, 0, -12) in Gamma, R J0 = J0 R, (R-I)\textasciicircum{}2 = 0 by program.\newline Given the return, bibliography and programs; one AI reader.\newline \path{status/REVIEW_S6_02_FAMILY.md}, Section 3, F.2.1\\
V035 & Claude Fable 5.2, effort high\newline 2026-11-12; F,D\newline FAM F.2.2 & correct. Transitivity with compact stabilizer on each sign region; density of rational conjugates of J0 in the null locus; the 12th-power step. Redone by hand.\newline Given the return, bibliography and programs; one AI reader.\newline \path{status/REVIEW_S6_02_FAMILY.md}, Section 3, F.2.2\\
V036 & Claude Fable 5.2, effort high\newline 2026-11-12; F,C\newline FAM F.2.3 & correct. q\_\allowbreak{}e = -6/\allowbreak{}D checked numerically at a random period point.\newline Given the return, bibliography and programs; one AI reader.\newline \path{status/REVIEW_S6_02_FAMILY.md}, Section 3, F.2.3\\
V037 & Claude Fable 5.2, effort high\newline 2026-11-12; F\newline FAM F.3.1 & correct. j(tau) = 1728 t is a holomorphic left inverse.\newline Given the return, bibliography and programs; one AI reader.\newline \path{status/REVIEW_S6_02_FAMILY.md}, Section 3, F.3.1\\
V038 & Claude Fable 5.2, effort high\newline 2026-11-12; F\newline FAM F.3.2 & correct. Read.\newline Given the return, bibliography and programs; one AI reader.\newline \path{status/REVIEW_S6_02_FAMILY.md}, Section 3, F.3.2\\
V039 & Claude Fable 5.2, effort high\newline 2026-11-12; F,D\newline FAM F.3.3 & correct. mu' = -tau'/\allowbreak{}3, beta' = 2 tau'/\allowbreak{}3 at z1 recomputed by hand.\newline Given the return, bibliography and programs; one AI reader.\newline \path{status/REVIEW_S6_02_FAMILY.md}, Section 3, F.3.3\\
V040 & Claude Fable 5.2, effort high\newline 2026-11-12; C\newline ab\_\allowbreak{}exact & identical. Python 3.13.9; sympy 1.13.1; mpmath 1.3.0. The review identifies the complete outputs as identical.\newline Rerun of the shipped program in a scratch copy.\newline \path{status/REVIEW_S6_02_FAMILY.md}, Section 2\\
V041 & Claude Fable 5.2, effort high\newline 2026-11-12; C\newline ab\_\allowbreak{}geometry & identical. Python 3.13.9; sympy 1.13.1; mpmath 1.3.0. The review identifies the complete outputs as identical.\newline Rerun of the shipped program in a scratch copy.\newline \path{status/REVIEW_S6_02_FAMILY.md}, Section 2\\
V042 & Claude Fable 5.2, effort high\newline 2026-11-12; C\newline ab\_\allowbreak{}deformation & identical. Python 3.13.9; sympy 1.13.1; mpmath 1.3.0. The review identifies the complete outputs as identical.\newline Rerun of the shipped program in a scratch copy.\newline \path{status/REVIEW_S6_02_FAMILY.md}, Section 2\\
V043 & Claude Fable 5.2, effort high\newline 2026-11-12; C\newline cd\_\allowbreak{}bounds & identical. Python 3.13.9; sympy 1.13.1; mpmath 1.3.0. The review identifies the complete outputs as identical.\newline Rerun of the shipped program in a scratch copy.\newline \path{status/REVIEW_S6_02_FAMILY.md}, Section 2\\
V044 & Claude Fable 5.2, effort high\newline 2026-11-12; C\newline cd\_\allowbreak{}exact & identical. Python 3.13.9; sympy 1.13.1; mpmath 1.3.0. The review identifies the complete outputs as identical.\newline Rerun of the shipped program in a scratch copy.\newline \path{status/REVIEW_S6_02_FAMILY.md}, Section 2\\
V045 & Claude Fable 5.2, effort high\newline 2026-11-12; C\newline cd\_\allowbreak{}local\_\allowbreak{}bound & identical. Python 3.13.9; sympy 1.13.1; mpmath 1.3.0. The review identifies the complete outputs as identical.\newline Rerun of the shipped program in a scratch copy.\newline \path{status/REVIEW_S6_02_FAMILY.md}, Section 2\\
V046 & Claude Fable 5.2, effort high\newline 2026-11-12; C\newline cd\_\allowbreak{}multiples\_\allowbreak{}initial & identical. Python 3.13.9; sympy 1.13.1; mpmath 1.3.0. The review identifies the complete outputs as identical.\newline Rerun of the shipped program in a scratch copy.\newline \path{status/REVIEW_S6_02_FAMILY.md}, Section 2\\
V047 & Claude Fable 5.2, effort high\newline 2026-11-12; C\newline cd\_\allowbreak{}periods & identical. Python 3.13.9; sympy 1.13.1; mpmath 1.3.0. The review identifies the complete outputs as identical.\newline Rerun of the shipped program in a scratch copy.\newline \path{status/REVIEW_S6_02_FAMILY.md}, Section 2\\
V048 & Claude Fable 5.2, effort high\newline 2026-11-12; C\newline cd\_\allowbreak{}special\_\allowbreak{}intervals & identical. Python 3.13.9; sympy 1.13.1; mpmath 1.3.0. The review identifies the complete outputs as identical.\newline Rerun of the shipped program in a scratch copy.\newline \path{status/REVIEW_S6_02_FAMILY.md}, Section 2\\
V049 & Claude Fable 5.2, effort high\newline 2026-11-12; C\newline cd\_\allowbreak{}torsion\_\allowbreak{}independent & identical. Python 3.13.9; sympy 1.13.1; mpmath 1.3.0. The review identifies the complete outputs as identical.\newline Rerun of the shipped program in a scratch copy.\newline \path{status/REVIEW_S6_02_FAMILY.md}, Section 2\\
V050 & Claude Fable 5.2, effort high\newline 2026-11-12; C\newline f\_\allowbreak{}exact & identical. Python 3.13.9; sympy 1.13.1; mpmath 1.3.0. The review identifies the complete outputs as identical.\newline Rerun of the shipped program in a scratch copy.\newline \path{status/REVIEW_S6_02_FAMILY.md}, Section 2\\
V051 & Claude Fable 5.2, effort high\newline 2026-11-12; C\newline cd\_\allowbreak{}numerics & identical apart from last platform digits. Two double-precision grid cross-checks differed only in their last digits.\newline Methods exactly as recorded by the review.\newline \path{status/REVIEW_S6_02_FAMILY.md}, Section 2\\
V052 & Claude Fable 5.2, effort high\newline 2026-11-12; C\newline f\_\allowbreak{}interval\_\allowbreak{}verify & identical. 19,911 leaves; exact area 23/\allowbreak{}40; 19,701 norm boxes; 36 local boxes; 174 outside the unit circle.\newline Methods exactly as recorded by the review.\newline \path{status/REVIEW_S6_02_FAMILY.md}, Section 2\\
V053 & Claude Fable 5.2, effort high\newline 2026-11-12; C\newline all 19,701 norm boxes & all recertified. Own \path{recertify_all.py} runner called the supplied certify routine; identical upper-bound strings; 73 seconds, 12 processes. Same mpmath interval library, not a new implementation.\newline Methods exactly as recorded by the review.\newline \path{status/REVIEW_S6_02_FAMILY.md}, Section 2\\
V054 & Claude Fable 5.2, effort high\newline 2026-11-12; C\newline independent exact lattice calculations & checked. \path{v1_exact.py} used the AC26 matrices for FAM A.2, A.3, B.2, B.4, F.1 and F.2.\newline Methods exactly as recorded by the review.\newline \path{status/REVIEW_S6_02_FAMILY.md}, Section 2\\
V055 & Claude Fable 5.2, effort high\newline 2026-11-12; C\newline independent period numerics & checked as numerical evidence. \path{v2_numeric.py}:\allowbreak{} 6,651 grid points, the modular arc and circles around i; independent mpmath theta evaluation.\newline Methods exactly as recorded by the review.\newline \path{status/REVIEW_S6_02_FAMILY.md}, Section 2\\
V056 & Claude Fable 5.2, effort high\newline 2026-11-12; C\newline independent elliptic arithmetic & checked. \path{v3_ec.py} recomputed the D.1 arithmetic and C.2.2 constants.\newline Methods exactly as recorded by the review.\newline \path{status/REVIEW_S6_02_FAMILY.md}, Section 2\\
V057 & Claude Fable 5.2, effort high\newline 2026-11-12; C\newline independent special values & checked as numerical evidence. \path{v4_closed_forms.py} evaluated the closed forms listed in Section 4.9; proofs were not supplied by that numerical check.\newline Methods exactly as recorded by the review.\newline \path{status/REVIEW_S6_02_FAMILY.md}, Section 2\\
V058 & Claude Opus 5.5, effort max\newline 2026-11-12; S\newline DGKXv1 three theorem statements and v2 contents & known labels confirmed. The Kuranishi, exact-period and tangent-cohomology statements were matched in v1; their absence in v2 was checked.\newline Both supplied TeX sources opened by the main session.\newline \path{status/REVIEW_S6_02_FAMILY.md}, added main-session note\\
V059 & Claude Fable 5.2, effort max\newline 2026-11-16; R,S\newline DGKXv1 and v2 scope & matched statements; proofs not read. Abstract and introduction; 17 statements in Sections 6--7; three Appendix C propositions; Appendix D opening; full term counts. For v2:\allowbreak{} abstract, introduction statements and section list.\newline Expert had the TeX and the status discussion.\newline \path{status/RESPONSE_EXPERT.md}, Section 1\\
V060 & Claude Fable 5.2, effort max\newline 2026-11-13, recorded 2026-11-16; D,C\newline ARI C.8.6 and C.8.7 & derived and numerically tested. Mahler-measure identity and two beta integrals derived by hand and tested to seven digits. This is not a claim of a fresh line-by-line read of the source proof.\newline Expert derivations reported in the supplied response.\newline \path{status/RESPONSE_EXPERT.md}, Section 4\\
V061 & Claude Fable 5.2, effort high\newline 2026-11-15; D\newline ARI C.4.2 & checks agree. Hand checks at t=1, 32/\allowbreak{}5 and 0; no code rerun.\newline Bounded arithmetic intake and its recorded hand checks.\newline \path{status/INTAKE_S6_05_FINAL.md}, Section 1; \path{INTAKE_FINAL.md}, Section 4\\
V062 & Claude Fable 5.2, effort high\newline 2026-11-15; F,D\newline ARI C.8.5 displayed inequalities & displayed inequalities agree. Every displayed inequality and the size of the three uniform bounds were checked. C.8.1--C.8.3 were not read line by line. Code and PDF were not read.\newline Given the source proof and intake notes.\newline \path{status/INTAKE_S6_05_FINAL.md}, Section 1\\
V063 & GPT-6 Astra, ChatGPT Pro mode\newline PUSH producing run; C\newline PUSH E.4 Arb certification & all boxes certified as reported. 19,701 boxes; largest upper enclosure below 1.999877571767410525. The intake reports this result and did not rerun the code. Detailed precision is in the source certificate block.\newline Separate interval library; same rational box endpoints.\newline \path{status/INTAKE_FINAL.md}, Section 1, item 4; PUSH E.4\\
V064 & Claude Fable 5.2, effort high\newline 2026-11-16; F\newline INV A1.3; direct images & no error found in the portion read. Proof read; A1.1--A1.2, which supply local freeness and direct images, were not read.\newline Bounded intake; precise read limits retained.\newline \path{status/INTAKE_S6_01.md}, opening and Section 2\\
V065 & Claude Fable 5.2, effort high\newline 2026-11-16; F\newline INV A2.1; period determinant identity & no error found in the portion read. Statement and full proof read.\newline Bounded intake; precise read limits retained.\newline \path{status/INTAKE_S6_01.md}, opening and Section 2\\
V066 & Claude Fable 5.2, effort high\newline 2026-11-16; F\newline INV A2.2; the fourth-root scalar & no error found in the portion read. Statement and full proof read; separate scalar computation is listed below.\newline Bounded intake; precise read limits retained.\newline \path{status/INTAKE_S6_01.md}, opening and Section 2\\
V067 & Claude Fable 5.2, effort high\newline 2026-11-16; F\newline INV A2.4; nonzero Kodaira--Spencer scalar & no error found in the portion read. Statement and full proof read.\newline Bounded intake; precise read limits retained.\newline \path{status/INTAKE_S6_01.md}, opening and Section 2\\
V068 & Claude Fable 5.2, effort high\newline 2026-11-16; F,D\newline INV A3.1; Hodge table & no error found in the portion read. Table recomputed from the six direct-image sheaves; Serre duality, Euler characteristic and the AC26 relation checked.\newline Bounded intake; precise read limits retained.\newline \path{status/INTAKE_S6_01.md}, opening and Section 2\\
V069 & Claude Fable 5.2, effort high\newline 2026-11-16; F\newline INV C1.1; Frolicher sequence & no error found in the portion read. Every arrow followed; no higher differential has both nonzero source and target; middle rows exact.\newline Bounded intake; precise read limits retained.\newline \path{status/INTAKE_S6_01.md}, opening and Section 2\\
V070 & Claude Fable 5.2, effort high\newline 2026-11-16; F,D\newline INV C2.1; Bott--Chern and Aeppli numbers & no error found in the portion read. First half of proof read; numbers recounted from the double-complex arrows. Remaining proof was not fully read.\newline Bounded intake; precise read limits retained.\newline \path{status/INTAKE_S6_01.md}, opening and Section 2\\
V071 & Claude Fable 5.2, effort high\newline 2026-11-16; F\newline INV C3.3; square-zero product & no error found in the portion read. Proof read and Leibniz argument followed; it uses C3.1, which was not read. Binding status remains A.\newline Bounded intake; precise read limits retained.\newline \path{status/INTAKE_S6_01.md}, opening and Section 2\\
V072 & Claude Fable 5.2, effort high\newline 2026-11-16; F,D\newline INV B2.1; higher-multiplicity canonical divisor & no error found in the portion read. Proof read; the two orders and the fibre degree checked. Higher-multiplicity Hodge inputs were not read.\newline Bounded intake; precise read limits retained.\newline \path{status/INTAKE_S6_01.md}, opening and Section 2\\
V073 & Claude Fable 5.2, effort high\newline 2026-11-16; F\newline INV B3.2; degree-n covers & no error found in the portion read. Proof read; the cusp residue mechanism and Betti agreement identified. B3.1 was not read.\newline Bounded intake; precise read limits retained.\newline \path{status/INTAKE_S6_01.md}, opening and Section 2\\
V074 & Claude Fable 5.2, effort high\newline 2026-11-16; F\newline INV B3.3; cover cohomology & no error found in the portion read. Statement and first half of proof read; not a complete proof check.\newline Bounded intake; precise read limits retained.\newline \path{status/INTAKE_S6_01.md}, opening and Section 2\\
V075 & Claude Fable 5.2, effort high\newline 2026-11-16; C\newline INV A2.2 scalar & 40 digits agree. \path{s601_scalar_check.py} reproduced 0.2366093140\allowbreak{}8078522468\allowbreak{}7705113450\allowbreak{}4496062943 and j second derivative -162 Gamma(1/\allowbreak{}4)\textasciicircum{}8/\allowbreak{}pi\textasciicircum{}4.\newline Independent numerical route using the local period expansion and j second derivative.\newline \path{status/INTAKE_S6_01.md}, Section 2\\
V076 & Claude Fable 5.2, effort high\newline 2026-11-16; F,D\newline PUSH A.3.1 & no gap found in the portion read. Full proof read; virtual Euler count redone as 24-(12+12)+0=0.\newline Bounded intake with earlier notes and source text.\newline \path{status/INTAKE_FINAL.md}, opening and Section 1\\
V077 & Claude Fable 5.2, effort high\newline 2026-11-16; F\newline PUSH A.3.2 & no gap found in the portion read. Full proof read; the specified candidate cannot smooth to the sphere fibre.\newline Bounded intake with earlier notes and source text.\newline \path{status/INTAKE_FINAL.md}, opening and Section 1\\
V078 & Claude Fable 5.2, effort high\newline 2026-11-16; F\newline PUSH D.1.1 & no gap found in the portion read. Full proof read; Albanese modular invariant natural boundary argument.\newline Bounded intake with earlier notes and source text.\newline \path{status/INTAKE_FINAL.md}, opening and Section 1\\
V079 & Claude Fable 5.2, effort high\newline 2026-11-16; F\newline PUSH D.1.2 & no gap found in the portion read. Full proof read; base change on the curve and Buchdahl input checked in the argument.\newline Bounded intake with earlier notes and source text.\newline \path{status/INTAKE_FINAL.md}, opening and Section 1\\
V080 & Claude Fable 5.2, effort high\newline 2026-11-16; R\newline PUSH D.1.3 & no gap found in the portion read. Opening read, not its whole proof.\newline Bounded intake with earlier notes and source text.\newline \path{status/INTAKE_FINAL.md}, opening and Section 1\\
V081 & Claude Fable 5.2, effort high\newline 2026-11-16; F\newline PUSH B.4.4 & no gap found in the portion read. Marked-map extension proof read.\newline Bounded intake with earlier notes and source text.\newline \path{status/INTAKE_FINAL.md}, opening and Section 1\\
V082 & Claude Fable 5.2, effort high\newline 2026-11-16; F\newline PUSH B.4.27 & no gap found in the portion read. Log-torus sheaf cohomology proof read.\newline Bounded intake with earlier notes and source text.\newline \path{status/INTAKE_FINAL.md}, opening and Section 1\\
V083 & Claude Fable 5.2, effort high\newline 2026-11-16; F\newline PUSH B.4.30 & no gap found in the portion read. Base modification and extension-of-differences proof read.\newline Bounded intake with earlier notes and source text.\newline \path{status/INTAKE_FINAL.md}, opening and Section 1\\
V084 & Claude Fable 5.2, effort high\newline 2026-11-16; F\newline PUSH B.4.31 & no gap found in the portion read. Global-meridian proof read.\newline Bounded intake with earlier notes and source text.\newline \path{status/INTAKE_FINAL.md}, opening and Section 1\\
V085 & Claude Fable 5.2, effort high\newline 2026-11-16; F\newline PUSH B.4.32 & no gap found in the portion read. Global-loop defect proof read.\newline Bounded intake with earlier notes and source text.\newline \path{status/INTAKE_FINAL.md}, opening and Section 1\\
V086 & Claude Fable 5.2, effort high\newline 2026-11-16; F\newline PUSH B.4.33 & no gap found in the portion read. Proof read; explicit objects B.4.9, B.4.18 and B.4.19 were not independently re-derived.\newline Bounded intake with earlier notes and source text.\newline \path{status/INTAKE_FINAL.md}, opening and Section 1\\
V087 & Claude Fable 5.2, effort high\newline 2026-11-16; F\newline PUSH C.3.21 & no gap found in the portion read. Actual delta-square proof read.\newline Bounded intake with earlier notes and source text.\newline \path{status/INTAKE_FINAL.md}, opening and Section 1\\
V088 & Claude Fable 5.2, effort high\newline 2026-11-16; F\newline PUSH C.3.22 & no gap found in the portion read. Signed-cap comparison proof read.\newline Bounded intake with earlier notes and source text.\newline \path{status/INTAKE_FINAL.md}, opening and Section 1\\
V089 & Claude Fable 5.2, effort high\newline 2026-11-16; F\newline PUSH C.3.23 & no gap found in the portion read. Equivariant integer proof read.\newline Bounded intake with earlier notes and source text.\newline \path{status/INTAKE_FINAL.md}, opening and Section 1\\
V090 & Claude Fable 5.2, effort high\newline 2026-11-16; F\newline PUSH C.3.24 & no gap found in the portion read. Monodromy-value proof read; C.3.10, C.3.14 and C.3.18 were not read in this intake.\newline Bounded intake with earlier notes and source text.\newline \path{status/INTAKE_FINAL.md}, opening and Section 1\\
V091 & Claude Fable 5.2, effort high\newline 2026-11-16; F\newline PUSH C.3.27 & no gap found in the portion read. Second determination proof read.\newline Bounded intake with earlier notes and source text.\newline \path{status/INTAKE_FINAL.md}, opening and Section 1\\
V092 & Claude Fable 5.2, effort high\newline 2026-11-16; F\newline PUSH E.2.21 & no gap found in the portion read. Disc-completion proof read.\newline Bounded intake with earlier notes and source text.\newline \path{status/INTAKE_FINAL.md}, opening and Section 1\\
V093 & Claude Fable 5.2, effort high\newline 2026-11-16; F\newline PUSH E.2.22 & no gap found in the portion read. Connected-base propagation proof read; E.2.5--E.2.20 remain unread inputs.\newline Bounded intake with earlier notes and source text.\newline \path{status/INTAKE_FINAL.md}, opening and Section 1\\
V094 & Claude Fable 5.2, effort high\newline 2026-11-16; F\newline PUSH E.2.24 & no gap found in the portion read. Continuous-deformation proof read; the Kuranishi theorem is an external input.\newline Bounded intake with earlier notes and source text.\newline \path{status/INTAKE_FINAL.md}, opening and Section 1\\
V095 & Claude Fable 5.2, effort high\newline 2026-11-16; F\newline MAPS A.3.1 & no conflict found. Full proof read. A.3.2 was read only as statement and setup; fixed-locus conclusions retain A status.\newline Bounded intake.\newline \path{status/INTAKE_FINAL.md}, opening and Sections 3,5\\
V096 & Claude Fable 5.2, effort high\newline 2026-11-16; F\newline ARI C.2.2 & agrees with the real-parameter criterion. Full proof read; compared with FAM B.4.2 and MAPS A.3.\newline Bounded intake.\newline \path{status/INTAKE_FINAL.md}, opening and Section 5\\
V097 & Claude Fable 5.2, effort high\newline 2026-11-16; D\newline CDP C5b at n=-1,0,12 & Euler characteristic agrees. The supplied response records the three values; INTAKE\_\allowbreak{}FINAL gives the n=0 tuple (1,1,1,0).\newline Hand substitution, not a full independent proof.\newline \path{status/RESPONSE_INTAKE.md}, Section 2; \path{INTAKE_FINAL.md}, Section 2\\
V098 & Claude Fable 5.2, effort high\newline 2026-11-16; R\newline PUSH E.2.1 and CDP E11 & same higher-multiplicity tangent tuple. Two runs, no second reader of these proofs; this qualifier remains attached to the T statement.\newline Cross-run comparison of separate producing sessions.\newline \path{status/INTAKE_FINAL.md}, Section 5 X3; \path{RESPONSE_INTAKE.md}, Section 2\\
V099 & Claude Opus 5.5, effort max\newline Recorded 2026-11-16; D\newline special norm and phase calculation & n=3 division case checked by intake. The intake records NOTES Section 33(a) as a valid derivation of the special norm and phases, and Section 34 as the exact nome-ratio calculation. No unseen notes are asserted to have been opened here.\newline Main-session calculation as evaluated by the intake.\newline \path{status/INTAKE_FINAL.md}, Section 5, NOTES Sections 33--34 discussion\\
V100 & GPT-6 Astra, ChatGPT Pro mode\newline 2026-11-16; S\newline AC26; Main Theorem items (1)--(5); Remark 3.12; Remark 6.4; Theorem 9.1; Section 10 & opened and matched. Construction on S6, algebraic dimension one, canonical bundle; elliptic-logarithm generator; classification left open in Remark 6.4; the source of the counterexample in the prescribed introduction paragraph. Checked in the supplied AI transcription, not in an original PDF.\newline Assembly with supplied source and review context; not an independent proof review.\newline \path{records/root/citation_checks.json}; \path{original/AC26_transcription.tex}\\
V101 & GPT-6 Astra, ChatGPT Pro mode\newline 2026-11-16; S\newline E26; Section 2; Lemma 2.1; Corollary 2.2; Definition 2.3; Proposition 2.4; Sections 3,5,6; Remark 8.1 and Conjecture 8.2 & opened and matched. Pointed cubic, Mordell-Weil generator and reference lift up to scale; inverse-translation convention; small-parameter compactness; lattice basis and affine data; exact correction and counterexample targets.\newline Assembly with supplied source and review context; not an independent proof review.\newline \path{records/root/citation_checks.json}; \path{papers/arXiv_2609.38442_Engel_src/S6.tex}\\
V102 & GPT-6 Astra, ChatGPT Pro mode\newline 2026-11-16; S\newline DGKXv1; Theorems 1.3,1.4; Theorem 7.1 and thm:\allowbreak{}all-maps-factor; p2:\allowbreak{}thm:\allowbreak{}full-automorphisms; p2:\allowbreak{}cor:\allowbreak{}bare-period; Appendix C & opened and matched. Standard-data exact 2Z identification, tangent tuple and Kuranishi germ, intrinsic bi-anticanonical pencil, unique factorization of every map to P1, full automorphism group and unevaluated punctured-disc radius. All-affine-data scope of FAM A.1.1 is not identical to the cited standard-data theorem.\newline Assembly with supplied source and review context; not an independent proof review.\newline \path{records/root/citation_checks.json}; \path{papers/arXiv_2610.00250v1_Du_Guo_Kusakabe_Xie.tex}\\
V103 & GPT-6 Astra, ChatGPT Pro mode\newline 2026-11-16; S\newline DGKXv2; Abstract, introduction, section input list; full-text term searches & opened and matched. Oka content retained; the v1 exact-period, tangent-cohomology, intrinsic-pencil and Kuranishi sections and statements are absent. No reason for this omission is inferred.\newline Assembly with supplied source and review context; not an independent proof review.\newline \path{records/root/citation_checks.json}; \path{papers/arXiv_2610.00250v2_Du_Guo_Kusakabe_Xie.tex}\\
V104 & GPT-6 Astra, ChatGPT Pro mode\newline 2026-11-16; S\newline VL26; Theorem 1.1; Proposition 4.7, label nonbih, and proof & opened and matched. The source's higher multiplicities and non-isotriviality argument; the factorization of maps to the base used for the stated-model proof. The full coordinate dictionary remains an announcement.\newline Assembly with supplied source and review context; not an independent proof review.\newline \path{records/root/citation_checks.json}; \path{papers/arXiv_2610.39102_Viaclovsky_log_transformations_src/Logarithmic_S6_October_28_2026_Final.tex}\\
V105 & GPT-6 Astra, ChatGPT Pro mode\newline 2026-11-16; S\newline OS91; Main Theorem table entry 49; Section 3 equation (3.2) & opened and matched. Root lattice E6 plus A1, determinant six, no torsion and Mordell-Weil lattice with generator height 1/\allowbreak{}6; height input for the IV*,III,I1 surface.\newline Assembly with supplied source and review context; not an independent proof review.\newline \path{records/root/citation_checks.json}; \url{https://rikkyo.repo.nii.ac.jp/record/10006/files/AA00610867_40-01_08.pdf}\\
V106 & GPT-6 Astra, ChatGPT Pro mode\newline 2026-11-16; S\newline Shi90; Theorem 8.6, equations (8.11)--(8.12), page 228; local contribution table (8.16), page 229 & opened and matched. Height pairing formula and local contributions 1/\allowbreak{}2 at III and 4/\allowbreak{}3 at IV*. The source-run citation to Table 8.9 is an incorrect locator; the matched table is (8.16).\newline Assembly with supplied source and review context; not an independent proof review.\newline \path{records/root/citation_checks.json}; \url{https://rikkyo.repo.nii.ac.jp/record/10027/files/AA00610867_39-02_05.pdf}\\
V107 & GPT-6 Astra, ChatGPT Pro mode\newline 2026-11-16; S\newline CDP20; Corollary 2.3 & opened and matched. The asserted algebraic dimension zero for a compact complex threefold homeomorphic to S6; used solely in the prescribed introduction paragraph.\newline Assembly with supplied source and review context; not an independent proof review.\newline \path{records/root/citation_checks.json}; \url{https://www-fourier.univ-grenoble-alpes.fr/~demailly/source_files/cdp/cdp-corrigendum-v5.pdf}\\
V108 & GPT-6 Astra, ChatGPT Pro mode\newline 2026-11-16; S\newline AC26; Theorem 3.4(ii)--(iv); Remark 3.12 & opened and matched. Finite mu values, period laws, invariant D and elliptic-logarithm interpretation; none of the theta norm, exact radius, integral, division or measure claims found by the recorded searches.\newline Assembly with supplied source and review context; not an independent proof review.\newline \path{records/periods/citation_checks.json}; \path{original/AC26_transcription.tex}\\
V109 & GPT-6 Astra, ChatGPT Pro mode\newline 2026-11-16; S\newline E26; Relation to the original manuscript; §2; Remark 8.1 & opened and matched. Geometric interpretation of the four periods and variation of the elliptic curves with linearization parameter; exact nomes and constants not present.\newline Assembly with supplied source and review context; not an independent proof review.\newline \path{records/periods/citation_checks.json}; \path{papers/arXiv_2609.38442_Engel_src/S6.tex}\\
V110 & GPT-6 Astra, ChatGPT Pro mode\newline 2026-11-16; S\newline DGKXv1; Theorem 1.4 (p2:\allowbreak{}thm:\allowbreak{}exact-period); Corollary p2:\allowbreak{}cor:\allowbreak{}bare-period; definition of M\_\allowbreak{}0 & opened and matched. Theorem 1.4 states bare and pair isomorphism iff parameter difference lies in 2Z; Corollary p2:\allowbreak{}cor:\allowbreak{}bare-period gives 0<|w|<exp(-pi M\_\allowbreak{}0); M\_\allowbreak{}0 is not evaluated.\newline Assembly with supplied source and review context; not an independent proof review.\newline \path{records/periods/citation_checks.json}; \path{papers/arXiv_2610.00250v1_Du_Guo_Kusakabe_Xie.tex}\\
V111 & GPT-6 Astra, ChatGPT Pro mode\newline 2026-11-16; S\newline EZ82; Unnumbered theorem p.399 and Fourier corollary p.400 & opened and matched. The zeros of Weierstrass wp are the elliptic integral branches with cusp constant log(5+2sqrt6)/\allowbreak{}(2pi i).\newline Assembly with supplied source and review context; not an independent proof review.\newline \path{records/periods/citation_checks.json}; \url{https://people.mpim-bonn.mpg.de/zagier/files/doi/10.1007/BF01453974/fulltext.pdf}\\
V112 & GPT-6 Astra, ChatGPT Pro mode\newline 2026-11-16; S\newline DI10; §2, equation (6), unnumbered Eichler--Zagier theorem; Theorem 1 & opened and matched. Zero formula and its branch normalization; local variable 1728/\allowbreak{}j is reciprocal to the family t.\newline Assembly with supplied source and review context; not an independent proof review.\newline \path{records/periods/citation_checks.json}; \url{https://www.math.ucla.edu/~wdduke/preprints/zeros.pdf}\\
V113 & GPT-6 Astra, ChatGPT Pro mode\newline 2026-11-16; S\newline DM20; Theorem 1.1; Corollary 1.2, preprint pp.2--3 & opened and matched. Adelically metrized height line and small-height Galois-orbit probability equidistribution for a section of nonzero generic canonical height over a number field.\newline Assembly with supplied source and review context; not an independent proof review.\newline \path{records/periods/citation_checks.json}; \url{https://arxiv.org/pdf/1701.07947}\\
V114 & GPT-6 Astra, ChatGPT Pro mode\newline 2026-11-16; S\newline Johansson2016; §§2--3 & opened and matched. Midpoint-radius representation and outward error bounds underpin the supplied Arb computation.\newline Assembly with supplied source and review context; not an independent proof review.\newline \path{records/periods/citation_checks.json}; \url{https://arxiv.org/pdf/1611.02831}\\
V115 & GPT-6 Astra, ChatGPT Pro mode\newline 2026-11-16; S\newline PythonFlint2026; arb constructor, abs\_\allowbreak{}upper and upper methods & opened and matched. Exact dyadic/\allowbreak{}rational constructor inputs, ball enclosures, outward upper endpoints.\newline Assembly with supplied source and review context; not an independent proof review.\newline \path{records/periods/citation_checks.json}; \url{https://python-flint.readthedocs.io/en/latest/arb.html}\\
V116 & GPT-6 Astra, ChatGPT Pro mode\newline 2026-11-16; S\newline DGKXv2; Abstract and Introduction, lines 134--270; whole-file searches for p2:\allowbreak{}thm:\allowbreak{}exact-period and p2:\allowbreak{}cor:\allowbreak{}bare-period & opened and matched. The exact-period theorem present in version 1 is omitted in version 2. No mathematical conclusion is inferred from this omission.\newline Assembly with supplied source and review context; not an independent proof review.\newline \path{records/periods/citation_checks.json}; \path{papers/arXiv_2610.00250v2_Du_Guo_Kusakabe_Xie.tex}\\
V117 & GPT-6 Astra, ChatGPT Pro mode\newline 2026-11-16; S\newline AC26; Supplied transcription:\allowbreak{} Definition 3.1 and Theorem 3.4(i), printed p. 14; Theorem 9.1, p. 64; Proposition 9.11, pp. 70--71; Lemma 9.12, pp. 72--74; Proposition 9.13, pp. 74--76; Remarks 1.3 and 9.21, pp. 5 and 81; Corollary 9.22 and Proposition 9.23, p. 81; Theorems 7.17 and 7.22; Lemma 10.3 and Theorem 10.5, pp. 87--89. & opened and matched. Opened and matched construction/\allowbreak{}period inputs, known structure direct images and Hodge constraints, canonical and topology credits, correction targets, and the memberwise algebraic-dimension input. No source prose reproduced. Verification concerns the supplied transcription, not a newly inspected original PDF.\newline Assembly with supplied source and review context; not an independent proof review.\newline \path{records/cohomology/citation_checks.json}; \url{https://alpo.ge/s6.pdf}\\
V118 & GPT-6 Astra, ChatGPT Pro mode\newline 2026-11-16; S\newline E26; Introduction, lines 215--277; proposition fundamental-group, lines 1215--1231; final invariant-Hodge-form remark, lines 1723--1731. & opened and matched. Opened the original supplied TeX; it credits AC26 for Dolbeault information and expressly omits those calculations. The construction comparison, standard-collar presentation and fibre-form background match the cited/\allowbreak{}credited scope.\newline Assembly with supplied source and review context; not an independent proof review.\newline \path{records/cohomology/citation_checks.json}; \url{https://arxiv.org/abs/2609.38442}\\
V119 & GPT-6 Astra, ChatGPT Pro mode\newline 2026-11-16; S\newline DGKXv1; Theorem 1.3 (p2:\allowbreak{}thm:\allowbreak{}kuranishi-aut), lines 310--327; Proposition 6.10 (p2:\allowbreak{}prop:\allowbreak{}structure-direct-images), lines 9151ff.; Theorem 6.12 (p2:\allowbreak{}thm:\allowbreak{}tangent-cohomology), lines 9445ff.; p2:\allowbreak{}prop:\allowbreak{}kuranishi-line, lines 9877ff.; p2:\allowbreak{}thm:\allowbreak{}intrinsic-fibration, lines 10018ff. & opened and matched. The exact tangent vector, smooth reduced Kuranishi line and miniversality are already present. The structure direct images and intrinsic-pencil inputs also match. The reviewed two DGKX lemma citations have been replaced by BS76, RRV71 and RR70 rather than retained as primary analytic references.\newline Assembly with supplied source and review context; not an independent proof review.\newline \path{records/cohomology/citation_checks.json}; \url{https://arxiv.org/abs/2610.00250v1}\\
V120 & GPT-6 Astra, ChatGPT Pro mode\newline 2026-11-16; S\newline VL26; Theorem 1.1 (t:\allowbreak{}main), source lines 210--244; Proposition 2.3 (logprop), lines 420--496. & opened and matched. Opened the actual source theorem and logarithmic-transformation proposition. They provide the sufficiently small nonzero construction domain, coprimality and admissible free cyclic quotients required in inv:\allowbreak{}B2.1.\newline Assembly with supplied source and review context; not an independent proof review.\newline \path{records/cohomology/citation_checks.json}; \url{https://arxiv.org/abs/2610.39102}\\
V121 & GPT-6 Astra, ChatGPT Pro mode\newline 2026-11-16; S\newline BS76; Chapter III, Theorem 4.10, pp. 132--133; Theorem 4.12(ii), pp. 134--135; web lines 5410--5508. & opened and matched. Original-book text reproduction opened. The first theorem supplies the analytic base-change criterion and comparison; the second gives cohomological flatness and local freeness from constant fibre-cohomology dimension for a flat coherent sheaf over a reduced base. Applied on the source complex-space charts, with canonical overlap comparison.\newline Assembly with supplied source and review context; not an independent proof review.\newline \path{records/cohomology/citation_checks.json}; \url{https://djvu.online/file/2hQiOnz8WuoQk}\\
V122 & GPT-6 Astra, ChatGPT Pro mode\newline 2026-11-16; S\newline RRV71; Original article introduction p. 261 and Section 2, pp. 264--266; original bibliographic identification Invent. Math. 13 (1971), 261--283, DOI 10.1007/\allowbreak{}BF01406078. & opened and matched. Full original-article text opened through a third-party reproduction. The proper analytic duality isomorphism and relative trace were matched. These support the source duality applications on smooth analytic atlases; RR70 identifies the Gorenstein dualizing line.\newline Assembly with supplied source and review context; not an independent proof review.\newline \path{records/cohomology/citation_checks.json}; \url{https://www.academia.edu/5435625/Dualit%C3%A9_relative_en_g%C3%A9om%C3%A9trie_analytique_complexe}\\
V123 & GPT-6 Astra, ChatGPT Pro mode\newline 2026-11-16; S\newline RR70; Section 3, Theorem 1, printed p. 81; Proposition 2 and complete-intersection paragraph, printed p. 91 (PDF page 16). & opened and matched. Original Numdam PDF opened. A Gorenstein analytic space has its dualizing complex represented by an invertible sheaf in the dimension shift; complete intersections are included. This is the dualizing-sheaf input used with the relative theorem.\newline Assembly with supplied source and review context; not an independent proof review.\newline \path{records/cohomology/citation_checks.json}; \url{https://www.numdam.org/item/10.1007/BF02684652.pdf}\\
V124 & GPT-6 Astra, ChatGPT Pro mode\newline 2026-11-16; S\newline Ugarte00; Proposition 3.1, pp. 175--176; Theorem 3.2, pp. 176--177; Corollary 3.3, p. 177. & opened and matched. Original PDF opened through the already retrieved primary-source reference after two fresh URL timeouts. Matched the conditional Hodge relations and the second- and third-page Frölicher constraints for a hypothetical complex S6. The paper cites these as background, not as an evaluation of this family's h12.\newline Assembly with supplied source and review context; not an independent proof review.\newline \path{records/cohomology/citation_checks.json}; \url{https://zaguan.unizar.es/record/165697/files/texto_completo.pdf?version=1}\\
V125 & GPT-6 Astra, ChatGPT Pro mode\newline 2026-11-16; S\newline McHugh18; Version 4, 15 November 2018; Section 7.1, printed pp. 24--25, including the table on p. 25; web lines 3342--3500. & opened and matched. Original PDF opened. The formulas express complex-S6 Bott--Chern dimensions using Dolbeault dimensions and the E2 term, with Aeppli obtained by the stated duality. Exact comparison scope is the background recorded in inv:\allowbreak{}C2.2; no certification of the whole paper is asserted.\newline Assembly with supplied source and review context; not an independent proof review.\newline \path{records/cohomology/citation_checks.json}; \url{https://arxiv.org/pdf/1708.03251}\\
V126 & GPT-6 Astra, ChatGPT Pro mode\newline 2026-11-16; S\newline Angella17; Version 2, November 2017; introduction, Figures 1--2, and Section 4, printed pp. 13--16. & opened and matched. Original PDF opened. Matched the conditional spectral-sequence/\allowbreak{}double-complex diagrams and Bott--Chern relations for complex S6. Used only for the precise historical discussion and credit, with the actual table proved from the stated INV input.\newline Assembly with supplied source and review context; not an independent proof review.\newline \path{records/cohomology/citation_checks.json}; \url{https://arxiv.org/pdf/1705.10518}\\
V127 & GPT-6 Astra, ChatGPT Pro mode\newline 2026-11-16; S\newline Stelzig21; Version 3, 8 March 2021; Theorem 3 (Theorem A in introduction, printed p. 2), Proposition 6(1), p. 16, Lemma 8, pp. 17--18. & opened and matched. Original PDF opened. The bounded-double-complex decomposition into squares and zigzags is valid for vector spaces without a finite-dimensional-component hypothesis. Proposition 6(1) identifies even-zigzag multiplicities with ranks of the spectral-sequence differentials; the BC/\allowbreak{}A counting interpretation was matched.\newline Assembly with supplied source and review context; not an independent proof review.\newline \path{records/cohomology/citation_checks.json}; \url{https://arxiv.org/pdf/1812.00865}\\
V128 & GPT-6 Astra, ChatGPT Pro mode\newline 2026-11-16; S\newline AC26; §9.7, Proposition 9.23 and Corollary 9.24, transcription lines4340--4390 & opened and matched. Aut\textasciicircum{}0=C*, its additive translation coordinate and effective period Z; the fixed double curve and semifree normal weights. Full Aut statement is credited separately to DGKXv1.\newline Assembly check with supplied proofs and review context; source opened and matched, not a new independent proof verification.\newline \path{records/geometry_ends/citation_checks.json}; \path{original/AC26_transcription.tex}\\
V129 & GPT-6 Astra, ChatGPT Pro mode\newline 2026-11-16; S\newline AC26; §9.1, Lemmas9.2--9.4, Proposition9.5 and Corollary9.8; transcription lines3274--3551 & opened and matched. Exact six-coefficient NS equation; five-function independence; indefiniteness; countable exceptional locus and generic NS=Z eta. The +12 wedge coefficient is distinguished from the complex-oriented numerical self-intersection −12.\newline Assembly check with supplied proofs and review context; source opened and matched, not a new independent proof verification.\newline \path{records/geometry_ends/citation_checks.json}; \path{original/AC26_transcription.tex}\\
V130 & GPT-6 Astra, ChatGPT Pro mode\newline 2026-11-16; S\newline AC26; §4.6 Proposition4.6, source PDF p.29; transcription lines1452--1470 & opened and matched. Finite dP6 normalization, opposite boundary-pair identifications, three double rational curves, two triple points and Euler number2.\newline Assembly check with supplied proofs and review context; source opened and matched, not a new independent proof verification.\newline \path{records/geometry_ends/citation_checks.json}; \path{original/AC26_transcription.tex}\\
V131 & GPT-6 Astra, ChatGPT Pro mode\newline 2026-11-16; S\newline E26; §2 Lemma2.1 and Corollary2.2; source lines330--401 & opened and matched. Height(P)=1/\allowbreak{}6; intersection inputs; the line O(P−O)⊗t6P*O(O−P) is the degree-one base pullback and its linearization section has divisor S\_\allowbreak{}infinity.\newline Assembly check with supplied proofs and review context; source opened and matched, not a new independent proof verification.\newline \path{records/geometry_ends/citation_checks.json}; \path{papers/arXiv_2609.38442_Engel_src/S6.tex}\\
V132 & GPT-6 Astra, ChatGPT Pro mode\newline 2026-11-16; S\newline E26; §8 Conjecture8.2, lines1673--1676, and Remarks8.1--8.4 & opened and matched. Exact as-worded conjecture for the specified nonnormal pinching; Remark8.4 asks Klingler's question for the indefinite (1,6) domain. No stated exact index/\allowbreak{}maximal-open answer.\newline Assembly check with supplied proofs and review context; source opened and matched, not a new independent proof verification.\newline \path{records/geometry_ends/citation_checks.json}; \path{papers/arXiv_2609.38442_Engel_src/S6.tex}\\
V133 & GPT-6 Astra, ChatGPT Pro mode\newline 2026-11-16; S\newline DGKXv1; Theorem p2:\allowbreak{}thm:\allowbreak{}full-automorphisms, source line11086, and proof & opened and matched. Aut(X)=Aut\_\allowbreak{}\{P1\}(X)=C*, including every component. The relevant precise theorem is the full-automorphisms theorem, rather than the exact-period theorem alone.\newline Assembly check with supplied proofs and review context; source opened and matched, not a new independent proof verification.\newline \path{records/geometry_ends/citation_checks.json}; \path{papers/arXiv_2610.00250v1_Du_Guo_Kusakabe_Xie.tex}\\
V134 & GPT-6 Astra, ChatGPT Pro mode\newline 2026-11-16; S\newline Schnell19; Theorem4.1 p.1; equation(4.3) p.2; effective-bundle quotient discussion pp.2--3; Theorem4.4 p.3 & opened and matched. Appell--Humbert correspondence, theta factor of automorphy, semipositive effective line bundles and positive integral Hermitian criterion for projectivity. Used in FAM D.1.6 and D.2.3.\newline Assembly check with supplied proofs and review context; source opened and matched, not a new independent proof verification.\newline \path{records/geometry_ends/citation_checks.json}; \url{https://www.math.stonybrook.edu/~cschnell/mat615/lectures/lecture4.pdf}\\
V135 & GPT-6 Astra, ChatGPT Pro mode\newline 2026-11-16; S\newline Buchdahl1999; Theorem11, printed p.297 & opened and matched. A compact complex surface with even first Betti number is Kähler. Applied to smooth connected surfaces with b1=2 in PUSH D.1.2.\newline Assembly check with supplied proofs and review context; source opened and matched, not a new independent proof verification.\newline \path{records/geometry_ends/citation_checks.json}; \url{https://www.numdam.org/article/AIF_1999__49_1_287_0.pdf}\\
V136 & GPT-6 Astra, ChatGPT Pro mode\newline 2026-11-16; S\newline Grauert1960D; §6 HauptsatzI, printedp.59; §7 Satz1, printedp.60 & opened and matched. Coherence of proper higher direct images and proper analytic image theorem. The proof does not use the originally stated closed exceptional locus of Satz2.\newline Assembly check with supplied proofs and review context; source opened and matched, not a new independent proof verification.\newline \path{records/geometry_ends/citation_checks.json}; \url{https://www.numdam.org/article/PMIHES_1960__5__5_0.pdf}\\
V137 & GPT-6 Astra, ChatGPT Pro mode\newline 2026-11-16; S\newline Grauert1963ErratumD; Item2, p.36 & opened and matched. Corrects the claimed closedness of the exceptional set in the 1960 §7 Satz2. PUSH D.1.2 instead proves the needed base change over a smooth curve directly.\newline Assembly check with supplied proofs and review context; source opened and matched, not a new independent proof verification.\newline \path{records/geometry_ends/citation_checks.json}; \url{https://www.numdam.org/item/PMIHES_1963__16__35_0.pdf}\\
V138 & GPT-6 Astra, ChatGPT Pro mode\newline 2026-11-16; S\newline Levy08; Lemmas3.1--3.2; §3.1, pp.588--590 & opened and matched. Analytic complexes give topological K-classes; the local complex comparison and stabilization/\allowbreak{}descent support the virtual tangent class used in PUSH A.3.1.\newline Assembly check with supplied proofs and review context; source opened and matched, not a new independent proof verification.\newline \path{records/geometry_ends/citation_checks.json}; \url{https://www.numdam.org/article/AIF_2008__58_2_571_0.pdf}\\
V139 & GPT-6 Astra, ChatGPT Pro mode\newline 2026-11-16; S\newline Loj64; §3 Theorems1--2, pp.463--464 & opened and matched. Compatible triangulation for the specified locally finite semianalytic sets in affine space and a countable real analytic manifold; applied to the compact radial inverse image and its singular/\allowbreak{}end strata.\newline Assembly check with supplied proofs and review context; source opened and matched, not a new independent proof verification.\newline \path{records/geometry_ends/citation_checks.json}; \url{https://www.numdam.org/item/ASNSP_1964_3_18_4_449_0.pdf}\\
V140 & GPT-6 Astra, ChatGPT Pro mode\newline 2026-11-16; S\newline FG10; §5.1 Theorem5.8 and Remark5.9, p.102 & opened and matched. Virtual Euler/\allowbreak{}Chern-number and smoothing result in the algebraic scheme setting. The paper cites this only as the algebraic counterpart, expressly not as a replacement for the analytic argument.\newline Assembly check with supplied proofs and review context; source opened and matched, not a new independent proof verification.\newline \path{records/geometry_ends/citation_checks.json}; \url{https://msp.org/gt/2010/14-1/gt-v14-n1-p02-s.pdf}\\
V141 & GPT-6 Astra, ChatGPT Pro mode\newline 2026-11-16; S\newline Ehresmann52; §1 Proposition and proof, pp.154--155 & opened and matched. A full-rank differentiable surjection with compact connected fibres has differentiable bundle structure; horizontal transport identifies the compact smooth fibres in PUSH D.1.2.\newline Assembly check with supplied proofs and review context; source opened and matched, not a new independent proof verification.\newline \path{records/geometry_ends/citation_checks.json}; \url{https://www.numdam.org/item/SB_1948-1951__1__153_0.pdf}\\
V142 & GPT-6 Astra, ChatGPT Pro mode\newline 2026-11-16; S\newline Meersseman2019E; §2.1 pp.883--884; §3.2 Theorem3.4 and Remark3.6 pp.898--899 & opened and matched. The smooth integrable-structure space has connected open local Kuranishi charts and a retraction to a finite-dimensional analytic space; the local product restricts to C-infinity operators. This is the mechanism cited for PUSH E.2.24; no global bounded-automorphism hypothesis is imported.\newline S-source match with supplied proof context; no independent F or D check asserted.\newline \path{records/geometry_ends/citation_checks.json}; \url{https://www.numdam.org/item/10.5802/jep.108.pdf}\\
\bottomrule\end{longtable}\endgroup
No human has checked any part of the mathematics. No formalization was supplied. No new execution of the missing mathematical programs is claimed. Gray is not a mathematical input; the public repository link was supplied but could not be fetched. Source comparisons and build checks are recorded separately from mathematical verification.

\begin{thebibliography}{99}
\bibitem[AC26]{AC26}
{Claude (effort not recorded)}; Alp{\"o}ge, Levent. \emph{The $(3,4,\infty)$ modular family of 2-tori, completed at its three special points, is a complex structure on $S^6$}. 2026. \url{https://alpo.ge/s6.pdf}. [verified against the supplied AI transcription: Main Theorem, Theorem 3.4, Remark 3.12, Proposition 4.6, Section 9 and Section 10 were opened and matched for construction and credit; the original PDF was not supplied to this assembly. S-check by GPT-6 Astra, ChatGPT Pro mode.]

\bibitem[E26]{E26}
Engel, Philip. \emph{Complex structures on $S^6$}. 2026. arXiv:2609.38442. \url{https://arxiv.org/abs/2609.38442}. [verified: supplied full TeX opened; Section 2, Lemma 2.1, Corollary 2.2, Definition 2.3, Sections 3--7, Remark 8.1, Conjecture 8.2 and Remark 8.4 matched. Construction, geometric dictionary and question source. S-check by GPT-6 Astra, ChatGPT Pro mode.]

\bibitem[DGKXv1]{DGKXv1}
Du, Yun-Heng; Guo, Bin; Kusakabe, Yuta; Xie, Song-Yan. \emph{Oka Geometry and Deformations of Torus-Fibred Complex Six-Spheres}. 2026. arXiv:2610.00250v1. \url{https://arxiv.org/abs/2610.00250v1}. [verified: supplied version 1 TeX opened; Theorems 1.3, 1.4, 7.1 and the companion all-maps-factor theorem, tangent-cohomology and full-automorphism theorems, bare-period corollary and Appendix C matched. Load-bearing known results. S-check by GPT-6 Astra, ChatGPT Pro mode.]

\bibitem[DGKXv2]{DGKXv2}
Du, Yun-Heng; Guo, Bin; Kusakabe, Yuta; Xie, Song-Yan. \emph{Oka Geometry of Complex Six-Spheres}. 2026. arXiv:2610.00250v2. \url{https://arxiv.org/abs/2610.00250v2}. [verified: supplied version 2 abstract, introduction and section list opened, with full-text term searches; the version 1 period-identification, tangent-cohomology and Kuranishi statements are absent. Scope comparison only. S-check by GPT-6 Astra, ChatGPT Pro mode.]

\bibitem[VL26]{VL26}
Viaclovsky, Jeff. \emph{Logarithmic Transformations on the Six-Sphere}. 2026. arXiv:2610.39102. \url{https://arxiv.org/abs/2610.39102}. [verified: supplied TeX opened; Theorem 1.1, translation data and Proposition 4.7 with proof matched. Higher-multiplicity setting and non-isotriviality comparison; the coordinate dictionary remains an announcement. S-check by GPT-6 Astra, ChatGPT Pro mode.]

\bibitem[OS91]{OS91}
Oguiso, Keiji; Shioda, Tetsuji. \emph{The Mordell--Weil lattice of a rational elliptic surface}. Commentarii Mathematici Universitatis Sancti Pauli 40 (1), 83--99. 1991. \url{https://rikkyo.repo.nii.ac.jp/record/10006/files/AA00610867_40-01_08.pdf}. [verified: original PDF opened; Main Theorem table entry 49 and Section 3 equation (3.2) matched: configuration IV*, III, I1, rank one, generator height 1/6. Load-bearing height credit. S-check by GPT-6 Astra, ChatGPT Pro mode.]

\bibitem[Shi90]{Shi90}
Shioda, Tetsuji. \emph{On the Mordell--Weil lattices}. Commentarii Mathematici Universitatis Sancti Pauli 39 (2), 211--240. 1990. doi:\url{10.14992/00009994}. \url{https://rikkyo.repo.nii.ac.jp/record/10027/files/AA00610867_39-02_05.pdf}. [verified: original PDF opened; Theorem 8.6, equations (8.11)--(8.12), and the local contribution table (8.16), pages 228--229, matched. The source-run locator Table 8.9 is replaced by (8.16). S-check by GPT-6 Astra, ChatGPT Pro mode.]

\bibitem[CDP20]{CDP20}
Campana, Fr{\'e}d{\'e}ric; Demailly, Jean-Pierre; Peternell, Thomas. \emph{The algebraic dimension of compact complex threefolds with vanishing second Betti number: corrigendum and addendum}. Compositio Mathematica 156, 679--696. 2020. doi:\url{10.1112/S0010437X19007802}. \url{https://www-fourier.univ-grenoble-alpes.fr/~demailly/source_files/cdp/cdp-corrigendum-v5.pdf}. [verified: author PDF opened and Corollary 2.3 matched. Cited solely in the prescribed introduction paragraph. S-check by GPT-6 Astra, ChatGPT Pro mode.]

\bibitem[Gray97]{Gray97}
Gray, Alfred. \emph{A property of a hypothetical complex structure on the six sphere}. Bollettino della Unione Matematica Italiana B 11 (2, suppl.), 251--255. 1997. \url{https://www.bdim.eu/item?id=BUMI_1997_7_11B_2_Suppl_251_0}. [unverified: the original article could not be opened; bibliographical identification was found through the opened Angella paper. Not used for any mathematical input. Retrieval attempt by GPT-6 Astra, ChatGPT Pro mode.]

\bibitem[Repo]{Repo}
{S6 complex-structure workbench}. \emph{Repository for the supplied S6 returns and announcement proofs}. 2026. \url{https://github.com/KokunoYumeto/s6-complex-structure-workbench}. [unverified public availability: the URL is supplied by the writing prompt, and retrieval failed. Statements were matched in the attached run TeX; repository pointers use only supplied run and statement identifiers. Not an independently opened proof source.]

\bibitem[EZ82]{EZ82}
Eichler, Martin; Zagier, Don. \emph{On the zeros of the {Weierstrass} $\wp$-function}. Mathematische Annalen 258, 399--407. 1982. doi:\url{10.1007/BF01453974}. \url{https://people.mpim-bonn.mpg.de/zagier/files/doi/10.1007/BF01453974/fulltext.pdf}. [verified: opened and matched by GPT-6 Astra, ChatGPT Pro mode, 16 November 2026; unnumbered theorem p.399 and Fourier corollary p.400. Role: source and credit for the zero formula, not for the global norm maximum.]

\bibitem[DI10]{DI10}
Duke, William; Imamo{\u g}lu, {\"O}zlem. \emph{The zeros of the {Weierstrass} $\wp$-function and hypergeometric series}. Mathematische Annalen 340 (4), 897--905. 2008. doi:\url{10.1007/s00208-007-0174-3}. \url{https://www.math.ucla.edu/~wdduke/preprints/zeros.pdf}. [verified: supplied PDF opened and matched by GPT-6 Astra, ChatGPT Pro mode, 16 November 2026; Section2, equation(6), unnumbered Eichler--Zagier theorem, and Theorem1. Role: zero-formula normalization and background attribution. Its local variable 1728/j is reciprocal to this paper's t.]

\bibitem[DM20]{DM20}
DeMarco, Laura; Mavraki, Niki Myrto. \emph{Variation of canonical height and equidistribution}. American Journal of Mathematics 142 (2), 443--473. 2020. arXiv:1701.07947v2. \url{https://arxiv.org/pdf/1701.07947}. [verified: primary PDF opened and matched by GPT-6 Astra, ChatGPT Pro mode, 16 November 2026; Theorem1.1 and Corollary1.2, pp.2--3. Role: mechanism of announcement ARI C.8.10; probability normalization is separately specified in that announcement.]

\bibitem[Johansson2016]{Johansson2016}
Johansson, Fredrik. \emph{Arb: Efficient Arbitrary-Precision Midpoint-Radius Interval Arithmetic}. 2016. arXiv:1611.02831. \url{https://arxiv.org/pdf/1611.02831}. [verified: primary PDF opened and matched by GPT-6 Astra, ChatGPT Pro mode, 16 November 2026; Sections2--3 describe the midpoint-radius representation and error bounds. Role: arithmetic trust base of the supplied PUSH E.4 certificate; not a claim of a current rerun.]

\bibitem[PythonFlint2026]{PythonFlint2026}
{python-flint developers}. \emph{python-flint 0.9.0 documentation: arb, real numbers}. 2026. \url{https://python-flint.readthedocs.io/en/latest/arb.html}. [verified: official documentation opened and matched by GPT-6 Astra, ChatGPT Pro mode, 16 November 2026; arb constructor and outward endpoint methods. Role: semantics of the supplied certificate's arithmetic implementation.]

\bibitem[BS76]{BS76}
B{\u a}nic{\u a}, Constantin; St{\u a}n{\u a}{\c s}il{\u a}, Octavian. \emph{Algebraic Methods in the Global Theory of Complex Spaces}. 1976. \url{https://djvu.online/file/2hQiOnz8WuoQk}. [verified: original-book text reproduction opened and Chapter III, Theorems 4.10 and 4.12, pp. 132--135, matched. Proper analytic base change, local freeness and fibre comparison under constant fibre-cohomology dimension. Replaces the reviewed DGKX lemma citations in FAM A.5.2. S-check by GPT-6 Astra, ChatGPT Pro mode.]

\bibitem[RRV71]{RRV71}
Ramis, Jean-Pierre; Ruget, Gabriel; Verdier, Jean-Louis. \emph{Dualit{\'e} relative en g{\'e}om{\'e}trie analytique complexe}. Inventiones Mathematicae 13, 261--283. 1971. doi:\url{10.1007/BF01406078}. \url{https://www.academia.edu/5435625/Dualit%C3%A9_relative_en_g%C3%A9om%C3%A9trie_analytique_complexe}. [verified: full text of the original article opened through a third-party reproduction; introduction, p. 261, and the relative trace and duality in Section 2, pp. 264--266, matched. Bibliographic data and DOI also matched in the supplied DGKXv1 bibliography. S-check by GPT-6 Astra, ChatGPT Pro mode.]

\bibitem[RR70]{RR70}
Ramis, Jean-Pierre; Ruget, Gabriel. \emph{Complexe dualisant et th{\'e}or{\`e}mes de dualit{\'e} en g{\'e}om{\'e}trie analytique complexe}. Publications Math{\'e}matiques de l'IH{\'E}S 38, 77--91. 1970. doi:\url{10.1007/BF02684652}. \url{https://www.numdam.org/item/10.1007/BF02684652.pdf}. [verified: original Numdam PDF opened; Section 3, Theorem 1, p. 81, and Proposition 2, p. 91, matched, including the invertible dualizing sheaf for a Gorenstein analytic space and the complete-intersection case. S-check by GPT-6 Astra, ChatGPT Pro mode.]

\bibitem[Ugarte00]{Ugarte00}
Ugarte, Luis. \emph{Hodge Numbers of a Hypothetical Complex Structure on the Six Sphere}. Geometriae Dedicata 81, 173--179. 2000. \url{https://zaguan.unizar.es/record/165697/files/texto_completo.pdf?version=1}. [verified: original PDF opened; Proposition 3.1, Theorem 3.2, pp. 175--177, and Corollary 3.3, p. 177, matched. Conditional Hodge and Frolicher constraints; not a proof of the evaluated middle Hodge number in this family. S-check by GPT-6 Astra, ChatGPT Pro mode.]

\bibitem[McHugh18]{McHugh18}
McHugh, Andrew. \emph{Bott-Chern-Aeppli, Dolbeault and Frolicher on Compact Complex 3-Folds}. 2018. arXiv:1708.03251v4. \url{https://arxiv.org/pdf/1708.03251}. [verified: original arXiv version 4, dated 15 November 2018, opened; Section 7.1, pp. 24--25, and its complex-S6 Bott--Chern table matched. Background credit for formulas in terms of Dolbeault and Frolicher data. S-check by GPT-6 Astra, ChatGPT Pro mode.]

\bibitem[Angella17]{Angella17}
Angella, Daniele. \emph{Hodge numbers of a hypothetical complex structure on $S^6$}. 2017. arXiv:1705.10518v2. \url{https://arxiv.org/pdf/1705.10518}. [verified: original arXiv version 2 opened; introduction, Figures 1--2, and Section 4, pp. 13--16, matched. Conditional spectral-sequence, double-complex and Bott--Chern background. S-check by GPT-6 Astra, ChatGPT Pro mode.]

\bibitem[Stelzig21]{Stelzig21}
Stelzig, Jonas. \emph{On the Structure of Double Complexes}. 2021. arXiv:1812.00865v3. \url{https://arxiv.org/pdf/1812.00865}. [verified: original arXiv version 3 opened; Theorem 3 (also Theorem A, p. 2), Proposition 6(1), p. 16, and Lemma 8 matched. The bounded-bidegree decomposition allows infinite-dimensional components. Load-bearing for INV C2.1. S-check by GPT-6 Astra, ChatGPT Pro mode.]

\bibitem[Schnell19]{Schnell19}
Schnell, Christian. \emph{Lecture 4: The {Appel--Humbert} theorem}. 2019. \url{https://www.math.stonybrook.edu/~cschnell/mat615/lectures/lecture4.pdf}. [verified; load-bearing; opened and matched by GPT-6 Astra, ChatGPT Pro mode, 16 November 2026; Theorems 4.1 and 4.4, equation (4.3), and effective-divisor/quotient discussion, pp.1--3]

\bibitem[Buchdahl1999]{Buchdahl1999}
Buchdahl, Nicholas. \emph{On compact Kähler surfaces}. Annales de l'Institut Fourier 49 (1), 287--302. 1999. doi:\url{10.5802/aif.1674}. \url{https://www.numdam.org/article/AIF_1999__49_1_287_0.pdf}. [verified; load-bearing; opened and matched by GPT-6 Astra, ChatGPT Pro mode, 16 November 2026; Theorem 11]

\bibitem[Grauert1960D]{Grauert1960D}
Grauert, Hans. \emph{Ein Theorem der analytischen Garbentheorie und die Modulräume komplexer Strukturen}. Publications Mathématiques de l'IHÉS 5, 5--64. 1960. doi:\url{10.1007/BF02684746}. \url{https://www.numdam.org/article/PMIHES_1960__5__5_0.pdf}. [verified; load-bearing; opened and matched by GPT-6 Astra, ChatGPT Pro mode, 16 November 2026; \S6 Hauptsatz I; \S7 Satz 1]

\bibitem[Grauert1963ErratumD]{Grauert1963ErratumD}
Grauert, Hans. \emph{Berichtigung: Ein Theorem der analytischen Garbentheorie und die Modulräume komplexer Strukturen}. Publications Mathématiques de l'IHÉS 16, 35--36. 1963. \url{https://www.numdam.org/item/PMIHES_1963__16__35_0.pdf}. [verified; load-bearing correction; opened and matched by GPT-6 Astra, ChatGPT Pro mode, 16 November 2026; item 2, p.36]

\bibitem[Levy08]{Levy08}
Levy, Roni N.. \emph{Riemann--Roch theorem for higher bivariant {K}-functors}. Annales de l'Institut Fourier 58 (2), 571--601. 2008. doi:\url{10.5802/aif.2361}. \url{https://www.numdam.org/article/AIF_2008__58_2_571_0.pdf}. [verified; load-bearing; opened and matched by GPT-6 Astra, ChatGPT Pro mode, 16 November 2026; Lemmas 3.1--3.2 and \S3.1, pp.588--590]

\bibitem[Loj64]{Loj64}
\L{}ojasiewicz, Stanis{\l}aw. \emph{Triangulation of semi-analytic sets}. Annali della Scuola Normale Superiore di Pisa, Classe di Scienze 18 (4), 449--474. 1964. \url{https://www.numdam.org/item/ASNSP_1964_3_18_4_449_0.pdf}. [verified; load-bearing; opened and matched by GPT-6 Astra, ChatGPT Pro mode, 16 November 2026; \S3, Theorems 1--2, pp.463--464]

\bibitem[FG10]{FG10}
Fantechi, Barbara; G{\"o}ttsche, Lothar. \emph{Riemann--Roch theorems and elliptic genus for virtually smooth schemes}. Geometry \& Topology 14 (1), 83--115. 2010. doi:\url{10.2140/gt.2010.14.83}. \url{https://msp.org/gt/2010/14-1/gt-v14-n1-p02-s.pdf}. [verified; background algebraic counterpart; opened and matched by GPT-6 Astra, ChatGPT Pro mode, 16 November 2026; Theorem 5.8 and Remark 5.9, p.102]

\bibitem[Ehresmann52]{Ehresmann52}
Ehresmann, Charles. \emph{Les connexions infinit\'esimales dans un espace fibr\'e diff\'erentiable}. S\'eminaire Bourbaki. 1952. \url{https://www.numdam.org/item/SB_1948-1951__1__153_0.pdf}. [verified; load-bearing; opened and matched by GPT-6 Astra, ChatGPT Pro mode, 16 November 2026; \S1, Proposition and proof, pp.154--155]

\bibitem[Meersseman2019E]{Meersseman2019E}
Meersseman, Laurent. \emph{The {Teichm\"uller} and {Riemann} moduli stacks}. Journal de l'\'Ecole polytechnique---Math\'ematiques 6, 879--945. 2019. doi:\url{10.5802/jep.108}. \url{https://www.numdam.org/item/10.5802/jep.108.pdf}. [verified; mechanism of announcement PUSH E.2.24; opened and matched by GPT-6 Astra, ChatGPT Pro mode, 16 November 2026; Section 2.1 pp.883--884 and Section 3.2 Theorem 3.4 and Remark 3.6 pp.898--899]

\end{thebibliography}


\end{document}
